% =====================================================================
%  H. Brøchner
%  Den græske Philosophies Historie i Grundrids.
%  (= Philosophiens Historie i Grundrids, 1. Deel.)
%  Kjøbenhavn: P. G. Philipsens Forlag, 1873.   274 pp.
%
%  "The History of Greek Philosophy in Outline" -- TRANSLATION (English).
%
%  Translated from transcription.tex (Danish) in this directory, the text
%  of record, as it stood at commit fac71a3 (2026-10-08).  If the
%  transcription is repaired, re-check the English at every changed
%  reading (TRANSLATION-PLAYBOOK.md §6a).  The companion volume is
%  texts/brochner/philosophiens-historie-2/translation.tex; the two share
%  conventions and a glossary (vol. 2 RESUME-NOTES.md, "TRANSLATION").
%
%  STRUCTURE mirrors the transcription 1:1: the same hand-built heads
%  (\begin{center} blocks, \divrule, \dblrule, \markboth), no \section,
%  and every \opage{N} of the Danish once and in order.  Footnotes are
%  marked *) and the counter is reset at each printed page, as in the
%  transcription.
%
%  CONVENTIONS:
%   * Greek is kept exactly as printed, in \gk{} (the print's cursive
%     Greek fount); a bracketed English gloss follows only where Brøchner
%     does not translate or paraphrase it himself.
%   * Latin, German and French are kept in the language and set upright,
%     as in the print.
%   * «...» become ``...''; ɔ: becomes "i.e."
%   * Letterspacing (\emph{} in the Danish) stays \emph{} (italic).
%   * Ancient names in their usual English form (Socrates, Plato,
%     Aristotle, Heraclitus, Anaximander, Epicurus, Plotinus ...).
%   * ERRATA: the Danish is diplomatic and applies none.  The English
%     follows Brøchner's corrections -- the five of this part's own
%     "Rettelser." leaf and the thirty of the list "1ste Deel." printed on
%     the errata leaf of the second part -- each with a % comment at the
%     site giving the printed reading.  Logged printer's defects are
%     translated in their evident sense, with a % comment at the site.
%
%  STATE: DRAFT COMPLETE (2026-10-08): Preface, Contents, the whole body
%  (pp. 1--274, with the schema on p. 10) and the errata note.  Made by 26
%  translator agents in six waves from one brief and glossary (RESUME-NOTES.md,
%  "TRANSLATION"); each part checked by script against the Danish (\opage,
%  footnotes, counter resets, \emph, paragraph counts, and the contents of
%  every \gk{} character for character -- identical but for the two Greek
%  errata, pp. 55 and 71).  Ancient names normalized to English forms by one
%  scripted pass.  The translators' odd spots in the Danish are listed in
%  RESUME-NOTES for the repair session; if the transcription is repaired at
%  any of them, re-check the English there (§6a).
% =====================================================================

\documentclass[12pt,a4paper,openany]{book}
\usepackage[utf8]{inputenc}
\usepackage[T1]{fontenc}
\usepackage{libertinus}
\usepackage{libertinust1math}
\usepackage{textalpha}
\usepackage[english]{babel}
\usepackage[protrusion=true,expansion=true]{microtype}
\usepackage{enumitem}
\usepackage{amsmath}
\usepackage{setspace}
\usepackage[margin=1.2in]{geometry}
\usepackage{fancyhdr}
\setlength{\headheight}{28pt}
\addtolength{\topmargin}{-13.5pt}
\usepackage{hyperref}
\hypersetup{pdftitle={The History of Greek Philosophy in Outline (1873)},
  pdfauthor={H. Brøchner}, hidelinks}

% Footnote marks as printed: *), reset at each printed page by
% \setcounter{footnote}{0}, exactly as in the transcription.
\renewcommand{\thefootnote}{%
  \ifcase\value{footnote}\or *)\or **)\or ***)\or \dag)\or \dag\dag)\else
  \arabic{footnote})\fi}

\newcommand{\divrule}{\begin{center}\rule{2cm}{0.4pt}\end{center}}
\newcommand{\dblrule}{\begin{center}%
  \rule{2cm}{0.4pt}\\[1.2pt]\rule{2cm}{0.4pt}\end{center}}

% Greek, in the print's cursive fount (see the transcription's preamble).
\newcommand{\gk}[1]{\textit{#1}}

\pagestyle{fancy}
\fancyhf{}
\fancyhead[LE,RO]{\thepage}
\fancyhead[RE]{\textit{The History of Greek Philosophy in Outline}}
\fancyhead[LO]{\textit{\leftmark}}
\fancypagestyle{plain}{\fancyhf{}\fancyhead[LE,RO]{\thepage}}

\setstretch{1.3}

\title{\textbf{The History of Greek Philosophy}\\[2pt]
  \large in\\[6pt]
  \textbf{\Large Outline}}
\author{Dr.\ H.\ Brøchner\\[4pt]
  \small Professor of Philosophy}
\date{Copenhagen: P.\ G.\ Philipsen, Publisher, 1873\\[18pt]
  \small Translated from the Danish}

% ---- original-page markers: the Danish printed page begins here ----
\usepackage{marginnote}
\newcommand{\opage}[1]{\marginnote{\footnotesize\textit{[#1]}}}

\begin{document}
\maketitle
\thispagestyle{empty}
\newpage

\chapter*{\emph{Preface.}}
\markboth{Preface}{Preface}

\divrule

This short outline of the history of philosophy has been written
chiefly to serve as the basis for the lectures on philosophical
propaedeutics which, under the new arrangement of the philosophical
course for the youngest students, I am to give this year; but it may
also be used by those who wish, by private study, to gain a survey of
the historical development of philosophy. I have endeavored to leave
out everything superfluous to a clear survey, and to include everything
that is necessary for one.

I have had to work out this book during a long and dangerous illness,
which often left me only little strength, and prevented me from going
back to the sources again as often as I could have wished. The
conditions under which I have had to work must serve as an excuse for
some of the book's shortcomings, to which I myself have not closed my
eyes. I hope, however, that the book may be of some use to students,
and in particular that it may give them a more reliable guide than
Schwegler's compendium, used for many years, easily and attractively
written, but in a high degree faulty.

Of the earlier works on the history of philosophy that have helped me
in working out this book, I name before all others \emph{Eduard
Zeller's} excellent, comprehensive and thorough work on Greek
philosophy, which is the fruit of a rare union of comprehensive and
reliable learning, acute and sound criticism, and philosophical insight,
and has brought light to countless points in the philosophy of
antiquity, and advanced its study more than any other earlier or
contemporary work. In this part of my book I have followed Zeller in the
whole division of the material, if partly with a different motivation,
and on many points in the presentation of particulars I have been able
to keep close to him.

From my earlier book, \emph{Bidrag til Opfattelsen af Philosophiens
historiske Udvikling}, various smaller portions have been taken up into
this one, partly unchanged, partly with slight reworking. The two books
nevertheless keep an independent position beside one another. The
earlier book sought the laws of the historical development of
philosophy, illustrated the laws it found by historical examples, and
through these gained the material for a characterization of the main
historical forms of philosophy. The present book is to give a
perspicuous, connected presentation of the whole development of
philosophy, articulated according to the main forms that the earlier
book distinguished.

The second part of this book, which contains the history of modern
philosophy, will appear toward the end of the year.

\textit{Copenhagen, May 30, 1873.}

{\raggedleft\textbf{H. Brøchner.}\hspace*{1.2cm}\par}

\divrule

\chapter*{Contents of the First Part.}
\markboth{Contents}{Contents}

\divrule

{\small [The page numbers are those of the Danish original, given in the
margin of this translation.]\par}
\medskip

\begingroup
\parindent=0pt \parfillskip=0pt \tolerance=9999
\newcommand{\indholdleader}{%
  \nobreak\leaders\hbox{\hspace{0.22em}.\hspace{0.22em}}%
  \hskip 0.8em plus 1fill\nobreak}
\newcommand{\indholdtop}[2]{\hangindent=2em\hangafter=1\noindent
  #1\indholdleader\mbox{\textbf{#2}}\par}
\newcommand{\indholdsub}[3]{\hangindent=3.6em\hangafter=1\noindent
  \hspace*{0.6em}\makebox[1.8em][r]{#1}\enspace #2\indholdleader\mbox{#3}\par}
\noindent\hfill\textbf{\small Page}\par
\indholdtop{\textbf{Introduction}}{1---10}
\indholdsub{1.}{The Concept of Philosophy; Its Peculiarity as a Science;
Its Relation to Consciousness}{1---3}
\indholdsub{2.}{The Historical Development of Philosophy; the Form of the Development}{3---5}
\indholdsub{3.}{Survey of the Historical Development of Philosophy}{5---10}
\indholdtop{\textbf{First Main Division: Greek Philosophy}}{11---258}
\indholdsub{1.}{The Historical Presuppositions for the Emergence of
Philosophy among the Greeks}{11---15}
\indholdsub{2.}{The Central Thought of Greek Philosophy; Its Division
as Conditioned by It}{15---21}
\indholdtop{\textbf{\emph{Chapter One:}} The \emph{First} Period of Greek
Philosophy: The Coming-to-Be of the Doctrine of Ideas}{21---65}
\indholdsub{1.}{Survey of the First Period}{21---23}
\indholdsub{2.}{The Old Ionians: Thales, Anaximander and Anaximenes}{23---26}
\indholdsub{3.}{The Pythagoreans}{26---32}
\indholdsub{4.}{The Eleatics}{32---41}
\indholdsub{5.}{Heraclitus}{41---45}
\indholdsub{6.}{Empedocles}{45---46}
\indholdsub{7.}{The Atomists}{46---49}
\indholdsub{8.}{Anaxagoras}{50---53}
\indholdsub{9.}{The Sophists}{53---63}
\indholdsub{10.}{The Character of the First Period in General}{63---65}
\indholdtop{\textbf{\emph{Chapter Two:}} Socrates}{65---81}
\indholdtop{\textbf{\emph{Chapter Three:}} The \emph{Second} Period of Greek
Philosophy: The Unfolding of the Doctrine of Ideas}{81---194}
\indholdsub{1.}{Survey of the Second Period}{81---85}
\indholdsub{2.}{The One-sided Socratics}{85---99}
\indholdsub{3.}{Plato}{99---133}
\indholdsub{4.}{The Older Academy}{133---134}
\indholdsub{5.}{Aristotle}{134---193}
\indholdsub{6.}{The Peripatetic School}{193---194}
\indholdtop{\textbf{\emph{Chapter Four:}} The \emph{Third} Period of Greek
Philosophy: The Dissolution of the Doctrine of Ideas}{195---254}
\indholdsub{1.}{Survey of the Third Period}{195---196}
\indholdsub{2.}{The General Character of the Third Period}{196---201}
\indholdsub{3.}{Stoicism}{201---214}
\indholdsub{4.}{Epicureanism}{214---223}
\indholdsub{5.}{The Skeptics and the Middle Academy}{223---229}
\indholdsub{6.}{Eclecticism}{229---234}
\indholdsub{7.}{The Dualistic Systems}{235---243}
\indholdsub{8.}{Neoplatonism}{243---254}
\indholdtop{\textbf{\emph{Chapter Five:}} The Abiding Significance of Greek
Philosophy; Its Relation to the Thought of Later Ages}{255---258}
\indholdtop{\textbf{\emph{Appendix:}} Philosophy in Christian Antiquity
and the Middle Ages}{258---274}
\endgroup

\divrule

\opage{1}
\begin{center}
{\Large\emph{Introduction.}}
\end{center}

\divrule

\begin{center}
\textbf{1.\enspace The Concept of Philosophy; Its Peculiarity as a Science;\\
Its Relation to Consciousness.}
\end{center}

\markboth{Introduction}{Introduction}
Within the common determination of scientific knowledge, philosophy
differs from the other sciences by the scope of its object, by the
peculiarity of its method, and by its relation of consciousness to its
knowledge. Whereas the particular sciences either, like zoology,
astronomy and empirical psychology, have as their object a definite,
bounded domain of natural or spiritual reality, or, like mathematics, a
single abstract determination of what is, philosophy has as its object
the whole concrete reality as a unity, in its determinateness by the
principle of existence. It is, according to its concept, \emph{thinking
determined by principle, carried through, and coherent}. Therefore it
does not, like the other sciences, take up fundamental concepts as
presuppositions that are not proved; but in going back to the ultimate
ground of existence and letting it show itself as that which is grounded
in itself, it thereby grounds the fundamental concepts of the individual
sciences, and, in seeing them as unfolding themselves out of the
principle, it determines their place in the realm of knowledge, their
relations, and the limit of their validity. And whereas the other
sciences in their knowing turn, as it were, outward, toward their
object, but do not bend thought back upon the very relation between
knowledge and object, do not make the possibility of a real and true
appropriation in thought of the object, the given, into a problem, and
therefore always lay themselves open to an attack from the doubt that
denies the possibility of that appro\opage{2}priation, philosophy in its
knowing undertakes, as it were, a doubling of consciousness, whereby a
consciousness of our consciousness is gained, a knowledge of our
knowledge, of its task and its fundamental relation. But this
potentiated consciousness, which with one glance sees itself, its
object, and its relation to the object, is gained definitively precisely
only by going back to that in which our thought and the existent meet,
to the principle that unites the fundamental oppositions of existence
and is at once the ground of spirit and of nature.

Such a universal knowledge as philosophical knowledge, characterized
according to its ideal goal, is grounded in the most original
fundamental relation of human consciousness itself, and is present as
an unconscious demand in the first elementary movements of the thinking
human mind. If we seek the reason why man, as he stands over against the
object in the relation of consciousness, is not content with the
phenomenon as it is immediately given, does not lose himself, as if
staring, in the object, but seeks to convert the given into something
determined by reflection, into something conceived, and thus to
dissolve it as given, we find the reason in this: that mind, which is
consciousness, is so only because it is self-consciousness, and is
driven by its unity to set the content of self-consciousness and of
consciousness in unity. But this can happen only as the self-conscious
mind finds itself again in what is the object of consciousness, and this
finding-again in turn becomes possible only as the object is transformed
into thought and what exists becomes, in this transformation, the image
of mind. And since the unity of mind contains this demand for an
abolition of the separation between the object of self-consciousness
and that of consciousness, it contains at the same time the demand that
the latter in its transformation form a unity, that the fragmentary and
scattered, and the contradiction inseparable from it, be lifted, and
that real existence, as explained in thought, close together into a
unity, united in the last instance by an all-grounding and thus for the
first time truly grounding principle. Only in such a
\opage{3}knowledge, in which every object of consciousness is
transformed into a determination of self-consciousness and the
foreignness of the object is lifted through the knowledge of its
essence, is the contradiction removed between the inner infinity of
self-consciousness and its limitation by a consciousness of something
foreign, between the subjective form and the objective content of
consciousness, between the I's knowledge of itself as a transparent
unity and as determined by something that is in the I and yet different
from the I. The drive of knowledge that stirs in the human mind from the
very first can be satisfied only by the annulment of these
contradictions, that is: by a \emph{universal} knowing, and this is one
with philosophical knowing, determined according to its ideal goal.
Philosophical knowledge thus issues with necessity from human nature,
from the drive that marks the essential peculiarity of that nature, and
it has its goal set by what lies originally in the determinateness of
that being.

\divrule

\begin{center}
\textbf{2.\enspace The Historical Development of Philosophy; the Form of the Development.}
\end{center}

In man's essence as mind lies the necessity of a historical
development, through which mind, as the free mind, acquires by its own
activity what was originally laid down in it as a gift, as a natural
endowment. Only what is mind's own deed is the real property of the free
mind; only through a free development does mind become itself. This
holds also of mind as knowing mind. Even if self-consciousness and
thought are the fundamental determination in the human, that which is
the distinctive mark of essence of the human, man must nevertheless
reach self-consciousness and thought through a development, and a series
of stages of development must be traversed by mankind before the
universal form of thinking that we designate as the philosophical can
find its beginning, from which a new, infinite development is then set
for it as law. Only when the concept of a unity of existence in its
principle \opage{4}comes forward for thought is the condition given for
the possibility of philosophy, of universal knowledge. But that concept
can come forward only where the human mind has won the consciousness of
its freedom. For without this inner liberation of mind, which includes
its instinctive confidence in its own power, thought would not be able
to set itself the universal task: ideally to master the totality of the
phenomenon in its infinitely changing manifoldness by finding its
grounding unity, and through the knowledge of this unity to reach that
relation to itself in its object which is precisely a relation of
freedom. Only on the presupposition of the freedom of mind, therefore,
can a philosophy be conceived; but the freedom of mind first comes to be
where its consciousness of being higher than the natural, of being
master over it, begins to dawn.

Thus the historical \emph{beginning} of philosophy is conditioned; and
already in the determination given above of the peculiarity of
philosophy there lies an intimation of the form of its historical
\emph{development}. Since philosophical knowledge strives to lead the
whole given reality back to its grounding principle, and is
philosophical precisely by this striving, philosophy stands from the
beginning in a necessary relation to the thought of the
\emph{Absolute}, to the thought of the \emph{principle} of the
\emph{totality} of existence. But since the principle is the principle
of an existence in which essential fundamental oppositions, distinct
fundamental forms, appear, it is, according to its concept, a full unity
of opposites, and must as such come to be for philosophical
consciousness in the course of the latter's deeper and deeper
penetration into the knowledge of the essence of existence. But in this
advancing knowledge the different positions that the fundamental
moments of the unity, which from the beginning of philosophy, at least
in the form of a dawning thought, are objects for consciousness, take up
toward one another precisely through that very advance, must with
necessity condition essentially different conceptions of the principle,
the Absolute, and these different conceptions must ground a
\opage{5}distinction of main forms of philosophy in its historical
development, and determine the fundamental character of these main
forms. The sequence of the historical main forms of philosophy must be
determined by the essential relation between the forms of existence:
spirit and nature, the fundamental moments in the thought of the
Absolute, and by what was stated above as the condition for the
coming-to-be of philosophy. A lawfulness grounded in essence will thus
determine the historical course of development in its great divisions,
and within each single main form of philosophy we shall see the same
lawfulness govern the development in its particulars. The principle of
the main form will show itself as its law, as its inner impelling and
moving energy of development, as that which conceals within itself, as
its conceptual moments, the principles of the individual systems it
comprises. In the principles of the individual systems the principle of
the whole unfolds itself; it grows, culminates and dissolves through
them, and, while governing them unconsciously to the thinkers, it
reveals the more comprehensive by which it is itself determined and
which governs its own unfolding and dissolution. Thus the single main
division is understood in its articulation from one central thought,
which governs the division as a whole in its unfolding, and by the same
fundamental thought it is fitted into the law-determined
\emph{total} development of human knowledge.

\divrule

\begin{center}
\textbf{3.\enspace Survey of the Historical Development of Philosophy.}
\end{center}

The historically given forms of philosophy gather essentially into two
distinctively formed main divisions: \emph{Greek philosophy and the
philosophy of modern times}. Greek philosophy is the beginning of
philosophy; for only in the Greek spirit's consciousness of itself as
free was the condition given for a knowledge philosophical in the
stricter sense, which had been lacking among the older civilized
peoples. In this first main form of philosophy thought expressed the
\opage{6}principate of the spiritual in the form in which this first had
to appear: in the form of the immediate unity of beauty of nature and of
the spirit that governs and ensouls nature. Since the dominion of spirit
reveals itself to intuition in the harmony of beauty, the unity is still
in the form of \emph{nature}, in the \emph{immediacy} of harmony. The
intuition of the universe as a world-work of art is the intuition in
which spirit finds rest; the universe becomes the cosmos ensouled by
ideas, the immediate unity of form and matter. But in this Greek
fundamental view there lay hidden an opposition which the Greek spirit
could not master, and which had to lead to the dissolution of the first
philosophical main form and to the development of a new form of
spirit's world-view. In the intuition of the world-work of art the union
of the opposites, of form and matter, spirit and nature, was a merely
immediate one; but since the union is only the immediate union of what
is not posited through one another, i.e.\ of what is originally distinct
in essence, the opposition is not in truth overcome and reconciled, but
breaks forth unavoidably as soon as concept-determining thought, in its
clear development, distinguishes what for intuition, and for thought
still resting in intuition, was unseparated in the unity of the work of
art. But when the separation between the fundamental opposites then
comes forward in the form of the concept, the thought that has
distinguished the opposites cannot unite them again so long as it does
not wholly abandon the ground of that world-view; and this it will not
be able to do as Greek thought, because as such it has a
nature-determined limitation of essence that binds it to that form of
unity: to the immediate unity of beauty, which expresses the peculiar
determinateness of essence of Greek thought, bound as it is to
intuition. It must therefore end in fruitless attempts, once the
consciousness of the discord, of the dualism of the opposites, has been
awakened, to find a unity again. The goal after which it strives without
being able to reach it must become: to assert that among the moments of
existence which in the Greek world-view had been conceived as the
highest and as the true \opage{7}reality, the spiritual, as the only
one, as all reality, and thus to develop a spiritualistic view of
existence, to which, however, only another spirit of the age could give
such carrying-through as is possible at all.

In opposition to the distinctively Greek world-view, which expressed
the immediate dominion of spirit over the natural ensouled by it, such a
spiritualistic world-view must conceive the spiritual, in its
distinction from the natural determined as essenceless, as alone valid.
This conception becomes the distinguishing mark of the span of time that
follows classical antiquity: of the Christian Middle Ages. But on the
basis of this nature-denying world-view, as will be shown more closely
later, no independent and distinctively formed philosophy could arise.
As nature was given up, thought lost its basis and its incitement to
development, that in which it has its confirmation; withdrawing into
pure spirituality and denying an essential side of the human, it got
only a contentless, merely negative ideality, from which no unfolding of
thought could issue. What has been called a Christian philosophy shows
itself, in its development among the Church Fathers and the
Scholastics, as a dependent and unprincipled external linking-together
of a content of thought taken up from the thought of antiquity and
belonging to a foreign, natural world-view, and of a Christian content
of faith, inaccessible to thought, incomprehensible: two elements which,
by virtue of their difference in essence, can never be united into a
real unity or a whole of thought. The ``Christian philosophy'' itself
becomes conscious, in Scholasticism, of the inner strife between its
heterogeneous constituents and their heterogeneous principles, and
through the acknowledgment of their heterogeneity pronounces judgment
on itself as philosophy.

Only at the close of the Middle Ages, when Scholasticism had dissolved
itself through its inner contradiction, and when the renewed
acquaintance with antiquity had again brought the spirit, prepared for
it by historical development, nearer to the
world-\opage{8}view that had been peculiar to Greek antiquity, did the
development of a new main form of philosophy become possible. The
removal from nature had made the thought of the Middle Ages unfruitful;
the rehabilitation of nature and of nature-grounded, autonomous human
thinking was the condition for the rebirth of philosophy, and the second
main form of philosophy, the philosophy of modern times, comes forward
independent and freely formed after thinking, through a renewed,
intimate attachment to antiquity, had in a transitional time gained
strength and content, and developed a wealth of forms which, with all
the energetic, youthfully fresh striving that asserts itself through
them, were nevertheless still bound by the past, and had not brought to
clarity an independently formed principle or been worked through to
inner consistency.

The characteristic peculiarity of the \emph{philosophy of modern times}
is its striving for a unity of thought, proceeding from the
presupposition of the absolute difference between the fundamental forms
of existence: nature and spirit. This peculiarity was conditioned by
what historically preceded it. By the view of life of the Middle Ages
the sharp dualistic separation had been posited between spirit, the one
true, and nature, the essenceless and untrue. The systems of the
transitional period, the age of the philosophical Renaissance, had
fought against this dualism, but, bound as they were to the fundamental
form of Greek thinking, and unconsciously influenced by the Middle Ages,
they had not been able to overcome it. They left this task unsolved to
modern philosophy, and with it the unconscious after-effect of the
presuppositions of the Middle Ages. The relation between the
fundamental moments could therefore, for modern philosophy, no longer
be, as with the Greeks, that of immediate harmony, which conceals the
discord present at bottom; the relation is from the beginning that of
the opposition between spirit, recognized in its essentiality, and
nature, restored in its essentiality. This opposition is to be brought
to unity, but the unity can, precisely because of the presupposition
lying in the \opage{9}starting point, become only a one-sided unity, one
in which one of the sides of the opposition has vanished, and the
remaining one, which seems to comprise everything, in reality gets only
the determination it had within the opposition. In the very kind of the
unity, then, the unreconciled opposition must show its influence. The
end point of the development must in this philosophical main form, too,
reveal the defect hidden in the starting point. Whether it ends in
spiritless nature (materialism) or in natureless spirit (Hegel), the
one-sidedness is the same. In both cases thought's striving for unity is
influenced and hampered by dualism. There is a striving toward the goal:
to unite the two positive fundamental opposites of existence into a
true unity, rich in content; but the thought of the unity, of the
highest, all-comprehending essentiality, is determined only under the
hidden influence of the dualistic presupposition, and the conception of
it therefore becomes one-sided and inconsistent. In both cases,
therefore, modern philosophy fails in the solution of its task, which
was: to overcome dualism; and in leaving this task unsolved it
intimates, albeit only in abstract and negative form, the goal for the
\emph{philosophy of the future}, which, according to the relations of
the concepts, becomes the last possible main form of philosophy, but in
its principle includes the condition for an unbounded development in
time.

In the first main form of philosophy the relation between the
fundamental opposites in existence was determined by immediate unity, in
the second by abstract opposition. After both conceptions have proved
impracticable, only one relation is still thinkable: that of the unity
of opposition. In order to solve its task, the philosophy of the future
must recognize the essential peculiarity of spirit and of nature, of the
ideal and of the real, and see their difference of essence as an
opposition, but as an opposition that is grounded in, and closed
together into unity in, the ideal-real principle determined as absolute
subjectivity. ---
\opage{10}Two main forms of philosophy thus lie before us historically
as wholly concluded, and by the fundamental problem they have posed but
not solved they point toward a third main form, whose development, which
by the abstract formula itself is intimated as conditioned by
philosophical thought's appropriating and assimilating to itself the
sciences that relate immediately to the real, belongs to the future, the
whole future.

Perspicuously, the fundamental relation between the three main forms,
in which the completeness of the division lies, can be stated in the
following schema, which shows the fundamental conception of the Absolute
in the different main forms, and the fundamental relation they posit
between the fundamental opposites of existence.

\begin{center}
1.\enspace\textbf{Greek Philosophy.}\\[2pt]
\emph{The Idea revealed in the cosmic work of art} (Kosmos).\\
(\emph{Immediate} unity of nature and spirit.)\\[6pt]
2.\enspace\textbf{The Philosophy of Modern Times.}\\[2pt]
$\overbrace{\hspace{0.67\linewidth}}$\\[2pt]
\begin{tabular}{@{}c@{}c@{}}
\makebox[0.5\textwidth]{\emph{Nature} = matter =} &
  \makebox[0.5\textwidth]{\emph{Spirit} (I --- absolute spirit) =}\\
\makebox[0.5\textwidth]{non-spirit.} &
  \makebox[0.5\textwidth]{non-nature = the non-material.}
\end{tabular}\\[2pt]
(\emph{Abstract opposition} of nature and spirit.)\\[6pt]
3.\enspace\textbf{The Philosophy of the Future.}\\[2pt]
\emph{Absolute subjectivity} as \emph{ideal-real} principle of
existence.\\
(\emph{Unity-in-opposition} of nature and spirit.)
\end{center}

\dblrule

\opage{11}
\begin{center}
{\Large First Main Division.}
\end{center}

\divrule

\begin{center}
{\Large\textbf{\emph{Greek Philosophy.}}}
\end{center}

\divrule

\begin{center}
\textbf{1.\enspace The Historical Presuppositions for the Emergence of Philosophy\\
among the Greeks.}
\end{center}

\markboth{Greek Philosophy}{Greek Philosophy}
When philosophy appears among the Greeks at the boundary between the 7th and
the 6th century before Christ, the Greek people had already run through many
stages of its history and, in its political and social conditions and in its
spiritual productions, had brought about the conditions that were necessary
to prepare a philosophical thinking.

Thales, the father of philosophy, is preceded, within the span of the Greek
people's history known to us, not only by the time of which we have the ideal
image in the Homeric songs, the people's heroic youth, in which myth and
history meet, just as gods and men meet in combat and in love; but also by the
period of the migrations of the Greek tribes, in which the peoples were
mingled, the old patriarchal forms of state were overthrown, and the Greek
peoples spread out beyond the borders of Greece proper and came into immediate
contact with the older culture of the Orient.

At the beginning of the 6th century we find among the Greeks not only a
material development, by which the external conditions of life were secured,
and conditions were thereby brought about under which thought could have
otium; but we find as well, \opage{12}during the transformation of the
constitutions, which had already in part led to republican forms of state, an
exertion of force and a free development of force in the sphere of civic life
which could not but hasten the development of the whole spiritual life; and we
find finally, through the change of the tribes' homelands and especially
through colonization, a ferment brought about in minds which was bound to
become of decisive significance for the life of consciousness. It is not in
Greece itself, but in its colonies on the Asiatic coast, that the first
philosophical attempts were made. This head start of the colonies was nothing
accidental; it was conditioned by the very breaking-away from the accustomed,
from what in the homeland had through long habit acquired, as it were, a
nimbus of holiness, such as tends to surround what is venerable by age and
inherited from the fathers, and by the reflection that was necessarily bound
to be awakened by the comparison between the native and the foreign, by the
adaptation of the old customs and forms to the new, altered conditions.

And while favorable conditions for the breakthrough of philosophical thinking
were thus brought about through the political conditions, it was more directly
prepared and supported through the development of the spiritual life of the
Greeks, and in particular through definite forms of the expression of
spiritual activity which aim in the same direction as philosophical inquiry
proper. As such we single out the oldest cosmological attempts, the religion
of the mysteries, and the incipient ethical reflection.

In the oldest cosmological attempts, such as they were made by Hesiodos,
Pherekydes, Epimenides, Akusilaos and the Orphics, we find as it were the
harbingers of a philosophy of nature, attempts which were bound to determine
the direction of thinking and to set it the task of leading the multiplicity
of natural existence back to its first ground. And since these attempts, which
agree in this comprehensive leading-back, differ from one another not only in
the details of execution but in the \opage{13}place they give to the most
perfect, the highest and the mightiest: as the conclusion of the world's
development or as its beginning, they thereby also bring essential problems
before comparative thinking.

In the religion of the mysteries there lie incitements for thinking in the
doctrine of the unity of the divine, though it is expressed in vague and
indeterminate forms, in the views of nature that were laid down in the sacred
legends, and in the belief in an immortality, chiefly in the form of the
transmigration of souls, which stood in connection with the conception of
man's moral vocation.

In the incipient ethical reflection among the Greeks, which meets us already
in Homer and Hesiodos, but acquires a more determinate and more developed form
in the older lyric poets and in the gnomic poets of the 6th century,
contemporary with incipient philosophy, and is concentrated in pregnant
sayings in the ``Seven Sages,'' we find not only in general a preparation of
thought for a sharp and sound conception of essential relations in human
existence, but we find a progressive development, through which a unified view
of life more and more begins to work its way forward, and the notion of the
divine as moral power is purified, and the thought of an ethical retribution
is clarified and refined.

But what thus shows itself as essential preparations and effective
incitements for philosophy is nevertheless, on the other hand, still
qualitatively distinct from it. Ethical reflection is, even in its clear,
sharp form of understanding, still the incipient, fragmentary reflection;
despite all its subtle observation, it does not consciously lead its
individual sayings back to general principles; still less does it ground them
theoretically by going back to the nature of man and to the essence of moral
activity in general, or seek to comprehend the ground of the moral. And the
moral teachings communicated in the mysteries, essentially wrapped in
religious forms of notion, have the same defect. \opage{14}As regards the
doctrine of the religion of the mysteries concerning the Godhead, and the
cosmogony's images of the coming-to-be of the world, identical with the
becoming of the gods, these do indeed contain germs of thought and incitements
to thought; but thought is still either only dawning and indeterminate, or it
is bound by imagination, governed by its power. In the religion of the
mysteries the undeveloped thought of unity only glimmers forth as an
indeterminate intimation, surrounded by the motley multiplicity of the forms of
imagination; and in the cosmological poems, which let existence in its genesis
unroll itself before the eye, it is not the principles of reality determined by
thought, real powers of nature, that come forward as working lawfully, but
creatures of imagination, which, in magical fusion with reality, borrow name
and shape from it; and it is not the essence of things that is sought, not
their natural grounds, not a comprehension of the law-determined real, but
only a childlike, intuitive notion of the coming-to-be of the all, formed in
the shape of narrative, through which thoughts and reflexes of experience may
indeed shine, but which essentially, in its uncertain double light, is in the
medium of the poetizing imagination, which transforms and shapes arbitrarily
according to the demands of the heart, and is for the imagination.

At the moment philosophy appears, seeking beyond the fragmentary toward the
whole, the sphere in which the mind, in its direction toward totality,
formerly moved is transformed. Philosophy separates out in principle the
fantastic, which becomes the organ of poetry and mythology; it is not content
with indeterminate intimations of thought, wrapped in the forms of
imagination, but moves in the element of clear thought, in which it will take
up the real as real. And as, seeking the true nature of the case, it asks
about the universal ground of this real, it directs itself toward the natural
(real) causes of things, toward the powers of reality. When this question
about the natural causes of things, causes valid in reality, is posed in
universal form, philosophy has begun. It comes into being when thought is
consciously directed toward the whole of what is, in order to explore
\opage{15}the natural genesis of things, and so to bring about a coherent,
universal knowledge, such a knowledge as can only come into being out of the
knowledge of the principle.

\divrule

\begin{center}
\textbf{2.\enspace The Central Thought of Greek Philosophy; Its Division as\\
Conditioned by It.}
\end{center}
% errata (Rettelser leaf), p. 15 l. 4: ";" for ":" after "Centraltanke"; printed "Centraltanke: dens Inddeling"

When we are to seek the thought that is to be determined as the central
thought of the Greek spirit, as the thought of Greece, we are led to find it
through a reflection that offers itself as if of its own accord. We find among
the Greek people, as the form for the intuition of the divine, as the ideal of
the human being, and as stamped in all the forms of life, as the fundamental
character pervading its whole existence: beauty, harmony. But what
characterizes the spiritual self of an entire people and designates its
formative principle, that philosophy, as the people's ideal self-knowledge at
its highest stage of development, must become conscious of and seek to
express, when the people has undisturbed developed an independent thinking,
essentially belonging to this people, and above all when the people has
developed as naturally and harmoniously as
% printed "harmonisk" (defect logged in the transcription)
the Greek people. And what is the formative, governing principle of human
existence as grasped by thought, that philosophy must, since the human mind
must posit itself as the goal of existence, strive to grasp as the determining
principle of all existence. Therefore Greek thinking at its height, in seeking
to grasp the whole of existence, had to grasp it from the standpoint of beauty
as the \emph{all-embracing world-artwork}, as \emph{cosmos}, and in this
highest intuition of the Greek spirit the thought of the \emph{Idea} had to
come into view as the thought of the ground of the form of beauty, of the
absolutely highest, of what truly is.

In analyzing the intuition that was transferred by the Greeks to the whole of
existence in its harmonious being, \opage{16}we are led to the distinction of
the fundamental moments whose inner relation it was the task of Greek thinking
to determine, and this relation contains what is normative for the
development of Greek philosophy and the germ of its dissolution.

If, in seeking to analyze the work of art, we take our point of departure from
the conception of it as the individual, limited work of art, and make the
intuition still more determinate by thinking of the work of art most
characteristic of the art of Greece: a plastic work of art, a statue, e.g.\ of
marble or bronze, then it appears that this work of art \emph{is} what it is
by its \emph{form}, that the matter has been absorbed into the form and is in
the work of art only as formed matter, so that it has there no independent
validity and being beside the form. It appears that the form, which expresses
the ideal significance, is not only what completes, what ensouls and
inspirits, but that it is what gives the work of art \emph{being} by giving it
the being of \emph{beauty}, i.e.\ \emph{true} being; for beauty is here truth
and reality. But as long as we think merely of the individual limited work of
art, the matter nevertheless constantly acquires a being that is relatively
% printed "bestandig" (defect logged in the transcription)
independent of the form, because it is relatively independent of the work of
art. The marble and the bronze also have a being as unformed material before
it is fashioned and formed by the artist, and the same form can find its
expression in both materials: the connection is not indissoluble. If, on the
other hand, we widen that intuition and exchange the notion of the individual
limited work of art for the notion of the all-embracing one, of the
world-artwork, of the universe as cosmos, then matter no longer acquires this
indifferent being; for the universe, for which the form of beauty is the
\emph{essential} form of existence, eternally comprehends in its unity all
matter, and form as totality. Matter in general then becomes, since it can no
longer find any place beside the all-embracing form, i.e.\ the form that
embraces and forms all matter, only something that is by virtue of the form,
of what ensouls the world-artwork. And this is to say in other words: then the
\opage{17}principle of form becomes the principle of being and of reality;
matter becomes what, torn loose from form, is a non-being, and acquires being
and reality only in unity with form. But form is the incorporeal, the ideal;
it is as it were of the kindred of spirit, and by form the intuiting Greek
thinking, for which the concept of form and that of the forming went together
into unity, understood what modern thinking understands by spirit, while
matter becomes the corresponding expression for nature, thought of as matter.
In the determination of form as the highest and ruling, the primacy of spirit
is thus pronounced, and the connection of form with spirit is stated expressly
by Greek philosophy itself. Where Aristotle reaches the concept of the highest
form, of form in its matter-free purity, the concluding point of his system,
at which Greek thinking as it were for a moment rises above itself: there form
changes its name and is called \gk{Νοῦς} [Mind], and this \gk{Νοῦς} is
determined as the energy of the divine thought that uninterruptedly thinks
itself, as the eternal \gk{νόησις τῆς νοήσεως}.

If, then, the world-artwork was the highest intuition of the Greek spirit,
then that one of the moments of the work of art which is the principle of
being becomes its highest content, and the expression in thought of this
moment becomes its highest thought. But the expression in thought of this
moment, which was grasped not by an abstract but by an intuiting thinking, was
the Idea, the intelligible primal form and ground of reality of sensuous
existence, independent of space and time and becoming.

The thought of the \emph{Idea} thus becomes the central thought which must
govern the development of Greek philosophy, since, through its relation to its
opposite: matter, which in the world-artwork is joined with it into
harmonious unity, it determines and characterizes the periods of the
development, and through the same relation conditions the insurmountable
limit of Greek thinking and makes its dissolution inevitable.

In the intuition of the work of art as the unity of form and matter there is,
as already indicated above, given an \emph{imme\opage{18}diate} union of two
moments, of which neither is posited by the other. Matter could not be
determined as posited, produced by form; for such a positing, which, given the
determination of the ideal as artistic \emph{form}, could not be seen as an
original relation of essence in the Idea, as necessary according to the very
concept of the ideal, would have had to be conceived by the Greeks as a
creating. But this notion was not only constantly foreign to the Greeks, even
to the Neoplatonists, but it had to be so by the nature of the fundamental
intuition, the work of art; for in the work of art matter, i.e.\ what
\emph{is to} be formed, is the \emph{presupposition} for form's being able to
come into activity. Form, as the form of the work of art, in its immediate
being for intuition, has immediately its matter, which it forms. And when the
unity is loosened by thought, and form as Idea separates itself from matter,
then it presupposes, --- not in the sense of time but in the sense of the
concept --- in its resting being outside matter, precisely matter as that to
which it, as the forming, \emph{is to} give form, even if in this separation
matter cannot be held fast by thought as a reality. But if matter thus could
not be thought as posited by form, it appears still more readily that form
could not be thought as posited by matter. For form is not only the higher and
ruling, but according to its concept immediately what is untouched by
coming-to-be and becoming, the eternal. The unity of the moments has therefore
not issued from the one's being grounded in the other; neither of them is the
producing that develops the other out of itself. The unity is the one that is
immediately present for intuition, but the immediate unity of what has not
shown itself as proceeding from one another is an \emph{immediate unity} of
something \emph{originally different}, which for thought must again reveal
this difference and be separated in unreconciled opposition. It is true that
according to the concept form is the ground of all true being; no being and
reality can be thought outside it; but since the form of beauty, in order to
be as forming form, in order to give itself reality as such, needs a
presupposed matter, \opage{19}this matter, which in its separateness from form
had to be comprehended as a \gk{μὴ ὂν}, nevertheless constantly acquires,
through this separation unavoidable for thought, a kind of being of its own,
of presupposed independence: it becomes the contradictory existing non-being.

Through the stated relation between form and matter the course of development
of Greek philosophy is now determined, and at the same time the problem is
indicated on which this form of thinking
% printed "lige" (defect logged in the transcription; site as placed in the transcriber's note)
was bound to break.

As Greek thinking seeks to find its distinctive world-view and to know the
highest principle of existence, and in this striving is guided by what is its
most distinctive essence, it must begin with \emph{matter} as that one of the
two moments which in its immediate being must be \emph{presupposed} in order
that form can act, and through which the forming must reveal itself, in which
it must be sought. It must begin by positing matter, the natural, as the
principle, but in such a way that from the beginning the material is as it
were transformed and betrays something ideal, and that thought, by the very
problems that arise when existence in its concretion is to be explained by the
material principle, must be driven forward to discover the other moment:
form, the \emph{Idea}, \emph{spirit}. When matter is to be conceived as
principle, it acquires the character of a universal, and in this a
determination of ideality lies hidden, since the determinateness, the
particularity, that conditions the actual being of matter as this matter must
vanish under that conception, and since the notion of a chaos as an existing
universal matter --- a notion that is insufficient to explain the becoming of
the manifold, i.e.\ of the manifoldly formed --- cannot be maintained. Matter
separated from all form, intuited as chaos, vanishes into a nothing; it
acquires only a semblance of being by being represented as the efflux of the
forms, thus still as having a last remnant of form in it. The determination of
the principle as matter will therefore unavoidably advance toward more ideal
determinations: principles such as water, air, unlimited and indeterminate
matter will be succeeded by principles such as \opage{20}number, being
(\gk{τὸ ὄν}), both conceived as the matter of which things consist, as the
substance of bodily existence; and as the problem of the coming-to-be of the
manifold from the principle, which is necessarily bound up with the problem of
the determination of the principle, comes forward and asserts itself as the
essential one, the concept of \emph{becoming} becomes that which, while it
appropriates unexplained the domain of the sensible, refines the concept of
being, which now seeks another domain, into the Idea. And as soon as the
concept of the Idea is found: the concept to which matter itself pointed as to
its completion, the Idea is posited as the true, essential reality, and the
task of thought becomes: to seek an expression for the unity of the two
moments in which this presupposition is held fast. But when this task comes
forward, it presupposes that thought has consciously separated the two
moments which for intuition are united in the unity of the work of art, and
the kind of this unity makes impossible, as was touched on above, a new union
by Greek thought, for which the immediacy of harmony was the only form of
unity. The immediacy of beauty does not tolerate the separating power of
thought; therefore thought brings discord into that on which it rests, and
through this discord this form of thought must perish. The moments, which in
their separation both compel a determination of being for themselves, stand
dualistically against each other, and Greek philosophy must end in the
fruitless endeavor, after the consciousness of the essential difference of the
moments has been clarified, to find once more a reconciliation.

What has been set forth implies that the development of Greek philosophy must
run through three main stages. The \emph{first} will be the stage at which the
principle of all was sought in nature, in matter; the \emph{second} the stage
at which the concept of the Idea had been found, and the unity in what is was
sought from the presupposition that the Idea was the true reality; the
\emph{third}, finally, will be the stage where the unity was sought from the
clearer or dimmer consciousness of the discord of the moments. If we seek
names for these three periods that can \opage{21}designate their relation to
the central element that is the moving power of the development, we can name
them \emph{the becoming of the doctrine of Ideas, the unfolding of the
doctrine of Ideas, and the dissolution of the doctrine of Ideas}.

\divrule

\begin{center}
{\Large\textbf{Chapter One.}}
\end{center}

\divrule

\begin{center}
{\Large The First Period of Greek Philosophy:\\
The Coming-to-Be of the Doctrine of Ideas.}
\end{center}

\divrule

\begin{center}
\textbf{1.\enspace Survey of the First Period.}
\end{center}

\markboth{The Coming-to-Be of the Doctrine of Ideas}{The Coming-to-Be of the Doctrine of Ideas}
The first period of Greek philosophy begins with the first question
of incipient philosophy, a question grounded in the striving of mind to
bring the content of consciousness in its manifoldness under the unity
of self-consciousness: with the question, what is that which becomes?
It asks about the abiding One in the changing, restless manifold. The
answers to this question advance rapidly: when the Ionians name a
sensuous matter, a determinate element, or the unlimited, indeterminate
material, the Pythagoreans name number, the measure of the material,
the Eleatics what is in its unity; and in these determinations, which
already bring with them a circle of essential concepts and of decisive
problems, they unconsciously intimate the three spheres into which, for
developed Greek thinking, existence divides itself.

In the question: \emph{what all is}, lies the second question:
\emph{how has this, which is what is in all, become the manifold}? ---
in other words: \emph{how has all come to be}, how is the manifoldness
in existence to be explained? This question comes forward as a
\opage{22}secondary question in those oldest philosophical attempts,
until the Eleatics determine being in such a way that by their
determination they cut off the question of becoming, without, however,
being able to banish becoming from sensuous notion, through which it,
as it were, again steals for itself a being. But against this being,
which cannot annihilate its adversary, Heraclitus then sets becoming
itself, the paraphrase of coming-to-be and of change, as the ground of
explanation of the manifoldness of existence; and this becoming,
intuited immediately as sensuous matter, as fire, becomes not merely
the principle by which, but that out of which, all is; it becomes force
and matter, and bears within itself, during the ``way downward'' and
the ``way upward'' of its transformation, as the only abiding and fixed
thing, what the Greek had to demand: the lawful, harmonious order of
change. But the ground of becoming was not explained, since becoming
was conceived not as what is comprehended by a being, grounded in a
being, but as its own presupposition, as the tautological paraphrase of
coming-to-be; and just as little were its rule and law explained,
because becoming, which is seen only as the becoming of the sensuous,
does not explain the indwelling measure, what abides in change. There
was therefore set for thought the new task: to demonstrate the ground
and the necessity of becoming; and as this task is taken up by
philosophers who know the doctrine of the Eleatics, it acquires the
more determinate form: to seek the ground of becoming in the underlying
itself without the concept of being being given up. And the attempt to
solve this task then becomes at the same time an attempt to find a
mediation between the fundamental concepts of the Eleatics and of
Heraclitus, which leads to a transformation of these. Empedocles, the
Atomists, and Anaxagoras form the group of philosophers that takes up
this task, and the last of them, in the attempt to solve it, touches on
the concept of that which is higher than the material: \gk{Νοῦς}
[Mind], which, however, he cannot yet with clarity distinguish
qualitatively from the material, or, while he ascribes to it the order
in the universe, maintain a form of activity for it other than the
production of mechanical motion. The attempts at mediation
\opage{23}failed, since they could neither maintain being nor explain
becoming; and they had to fail because the opposed principles were to
meet on the same domain of existence: within sensuous existence. And
as the attempts at mediation failed, the two principles had again to
separate themselves abstractly from one another; and in the carrying
through of them in this sharp separation, in which the Eleatics'
principle of unity, by making being determinationless, hence equal to
nothing, renders a knowledge of it impossible, while Heraclitus's
becoming leads to the same result, the annihilation of knowledge, in
that it lets all abiding determinateness vanish, the oldest form of
Greek philosophy is dissolved and a higher form of development is
prepared. This dissolving activity, which prepares indirectly, through
the negative, falls to Sophistic, which, as a philosophical phenomenon,
represents the dissolution of the older philosophy through the
one-sided holding fast and dialectical carrying through of its opposed
fundamental concepts.

The many different philosophical attempts within the first period thus
gather themselves into four groups: one whose fundamental thought is
being; one whose fundamental thought is becoming; one that seeks to
mediate between these opposites and to bring them into agreement with
one another; and one that, through these fundamental thoughts taken as
abstract opposites, lets this stage of the development of philosophy
dissolve itself.

\divrule

\begin{center}
\textbf{2.\enspace The Old Ionians: Thales, Anaximander and Anaximenes.}
\end{center}

a. \emph{Thales} of Miletus (born c. 640 B.C.) is, according to the
testimony of antiquity, especially of Aristotle, the father of
philosophy. He founds philosophical thinking in that, in his striving
to find the ground and source of all that is, he seeks it in the world
of nature, of reality, not, like the mythologists, in that of the
arbitrarily transforming imagination. He posits \emph{water as that out
of which all comes to be and to which all \opage{24}returns}. That is
all we know with certainty of his philosophy. How he grounded his
doctrine, the sources contain nothing certain about; but it sounds
probable, and does not go beyond the stage of development of his age,
when Aristotle supposes that he ascribed so great a significance to
water on account of the importance that it (as the moist) has for the
organic in nutrition and generation. It is likewise probable that he,
in a childlike, anthropomorphizing manner, intuited his primal matter as
living (ensouled) and active. Concerning the manner in which the
individual things come to be out of the primal matter he seems to have
given no determinations. But by that single proposition about water as
the primal matter Thales gave mind another position toward the object
of knowledge; he broke the path for the philosophizing thought that
acknowledges an independent truth, and this thought now advances
swiftly and energetically.

b. \emph{Anaximander of Miletus} (born c. 611 B.C., flourished toward
the middle of the 6th century. Writing: On Nature). Already in this
first successor of Thales the progress of philosophical thinking is
conspicuous. When Thales gave water as that out of which all comes to
be and to which all goes back, water was thereby described as the
element, as the origin and principle of things; but the expression in
thought for the relation was first formed when Anaximander found the
determination \gk{ἀρχή} and thereby fixed the thought in its own form.
And whereas Thales had posited as the underlying a single matter, given
to immediate sensation, which keeps its name as this single one
although it appears as principle, as what is in the manifold and
different, Anaximander determines as the principle matter \emph{as
\gk{τὸ ἄπειρον}}. By this determination is denoted first the boundless,
matter being seen as an inexhaustible source of coming-to-be; but to the
meaning of the boundless there is then joined, as the element is seen
as that which first \emph{is to} become the determinate matters, the
meaning of the determinationless,
\opage{25}\setcounter{footnote}{0}formless. \gk{Τὸ ἄπειρον} thus
becomes, as it were, the translation into thought of the mythical
chaos. Anaximander determines it as immortal and imperishable, as
infinite, as living, i.e.\ as having the moving force (an \gk{ἀΐδιος κίνησις} [eternal motion]) within itself, as encompassing all and
governing all, and as that from which things are \emph{separated out}
(\gk{ἐκκρίνεσθαι, ἀποκρίνεσθαι}) in accordance with its constant
motion. In the indeterminateness of this expression is mirrored the
indeterminateness in the conception of the relation of the unlimited
to the determinate (cf. the expressions: \gk{μίγμα} [mixture] in
Aristotle, \gk{μία φύσις ἀόριστος} [one indeterminate nature] in Theophrastus, \gk{τὸ παρὰ τὰ στοιχεῖα} [that which is beside the elements] in Simplicius). As the
first oppositions separated out are determined the warm and the cold;
as their first product the fluid; and from its further transformations
is then explained, with bold imagination, the coming-to-be of the
universe enclosing manifold forms of existence, and the formation of an
infinite series of successive worlds, all of which pay the debt of
individual existence by returning, according to the lawfully determined
order of time, to that out of which they have come to be.

c. \emph{Anaximenes of Miletus} (lived in the latter half of the 6th
century. Writing: On Nature). In Anaximenes we find, in comparison with
his predecessor, an apparent step backward, since no longer
indeterminate matter in its generality, but once again a single matter,
\emph{air}, is posited as the principle. In the difficulty of making
intuitable to oneself the manifoldness of the determinate matters as
coming to be out of the indeterminate primal matter, and in the
vagueness of Anaximander's determination of the relation, the motive
could lie for returning to the determinate; and air could then for
several reasons seem suited to be principle. It could do so on account
of its significance for life as the condition of breathing, since the
universe too was conceived as a great living thing\footnote{Cf.
\gk{οἷον ἡ ψυχὴ ἡ ἡμετέρα ἀήρ οὖσα συγχρατεῖ ἡμᾶς, καὶ ὅλον τὸν κόσμον πνεῦμα καὶ ἀὴρ περιέχει}
[As our soul, being air, holds us together, so breath and air encompass
the whole cosmos]. Stob. Ecl. phys. p. 296.}; on account of its
manifold, essential effects, which
\opage{26}\setcounter{footnote}{0}Anaximander had already emphasized;
on account of its constant motion and change, which indicated its inner
life; and finally on account of its apparent infinity. The most
important of the properties Anaximander ascribed to the unlimited,
Anaximenes could find again in air. And out of this determinate element
the manifold forms of existence could now be conceived more
determinately and more intuitably as coming to be through a
condensation (\gk{πύκνωσις}) and a rarefaction (\gk{μάνωσις, ἀραίωσις})\footnote{By rarefaction air was supposed to be transformed
into fire (= the finest air), by condensation successively into wind,
clouds, water, earth, and stones.}, and on the basis of this notion,
and with the use of few and meager observations, imagination could then
in turn support thought in working out, in a manner similar to that of
Anaximander, a whole picture of the world.

d. The various determinations of the principle by the old Ionian
philosophers all keep, as the foregoing shows, within the bounds of the
purely material. Their world-view is a \emph{naive materialism}; but
through it there runs, while mind still hides itself from thought and
is only instinctively felt, in the very endeavor to find the principle
of \emph{unity}, a movement that is opposed to materialism; for in the
principle of unity for the sensuous something ideal is unconsciously
sought. The next step of philosophy also shows that its direction is
moving away from materialism. This next step forward, which gives the
principle a freer, purer form, is taken in the Pythagorean philosophy.

\divrule

\begin{center}
\textbf{3.\enspace The Pythagorean Philosophy.}
\end{center}

In stating the fundamental features of the Pythagorean doctrine, we
speak of the Pythagoreans as a collective unity; for the historical
reports are indeed sufficient to assure us of Pythagoras's historical
existence, but not to separate the philosophical doctrine that belongs
to him from the forms it \opage{27}successively assumed among his
successors. Only internal grounds can here lead onto the trail of the
probable order of the development, but without definite periods of time
being distinguishable.

In the case of the Pythagoreans we again, in accordance with the
limitation of our task, take account only of the purely philosophical
element, not of the political and moral-religious communal life,
ordered in definite, strict forms, through which the Pythagorean league
acquired so great an influence, which was the original interest, and
which formed, as it were, the home of the theoretical development.

Pythagoras was, according to the most probable traditions, born on
Samos about 582 B.C., came toward the middle of the century to Italy,
where in Croton he founded his aristocratic political and
ethical-religious league, and where he died not long after 510 B.C. He
is thus contemporary with Anaximenes. Among the older Pythagoreans are
named Simmias and Kebes, Timaeos of Locri, Archytas of Tarentum, the
Lucanian Okellos, and Socrates's contemporary Philolaos, the first of
the Pythagoreans to set down the school's philosophical doctrine in a
writing, of which genuine fragments are preserved, mixed with spurious
ones from a later time. The surest source of knowledge of the
Pythagoreans is for us Aristotle.

Concerning the relation of the Pythagoreans to the earlier philosophers
Aristotle says: \gk{οἱ μὲν πολλοὶ καὶ οἱ πρότεροι τὴν οὐσίαν καὶ τὸ ὄν ᾤοντο τὸ σῶμα εἶναι}
. . . . \gk{οἱ δὲ ὕστερον καὶ σοφώτεροι τοὺτων εἶναι δόξαντες}
\emph{\gk{τοῦς ἀριθμούς}} [The many and the earlier thought that
substance and being is body . . . . but those who came later and were
reputed to be wiser than these, the numbers]. No longer, then, matter
as such, but the measure of the material is now posited as the
principle of things; but as number is posited as the principle, it is
understood in such a way that it is the very substance (substratum) of
things, their essence, not merely a form and determination in them.
Things themselves, says Aristotle, they determined as numbers.

To this determination of the principle, which at first glance must
seem strange and peculiar to us, but which acquired an essential and
lasting significance for the development of Greek philosophy,
\opage{28}\setcounter{footnote}{0}the Pythagoreans came, as it were,
through a presentiment of the idea, of whose concept number is the
first, naive, undeveloped expression. Greek thinking, having begun to
seek the grounds of existence, fastens, as it were instinctively, on
% errata, p. 28: printed "Grund", corrected "Grunde"
the form of beauty; it then sees in things lawfulness, order,
determinateness, abiding relations; it divines in these relations the
essential; it seeks an expression for them and finds it most nearly and
naturally in number: the measure of the material, the condition of all
being in space, of all lawfulness in the phenomena; and this number it
then conceives, after the manner of childlike thinking, as the matter
of the things in which it is seen, as their substance. Precisely as
number, form, too, still adheres, as it were, fast to the material, and
cannot yet free itself from it as idea.

We now consider, after having provisionally made clear the significance
of the principle of number: 1) the probable development of the
principle; 2) its application to the concrete; 3) its relation to the
moral and religious teachings of the Pythagoreans.

1) In the reports of the principle of the Pythagoreans it is said
sometimes that they determined number as such, sometimes the elements
of number (\gk{στοιχεῖα} or \gk{ἀρχαὶ}), sometimes harmony; and as the
elements of number are given partly the odd and the even, partly limit
(or the limiting: \gk{τό πέραινον}\footnote{As synonymous with this,
the limited (\gk{τὸ πεπερασμένον}) is also used, which, however, from
another point of view has also been conceived as the product of \gk{τὸ ἕν} [the one] and \gk{τὸ ἄπειρον}.}) and the unlimited. These different
expressions point to a successive development, which on internal
grounds may be assumed to have been the following. The oldest
determination of the principle, deriving from Pythagoras himself, was
presumably simply \emph{number} as arithmetical number. But the
observation of oppositions in nature --- and such had already been
emphasized by Anaximander, to whom Pythagoras is said to have stood in
the relation of a disciple, --- had to lead to seeking something
corresponding in number, and it lay near at hand to find it in the two
kinds of number, the even and the odd, which through a playful
reflection were brought into connection with the \opage{29}less perfect
and the more perfect. The odd is supposed to correspond to the more
perfect, because, unlike the even, it resists division in two, and
because it does not change its nature when the even is added to it,
whereas it itself, added to the even, changes this into an odd. The odd
and the even were no doubt originally conceived as the \emph{kinds} of
number, only later as the \emph{elements} of numbers. But as the odd
and the even were posited as the universal principles, there had to
arise, since the determination of number as principle had originally
proceeded from the intuition of form, a need to exchange these more
symbolic designations for such as pointed directly to the fundamental
relation of the formed, such as the concepts of limit (the limiting) and
the unlimited. These new determinations were then, if somewhat
artificially, combined with the former ones; and as the opposites, in
the new formal determinateness, now showed themselves as connected,
their union was sought in harmony. For since, in the formed actuality
of things, the unlimited could not be thought of, as the even or the
odd could, as subsisting by itself alone,
% errata, p. 29: printed "det Ulige", corrected "det Lige eller det Ulige"
it had in each single thing to be posited as connected with limit, and
the bond of this connection was precisely harmony: ``the unity of the
manifold and the agreement of the conflicting.'' At this point an
attempt could now be made to set up a table of opposites, in which the
fundamental opposition between limit and unlimited, the formed and the
formless, is mirrored, and such a table is reported in the old sources,
without, however, its author being certain. The 10 pairs of opposites
are: 1) Limit and Unlimited; 2) Odd and Even; 3) Unity and Plurality;
4) Right and Left; 5) Male and Female; 6) Resting and Moved; 7) Straight
and Curved; 8) Light and Darkness; 9) Good and Evil; 10) Square and
Rectangle. --- This table is characteristic in the kind of opposites it
puts together, in the lack of distinction between the mental and the
bodily that shows itself in it.

At the same time the significance of number for knowledge could now
also be emphasized, and it is expressed in Philolaos's statement: ``All
% printed "Plilolaos's" (defect logged in the transcription)
\opage{30}that is known contains number within it, since without number
nothing can be thought or known, but all is indistinct and unclear,
whereas number, in harmoniously joining things together with the soul,
makes it possible that all things can be known and determined according
to their mutual relations''; --- and: ``The nature of number is
lawgiving and governing, and is a teacher of all that is disputed and
of all that is not known.''

2) The working out of the doctrine of numbers and the application of
the principle of number to the concrete had, by the nature of the
case, soon to lead to the arbitrary and fantastic. Already in the
derivation of the arithmetical numbers from the principles much
fantastic speculation asserts itself, as in the determination of the
sacred significance and mutual relation of the number four and the
number ten. And yet here, as in the derivation of the ``geometrical
numbers'': of the point as corresponding to unity, of the line as
corresponding to the number two, of the plane as corresponding to the
number three, and of the solid as corresponding to the number four,
one still moves in the domain of abstract determinations. But when the
numbers are to be applied to concrete objects of nature (as when the
regular stereometric figures derived from the numbers are brought
together with the elements: the cube with earth, the tetrahedron with
fire, the octahedron with air, the icosahedron with water, and the
dodecahedron with the fifth, all-encompassing element) or to spiritual
relations (as when the number four and the number nine, as \gk{ἀριθμός ἰσάκις ἴσος},
as square numbers, represent justice with its like for like, or the
number five, as the first union of an even and an odd number,
represents marriage, the union of the female and the male), then the
% printed "Mandliges" (defect logged in the transcription)
arbitrarily symbolizing comes forward ever more clearly.

More important and more significant than the application of the
principle of number to the single concrete thing is the application of
the principle in its generality to the whole of the universe, which is
intuited as a \gk{Κόσμος} [ordered world] harmoniously formed through
numerical relations. The designation of the universe as \gk{Κύσμος}
% printed "Κύσμος" (defect logged in the transcription)
has its origin precisely with the Pythagoreans. They conceived the
formation of the world and the world-\opage{31}order in the following
manner. First there formed itself in the middle of the world-all a
central fire, figuratively designated as the One, the Mother of the
Gods, Hestia, the Citadel of Zeus, in which the limiting, which here,
as principle of the existence that is coming to be, must step forth
from its opposite, is intuited bodily. From it the nearest parts of the
unlimited --- which at once designates infinite space (infinite
extension) and infinite matter --- were attracted and, through the
attraction, limited, until, by a constantly further spreading of this
effect, the structure of the world came to completion. The structure of
the world was thought of as spherical, as a \gk{σφαῖρα}; its middle is
the central fire, and its outermost limit is formed in turn by a
peripheral fire, corresponding to the fire in the middle of the
universe. In ten circles the heavenly bodies move about the central
fire, among them the earth as the next-to-innermost, and, as the
innermost, an \gk{ἀντίχθων} [counter-earth], invisible to us, which
completes the sacred number ten. The central fire is \gk{πρῶτην φύσει}
[first by nature], it is \gk{συνοχὴ} [what holds together] and
\gk{μέτρον φύσεως} [measure of nature], and by the latter
determinations the Pythagoreans no doubt indicated that a continuing,
intervening influence, through an outflow of warmth and life-force
(soul), proceeds from it. In the regular ordering of the heavenly
bodies the Pythagoreans found the proof of their divine nature. A music
of the spheres was, according to the doctrine of some Pythagoreans,
supposed to be produced by the distance of the heavenly spheres from
one another being precisely such as corresponded to the ratio of length
between the strings that give mutually harmonious tones; for from such
tones, since the distances of the heavenly bodies from one another
correspond exactly to the intervals of the musical octave, a chord was
then to be formed during the revolution of the spheres about the
center, even though it is not heard by us, who always live in this
harmony of the world-all. --- Beyond the peripheral fire, which as the
enclosing bond encompasses the world-all, the Pythagoreans set the
unlimited (\gk{τὸ ἄπειρον}) or the unlimited air, from which the
world-all, as a \gk{ζῷον} [living being], draws its breath, and from
which time (the successively infinite) and the void also enter the
world. \opage{32}Within the universe so formed, in which a distinction
is sometimes drawn between \gk{Ὄλυμπος}, the peripheral fire,
\gk{Κόσμος}, the region of the heavenly spheres from the fixed stars to
the moon, and \gk{Οὐρανός}, the region below the moon, the Pythagoreans
perhaps conceived the various spheres of reality as forming a graded
series with respect to perfection.

3) The moral and religious teachings that derive from the Pythagorean
society are set in only a vague and abstract connection with the
principle of the Pythagorean philosophy. The ethical determinations are
either immediately placed alongside the principle, as when it is said:
virtue is harmony, or they are, in the arbitrarily symbolizing manner
touched on above, paralleled with the individual numbers. A properly
philosophical determination of ethical concepts the Pythagoreans do not
give. Their ethics is a popular ethics with a religious stamp. Their
doctrine of the Godhead is not set in connection with the doctrine of
the principles, and in the religious teachings there appears, despite
the belief in the gods being purer than in the popular religion, much
superstition and much fantasizing, e.g.\ in the doctrine of the daemons
and of the transmigration of souls. ---

The philosophy of the Pythagoreans is thus a philosophy of nature
which, like that of the Ionians, has as its essential task to find the
element of things. By determining number as this, it takes a step
forward toward the ideal, yet in such a way that it does not leave the
bounds of matter, which also shows itself in the materialistic
conception of the psychical (the soul = \gk{τὰ ἐν τῷ ἀέρι ξύσματα}
[the motes in the air]; the soul bound in the body by the veins). A
step further toward a more ideal conception is now taken by the
Eleatics.

\divrule

\begin{center}
\textbf{4.\enspace The Eleatics.}
\end{center}

While the Ionians had determined the material, and the Pythagoreans number, as the principle of existence, the Eleatics determine it
\opage{33}as what is, which as unity excludes becoming and
multiplicity, and which, in being posited as one with thought,
intimates the ideal still more determinately than the principle of the Pythagoreans.
But since the Eleatics, like the Pythagoreans, seek being as
the essence of the natural, corporeal existent, it is determined
for them as the intuited existent, which receives the determination of space, and in the principle, which by its fundamental determination points
toward the ideal, we nevertheless obtain only the substance of corporeal
existence. Being (\gk{τὸ εἶναι}), which becomes what is
(\gk{τὸ ἐόν}), is conceived as matter, and since thought is identified
with being, thought itself is conceived corporeally.

The philosophy of the Eleatics passes through a threefold stage of development. First the doctrine of unity is founded by \emph{Xenophanes}, still with a theological coloring; then it is developed in its
pure conceptual form by \emph{Parmenides}, the real representative of the school; and thereafter the doctrine of being, which in its abstraction permits no development, is defended indirectly by \emph{Zeno}
through the dialectical demonstration of the insoluble contradictions
of the opposite view, while \emph{Melissus} defends it directly,
in connection with what is commonly assumed, but with less
acuteness and stringency.

\textit{a.} \emph{Xenophanes} of Colophon (born c. 569 B.C.;
lived in his later years at Elea in Lower Italy, and died c. 477). As
Xenophanes seeks the ground of existence, he finds it in the
one, unbecome, unchangeable, thinking divine power: All
is one, and this is the Godhead. This notion of the Godhead is still held, as Aristotle already remarks, in
an indeterminateness that leaves it unclear whether by the unity of the divine he has thought of a unity grounded in the concept
(as Parmenides), or of a unity grounded in the infinite continuity of extension. Xenophanes has merely, ``in
looking at the whole universe,'' designated as God the unity whose notion
took shape in him; but he has not more precisely determined
the essence of the divine unity. Nor is the relation
between God and the world yet clear in him. On the one
\opage{34}hand the Godhead, which as the One becomes the world's
immanent ground, merges with the world, so that its determinations
can be transferred immediately to the latter, and God and the world
stand related to each other; as essence and phenomenon they are inseparable.
% errata, p. 34: printed "forholde sig som", corrected "forholde sig; som"
But the conception of the one existent as Godhead leads, on the
other hand, to an involuntary personification and thereby to a
distinction. When it is said of the Godhead that it is wholly
seeing, wholly thinking, etc., a separation has thereby involuntarily
been made between the seeing, thinking, hearing
God and what he sees, thinks and hears. Thereby it becomes
possible for the manifold and changeable in the phenomenon
still to retain a semblance of reality beside the One
and unchangeable. Only when the \emph{notion of the Godhead} is replaced by the \emph{concept of being} (in Parmenides)
can the relation between the One and the many be sharply
determined.

Xenophanes becomes of great significance for the position of Greek philosophy toward
the Greek popular religion through his sharp
polemic against the anthropomorphic and polytheistic notions of the gods which, through the poets, had gained established authority in the common consciousness. Proceeding from the notion of the Godhead as
the one, unchangeable, all-governing, which is wholly seeing, wholly
thinking, wholly hearing, unlike man in form and thought,
and which without toil governs all things by the power of its thinking,
he censures the poets because in their tales of the passions
and battles of the gods they had ascribed to them everything shameful and
unworthy; and he seeks the ground of the anthropomorphisms of popular religion in man's urge to form the gods after
his own image, an urge that would show itself in a corresponding form in
the animals, if they had man's ability to form images of the gods.

\textit{b.} \emph{Parmenides} of Elea (born in the next-to-last decade of the 6th century), who was admired by Greek antiquity for
his depth and greatness, gives Xenophanes's doctrine of unity its
pure metaphysical form, in that he abandons the theological conception of unity and grounds the unity of all that is on
\opage{35}the very concept of being, from which he derives the marks of being (\gk{σήματα}). His starting point is the proposition that
only what is can have being, whereas what is not,
as the opposite of being, can neither be nor be thought. His
charge against the common notion is that, by letting what is not have being, it moves
in the contradiction
of positing what is and what is not at once as
one and as different. From the concept of being, which is thus conceived in sharp distinction from its opposite, is derived first the
mark that what is must be unbecome and imperishable. It cannot have come to be, because it would then either
have had to come to be from what is not, which precisely as
such can produce nothing, or from what is, which is, after all,
itself. Nor can it pass away, because what is cannot
become something other than itself: if being ever were
not, what is would indeed not-be. --- Next it is shown
that what is must be untouched by differences of time. It cannot
be said of it that it has been or will be; for thereby it would be said that it no longer was, or that it was not
yet, thus in both cases a not-being that is incompatible with what is. Further, what is is an indivisible unity, since there is nowhere anything different from it
that could separate its parts; it is in the same place, i.e.\ without
spatial distinctions, as a pure continuum. It is unchangeable, like itself and immovable. It is limited; for
otherwise it would be unfinished and deficient. It is, finally, one with thought; for there is nothing outside what is;
all thought is thought of being. What is, then, is the All as
unity, without coming-to-be and perishing, without change of
place and form, an undivided, homogeneous whole, equally
perfect at every point, which Parmenides can therefore designate as
a well-rounded sphere.

What must strike us in this derivation of the marks of being,
the first properly dialectical determination of a concept that we meet
in the history of philosophy, is the
\opage{36}immense boldness with which thought, relying on itself,
defies the impression of the multiplicity and change given to sensation,
and over against it holds fast to the infinite
unity and unchangeableness of being. And this thought, which thus
confidently rejects the testimony of sensation, points, in that it
affirms the eternal, unchangeable unity of what is, its being untouched by the division of space and time, and its unity
with thought, beyond the limit of the material to the region of the idea.
The true concept of being, the idea, is as it were on the point of breaking
forth. Yet it becomes clear that at this standpoint being is not yet conceived and held fast in an ideal sense
and as a purely metaphysical concept, but, despite those determinations reminiscent of ideality, as what is in the sense
of the substance of the corporeal. What is is determined by
space, is the space-filling (\gk{τὸ πλέον}), and Parmenides
as yet has no thought at all of a non-spatial being, and a
separation of the corporeal and the incorporeal is foreign to him, indeed even directly incompatible with his whole standpoint.
When he posits thought in unity with being, this
unity is not to be understood in the spirit of modern idealism, but
in such a way that thought falls under being, itself becomes
corporeal.

Through Parmenides's answer to the question of
the ground of what exists, and his determination of being as
excluding all not-being, and with it all multiplicity
and all becoming, a stage within the first period of Greek philosophy is concluded, and at its conclusion
thought stands facing a problem that cannot be solved from its presuppositions.
It is rational thinking (\gk{λόγος}) that
pronounces the truth: that only what is, is; the false opinion,
that what is not can have being, stems from the deception of the senses,
which delude us with the manifold and the changeable: that whose being includes the not-being that lies in difference,
and that whose being alternates with not-being. But
precisely this existence of what is not in consciousness becomes
\opage{37}inexplicable for Parmenides. As existing for consciousness,
between whose forms the doctrine of unity can establish no distinction in principle, not-being, which has been driven back
as it were to its last entrenchment, asserts for itself the being that was posited
as impossible; and since what is not cannot be expelled
from the sensuous notion, the question of the multiplicity and becoming inseparable from it must come forward again.
Parmenides too, in his didactic poem, after the true doctrine
of being is concluded, has a doctrine of the world of appearance follow,
which, in hypothetically ascribing validity to it, he explains by two opposed principles: the light and warm and the
cold and dark, which in their relation to each other form an
analogy to the opposition between being and not-being and
by their various mixture are supposed to produce the multiplicity of things. But from his principle of being the very concept
of appearance is inexplicable. The world of appearance emerges, uncomprehended,
and asserts a being for itself; and although it is dialectically dissolved by the Eleatics through the demonstration of the contradictions it
harbors within it, it constantly returns, as if bewitching
their thought.

\textit{c.} \emph{Zeno} of Elea (Parmenides's disciple, born c. 485)
championed the Eleatic doctrine in a
writing composed in dialogue form, in which he shows the contradictions in which the opposite
assumption entangles itself, and from its inner contradiction concludes
to its impossibility. His attack is directed at the two points
in which the common notion's confusion of being
and not-being appears most pregnantly: against the notions
of \emph{plurality} (multiplicity) and of \emph{motion} (the most general
expression of becoming). He uses the concept \gk{τὸ ἄπειρον}, the
infinite, the opposite of the \gk{πέρας} [limit] that was the determination of the one
existent (cf.\ Parmenides), to bring contradiction into those notions, which point to this \gk{ἄπειρον}.
In its two forms, as the infinite in the direction of increase
and as the infinite in the direction of division, \gk{τὸ ἄπειρον}
brings contradiction into those notions by becoming what is incommensurable with the
\opage{38}\emph{determinate} magnitude, and thereby what
constantly negates this magnitude, which is the condition of the
reality of multiplicity and motion. The infinite
thus becomes the instrument of \emph{negative} dialectic.

Thus Zeno argues against the concept of \emph{plurality},
saying that what is, if it were a manifold,
would have to be at once infinitely small and infinitely large.
It would have to be infinitely small; for every plurality consists
of units, but the unit is a real unit only as indivisible,
and the indivisible is without magnitude, since all magnitude can be divided.
Everything that is a unit or consists of units must therefore
be magnitudeless, i.e.\ infinitely small. But equally what is
would, as plurality, have to be infinitely large; for in order to be,
the manifold would have to have a magnitude, and in order to be as
manifold, as plurality, it would have to be separated. But that
which separates the manifold also has a magnitude, and so likewise does that
which separates these new magnitudes from the earlier ones,
and thus we get infinitely many magnitudes and by their
summation something infinitely large.

He demonstrates a similar contradiction when \emph{number} is applied to plurality. The manifold must be limited in
number; for its number is equal to the number of the many things. But it
must also be unlimited, since what separates, and what
in turn separates from what separates, etc., is counted in to
infinity.

Against the reality of \emph{motion} Zeno sets up 4 proofs.
1) Every motion from one point to another becomes impossible,
because the moving thing, in order to reach the goal, would first have to traverse the infinitely many parts of space that the distance includes, but
the infinitely many parts of space cannot be traversed in a
finite time. 2) If one assumes motion, the contradiction arises that the faster-moving cannot overtake the slower, which has a head start; for in the time the faster
covers the head start, the slower gains a new, smaller
head start, and while this in turn is covered by the former, again
\opage{39}a smaller one, which can be divided to infinity but never vanishes. 3) The arrow that is imagined to fly is at rest; for in each
instant, i.e.\ in each infinitely small part of time, it is in one space;
but what is in one space is at rest. And if it is at rest in each
instant, it is always at rest; for time is the sum of
the instants. 4) What moves with the same speed would,
according as its motion is measured against something at rest or
something moving, be able to cover unequal spaces in the same time, the
same motion thus being unequal in magnitude to itself (motion thus having in general an infinite relativity).

By his proof against the fundamental determinations of the common notion, Zeno founded the \emph{negative} form
of dialectic, which later played so great a role in
Greek philosophy. For him, however, this negative dialectic was not, as later, an end, but only a means to uphold
a positive conviction. His dialectic is, as Greek
dialectic always remained, a dialectic of \emph{distinction}, which holds fast the sharp distinction between the opposites, and precisely
thereby seeks to uphold them in their truth. Therefore the infinite, too, is
sharply distinguished from the determinate, finite, not
seen in that unity with it which is the condition of all real life, all becoming and all motion.

\textit{d.} \emph{Melissus} of Samos (flourished about the 84th Olympiad;
% printed "4." (defect logged in the transcription)
a writing in prose on what is) conducts \emph{direct} proofs of
the doctrine of the unity of being, making use of the points of attachment that
are given in the representatives of the common conception.
The proofs concern the determinations: eternity, infinity, unity and unchangeableness. For \emph{eternity} the proof is conducted essentially as in Parmenides. \emph{Infinity} is derived from eternity, since that which, as unbecome and imperishable, has neither
beginning nor end is the infinite. (For this
inference from infinity in time to infinity in space Melissus
is censured. The connection of thought in him seems to be this:
what has a beginning is not whole all at once;
what has no beginning is whole all at once, indi\opage{40}visible,
and that is only the infinite). From infinity follows
in turn \emph{unity}, since several existents would mutually limit
one another. They could also only be thought by means of a void
that separated them, but the void is something that is not. From
unity follows, finally, the \emph{unchangeableness}
and \emph{immobility} of what is. What is is unchangeable, since every change would make it a plurality. It is
immovable because there is no void outside it in which it
could move, and because as unity it cannot move
within itself, which would involve a division. The testimony
of the senses to a plurality has no reliability, because they
show us things as changing, thus not as being, but as passing over into not-being or coming to be
from not-being. What is, is for Melissus the one continuous substance, and when he designates what is as
incorporeal, only the divisibility that attaches to
corporeality (\gk{σῶμα, πάχος}) is thereby to be excluded: the determination of extension is not cancelled; it is held fast precisely in the determination:
the unlimited, and in this determination Melissus seeks
a guarantee of unity, since the unlimited, in his
view, excludes division.

\divrule

The significance of Eleatic philosophy for the development of Greek
thought we must, after the foregoing, place
partly in its energetic breaking away from the sensuous, whereby
a great step is taken toward the idea, as we saw intimated in
the determination of the principle of being; partly in the clear and sharp
distinction of such fundamental concepts as what is and what
is not, in the incipient analysis of the marks of concepts,
and in the germs of a dialectic both as positively grounding and as negatively dissolving; partly, finally, in
the problems that the representatives of the doctrine of unity, and in particular Zeno,
set in motion. We shall see in what follows how
essential an influence the Eleatics exercised on the subsequent
systems even in the culminating period of Greek philosophy.
\opage{41}With the preceding philosophical attempts Eleatic philosophy shares the character of \emph{natural philosophy}. The
dialectical tendency is a \emph{consequence} of the way in
which the world-principle is determined; the Eleatics' propositions about
true rational knowledge and untrue sensation are
\emph{derived} from the metaphysical and physical propositions; \gk{λόγος}
and \gk{δόξα} are distinguished by their content and are both traced back (by Parmenides,
in the second part of his poem) to the mixture of the kinds of matter without any distinction in principle being able to be made. A theory
of dialectic the Eleatics do not give. ---

The problem at which the thinking of the Eleatics had to
come to a halt has been pointed out above. Since the question of the ground of things
has been answered in such a way that the determination of the principle as
the One, the existent, excludes becoming, but this becoming, denied in the principle,
nevertheless cannot be expelled from the sensuous
notion, just as it is not forgotten by the thought that had gone back from the multiplicity of existence to the one existent,
becoming asserts its validity as that which cannot be
annihilated, and the question of the coming-to-be of what exists
now presses forward anew and now becomes the \emph{primary} problem. This turn in thought is represented by Heraclitus.

\divrule

\begin{center}
\textbf{5.\enspace Heraclitus.}
\end{center}

\emph{Heraclitus} of Ephesus (flourished at the close of the 6th
century; a writing on nature, obscure in its form, rich in imagery and
gnomic, often poetic and grand) fixed his gaze from the very beginning
on the current in existence by which everything comes to be and
vanishes again. Not resting being, but living production captivates
his thought. \gk{Πάντα ῥεῖ} is his point of view on existence. ``All
is in flux,'' i.e.\ in uninterrupted motion, in unceasing becoming. The
older Ionians and the Pythagoreans had also admitted a coming-to-be,
but had not let the principle be absorbed into \opage{42}becoming,
become one with it; but in Heraclitus the principle is immediately
conceived in such a way that it coincides with becoming.

This fundamental thought of the restless flux of existence Heraclitus
now sets thetically against the one unchangeable being of the Eleatics;
but it forms at the same time an opposition to the view of ordinary
notion, which makes room for a relatively abiding and tranquil being
and thus ascribes to things a relative unchangedness, whereas according
to Heraclitus's view nothing has subsistence and persistence for a
single moment, but everything passes uninterruptedly over into its
opposite, comes to be out of everything, becomes everything.

This restless life of existence Heraclitus intuits as the primal fire,
which eternally transforms itself into things and lets things go back
into itself. Fire is not, with him, the symbol of a metaphysical
concept of becoming; rather, the notion of the process of coming-to-be,
of becoming, coincides immediately with the intuition of fire as the
living, mobile element that transforms itself and everything else.
Fire, which exists only as unceasingly transforming itself, produces
all things through the double motion of its transformation, on the way
downward, which leads fire (a name that also designates the ethereal
air) over into water and earth, and on the way upward, which leads back
again to fire, inasmuch as a constant double current takes place: a
development and combination of opposites. And in this production fire
is itself the power that governs the production and leads what is
produced back again to itself. As it unfolds opposites, which is
precisely the essence of production (\gk{πόλεμος πάντων πατήρ} [war is
the father of all]), it joins these together again into unity, and out
of the very transition and separation it calls forth a return, a union,
so that the motion of becoming forms a cycle in which one becomes all
and all one, and which is governed by harmony. Harmony expresses partly
the harmonious joining together of the opposites in what exists (``in
separating from itself, the one comes together with itself, like the
harmony of the bow and the lyre''; --- ``from what is at variance comes
the most beautiful harmony''), partly the lawful \opage{43}relation in
% printed defect at the head of p. 43, site not specified in the note (defect logged in the transcription)
the transition between them. Harmony is designated by Heraclitus as
the divine law, as the world-governing wisdom, as Zeus and as the
Godhead, and is conceived in unity with the principle itself. In this
principle, which is accessible only to rational knowledge, are thus
comprised the world-forming force, matter and law; production,
intuited as a \gk{πῦρ
ἀείζωον} [ever-living fire], becomes not merely the principle through
which, but also that out of which, everything is, and at the same time
it bears within itself the lawful, harmonious order of change. It is
characteristic of Greek thought that even the philosopher who lets
existence flow in restless change nevertheless posits a measure, a
law-determined relation in becoming, and posits it as the only fixed
and abiding thing in eternally moving existence. But this abiding
element in existence forms an opposition to the Eleatic existent unity,
in that it is a unity in change, and in change the manifold individual
acquires significance as a member in a current which just as fully
gives it existence as it cancels this existence. As the cycle takes
place in the individual according to the world-law of strife, so it
takes place in the large. In great periods everything returns to the
primal fire, from which a new world-formation then proceeds again. The
primal fire is kindled and quenched according to law-determined
measures. As in restless play the Godhead builds worlds without end,
and lets them return again into itself.

The divine law, which governs existence as the universal order, is to
be the law for the individual's thinking and acting. The soul is of the
nature of fire (it is a \gk{ξηρὰ ἀναθυμίασις} [dry exhalation]); and
it draws its nourishment from the rational fire diffused through the
whole world, which streams into it through breathing and the open
senses. Man is to follow in his knowledge the universal reason
(\gk{λόγος ξυνός}=\gk{κοινός}), which is apprehended by thought, not his
particular opinions (his \gk{ἰδία φρόνησις}), which rest on deceptions
of the senses and imaginings. And likewise he is to follow in his
action the universal law, the law of the state and the law of the
universe; he is to flee self-will and willfulness and all \gk{ὕβρις}
[insolence]. \opage{44}It is wisdom to obey the law of nature in one's
speech and in one's action; in submission to it
% printed "i Underkastelsen under den" with a defect between "Underkastelsen" and "under" (defect logged in the transcription)
is reached that contentment of mind (\gk{εὐαρέστησις}) which is the
highest aim of life.

It is Heraclitus's merit, in bringing forward the problem of
coming-to-be as the primary one, to have held fast the intuition of
becoming in its sharpness and to have conceived the process of
coming-to-be as the restless production and joining together of
opposites, as this outflow which at every moment is at the same time a
backflow. The restlessly pulsing life of existence he grasped in a
lively intuition. But while Heraclitus does demonstrate becoming as a
fact, and asserts an eternal lawfulness in it, he does not explain the
ground of becoming or this lawfulness. He cannot do so, because
becoming does not for him become what is comprehended by an ideal being
and grounded in such a being, but stands as its own presupposition, its
own subject, and, as the becoming of the sensible --- for the principle
was, after all, a sensible matter --- does not explain the abiding
measure, the dominion of the transformed element over the
transformation. The consequence of this must be, on the one hand, that
where the principle is held fast, as in H.'s school, everything abiding
must vanish, and a becoming that dissolves itself and all thought, a
becoming out of which nothing becomes, must alone remain; --- on the
other hand, that a new task must present itself to thought: to
demonstrate the ground and necessity of becoming; and as this task is
taken up by thinkers who know the doctrine of the Eleatics, it
acquires, as already indicated above (p. 22), the more determinate
form: to seek the source of becoming in the underlying itself, without
the concept of what is being lost thereby. The attempt to solve this
task is then at the same time to be regarded as a mediation between the
fundamental determinations of the Eleatics and of Heraclitus, each
insufficient on its own, being and becoming, a mediation in the course of which
% errata, p. 44: printed "hvilket", corrected "hvilken"
these are transformed. Being, which is to leave a place open for
becoming, can no longer be the one, undivided, unmoved being; plurality
and motion must enter into it, but the essence of being must be upheld
by each of the \opage{45}originally existent things remaining
qualitatively unchanged what it was. And precisely on account of this
conception, becoming must be determined otherwise than in Heraclitus.
Becoming is no longer a transformation, an inner qualitative change, a
transition to an other, to the opposite; it now concerns only the
external relation to an other of what is inwardly unchangeable; in
other words: it is reduced to a mechanical change of place of
unchangeable, ungenerated elements. And for this change of place, for
this combining and separating of what originally exists as a plurality,
a ground is now sought in the very nature of the principles.

The philosophers who represent this conception form the third group in
the oldest period of Greek philosophy; they are: \emph{Empedocles},
\emph{the Atomists} and \emph{Anaxagoras}.

\divrule

\begin{center}
\textbf{6.\enspace Empedocles.}
\end{center}

\emph{Empedocles} of Akragas (lived c. 492---432 B.C.; didactic poem
on nature) is the first who distinguishes the moving force from matter,
with which it had been joined by the older Ionians, the Pythagoreans
and Heraclitus, through childlike anthropomorphizing or through
conscious reflection. As the unchangeable primal constituents of what
is, he posits the four elements (\gk{ριζώματα} [roots]): fire, air,
water and earth, and as principles of motion, of combination (mixture,
\gk{μῖξις}) and separation (\gk{διάλλαξις}), two mythical figures:
Friendship (Love) and Strife (Hate): \gk{Φιλότης} and \gk{Νεῖκος},
which, however, he still conceives as mixed in a material way with the
elements. That the materials are not to be transformed, but only to
change their external relations (\gk{φύσις}, coming-to-be, is an empty
name), conditions in Empedocles the distinction of the two kinds of
principles (the material and the moving); for in the unchangeable
original elements he could not conceive a drive toward a combination or
separation that did not touch their essence. Therefore the moving force
had to be sought outside them, and it was presumably posited as twofold
so that it should not, by combining opposite \opage{46}activities, as it
were change its nature and itself prove subject to transformation. More
precisely, Empedocles conceived the elements as originally mixed, as
being in undifferentiated tranquil unity in the \gk{σφαῖρος} [Sphere],
governed by \gk{φιλότης}, until \gk{νεῖκος} penetrated from the
periphery and brought about the differentiation, which, having reached
complete separation, gave place again to a successively advancing union
effected by \gk{φιλότης}, in the course of whose advance the individual
forms of existence and the order of the universe come to be, only to
cease again when the original unity is completely restored, but then to
be called into life anew through the same cycle without end.

In Empedocles's doctrine of the coming-to-be of natural things there
are, beside much that is fantastic and fabulous, interesting
anticipations of the theories of later times on the successive
development of the higher forms of existence from the lower, and on the
as it were groping attempts at organic formations, in which the
monstrous and the normal arise side by side, but only the latter
achieves subsistence, because only it has the capacity to maintain
itself. In man the principles are mixed as in the rest of existence; we
know the elements of things (the material and the moving) by what
corresponds to them in us. ``Like is known by like.'' Effluences from
things, as they meet with corresponding ones from the senses or touch
corresponding particles of matter in the organs, produce our
% errata, p. 46: printed "Opgaverne", corrected "Organerne"
sensations.

\divrule

\begin{center}
\textbf{7.\enspace The Atomists.}
\end{center}

The atomic doctrine is founded by \emph{Leucippus} of Abdera (probably
contemporary with Empedocles and Anaxagoras), but finds its proper
representative in \emph{Democritus} of Abdera (c. 460---360), author of
numerous writings, mostly on natural philosophy, of which his
\gk{μέγας διάκοσμος} [Great World-Order] was the most famous, but which,
except for a few fragments, are lost. By the many-sidedness of his
learning, which he had gathered on long journeys in the Orient and
Egypt, \opage{47}Democritus recalls Aristotle, who mentions him with
respect and frequently makes use of him, whereas Plato ignores him.

As principles the Atomists, since their task is to explain the
multiplicity and change in the real that is given as \emph{fact},
determine the full and the empty (the atoms and empty space), which
become the physical expression for the metaphysical: what is and what
is not, to both of which being is equally ascribed. (``\gk{Τὸ δὲν}
[the thing] is no more than \gk{τὸ μηδέν} [nothing].'' Whereas the
Eleatics denied the fact by virtue of the concept, the Atomists deny
the consequence of the concept, that non-being cannot be, in favor of
the fact.) The full is the primal bodies, which, since a division to
infinity would cancel magnitude, are indivisible, atoms. The atoms, as
the elements of the infinitely manifold, are infinite in number and, as
the expression of what is, ungenerated and imperishable, and
insusceptible of inner change. They differ from one another not by
qualitative differences, but only by differences in shape (\gk{ἰδέα},
\gk{σχῆμα}) and size and in the weight corresponding to size; and from
these differences and from the difference in the order and position of
the combined atoms are derived all concrete differences in what exists,
which appear to our senses as qualitative differences. The unlimited
empty surrounds and separates the atoms and is the condition of motion,
and precisely as such a necessary existent. By the eternal motion,
which in its original form is a vertical motion of falling, but which,
on account of the atoms' difference in weight and the consequent
difference in speed of fall, leads to collisions between the atoms and,
through the atoms' rebounding, to a vortex motion that constantly
spreads further, the composition of the atoms into concrete existences
is to be explained. Through the vortex motion what is alike in weight
and shape is first gathered in the same place by a natural necessity,
but along with this, through the shaking together, unlike atoms can
then be united into lasting combinations, whereby are formed not merely
the individual composite things, but countless coordinate worlds,
\opage{48}subject to a constant alternation, all governed by the same
\gk{Ἀνάγκη} [Necessity], not steered by a thought and a purpose.

The soul of man consists of the finest, smoothest and roundest atoms,
the fire-atoms, which are diffused through the whole body, but in such
a way that the different activities of the soul have their seat in
different organs. We draw the soul-atoms into ourselves by inhaling and
expel them by exhaling. Everywhere in existence there is soul and
reason. As atoms stream out from things, forming images and reaching
the soul through our senses, sensation arises, which deceives in that
it declares qualitative differences (such as warm and cold, sweet and
bitter) to exist objectively, whereas in reality only the qualityless
atoms and the empty have being. The true, genuine knowledge, the
knowledge of this, is won by thinking, which infers from phenomena to
what is hidden and is conditioned by the right temperature in the soul.
How this thinking, which cannot be independent of sensuous feeling,
obtains its peculiar capacity, Democritus did not reflect on more
closely. --- As the highest good Democritus posits happiness, which
consists not in external goods but in the tranquility and harmony of
the soul, and is conditioned by moderation. Not the external action but
the disposition determines, in his view, the moral character. For the
rest, he does not lead his ethical propositions back to the general
principles, although in their character they are determined by his
theoretical propositions on the unchangeable order of nature and on the
higher worth of what is hidden behind the phenomenon.

In comparison with Empedocles the Atomists have made an essential
advance, in that the principles are led back to a more universal
determinateness and to a more conceptual opposition, and in that their
relation loses the unclarity and arbitrariness it retained so long as
the moving forces, though distinguished from the material, were
nevertheless represented as mixed with it in a material way, and so
long as they were conceived in a mythological, anthropomorphizing form.

\opage{49}Peculiar to the ancient atomic doctrine, in comparison with
atomic theory as it appears in modern science, is its metaphysical
character. It is not through experience and through the analysis of
the concrete phenomena of nature, but through speculation, that
philosophy is originally led to set up the doctrine of atoms. Only the
quite general problem, to explain the possibility of becoming and of
the multiplicity of existence, was brought forward by the intuition of
real existence; but in this general form the problem is solved by way
of speculation alone, and precisely the generality of the problem must
lead onto this way. The atoms and empty space (the full and the empty)
are indeed physically named determinations, and the explanation of the
coming-to-be of things by the motion of the atoms rests on a physical
intuition: the intuition of space and of what moves in space; but
behind the physical names there hide only abstract metaphysical
concepts: the concepts on whose relation earlier thinkers had already
pondered. It is the concepts of being and non-being that are translated
into physical determinations as they are intuited as the atoms and
empty space, as the existent full and the existent empty; but in this
translation they still retain their earlier content. That the atoms
and the empty are the explanatory ground of being for what exists
expresses precisely the same as the Atomists' paradox: what is is no
more at all than what is not. Existence is explained by the atomic
doctrine not physically, but metaphysically; the \emph{concept} of
multiplicity (plurality) is explained by the \emph{concept} of atoms,
the \emph{concept} of becoming (motion) is explained by the relation of
% errata, p. 49: printed "Vorden (Flerhed)", corrected "Vorden (Bevægelse)"
what is and what is not; and of an exact determination of the phenomena
there is no question. The metaphysical stamp that the Greek conception
of nature retains throughout, the conception of nature has also in the
ancient atomic doctrine, even though this doctrine contains incentives
for a later time toward an exactly executed, consistent mechanical
theory of nature.

\divrule

\opage{50}
\begin{center}
\textbf{8.\enspace Anaxagoras.}
\end{center}

\emph{Anaxagoras} of Clazomenae (born c. 500, lived long in Athens as
the friend of Perikles, until, accused of atheism, he had to leave the
city; died at the age of 72; writing \gk{περὶ φύσεως} [On Nature]) was
older than Empedocles and perhaps also than Leucippus, but seems to
have come forward as an author later. He seems at least to have known
Empedocles's doctrine. As the one who came forward later, and as the
one who at least names a higher principle than all the others, he
receives his place here as concluding the positive philosophical
attempts of the first period.

As principles of being Anaxagoras posits a multiplicity of primal
matters infinite in number and kind, which are not, like the atoms,
merely quantitatively different, but are qualitatively distinct from
one another and are the seeds of the concrete existent, whose smallest
parts they are. He calls them \gk{σπέρματα}, as that from which things
come to be, or simply \gk{χρήματα} [things]. The later designation,
the homoeomeries, derives from the Aristotelian terminology, according
% printed "bidrører" (defect logged in the transcription)
to which \gk{τὰ ὁμοιομερῆ} designate such things as have homogeneous
parts (the examples given are gold, silver, bone, flesh, marrow,
blood), and at the same time the parts themselves, since these are
represented as consisting in turn of homogeneous parts. In the
beginning all these \gk{σπέρματα} were chaotically mixed together, so
that the quality of no matter was visible; but after they had been in
this resting mixture for an infinite span of time, a separating-out
began through a vortex motion, which spread from a single point and
first separated the elementary opposites, fire and air, water and
earth, each of which is a mixture of many heterogeneous parts, and
then, as the separation advanced, separated out from the elements the
homogeneous small parts, through whose combination the concrete things,
consisting essentially of like parts, arose. Yet there is not yet a
complete separation in things, but in all of them, with the exception
of Mind, constituents of all other things are mixed in, if only in
small measure. The origin of that motion by which the combination and
the separation in which \opage{51}coming-to-be and perishing consist
are effected, Anaxagoras cannot seek in the sensibly existent itself,
to which motion is indifferent, nor in the empty, whose existence he
finds contradictory and therefore denies; and still less can he seek
there the ground of the order and beauty in what has come to be. He is
therefore compelled to seek it outside the sensible, in the
world-moving and world-ordering divine \gk{Νοῦς} [Mind], and by
positing \gk{Νοῦς} as what governs the world he names the fundamental
thought of Greek philosophy, and gives expression to the Greek thought
of the sovereignty of mind. \gk{Νοῦς} has power and knowledge and is,
in its unmixed purity, distinct from things and independent of them.
But despite these predicates, which Anaxagoras ascribes to it, and
which make it clear that he strives to hold it fast as distinct from
the material, as incorporeal, he is nevertheless able to separate its
essence from matter only in a quantitative way. It is ``the finest and
purest of all things'' (the purest, in that it is unmixed, whereas all
things are relatively mixed). But these determinations, by which for
Anaxagoras its independence, power and knowledge are conditioned,
indicate precisely only a quantitative difference. And when Anaxagoras
is to determine Mind's form of activity, he ascribes to it only the
production of a mechanical motion, of the vortex motion, the intuition
of which is taken from the revolution of the vault of heaven; and in
the explanation of the particular in existence he lets it again recede
behind the mechanical causes, which Democritus alone holds fast, in
that he consistently denies \gk{Νοῦς}.

Like so many of the philosophers mentioned earlier, Anaxagoras must
declare the senses insufficient to know what is in its truth; they are
not fine enough to distinguish the constituents of things. It is the
\gk{Νοῦς} in us, which is of like kind with the divine, that knows the
truth, and the knowledge of the truth is man's highest task. ``Man
exists in order to contemplate the heavens and the order that reigns
in the whole universe.''

\divrule

\opage{52}None of the various philosophical attempts that we have
considered as mediating between the Eleatic and the Heraclitean
fundamental determination has succeeded in solving the task they set
themselves: in reality they have been able neither to explain becoming
nor to uphold being. Becoming, which is now reduced to motion and
change of place, is not made comprehensible by the being that is
posited, but comes forward by a postulate. In Empedocles and Anaxagoras
it is derived not from the principles that represent being, but from
those which in reality are only a paraphrase of motion itself, and
which at the same time, as they involuntarily assume the shape of
material being, thereby point beyond themselves anew in order that the
origin of motion may be explained. In the Atomists motion is posited as
that which, as something without beginning, has no \gk{διὰ τί}
[wherefore], any more than the being of the atoms and their properties.
Empty space, non-being physically intuited, is named as the condition
of motion; but even if all motion, as change, is conditioned by a
non-being, this non-being must stand in an inner relation to being, not
be what is absolutely excluded from being; but empty space, the
intuited non-existent, has precisely no such relation to what is: the
atoms. The Atomists represent the atoms as originally falling in empty
space, but this notion of a falling presupposes that the falling atoms
are in a place other than that toward which they strive by their
nature, that they have been carried away from their ``natural place'';
that is, a motion prior to motion is presupposed, which thus remains
unexplained. --- And while motion thus stands unexplained, being is not
upheld. The Atomists, in positing what is not as the existent empty,
join the determination of non-being together with that of being. All
the thinkers named assert a plurality, in part an infinite plurality,
in being; but thereby the non-being of difference enters into the
existents, each of which is what it is only by not being what is
different. \opage{53}And as change is recognized as motion, the
non-being of change maintains itself in motion.

These defects in the attempts at mediation were unavoidable, because
the opposed principles were to meet within the same domain of
existence: within \emph{sensible} existence, where being, when it
cannot keep itself free of non-being --- (which is impossible, because
being, if all non-being were to be excluded from it, would itself, as
% printed "udelukkes fra, den, selv," (defect logged in the transcription)
indeterminate, become a non-being) --- cannot maintain itself through
it, and where becoming, which by its non-being disturbs being, for that
very reason cannot find a grounding in such a vanishing being, any more
than it could find one in a being that separated all non-being from
itself. But as the attempts at mediation fail, the principles must
again separate abstractly from one another, and the carrying-through of
this sharp separation had to lead, as was shown above in the survey of
this period, to the dissolution of the oldest Greek philosophy. This
process of dissolution is completed by the Sophists.

\divrule

\begin{center}
\textbf{9.\enspace The Sophists.}
\end{center}

In using the name Sophistic as a common designation for a philosophical
phenomenon, we do not thereby designate a closed and fully formed
philosophical school, but a tendency marked by the same essential
peculiarity, which, from the middle of the 5th century to somewhat into the
4th, asserts itself where philosophical endeavors come forward as a moment
within a phase of spiritual development that comprises all the spheres of
spiritual life and therefore has a general significance for the history of
culture.

In order to characterize Sophistic as a philosophical phenomenon, and in
particular to determine the relation between the philosophical element in
Sophistic and the general element belonging to the history of culture, we
shall proceed from what is reported about the most \opage{54}significant
individual Sophists, so as then to seek what is common and pervasive.

\emph{Gorgias} of Leontini (born c. 483, died c. 375, famous as a rhetor),
in his work \gk{περὶ φύσεως ἤ τοῦ μὴ ὄντος} [On Nature, or On What Is Not],
relying on the dialectical arguments of the Eleatics, gives proofs 1) that
nothing can be; 2) that even if something were, it still could not be known;
and 3) that even if it were and could be known, it still could not, by being
uttered, be communicated to others. For the first thesis the proof is
conducted thus: it is shown that just as little as what is not can be, since
this immediately contains a contradiction, just so little can what is be,
because all the opposed predicates, as leading to contradictions, must be
denied of it, so that its being cannot be held fast, since it becomes devoid
of determination, but being devoid of determination is equal to nothing, to
not-being. What is cannot be ungenerated; for then, as having no beginning,
it would have to be infinite, but the infinite cannot be anywhere, neither in
another, since then it would not be infinite, nor in itself, since it cannot
enclose itself; but what is nowhere is not. Nor can it have come to be; for,
as the Eleatics already showed, it can have come to be neither from what is
not nor from what is. Likewise what is cannot be one; for real unity would
have to be without magnitude, but what is without magnitude is a nothing. Nor
can it be a plurality; for plurality is thinkable only as a plurality of
unities, but these are not. And what holds of these opposed predicates holds
also of predicates such as moved and unmoved, etc.
--- The second thesis \emph{Gorgias} proves by this: that if what is were
identical with what is thought, then everything we think would have to have
actuality, and there could be no false notion, i.e.\ notion of something that
is not as actual. But if what is and what is thought are two different
things, then the former is not thought, i.e.\ not known. --- For the third
thesis, finally, he conducts the proof thus: even if there were
\opage{55} \setcounter{footnote}{0}
knowledge, it could not be communicated, because the sign is different from
what is signified, so that the word, e.g.\ the one that designates color,
cannot therefore produce the intuition of color. Nor is it possible that the
hearer can think the same thing by the words as the speaker, since the same
cannot be in different persons. And even if it were, it would still have to
appear to them as different, since they are different persons and in
different places. (By the last thesis Gorgias thus cancels all
\emph{universally} valid knowledge, just as by the preceding one he cancels
all \emph{objectively} valid knowledge.)

\emph{Protagoras} of Abdera (c. 490 --- c. 415, teacher of the art of
speaking and grammarian, the first who came forward as a Sophist, author of
various writings, among them \gk{τέχνη ἐριστικῶν} [Art of Eristics], and
\gk{ἀντιλογίαι} [Contradictions]) takes his point of departure from the
Heraclitean view of the constant change of things and the opposed streaming
of motions, and in applying it to the apprehending subject he comes to deny a
universally valid truth, and to let the expression for the momentary relation
between the thing and the sensing subject be the only truth. According to
Plato's presentation (in the Theætetos)
% printed "(i Theætetos" with the parenthesis never closed (defect logged in the transcription)
he reaches this result in the following way. Everything is in constant
motion. The motions can, in their infinite multiplicity, all be traced back
to two main classes: acting and suffering, (active and passive motion); but
both of these express a relation, and the properties of things, which arise
through motion, are therefore something relative and something merely
becoming. This holds also of our notions of things (\gk{τὸ αἰσθητόν}
% errata, p. 55: printed "αἰσθητίν", corrected "τὸ αἰσθητόν"; the \gk here is changed to the corrected reading (the transcription reads "αἰσθητόν")
[the perceived] and \gk{αἴσθησις} [perception] arise simultaneously), which
express only the momentary sensory feelings produced by the momentary contact
between things as acting and our sense organs as suffering. What is
represented is only in relation to us, determined by our changing
determinateness, our changing state. ``Man is the measure of
things''\footnote{\gk{πάντων χρημάτων μέτρον ἄνθρωπος}. \gk{τῶν μὲν ὄντων ὡς ἔστι, τῶν δὲ οὐκ ὄντων ὡς οὐκ ἔστιν} [of the things that are, that they are; of the things that are not, that they are not].},
and this means, since man is \opage{56}here understood only as that which
senses: there is no objective, universally valid truth, but only a subjective
semblance of truth, the individual, momentary opinion. As something seems to
someone, so it is for him; there is no other truth.

The somewhat later Sophist \emph{Euthydemos} uses both Eleatic and
Heraclitean points of departure, and from both arrives at consequences that
equally cancel all objective truth and all knowledge. From the thought of a
becoming that ceaselessly transforms all being he arrives at the proposition:
\gk{πὰντα πᾶσιν ὁμοίως εἶναι
καὶ ἀεί} [all things belong to all alike and always]; from the Eleatics'
doctrine of the unity of thought with being, at the proposition:
\gk{ψευδῆ λέγειν οὐκ ἔστιν οὐδὲ δοξάζειν} [it is not possible to say or to
believe what is false].

When objective truth and the possibility of objective knowledge are thus
denied, the task that remains for thought can only be, partly, the negating
activity itself, by which, through a merely dissolving dialectic, it
annihilates at every moment everything objective, partly the activity by
which the individual, for whom only the vanishing semblance of truth remains,
seeks, in the interest of egoistic ends, to give this semblance momentary
validity for others. In the Sophists, therefore, \emph{eristic} and
\emph{rhetoric} take the place of the knowledge of truth.

Most of the Sophists came forward as rhetors and as teachers of the art of
speaking. This is the Sophists' proper outward distinguishing mark. They come
forward as teachers of youth through imparting skill in speaking, eloquence,
and the art of the word in debate and discussion. In their own speeches the
epideictic is strongly prominent, and they display their virtuosity in
extemporized lectures on all subjects, in which their polymathy stands them
in good stead, but in which also a superficial reasoning and a hair-splitting
arguing for and against everything becomes the prominent thing. Rhetoric
becomes in their hands an art of persuasion resting on a calculation of the
prejudices and passions of the listeners, whose task is: to contrive for
others the semblance of truth and right that serves the
\opage{57}individual's contingent aim. In popular opinion the Sophist stood
as the one who understood \gk{τὸν ἥττω λόγον κρείττω ποιεῖν} [to make the
weaker argument the stronger].

The character of eristic as a dialectical art of disputation is designated by
Protagoras's proposition: \gk{δύο λόγους εἶναι περὶ παντὸς πράγματος ἀντικειμένους ἀλλήλοις}
[that there are two arguments about every matter, opposed to one another]. As
a dialectical art it becomes a formal virtuosity of argumentation, which has
its strength in the hair-splitting and the ambiguous, its aim in semblance,
its law in the individual. Its general means are the absolute distinction
between opposites (as when someone who says that he knows something is forced
to the admission that he must then know everything, because, when he knows
something, he is a knower, but the knower cannot in any respect be a
non-knower), and the use of the ambiguity of grammatical expression (e.g.\ in
the inference in Plato's Euthydemos: you have a dog and the dog has puppies,
\gk{οὐκοῦν πατὴρ ὢν σός ἐστιν,
ὥστε σὸς πατὴρ γίγνεται} [so, being a father, he is yours, so that he becomes
your father]).

In the domain of the practical the theoretical negativity was bound to have a
far-reaching effect, but the consequences develop here only gradually, as
Plato clearly lets appear in his portrayals of the Sophists. Originally the
Sophists came forward as teachers of virtue, and the concept of virtue
(\gk{ἀρετὴ,
ἀνθρωπίνη καὶ πολιτική} [human and political]) they understood in accordance
with the old usage of speech, which included in it at once all practically
useful skills and the capability and uprightness of character. Neither
Protagoras, Gorgias nor Prodikos, the most eminent in the older generation of
Sophists, placed themselves
% printed "sig sig" (defect logged in the transcription)
in principled opposition to the people's ethical view of life, or made the
distinction, which later became common, between natural and positive right.
But in the very existence of Sophistic lies a breaking loose from the ethical
tradition. In that a separate instruction in virtue is given, the influence of
the political public spirit and of attachment to the eminent men in the state
is posited as insufficient; theoretical education for civic life is demanded
instead of a merely practical upbringing, and a \opage{58}consciousness is
demanded of what virtue consists in, what is right and wrong and why it is so;
and thereby the old relation of piety is broken.

Now when the theoretical result is applied to the domain of the practical,
there can be no universally valid law in it, but the consequence must be that
everyone must be said to have a natural right to follow his arbitrary
opinion, and, when he has the power for it, to free himself from the
unjustified compulsion of positive law. The determining factor, the measure
of things, must in the practical domain too become the individual, but the
individual in its contingency, as arbitrary, egoistic and sensuous, and the
good can be understood only as what is useful, i.e.\ satisfying, to this
individual, not as something universally binding and universally valid.
Hippias of Elis, a younger contemporary of Protagoras, is the first to state,
with regard to what is practically valid, the opposition between natural and
positive right, and he disputes the binding character of the laws, because
they change so frequently and are so different in the different states,
whereas only what is universally accepted could hold as divine or natural
right. This principle is now soon stated generally. The laws are regarded as
arbitrary determinations that aim only at the advantage of those in power.
Natural right is the right of the stronger, and he who has strength need not
let himself be bound by those laws. Among the arbitrary legal determinations
were counted also the religious doctrines. The Godhead can least of all be an
object of knowledge; its worship is due either to a statesmanlike calculation
of what can call forth obedience in the citizens (Kritias), or to the
deification of what is useful in nature, such as bread, wine and fire
(Prodikos). Protagoras was already condemned as an ``atheist.''

From the foregoing we shall be able to extract the fuller determination of
what in general characterizes Sophistic in a philosophical respect. Sophistic
has shown itself, from the different points of departure, constantly to
arrive at the same negative, destructive result. It gives thinking a new
direction, away from the knowledge of nature,
\opage{59} \setcounter{footnote}{0}
which is made uncertain by thought, toward the subject, powerful through its
thought, and its interests. Not speculation on the objective in the world of
things, but reflection on the human, and in particular on the practical
interests of life, steps into the foreground. The subject, with the
consciousness of its mental power, demands to test everything, and to grant
it validity only when it has stood this test before reflection. It is the
right of subjectivity over against the objectively given that asserts itself
here. But the subjective consciousness that would test everything given is
the still abstract and formal one, which, in turning reflectively toward
itself, has not yet found in itself the point where thought, as it were,
meets the universal, conceptually determined essence of things, and thereby
again acquires objective validity. The sophistic consciousness can
therefore, in its relation to the given, assert its superiority not by
putting a higher content of knowledge in its place, but only by a breaking
loose from the given in the sphere of thought and of life, by a destructive
reasoning that dissolves the earlier content of knowledge, and by then,
starting from formal and abstract subjectivity, arbitrarily positing its own
determinations. These are then maintained by the subjective facility of
thought, which now comes to serve not the power of truth but the interests of
the individual. Plato draws a parallel between the relation of sophistic and
rhetoric to statesmanship and justice and the relation of the art of
adornment and the art of cookery to gymnastics and medicine.

How the philosophical element in Sophistic relates to the general foundation
in the history of culture will now also become evident. Philosophy is for the
Sophists not a primitive interest, but only a means for an aim proceeding
from a more general movement. None of the Sophists is in the stricter sense a
disciple of the older schools; they take up, partly
eclectically,\footnote{Thus, e.g., Gorgias uses Empedocles's propositions
about effluences from things that bring about the sensations.} motifs from
them, and use them in the interest of \opage{60}practical-egoistic ends. The
philosophical element is in Sophistic only something secondary. The Sophists
use the philosophical partly to give reasoning a greater semblance of
thoroughness by procuring for it a more general foundation, partly, through
the negative, to make room for the arbitrary determining of subjectivity.
Nor can the peculiarity of Sophistic in a philosophical respect be understood
from its philosophical presuppositions alone; it points back to a more
comprehensive movement that first properly explains it in its completeness.
From the philosophical presuppositions alone one would thus arrive only at a
skepticism with respect to the knowledge of the objective, not at that
assertion of contingent, egoistic subjectivity which characterizes Sophistic.
If, on the other hand, one seeks in the more comprehensive movement the
explanation of the assertion of subjectivity, and that as still abstract
subjectivity, and lets this become the point of departure, then it is easily
understood that this stage will use the consequences of the earlier
philosophical development that lead to a dissolution of philosophical
knowledge and of objective knowledge in general. For everything that would
establish itself as law can then be made ineffective by the destructive
dialectic, and by it room can be made for the arbitrariness of subjectivity.
--- Just as little does the change in the direction of thought that takes
place in Sophistic, as it turns essentially from the contemplation of nature
toward the sphere of the human, find its sufficient explanation in the
earlier philosophy alone. That the consciousness of mind as the highest
became clearer and clearer does not explain it; for we see in Anaxagoras, who
precisely recognizes as the highest the \gk{Νοῦς} [Mind] that is essential in
the Godhead and in man, that man's task is nevertheless determined as this:
to contemplate the universe and understand its eternal order. And the
Sophists' denial of a true knowledge contains by itself alone no ground of
explanation either, but would, since it would definitively apply just as much
to the human being and \opage{61}human affairs as to external nature,
consistently have to lead to thought's concerning itself equally little with
both spheres and giving up all philosophizing. But here too the given
phenomenon becomes explicable when account is taken of the more comprehensive
spiritual movement.

The general spiritual movement of which Sophistic, as a philosophical
phenomenon, is a peculiar expression, we can characterize in general as the
breakthrough of reflection into the common consciousness. This breakthrough
signifies that among \emph{that} Greek people with whom Greek philosophy, in
finding its fundamental thought, gets its proper home, the first main stage
of the life of the state has been left behind and a new one has begun. The
stage has been passed at which consciousness rested with secure immediacy in
what was handed down, in inherited mores and customs, opinions and
convictions, and spirit has now followed the impulse, originally lying in the
essence of self-consciousness, to make all this content of consciousness
transparent, to seek grounds for what previously asserted itself immediately
by authority, and by the grounding to justify it before the self as something
in accord with it. Such a reflection must come to light in every cultured
people at a definite stage of development, with the same necessity with which
it had previously to awaken in individual persons. But in its first
appearance reflection is not able to penetrate to the essence and deepest
ground of things; it moves around its object, as if attacking it at many
points, seeking, not the ground of the matter, but grounds for and against,
and in this seeking finding the final decision in the not yet clarified
subject with its interests, opinions, desires and passions, in the subject
that now feels itself the measure of things, and, in the absence of anything
universally valid and universally binding, makes its opinion truth, its will
law. It is thus this reflection that forms the general foundation of the
sophistic form of thinking, and its development was, in the time immediately
preceding the appearance of the Sophists, \opage{62}prepared in the common
consciousness by the historical events and the political conditions
influenced by them, in particular the development of democracy with the
striving for individual freedom and power that it awakened, by the rising
level of education, and by the influence of art and poetry on the spirit of
the age. The earlier philosophy itself made a contribution to such a
development by the separation it made between nature and convention, between
the knowledge of reason and the semblance of sense, by its polemic against
popular belief, by the development of the skill of thinking, and by its early
begun, if unclear, striving to posit the thinking mind as the highest. And
while the earlier philosophy thereby supplied a moment to the general
movement in the direction of the freedom of subjective thought, it prepared,
as touched on above, the negative character of Sophistic through the seeds of
dissolution that lay in the old philosophy of nature's one-sided limitation
to the domain of the sensuous, in the incompatibility of its fundamental
concepts, and in its lack of a consistent, principled separation of sensation
and thinking.

When, after what has been developed, we are to pass judgment on the
significance and justification of Sophistic, the judgment will be essentially
nuanced according as our point of view is that of the history of culture or
that of philosophy. In respect of the history of culture Sophistic has
exercised an essential, forward-driving influence on the development of
spiritual life. It has spread education and enlightenment, made science
common property, and spread a wealth of learning. It has striven to raise
the practical activity of life from the quasi-instinctive to the clearly
conscious. It has, in developing reflection, given it the courage to try
itself on everything. Finally, it has contributed to forming the Greek
language into a pliant and supple instrument for thought. These merits of
Sophistic must not be forgotten over its defects: over the superficiality and
boastfulness of so many Sophists, over the empty facility of reasoning that
so frequently took the place of a more serious search, and over
\opage{63}the egoism and self-will that in the practical and political domain
became the consequence of the theoretical presuppositions. --- In a
philosophical respect Sophistic has a lesser, almost merely negative and
preparatory significance. It does indeed assert the higher right of
subjectivity, but in an imperfect and superficial form, since the subject is
still understood by the Sophists only as sensuously determined in its knowing
and willing. It prepares a higher form of philosophy essentially indirectly
and negatively, by, as it were, making free room for the new thought through
its dissolution of the earlier thought. Only in particulars, through
contributions it has made to the theory of knowledge (Protagoras) and to
dialectic, and through its investigations of language, has it influenced the
development of thought more directly; but here too it retains a stamp of the
superficial.

The standpoint of Sophistic must be transcended, because in its principle of
subjectivity it does indeed express the higher principle of a new age, but
only in an untrue form; because the subjectivity that it asserts is one that
in its thinking is formal and abstractly destructive, and that, when it is to
posit a positive content, must fetch this from the sensuous side of the
subject, from the isolated individual subject, sensuous and egoistic as such,
in its contingency. It is transcended by there being found, in thinking
itself, whose power in Sophistic showed itself through the dissolution of
earlier conviction, a deeper, universally valid foundation for science and
ethical life. To this source of a new content of knowledge, higher than that
which the oldest Greek philosophy possessed, \emph{Socrates} points back, and
in doing so he opens a new possibility to the philosophy that had been
dissolved in doubt.

\divrule

\begin{center}
\textbf{10.\enspace The Character of the First Period in General.}
\end{center}
% errata (Rettelser leaf), p. 63: the words "Tilbageblik paa Udviklingen" deleted; printed "10. Tilbageblik paa Udviklingen. Den første Periodes Charakteer i Almindelighed."

If we look back on all the philosophical attempts of a positive character
that we have come to know in the foregoing, then \opage{64}it appears that
they all had in common the fundamental peculiarity that they sought the
principle of what is in matter, in nature. In this fundamental peculiarity
the essential distinguishing marks of this oldest philosophy lie enclosed. It
had essentially to confine itself to being philosophy of nature, physics, and
the ethical and dialectical determinations we come across could come forward
only as consequences of the physical theory or as unworked and isolated
material taken up from notion. Next, although an endeavor to secure for
thought a higher dignity early revealed itself, no sharp distinction could be
made in this period between the mental and the bodily, the non-sensuous and
the sensuous, and even where, at the close of the period, the concept of mind
began to clarify itself out of matter, the world-forming mind could be
understood only as a mechanically moving natural force and depicted
half-sensuously as the finest and purest matter. Finally, thought, precisely
because it had not yet distinguished itself with determinateness from matter,
could not turn reflectively toward itself in order to know the conditions of
knowledge; but it turned immediately toward the object without doubting the
truth of rational knowledge, and the statements about the nature of knowledge
that prepare the later philosophy of concepts are consequences of the
investigations into the objective constitution of things, not the foundation
for this investigation.

But although the oldest Greek philosophy thus, as it were, stops at the
boundary of the world of mind, we have nevertheless designated it as
\emph{the becoming of the doctrine of Ideas}, because what is the goal of the
development within Greek philosophy asserts itself, unconsciously, in the
first attempts of the still undeveloped thinking, and because it is pointed
to through their fundamental concepts and problems, and through the kind and
peculiar character of the thinking. The determinations given above of the
principles of being among the old Ionians, among the Pythagoreans and among
the Eleatics form an analogy to what Plato distinguishes as the spheres of
the material, of the mathematical and of the Ideas.
\opage{65}The two fundamental concepts, being and becoming, which govern the
old philosophy and around which its movement of thought turns, are, in a
deeper conception, precisely the fundamental concepts of the developed
doctrine of Ideas. The endeavor to bring them into harmony is the same that
repeats itself in a higher form in the doctrine of Ideas, and repeats itself
with the same lack of result. The principle of subjectivity that in Sophistic
acted as dissolving the older forms of thought is, in its imperfect form, the
preparation for the more developed one that is the condition for the thought
of the Idea being able to become clear, for the Idea being posited as the
essence of existence, and for a philosophy determined by this fundamental
thought developing. Finally, the kind and basic stamp of the thinking and the
form of the dialectic as a dialectic of distinction accord with the one
conditioned by the intuition of the principle of being as the hypostatized
ideal archetype, as the form from which negation is excluded.

\divrule

\begin{center}
{\Large\textbf{Chapter Two.}}
\end{center}

\divrule

\begin{center}
\textbf{Socrates.}
\end{center}

\markboth{Socrates}{Socrates}
Socrates, son of the sculptor Sophroniscus and the midwife
Phaenarete, was born in Athens about 470, and died, condemned for
the introduction of new divinities and the corruption of the youth, in 399.
From the sculptor's work, which he practiced in his youth,
he turned in his maturer years to the work he regarded as his divine calling: to awaken his fellow citizens, especially the young, to self-knowledge and to ethical action,
by preventing them, through his dialectic, from settling into
repose in an imagined knowledge and in self-deception. Of his mental
development Plato gives in his Phædon a picture that is, in its main features,
no doubt historical. From an occupation with the writings of the older natural
\opage{66}philosophers Socrates, unsatisfied, had turned
to the study of Anaxagoras's treatise, in which he hoped to
find an explanation of things from the point of view of purposiveness and by a divine reason. But here too he found
in the concrete only mechanical natural causes, which did not
satisfy his teleological-religious outlook, and he
now abandoned the consideration of natural things, which elude our
knowledge, in order to turn to the consideration of what lies
nearest to us: the human, and to ethical action. Of the
older philosophers Socrates no doubt came into personal contact with Parmenides, perhaps also with Anaxagoras. Of
the Sophists he was in contact with many, and his ethical
standpoint, which has the general basis of subjectivity in common with sophistry, developed in its peculiarity precisely through its opposition to sophistry.

When we assign Socrates a special place, and refer him
neither to the first nor to the second period of Greek
philosophy, although it was he who positively prepared the development in the second period, by demonstrating, in thought, which in
sophistry had been only destructive, the deeper foundation of science and morality, we do
so because he used his dialectic only as a serving moment for the ethical and because his thought ends in the negation of knowledge, in ignorance. Socrates places himself outside speculation; he becomes essentially an ethical-religious personality,
not a philosopher, because he \emph{will not} be one. But he gets
a place here because by his method, by his demand for
knowledge and by his conception of this knowledge he has occasioned a speculative philosophy and has become the incitement
% errata, p. 66: printed "med", corrected "ved"
for the whole later philosophical development among the Greeks.

In seeking to determine what is peculiar to Socrates,
we at once come upon a difficulty that might seem insuperable. Just as a multitude
of schools attached themselves to Socrates, all of which would express the Socratic, and yet went
in the most diverse directions, so, in the reports of
\opage{67}Socrates's ``doctrine'', the accounts diverge strongly, and comprise standpoints ranging from that of the right-thinking Sophist to that of the speculative philosopher proclaiming the doctrine of Ideas. But in the divergent, incomplete, trivializing or idealizing conceptions of Socrates there are certain points in which
the conceptions touch one another, and from which it becomes
possible to bring forth the figure of the real historical Socrates,
to clarify his standpoint, and to find the coherence in
the ethical view that corresponds to it, and that on closer consideration shows itself in its peculiarly Greek character.

It is reported of Socrates that, in as it were bringing
philosophy down from heaven to earth, that is, leading it back from the speculation on nature to self-knowledge, he occupied himself exclusively with the ethical. It is reported that he
determined virtue as knowledge. It is reported that he was the first
who, through induction from the particular, brought forth the universal
concepts, which were fixed in definition. And it is reported, finally, that he declared himself ignorant.

These statements are as it were the sum of what provisionally emerges, from the points of contact of the various reports, as a foothold for the conception; but it is characteristic that what
the various reports essentially agree on, taken
immediately, is something inconsistent in itself, something contradictory. The contradiction shows itself at once when one puts the statements together,
but precisely through the resolution of the contradiction we arrive at the conception of Socrates and at the understanding of the misunderstandings.

For if we first put together the two statements about Socrates's
method of induction and about his exclusive occupation with
the ethical, a contradiction at once comes in; for
when Socrates, as Xenophon puts it, constantly investigated with
those around him \gk{τί ἕκαστον εἴη τῶν ὄντων} [what each of the things that are might be], the
interest becomes a logical, not an ethical one. But on closer investigation of how matters stand with the result of Socrates's search for the concept of things and with his
\opage{68}definitions, it turns out that his inductive method --- more precisely: his method of concept-forming abstraction, --- either
leads straightforwardly to a negative result, since the attempted universal
determinations of the matter are overturned again by contradicting cases, or that it only apparently leads to a positive one,
namely to one that has already been shown untenable by
the preceding dialectical investigations (e.g.\ where the
good is determined as the useful, cf.\ the investigation in Plato's Laches). And precisely this seems to indicate that Socrates's
induction and definition, in general his logical activity of thought, must be conceived not as an end, but as something
that in one way or another serves as an instrument for something
else, and this something else would then most naturally be sought in the
ethical, as that on which Socrates's interest is, after all, constantly
said to have essentially rested.

That the definitive conceptual result of Socrates's inductions was always negative, and that therefore the \emph{final} goal
of the thought that ceaselessly returned to the same
movement could not be the logical one, a goal of knowledge, finds
its confirmation in what has been handed down: that Socrates declared himself
ignorant. For when, as tradition says, he
posited the concept (\gk{τί ἐστι}) as the true expression of the essential, and
would give the determination of the concept in definition, then,
had he really arrived at a positive result, he could not, however
much he ``saw the concept in its relation to the
subjective life'', without self-contradiction call himself ignorant;
for the content of the definition would precisely be a real
knowledge.

But if we now wish to hold fast to the indication that Socrates's logical interest served the ethical, a difficulty
again comes forward, since he, the ignorant one, determined
virtue as knowledge. That virtue for Socrates must have been an
actuality, no one will deny; for the opposite would not only
presuppose a wholly different standpoint from the Greek, but
conflict with all the reports of his personal life. But if
\opage{69}virtue was an actuality, its synonymous determination, knowledge, must be so too. Knowledge must be a reality, but
it becomes so only as knowledge of something, and whence should
Socrates's knowledge, which was only negative, a knowledge of being
ignorant, get this something? Of what is virtue a knowledge?
If the answer is: of the good, but the good as the ethically
good can in turn be determined only by knowledge, then we move
in a circle that we cannot get out of. But Socrates, if he was an ethicist, had to get out of it; for in order that
action may be reached, knowledge must have a content; its
object must have determinateness: form, as a Greek would
say. Now the historical records have Socrates determine
the content of the good as what accords with the laws
(just as he demands that one worship the gods \gk{νόμῳ πόλεως} [by the law of the city]),
and in Xenophon he is represented as determining the good
by the useful, hence eudaemonistically. But the dialectic in
the useful at once asserts itself: the concept shows its relativity, and the question, what is the useful? cannot be answered in positive form by him whose knowledge is ignorance, and an
% printed "sin.Relativitet" (defect logged in the transcription)
answer in this form is necessary if there is to be action,
and life is not to be absorbed in asceticism. If anything
certain is to remain in the determination, an intermediate determination that supports it is at any rate required. The same is the case when
the good is determined as what accords with the laws;
for the opposition between the laws among themselves, and between the laws and a knowledge which, in negative form in man, in positive form
in the divine, was higher than the laws, stood clearly before
Socrates. Above the contradictory laws of the states stand \gk{νόμοι ἄγραφοι} [unwritten laws],
but who is to read their unwritten script?

The difficulties and contradictions that have appeared can
hardly be removed by any other way than by conceiving Socrates
as an \emph{ethical-religious} personality, and by seeing in the ethical-religious the object to which his interest was
consciously confined, and which his dialectic, which must be conceived as purely negative, served.

\opage{70}It was touched on above that the logical interest was not
what in him was ultimately the impelling force in his
dialectic. Although he was the first who, as Aristotle
puts it, brought forth \gk{τούς τ' ἐπακτικούς λόγους καὶ τὸ ὁρίζεσθαι καθόλου}
(method of induction and definition), yet
his interest was not the logical in the method as such,
in the way that, e.g.\ Plato and still more Aristotle make
it the object of investigation; nor was it a positive result of the method: the determination of the concept; for we have
seen that such a result was illusory, and Socrates was,
as will appear, conscious that it \emph{had to} be so. But
if dialectic in Socrates could not definitively be an end, because the final goal must always be something positive, and even
in the Skeptics is the undisturbed repose of the individual in himself; if
instead, in its peculiarity, it had to be a serving means, then we must conceive it as the power that, through the dissolution of everything
sensuous and finite, brings forth something infinite, and this infinite, which Socrates would bring to breakthrough, was the as yet provisionally negative infinity of the self,
the point of departure for the ethical, for self-determination by virtue
of freedom and by virtue of thought. Precisely through the consciousness of ignorance the individual is to be led to the ethical, and dialectic, in which the thinking subject asserts its infinity
in the form of negativity, constantly keeps that consciousness awake, at the same time as it keeps the energy of infinity alive in the self. But if the negative, dissolving
dialectic could as it were clear the way back to that point of departure of the ethical, then --- as will be established
below --- a new moment was required at this standpoint:
the religious relation, in order to come out again from this invisible point, where thinking self-consciousness closes together with itself in freeing itself from all finitude, in
using its negative infinity for this liberation, and
in order to give this self-consciousness, which had as yet only
reached the infinity of negativity, a content of knowledge again. So\opage{71}crates could not, by the way of thought, by thought itself, give the abstract infinity of self-consciousness concrete determinateness. He
had to seek the filling of this infinity through something other
than his subjective thought. This other is the divine
thought, which for Socrates bore witness to itself in the purposiveness of existence, and is conceived by him as the \gk{φρόνησις} [wisdom] that is
\gk{τὸν ὅλον κόσμον συντάττουσα καὶ συνέχουσα} [ordering and holding together the whole cosmos].
% errata, p. 71: printed "συντάττων καὶ συνέχων", corrected "συντάττουσα καὶ συνέχουσα" (the Greek here is changed to the corrected reading)
This divine thought is in unity with the human, is its essence,
but in this unity it nevertheless becomes, although something ideal in the subject, as the object of an inner intuition, something separated from the subject's thought and personality, and speaks as a divine oracle placed in the subject's interior. And since
Socrates, through the religious relation to the divine,
can intuit the given as that in which there is a divine
order (cf.\ his teleological conception), he has, when
the oracle in his interior, the demonic voice (\gk{τὸ δαιμόνιον}),
speaks in the negative form of infinity, in its negative, warning pronouncement a corrective in the relation to this given. Where
the demonic voice does not speak, he confidently follows the law
as something divinely sanctioned. Unless the religious
relation becomes an intermediate determination, it is incomprehensible that Socrates, despite his dissolving dialectic, can again set up as the norm for
personal action the lawful, which in its particularity had already been made to waver by dialectic, and
the useful, which had been subjected to a consuming skepticism. Only
when the relation to the divine is drawn in can
the law, as something divinely ordained, again acquire firmness,
and the useful, positively determined, acquire a security in the midst of
insecurity, become an object of certainty in the midst of ignorance. For under the determination of the useful
Socrates can then bring in the relations that the general ethical
consciousness sanctions, and which he can conceive from the same
point of view as the lawful, namely as ordained by
the Godhead. The concept of the useful then acquires a more concrete content than it had so long as it comprised only the
\opage{72}consciousness of ignorance, ethical self-knowledge, as the negative condition for ethical action, and autarky, the negative expression of ethical
freedom, together with whatever can fall under
the point of view of means to its acquisition. As the expression of
this negative the useful is known by knowledge, just as ignorance is known, and just as the fundamental relation conditioning ignorance (see below) is known in its negative determinateness.

To that intermediate determination by the divine the
historical sources everywhere have Socrates point back when he speaks
of his life's task and of what definitively determines his
work. ---

The essence and significance of the Socratic standpoint, thus determined in its fundamental features, will become clearer through a comparison with two of its historical presuppositions: with the Sophists and Anaxagoras.

With the Sophists Socrates has this in common, that for him too the
essential interest of thought becomes not nature but
the human, and that he, like them, asserts the principle of subjectivity. But a decisive difference separates him from them, and this difference, beside which the likeness recedes, is grounded in the fundamental difference in the determination of subjectivity, and in the kind of thinking through which the energy of subjectivity expresses itself. In the Sophists subjectivity had asserted its negativity
through the dissolving dialectic, but for them the contingent
individuality in its finitude had still become the decisive thing: the measure of things; their dialectic turned dissolving and
destructive against the whole of objectivity, but not infinitizing and liberating against thought and the subject itself. The Sophists
were not too much dialecticians, but too little: they stopped
at something finite, at something positive which had only the illusory sanction of opinion and
desire, and did not bring to light the infinity from which something universally valid, truly
positive, was to be grounded. Precisely this is Socrates's task.
His dialectic, which has the infinite negativity of irony as its back\opage{73}ground, disturbs everything finite, all finite positivity, even that
posited by the individual himself, and reveals in this thoroughgoing
dissolving the inner infinity of self-consciousness. The goal of dialectic
is for him not to destroy and dissolve: through
the dissolution of the finite determinations, through the confounding of the content of opinion, of imagined knowledge, he lets
thought seek something fixed and abiding: the universal essence of things, expressed in the concept, which is as it were the formed infinite. In
striving toward this goal, it uses that method
which was indicated above: the method of induction, which takes
its point of departure from the particular in order, through a testing
comparison of it, to find and determine the pervading universal, the essence of the matter independent of the alternation of particulars and of the contingent individual's conception. But when the concept that is sought, and sought with the consciousness that thought would in
it grasp the essence of things, ``that which they are'', flees
from thought, which possesses only the infinity of the negative, without
being able from this, by itself, to win a new content:
a conceptual content, then the thinker, who cannot
come to rest in the finite and contingent, is forced over into the knowledge of his ignorance, and from this again over into the ethical-religious.
Here the subject, with the consciousness of ignorance
as its presupposition, finds a new content for its knowledge, secured by
trust in the divine, and by the corrective that the
demonic voice applies within. This voice,
which for Socrates sounds as a divine oracle, is, psychologically seen, an expression of the ethical tact of the subject purified and freed through the negative
infinity of irony, of the involuntary reaction, unconscious in its coming-to-be, of ethical feeling against the particular untrue and unethical (inharmonious) action.

But that Socrates, although knowledge for him ends in ignorance, can nevertheless set a knowledge of the concept of things as
the goal of thought and determine virtue as knowledge, becomes explicable through the other historical presupposition mentioned above:
% printed "Forndsætning" (defect logged in the transcription)
\opage{74}Anaxagoras's doctrine, and ultimately through a regard to the
Greek fundamental view.

In the Socratic determination of virtue as knowledge we must
be reminded of the Anaxagorean philosophy's doctrine of \gk{Νοῦς}, whose
significance for Socrates is not diminished by his having found himself
unsatisfied with the Anaxagorean theory in its execution.
For Anaxagoras \gk{Νοῦς}, the divine thought, was the highest power of reality; for Socrates likewise, although this
view is not stated directly, nor becomes the point of departure for
a dogmatic doctrine, the ordering and governing principle is the
divine thought (reason, \gk{φρόνησις}, \gk{γνώμη}), which in man, according to his essential determination, is as knowledge, i.e.\ as a
knowing of what is abiding and universal in things, and this conception forms as it were the translucent background of
his whole view. That this is so is easily
shown. When Socrates determines virtue as knowledge, there lies
herein first, since virtue according to the Greek conception, which we see
consciously carried through in Plato, is the working of a being according to its
peculiar, undisturbed nature, that knowledge is man's
essence. But the knowing thought, which is man's essence,
becomes at the same time the essence of the objective, and, as divine thought,
what rules in existence. This shows itself partly in that
thought in Socrates is not a thought torn loose from objectivity,
but a thought that posits the objective as its
object in order to express its essence in the concept formed by thought; partly it shows itself through those statements of Socrates about the gods which, even if they are not to be taken in a strictly
dogmatic sense, nevertheless indicate a relation of God to the world
like that of the soul to the body, while there is a likeness of essence
between the human soul and the divine (Mem.
IV, 3; cf.\ I, 4). Now indeed the knowledge eternally actual in the Godhead becomes for man, in uncertainty, ignorance, but
there is nevertheless in Socrates no skepticism about the reality of knowledge in
and for itself, about its essential relation to human nature;
it is intimated (cf.\ with the passages in Plato, especially in the Phæ\opage{75}don, the words that in the Kyropædie VIII, 7, 20 are put into the mouth of the dying
Cyrus) that in death it again becomes actual through recollection, which is no longer hampered by anything obscuring,
and through the religious it again acquires a positive relation
to objectivity. It is, then, when knowledge in and for itself has
reality, but as human knowledge becomes ignorance,
existence with its relativity and opposition that prevents thought from penetrating reality masteringly,
and precisely thereby leads the existing one from speculation to action; but essentially thought is knowledge, one with its object,
which is to be grasped in the concept. Knowledge is in Socrates, in personal form, the same as the Idea is in Plato.

Thus Socratic ignorance, while it is
the object of knowledge, has a background of knowledge, which nevertheless never gains shape in a comprehension of the concrete. The apparent contradiction in such a knowledge is removed when we go back to what became, for Greek thought, the fundamental relation in
existence: the relation between the moments of the world's work of art,
matter and form (cf.\ above p. 16 f.), and through this the whole Socratic standpoint at the same time receives its definitive explanation.
For in that fundamental relation there is enclosed a twofold
consequence with regard to human thought. When
that relation on the one side shows itself to us in such a way that
form, which realizes \emph{itself} in matter, is the true and sole
existent, then from this point of view form is seen as what rules unconditionally, as that which alone determines the essence of what is formed,
and therein lies, by virtue of the kinship of form with \gk{Νοῦς},
the consequence that thought (the concept) is the essence of what is,
and with regard to human thought, that this, as
what rules in the human being and determines the human being, has, by virtue of its unity with the divine \gk{Νοῦς},
the possibility of penetrating, as knowing, objectivity, whose essence is thought itself. --- But when the relation on
the other side shows itself in such a way that matter is seen, as it were in a different light, as the con\opage{76}dition, \emph{not} posited by form, for the realization of form; then matter acquires a kind of
independence of form, a kind of independent working. And
this relation of independence then leads, within the sphere of the relative
(cf.\ Plato's and Aristotle's physics) --- in the totality,
namely, the relation is undisturbed and the divine knowledge
therefore an undisturbed knowledge (cf.\ Xenoph. Mem. I, 1, 8; IV, 6,
7; Plat. Apol. 20, 23) --- to a negative relation, to a relation of opposition between matter and form, and the human
mind is seen, by virtue of this, although it is by form essentially
determined as knowledge, since it does not freely master the heterogeneous matter but is obscured by it, as that whose essential knowledge during the existence of the individual, i.e.\ during the
conjunction of form with what is heterogeneous, takes on the form of ignorance. Precisely in the difference between these two consequences is the difference
given between Socrates on the one side, and Plato
and the whole of Greek speculation on the other. Socrates, who
in his individuality as it were symbolically expresses the disharmony
between mind and the material, stops at ignorance,
gives up speculation; but since the ideal, form, has the significance of reality for him as essential determination and goal, he does not lose
himself in a theoretical skepticism, but strives to actualize the ideal through action, and the ethical-religious
becomes his definitive interest. Plato holds fast, even if the
other relation also comes to light in him, as also in Aristotle, though only to be easily dismissed, to the essential determination knowledge as a determination of reality, and knowledge becomes
for him speculative knowledge, whose object is the Idea; his
interest becomes philosophy.

Socrates, and he alone, represents wholly and consistently a standpoint of mind that lies as a consequence in the
Greek fundamental view and is completely explained by it.
Therefore the Greek background shines through everywhere in him, even where he seems to depart most decisively from
his time and his people, and where he touches the fundamental ethical thought of a distant
future. A skeptical standpoint he avoids
\opage{77}by the fact that the ideal, form, has not vanished for his consciousness in ignorance, but has validity and actuality
for it; so that the essential knowledge, the disposition for knowledge, remains
in the midst of ignorance. But form, which from his point of view
shows itself to thought in its original being \emph{beside} the
presupposed matter (cf.\ p. 18 f.) --- thus it must here be seen that
the two factors, despite the preponderance of form, are equally original ---
becomes negatively infinite, since in the separation it is that which
forms nothing, the formlessly infinite form. The knowledge
that remains, despite ignorance, thus becomes a knowledge of a
\emph{formal} relation; but a knowledge which, owing to the determinateness of this relation, \emph{cannot} be filled out in the knowledge of the
actual, and, when it turns toward this, gets only ignorance as its sure content.

With the Socratic fundamental view of knowledge in
its relation to existence set forth here (which finds an analogy in Greek
poets: in the conception of death as reconciling), the
essential determinations of the Socratic ethics can be placed in close connection.

Since knowledge is the expression of essence, knowledge is according to its concept
what rules (cf.\ Plato's Protagoras, 352, C f.), over which passion and sensuousness have no power. Knowledge and action are inseparable; no one is voluntarily evil. Sin is ignorance; but ignorance, whose \emph{first} possibility becomes a riddle
that Socrates does not pursue, is like a jesting mask, vanishing with existence, which, as what is personally unmerited,
takes the impersonal guilt with it. The mind becomes free of
guilt in being loosened from what obscures. --- As knowledge,
virtue is furthermore one: there arises neither a separation between the different forms of manifestation of virtue, nor between the virtues of the different individuals; for that which should give virtue
diversity, the natural, whether it is posited as the object of action or as the natural determinateness of the individual, is what is
separated from form (the ideal), which knowledge represents,
and is, over against knowledge, in its concept something subordinate, which
\opage{78}\setcounter{footnote}{0}cannot bring separation into what has unity as its essence.
--- Since the individual is essentially knowing, it is essentially complete, and a twofold demand arises upon it. Over against the natural it is to be \gk{αὐτάρκης} [self-sufficient], to realize its
ethical independence in freedom from dependence on everything external,
not in an ascetic manner, but by mental freedom in being able to
enjoy and to do without, so that this autarky becomes the expression, corresponding to the conditions of existence, of the positive: that the knower is what rules, since the principle of knowledge, rational
thought, is the essence of all. Being the ruler is expressed in existence by not letting oneself be ruled by anything else\footnote{Through passion and sensuousness we become unfree (\gk{ἀνδραποδώδεις}),
and lose, together with the right, clear insight, that which is the most essential
condition of happiness, and which in the midst of privation preserves good fortune
in preserving the freedom of the soul. Moderation and self-mastery are
an unconditional demand; they are \gk{ἀρετῆς κρηπὶς} [the foundation of virtue], that which gives it
firmness, that without which the insight won is obscured again in
the individual.}. --- Over against other individuals the individual is to be itself
independent and to acknowledge them as independent. This consequence is necessary whether the individuals are conceived as
ignorant or as essentially knowing. In the first case
the one cannot teach the other anything because he himself knows nothing;
in the latter he cannot because the other essentially knows everything.
The task of the one individual in relation to the other can only
be to give it occasion for self-reflection, to \emph{awaken}
it to ethical self-knowledge. It is to \emph{help} the other individual to the consciousness of its ignorance, but in eliciting this consciousness it is presupposed that the other individual
has in itself what furnishes the standard for ignorance, and
therefore essentially comes to consciousness by itself. That
awakening takes place as the imagined knowledge, which is ironically acknowledged by him who confesses his own ignorance, is, by
being tested, by being forced to give an account of itself, brought to
waver through dialectic, as the in\opage{79}\setcounter{footnote}{0}dividual wrapped in illusions is as it were stripped, until it finds itself naked before
itself, and now stands where the leap can be made, recollecting,
to take itself back into the Idea, into the essential relation, unless this
recollection, as by Socrates himself, is ethically broken off by action through the consciousness of the conditions and claims of existence.
That bringing to self-knowledge, which for Socrates was a life's task, is what he, with a playful reminder of his mother's work, calls his maieutic, his midwife's art. It
is the help to give birth to \emph{oneself}: the rebirth that Greece
knew. The relation between the individuals who have already attained ethical consciousness becomes friendship, which, in Socrates's conception, presupposes the autarky of the friends and their equality in being good, and whose essential form of manifestation is
the mutual, continued prompting to self-knowledge and self-development through the dialectic of the ethically awakening ``conversation''.

The mutual independence and the essential
completeness of individuals lead to the same consequence as the fact that knowledge in every single existing individual turns into ignorance, namely to this: that virtue \emph{cannot} be taught, that is, taught
\emph{by others}. Socrates says expressly that no teacher of virtue could
be found, just as he himself declares that he has not
been anyone's teacher (Plato, Apol. 33, A f.), and how indeed should a teacher be found among the ignorant?
When it is said in some reports about Socrates that he
taught the opposite, this is partly only a false inference
from the Socratic statement that virtue is \gk{ἐπιστήμη} (\gk{σοφία})\footnote{In Aristotle, being communicable by teaching is synonymous with
being capable of being grounded. For him, then, it is necessity and
universal validity that make something teachable (cf.\ Plato's Timæos 51, E).}, partly
a confusion of learning from oneself through self-reflection
(``recollection'') and practice with learning from others. If, then,
virtue is to come to be in the individual, the individual, who is posited as essentially knowing, is referred to himself, to call the essential knowledge to life through self-reflection in the negative and
\opage{80}formal shape in which it can be had, and in which it
can be the foundation for the negative form of the ethical, in order then,
as it turns into ignorance in relation to the concrete, to seek a knowledge that corresponds to the positively ethical,
through the religious. That self-reflection is in Socrates
the counterpart of the Platonic recollection, through which
knowledge comes to be as positive knowledge of the Idea.

These are the general, essentially formal,
determinations of the Socratic ethics. If the ethical consciousness was to have a
concrete content, this could not, according to what was developed above about the fundamental relation of thought, come about through an unfolding of \emph{thought}, and thought could find the way only inward, toward
the infinity of self-consciousness, the point of departure of the ethical,
but not outward, from infinity to the unfolding of the ethical in action: in this movement knowledge became ignorance for the existing
individual. The infinity of the ideal (cf.\
p. 76 f.) could reveal itself to the mind only in its negative
working as moving power in the dialectic that consumed
the finite content of opinion, or in ecstasy, in which the self
loses itself intuiting in the projected infinity, which, captivating, spellbinds consciousness into immobility (cf.\ the accounts in Plato's Symposion of Socrates's ecstatic states),
but does not gain shape and unfold itself. In the intuition of
the infinite, finite knowledge is negated; the infinite
becomes the shapeless, ununfolded mighty. In order to arrive at
action full of content, ethical action, the religious relation therefore became indispensable at this standpoint, and through it we saw the
good acquire concretion for Socrates.

\divrule

In Socrates, then, we have the supplement to the speculative direction of Greek philosophy. In his ``ignorance,'' which is yet only as it were an inhibited knowledge, a decisive consequence of the Greek fundamental view comes forth, and it unfolds itself in a circle of ethical determinations. From Socrates \opage{81}a new and higher philosophical development takes its point of departure, and the impulse he has given (cf.\ p. 66) works through Plato and the imperfect Socratics right down to the close of Greek philosophy. The most typical expression of Greek philosophy, the Platonic philosophy, attaches itself immediately to Socrates. When the Greek spirit, which by its nature could not permanently content itself with making the movement of irony toward ignorance, in the urge to hold fast also the harmony of finite human thought with what is, a harmony that had been broken by the Socratic conception, once more, proceeding from the intuition of the immediate unity, postulated the likeness of \emph{reality} to \emph{essence}, the Socratic essential knowledge cast off the form of ignorance and came forth as the existing thinker's knowledge of the Idea. The concept, which in Socrates was posited as the unattained goal of thought and was sought as a regulator for the life of the individual, is now, while it is seen as the essence and truth of thought and as the norm of all human life, conceived as the essence of the objective world \emph{accessible to thought} and as \emph{reality being in and for itself}. Thus the Platonic doctrine of Ideas emerges from the Socratic attempted formation of concepts; the ethical standpoint is transformed into a speculative one, and this knows in the standpoint from which it has emerged what is akin to it in essence.

\divrule

\begin{center}
{\Large\textbf{Chapter Three.}}
\end{center}

\divrule

\begin{center}
{\Large The Second Period of Greek Philosophy: The Unfolding\\
of the Doctrine of Ideas.}
\end{center}

\divrule

\begin{center}
\textbf{1.\enspace Survey of the Second Period.}
\end{center}

\markboth{The Unfolding of the Doctrine of Ideas}{The Unfolding of the Doctrine of Ideas}
To the second period of Greek philosophy we reckon the so-called \emph{imperfect} or \emph{one-sided Socratics}: \opage{82}the philosophers of the \emph{Cynic}, \emph{Megarian} and \emph{Cyrenaic} school, together with the two philosophers who in union represent Greek thinking in its highest unfolding, \emph{Plato} and \emph{Aristotle}, and \emph{their schools}, the Academic and the Peripatetic, in their first generations.

The imperfect Socratics retain from their master the interest in the ethical and the determination of virtue as knowledge (insight), and this Socratic element they all seek to fuse with one of the main directions of the older philosophy. In this attempt there appears the impossibility of calling forth a fruitful development by a one-sided combination of the new principle of knowledge with one of the fundamental concepts of the earlier philosophy, and the necessity, in order that such a development might take place, of granting each of the fundamental thoughts of the older philosophy its place. The task of finding a new fruitful principle for philosophy, and of forming a system that comprehends in itself the essential content of the whole earlier philosophy and in its richer world-view assigns both of that philosophy's opposed principles their domain, is first solved by Plato. He is the one who actualizes the new higher form of thinking and carries out what Socrates had intimated and prepared, and his edifice of thought finds its necessary supplement in the Aristotelian.

If we characterize the second period by its main phenomena, its opposition to the first shows itself in this, that the principle of being and of reality is no longer sought in the world of matter, but in that of thought, of spirit, in the concept that expresses the essence of things and is in and for itself: in the Idea, the form. The principate of spirit which the older philosophy divined, which the words of Anaxagoras, only half understood by himself, pronounced, which the Sophists asserted negatively, Socrates positively but outside the domain of philosophy, is now, within that domain, pronounced as the solution of the riddle of existence. Thought, which has turned from nature to spirit, from the contemplation of what objectively is to the contemplation of its knowledge of it, now investigates the conditions of knowledge, and as this investigation \opage{83}leads to positing the content of the concept, the Idea, true reality, as the sole object of real knowledge, and to keeping, in a conceptual distinction, the eternal essence of the Idea, of the spiritual, of what in truth is, clearly apart from the sensible and changeable, new philosophical sciences, conditioned by the knowledge of the Idea, can take shape: dialectic and ethics, and these gain precedence over physics, which is transformed by its relation to the Idea.

It is Plato and Aristotle who create this culminating form of Greek thinking; they essentially fill the second period, and beside them the one-sided Socratics, who make only powerless attempts, step into the shadow, and gain new significance only in the span of time after Aristotle, by giving impulses to Stoicism, Epicureanism and Skepticism. Plato and Aristotle supplement each other, in that their principles express two sides of one and the same thought, the fundamental thought of Greek philosophy, and in that the two formulas they set up for the relation between the fundamental moments of existence point back, in their difference, to the same fundamental view, of which they are the only fully Greek expressions, and which together enclose the whole essential content of Greek philosophy.

In the genesis of the concept of the Idea lies the motive for the twofold determination of the essence of the Idea and of its relation to matter. As thought, which finds the Idea, finds itself again in the Idea, and now, giving up matter as principle, determines the Idea as what truly is, as that which gives existence being, as the eternal that is not touched by becoming, there remains at this stage, at this first revelation of the Idea, no other determination for the moment distinct from the Idea than the non-existent. The Idea separates itself sharply from this in essence, and a gaping chasm divides the spiritual, the non-sensible object of knowledge, from the material and, as such, unknowable. But this material element, which is the foundation of the sensible world, thought cannot, even if, as it were through a subtraction of the being that the sensible world owes to the Idea from \opage{84}the sensible world's being, it determines it as a non-existent, yet hold fast in this determination; it cannot press outside being. When finite existence is to be comprehended, the material element asserts its place as a co-determining factor in it, and now the contradiction comes forth that the non-existent nevertheless receives a kind of being and activity, and that an unexplained world of phenomena, with an imperfect being which yet cannot be denied the character of being, comes, seen in its sensuous determinateness, to separate itself from the world of Ideas with its perfect and complete being; while in the intuition of the cosmos the whole of real existence is seen in harmonious unity with the world of Ideas. These difficulties are the ones at which Plato stops, and they lead Aristotle to set up another formula for the relation, one that lets both the Idea and its opposite, both form and matter, as he names them, be conceived under the determination of being, but within this concept makes a distinction between a being according to actuality (\gk{ἐνεργείᾳ}) and a being according to potentiality (\gk{δυνάμει}). Thereby the original connection of essence is designated. Matter is the potentiality for that which in the form has become unfolded, true actuality, and matter desires the form, strives toward it as toward its essence, while conversely the form is the indwelling form which as entelechy works in matter, as principle of development leads it to actuality, and in the whole of the world's existence works as the divine principle of unity of the harmonious universe. But in Aristotle too the system cannot be concluded, the ultimate ground of energy cannot be determined, without the form in its purity, as divine \gk{Νοῦς} [mind] (and as active \gk{Νοῦς} in man), separating itself from all matter and from all motion, and through this separation, at the point where essence appears most purely, dualism again gains dominion.

Through these two conceptions of the Idea in its relation to the material the doctrine of Ideas unfolds itself, and only these two conceptions become possible where the relation of the opposites is to be determined from the standpoint of \gk{Κόσμος}, of the harmonious unity of the equally original \opage{85}moments of existence under the principate of the Idea. Either, from this standpoint, the ideal becomes the only thing that gives existence reality and truth, or it becomes that which alone is an existent in the highest sense, but by its nature stands in a necessary relation to matter, as to what is different in form of existence, alike in kinship of essence. But whereas the first conception, through its absolute opposition between the existent, the Idea, and the non-existent, matter, had to lead to the existence of the sensible world becoming inexplicable and the being of matter contradictory, the second, which, by conceiving both opposites under the determination of being, avoids the inner contradiction of the former, and, by seeking the ideal in external reality itself, can acknowledge this reality in its essentiality, shows a twofold defect grounded in the very essence of Greek thinking: that the activity of form in matter, by which the unity was to reveal itself and the concretion of existence to be explained in its coming-to-be, nevertheless remains uncomprehended, and that with this the unity of principle, energetically striven for and as it were touched, dissolves again into dualism. But since the Aristotelian conception of the relation of the principles, which would unite the opposites separated from one another in a unity of thought under the concept of being, and thereby marks the highest striving of Greek philosophy, is the last many-sidedly Greek one, because it is the only genuinely Greek one beside Plato's, the philosophy that arises after Aristotle, and that with clearer or dimmer consciousness of the rupture strives after a unity which sacrifices one of the opposites, and therefore does not in truth reconcile the opposites, becomes the period of dissolution of Greek philosophy.

\divrule

\begin{center}
\textbf{2.\enspace The One-sided Socratics.}
\end{center}

What is common to these has been stated above. It was the interest in the ethical, the determination of virtue as knowledge (insight), and the endeavor to fuse this \opage{86}Socratic element with one of the main directions of the older philosophy. Their difference is grounded in this, that while the \emph{Cyrenaic} philosophy takes up an element from the \emph{Heraclitean} direction, in the shape in which it had been prepared by \emph{Protagoras}, the two others take up an \emph{Eleatic} one, but in such a way that in the \emph{Megarian} school the main weight falls on the dialectical, in the \emph{Cynic} on the practical. The former accentuates the \emph{knowledge} that in itself is virtue, the latter the \emph{virtue} that, one in essence with knowledge, expresses itself in \emph{action}. --- We now consider the individual schools briefly.

\begin{center}
a) \emph{The Megarian School.}
\end{center}

The \emph{Megarian} school had its name from its founder \emph{Euclid's} native city, \emph{Megara}. Euclid had as a young man become acquainted with the Eleatic doctrine; later he became one of Socrates's most devoted disciples, and tradition tells of him that when, during a war with Megara, the Megarians had been forbidden on pain of death to come to Athens, he stole into the city disguised, by night, merely to hear Socrates. He was present at Socrates's death, and after this catastrophe several of Socrates's disciples fled to his home in Megara, among them Plato, on whom Euclid had no small influence. Euclid seems to have lived and presided over the school he founded for some 20 years after Socrates's death. Of his writings none are preserved.

Among the philosophers of this school the best known are \emph{Eubulides}, \emph{Diodoros Kronos} and \emph{Alexinus}. \emph{Menedemos}, who came out of the \emph{Elian} school founded by Socrates's disciple \emph{Phædon} of Elis, which essentially agreed with the Megarian in dialectical character, fused the doctrine of the two schools in that of the \emph{Eretrian}, while the Megarian direction had, in \emph{Stilpon} (teacher of the Stoic Zeno; he appeared in Athens c. 320), united itself with Cynicism. For the rest, the reports we have of the history of the Megarian school are few and disconnected.
\opage{87}Since Euclid combines the Socratic determination of the good as the highest object of knowledge with the Eleatics' doctrine of the one being known by reason, he determines the good as that which is one, even though it is named with many names, and as that which is unchangeably like itself, and he acknowledges only one virtue: knowledge of the good, inasmuch as what are designated as different virtues are only different names for the same virtue. In him the concept of the good is transformed from an ethical concept into a metaphysical one, and coincides with being. What is opposed to the good thus becomes a non-existent (\gk{τὰ ἀντικείμενα τῷ ἀγαθῷ ἀνῄρει μὴ εἶναι
φάσκων}). But the same lack of content and development that became the stumbling block for the Eleatic doctrine had now to become so for the Megarian as well. The one did, it is true, apparently receive closer determinations through the statement that the good was named in many ways, was called now insight (\gk{φρόνησις}), now deity, now reason, etc.; but precisely because the difference was determined only as a difference of names, the content vanished again. It does not seem, however, that the first generation of the Megarians strictly held fast to the exclusive unity; at least in one place in Plato's ``Sophist,'' where the Megarians are doubtless meant, there is talk of \gk{νοητὰ ἄττα καὶ ἀσώματα
εἴδη} [certain intelligible and incorporeal forms] in the plural, and these must then presumably be thought of as standing in the same relation to the one as the plurality of thought concepts to the unity of thought. Yet the connection of this determination with the original point of departure is not clear.

Since the Megarians brought the concept of the good into connection with the Eleatic being, the ethical \emph{could} not unfold itself; for the good in its unity was without determination, and the possibility of a determinateness was excluded by the fact that the manifold, here ethically conceived as that which is opposed to the one good, was posited as a non-existent. Here we see again the significance of Socrates's being an ethical-religious personality. He could, through the religious, lead knowledge back into action; among the Megarians knowledge was determined only, so to speak, retrogradely; it
% printed "retograd" (defect logged in the transcription)
\opage{88}ended in abstraction. The ethical interest must therefore soon be exhausted and give place to the dialectical. It is also essentially dialectical determinations that have been handed down to us as stemming from the Megarians. It is already a dialectical determination that the opposite of the good, of the one, is posited as a non-existent. From the conception of the fundamental concept there follows as a consequence the denial of a connection between the existent and the non-existent, of anything merely possible (i.e.\ of an existent that includes a non-being). Only what is necessary is actual, and only the actual is possible. But with the concept of possibility motion is eo ipso abolished, although it was attacked on particular grounds by the later Megarians, notably by Diodoros Kronos, who determined the Eleatic polemic against it more closely with regard to the doctrine of the Atomists. --- On the abstraction in the conception of the existent there also rests the protest against asserting predicates of the subject, and against applying the universal concept to the individual things designated by it. The abstraction of the Megarian dialectic soon had to lead over into the eristic. Various captious arguments are enumerated, which are attributed especially to Eubulides and Alexinus. Among these the best known are ``the Concealed,'' ``the Liar,'' ``the Horned,'' ``Sorites'' and ``the Bald.'' They can be traced back partly to a verbal subtlety (e.g.\ the Horned: What you have not lost, you have. You have not lost horns. Therefore you have horns); partly to the opposition between thought and sensation, \gk{λόγος} and \gk{αἴσθησις}, (e.g.\ the Veiled, or Elektra: Does Elektra know Orest, who stands before her? Yes! he is, after all, her brother. But in his disguise she does not recognize him. Therefore she both knows and does not know him at the same time); partly to the dialectic in the quantitative (e.g.\ Sorites: One grain makes no heap [\gk{σῶρος}], nor do two or three, etc., but at last, by the addition of one grain, a heap does come to be; and in the same way a man becomes bald by losing one hair). --- In Stilpon the abstraction in the Megarian \opage{89}philosophy receives a peculiar practical expression, as he posits pure apathy, insensibility to pleasure and pain, as the ethical end.

\begin{center}
b) \emph{The Cynic School.}
\end{center}

The founder of this school is \emph{Antisthenes} (born c. 444), who,
after having earlier heard Gorgias, joined Socrates at an already riper
age, attached himself to him with the greatest inwardness, and was present at his death.
He seems to have lived about 30 years more after Socrates's
death. The name of the school which he founded after his master's death
derives from its meeting place, the gymnasium Kynosarges; perhaps, however, the school's name also contains an allusion to
the way of life of its adherents. Of Antisthenes's writings nothing is
preserved. --- Of the later Cynics, \emph{Diogenes} of \emph{Sinope}
and \emph{Krates} of Thebes, who was also the teacher of the Stoic Zeno,
are the best known. After having been absorbed into Stoicism,
which gives it a nobler form, Cynicism reappears in
the first century after Christ, but without acquiring
philosophical significance.

\emph{Antisthenes} holds fast to the Socratic determination: that
virtue is knowledge (insight, \gk{φρόνησις}), and that it is a unity,
and he conceives virtue, which alone is our property, not merely
as the highest, but as the only good, as necessarily
bringing happiness with it, and as inalienable once it
has been acquired. His statement that virtue is something that can
be taught (is a \gk{διδακτόν}) seems to be understood thus, that
what can be taught is the negative: \gk{τὰ κακὰ ἀπομαθεῖν} [to unlearn what is bad], and
Antisthenes hardly made any sharper distinction between learning from others and teaching oneself by practice. --- When he
is to determine the content of virtuous action, the contentlessness of the knowledge
identified with virtue, which he does not remove
in the Socratic manner but precisely sharpens by his Eleatic-sophistic
dialectic, compels him to determine virtue negatively.
It becomes independence, autarky, a sensuousness-negating
\opage{90}self-mastery. For the agent who cannot give
his action concrete content can assert himself through action only
in an abstract way, i.e.\ only by making himself count in his
negative independence. Cynic autarky becomes essentially different from the Socratic. Socrates could, in the midst of his
autarky, enjoy; the Cynic could not: he feared enjoyment; the only interest for him became the purely abstract
virtue, and this virtue, which alone sufficed for happiness, was,
even though it be brought into the closest relation to insight
(knowledge), not to need many concepts (\gk{λόγοι}) or items of learning,
but above all Socratic strength of soul. Here too the divergence between the Megarians and the Cynics shows itself. The former turn off in
the direction of thought, the latter in the direction of action, but both
remain in the abstract.

Since being virtuous is determined negatively, as independence, this determination is applied most immediately to the relation to the sensuous desires and the sensuous needs. Pleasure must then, when it becomes an aim, be determined
as an evil, because it holds us fast in the unfreedom of desire and need. ``Rather mad than feel pleasure,'' says
Antisthenes. But if pleasure is an evil, toil, exertion (\gk{πόνος}), is a good, because it strengthens ethical self-mastery. What does not fall under the determination of the
good and the evil is an \emph{Adiaphoron}, an indifferent thing. But
this concept of the indifferent acquires a wide extension,
and with it the significance of autarky is widened and modified.
Among the adiaphora the Cynics had consistently to count all the
objective, concrete forms of family life and the life of the state, which
precisely on account of their concretion could not be brought into
connection with the abstract principle of virtue. Not the laws of the state,
but those of virtue, were to be the norm. And while the Cynics freed themselves from the laws of the state, they also freed themselves
from popular belief. Not the many popular gods, but only the one
Godhead has truth; only virtue is true worship of God; the
wise man is to be independent of those notions. The wise man,
\opage{91}i.e.\ he who has attained autarky, becomes the ethical norm for
himself. But since the decision in the last instance in all individual
cases is to proceed from the individual who, precisely by one-sidedly
throwing all the weight on action, cut himself off from the moral
purification that knowledge and scientific culture give,
it is after all once more the contingency of individuality that
becomes decisive, and it depends entirely on it whether
the ethical is not to be transformed into its opposite. This in
reality became the case with the later Cynics. From the polemic against pleasure they came to an opposition against what,
as the concern of tact and of the sense of beauty, occupies so large
a place in Greek morality, and while independence of
the laws of the state led to a contesting of them and to the preaching
of a hazy cosmopolitanism, the concept of the indifferent was widened
in such a way that for the independent wise man it even factually took in
pleasure itself, and that in its most unbeautiful form, stripped of modesty and restraint. The wise man's independence
at last turns over into dependence on that which was to be subdued, and
on the opinions of others, his virtue into egoism, vanity, and
coarseness. --- Cynicism ends, in practice, in the individualism which,
as a principle, had from the beginning been the distinctive feature of the Cyrenaic school.

The logical determinations and dialectical propositions that
are traced back to Antisthenes all seem to aim at showing, in
the interest of action, the impossibility of theory. He does indeed determine the essence of definition as stating the
concept of things, what they are or were (\gk{τὸ τί ἦν ἢ ἔστιν}), and defines
knowledge as correct opinion accompanied by a conceptual account (\gk{δόξα ἀληθὴς μετὰ λόγου}). But since he, in
the manner of the Sophists, one-sidedly holds fast to the unity of what is with
itself, he denies the validity of every judgment that determines the subject by a predicate different from it, so that
only tautological judgments, e.g.\ Man is man,
% printed "Mennesket er Menneske det Gode er godt" with no comma after "Menneske" (defect logged in the transcription)
the good is good, remain. Thereby, however, a genuine definition is made impossible, and Antisthenes also denies
\opage{92}the possibility of one with respect to the simple, which
can only be named and compared with something else, while he
admits, with respect to the composite, an explanation
through an enumeration of its constituents according to their real
connection. Of the simple, the elements of the composite,
there is only a correct notion, no knowledge. --- From the determination of the relation of the thing to its proper designation
follows as a consequence Antisthenes's paradox that contradiction
is impossible. For when the selfsame thing is spoken of, then,
since it can have only one proper designation, only
the same can be said of it; if, on the other hand, different
things are spoken of, then opposite statements contain no contradiction. Against
Plato's doctrine of Ideas Antisthenes polemicized without the motive
emerging clearly from what has been handed down. His endeavor
to lead the individual from theory to action may in any
case have been a contributing factor.

In philosophical respect the Cynic school acquires only
slight significance. Its essential significance lies in the
influence which, through a strict ethical demand and through a merciless combating of follies and vices, it has
exerted on its effeminate contemporaries, for whose moral misery the Cynics, despite all their proud self-enclosure, nourished a deep
compassion. But the ethical life the Cynics could awaken in
others could still not acquire a different character from the one it had
in themselves, and in referring the individual to himself, to
self-enclosure in autarky, they abandoned him to the same
consequences from whose power they themselves had not been able to escape.

\begin{center}
c) \emph{The Cyrenaic School.}
\end{center}

\emph{Aristippus} of Cyrene (according to doubtful reports born
435, died after 356) founded this school, in which the Socratic principle of knowledge seeks the filling-out that Socrates gave
it through the religious, through the sensuous determination
in us, through the pleasure which knowledge, insight, is to shape
\opage{93}\setcounter{footnote}{0}harmoniously. \emph{Aristippus} seems for a long series of years to have
belonged to Socrates's circle, but already earlier to
have occupied himself with Protagoras's doctrine, through which
he was influenced by Heracliteanism. After Socrates's death he seems
to have appeared in various places as a teacher for
pay. Aristotle calls him a sophist. --- Of the members of the Cyrenaic school we name, besides Aristippus's daughter
\emph{Arete} and her son \emph{Aristippus the Younger}, the three
men in whom the consequences of hedonism develop:
\emph{Theodoros}, \emph{Hegesias Peisithanatos}, and \emph{Annikeris}.

As the goal of man's striving Aristippus determines
--- seizing the eudaemonistic element that in Socrates
could find a place in the reflection on the ethical, ---
\emph{pleasure} (\gk{ἡδονὴ}), as \emph{positive} pleasure, not as freedom from
pain, as the \emph{particular} pleasure\footnote{--- --- \gk{εἶναι τε τήν μερικὴν ἡδονὴν δι' αὑτὴν αἱρετήν· τὴν δ' εὐδαιμονίαν οὐ δι' αὑτήν, ἀλλὰ διὰ τὰς κατὰ μέρος ἡδονάς.} [and that particular pleasure is choiceworthy for its own sake, but happiness not for its own sake, but for the sake of the particular pleasures.]}, the pleasure of the present
moment, not as a general condition of life: happiness (\gk{εὐδαιμονία}),
% printed "τήν" with an acute in the note, l. 1 (defect logged in the transcription)
and as \emph{bodily} pleasure. Pleasure \emph{as such} is a good,
and between the different pleasures there is only a difference of degree.
That pleasure must become the goal of life is supported by a reasoning that is much reminiscent of Protagoras. All our experience
is only a consciousness of our affections; it restricts
itself to these, without our being able to state anything about the affecting
object and its properties and without our being able to decide
whether our individual being-affected by it agrees with that of other
individuals. Only in our individual feeling can we therefore seek the norm for our action. The goal must be the
feeling that appeals to us most, and that is the feeling of pleasure.
Pleasure is the feeling of a gentle motion of the soul (\gk{λεία}
\gk{κίνησις}), unpleasure the affection that arises from an uneven and
violent motion (\gk{τραχεῖα κίνησις}). The complete
rest of the soul is neither pleasure nor unpleasure.

But when Aristippus determines the goal by pleasure, and
\opage{94}seeks in the stated manner to ground this doctrine, he
demands, for the realization of this goal of pleasure, and for
the determining of the form of enjoyment, what points back to the
Socratic: \emph{insight}. Insight is to extract from every relation of life
the pleasure it contains; it is to be able, through consciousness
of the consequences of pleasure, to choose rightly among the different pleasures,
equal in and for themselves, that offer themselves; it is to free from false
notions, from prejudices and fancies; it is to remove
longing for what has vanished and desire for what is to come; it is to procure freedom in enjoyment itself (\gk{ἔχω οὐκ}
\gk{ἔχομαι} [I possess, I am not possessed]) and that independence of consciousness which in every situation gives us contentment in the present. It is, in
a word, to make possible an \emph{art of enjoyment}, which makes it
possible for us, with the freedom and security of consciousness, to give ourselves over
to the enjoyment of the moment, and thus to extract the greatest
possible pleasure from life. Pleasure as the goal and insight as
the essential condition are thus placed in the closest
connection, and philosophy is designated as the way that leads
to the truly human life: the life of the wise man, which includes the greatest possible pleasure. But the relation of pleasure and insight
nonetheless at once becomes uncertain as soon as one recalls that
pleasure is expressly posited as the pleasure of the moment, because only the
present belongs to us, and because pleasure, as the ``gentle motion'' that takes place in the one who feels, as sensuously
determined, is in its arising something vanishing, which must be seized
in the moment, and can indeed be preserved as a ``pleasure of the soul,'' but
then only as a fainter recollection, as an echo of the
proper, i.e.\ bodily, pleasure. For when pleasure, the goal of life,
is a matter of the moment, the one who enjoys, who relates himself
to this goal, must also live in the moment, concentrate himself in it;
and precisely herein lies the difficulty for the determination of the relation of insight to pleasure. Insight would not, at this standpoint, be able to procure
pleasure in a positive way, for
then one would have to go beyond life in the moment, and
the standpoint would be abandoned. The task would then most nearly
\opage{95}become that of warding off what disturbs, and then letting the
pleasure that the course of life brings with it fill the moment unhindered. Under this point of view the function of insight, as it was stated above, would then have to be understood.
But here too insoluble difficulties arise. It cannot
be made clear how insight, whose content could \emph{only}
be the consciousness of pleasure as the goal of life, and as
being in the moment, should be able to give consciousness independence of the contingent filling of the moment, which can also
be a filling merely with pain. Pain must be
disturbing to happiness, since consciousness may not abstract
from the moment and does not have a higher content of life present
in itself. And precisely because consciousness does not have a content that
is higher than enjoyment, neither can it be made clear how insight should be able to give a freedom in the enjoyment
that is to fill the human mind bound to the moment.
Nor can it become clear how insight can choose
between the pleasures according to regard for their consequences; for on the one hand
every pleasure as pleasure is a good; there is no difference of kind
between them, and even the difference of degree vanishes when life's
dissolution into discrete moments takes away the standard of comparison;
on the other hand regard for the consequences points beyond the moment,
which conflicts with the presuppositions. There then remained only
this: to push aside everything that, in the shape of disturbing
notions, envy, longing, and desire, intruded
upon the enjoyment of the moment, and to let insight bring it about that the concentration in the moment became complete, that past and future ceased to be for the consciousness of the one who enjoys. But this
most extreme consequence, which would make the cultivation of insight impossible, since no room would be left for such cultivation in the life destined for
enjoyment, would directly contradict the essence of insight
as consciousness; for consciousness contains the universal, it
cannot be bound to the moment and to the pleasure of the moment, but
its essence constantly points out beyond it. That demand for
a joining of pleasure in Aristippus's conception and in\opage{96}sight thus proves indeed to be a demand that cannot be avoided, because the one who enjoys is a man, but at the same time a
demand that cannot be fulfilled. The two moments, which in
their essence are heterogeneous, must again separate from one another,
and the standpoint must run through a development in which the
inner contradiction reveals itself, and that mind asserts itself from which this ethical standpoint, which would seize
man in the midst of existence, as sensuous, and as consciously
enjoying this sensuousness, and as freed in it by insight, has emerged as a one-sided and yet relatively justified
form.

For Aristippus the individual stands, as a living work of art, in the midst of the world of sensuousness, in its motion and
flux, and its task is what Aristippus, to whom it was
given \gk{χλανίδα φορεῖν καὶ ῥάκος} [to wear both the fine cloak and the rag], attempted in his own life with brilliant
virtuosity: artistically to enjoy the pleasure-filled moment.
But we saw that in his task the contradiction showed itself
that thought was to hold fast the moment, concentrate itself in the moment, while its nature is precisely not to be in the moment,
and, precisely as it consciously sinks itself into it, to reach out beyond it. And the content of the moment always remains the contingent,
what is foreign to thought. But thought is the mightiest,
the norm of the human mind; it must therefore, when the demanded
unity of pleasure and insight is dissolved, assert its power
% printed "oplöses" with a German ö (defect logged in the transcription)
by displacing the other element. This, of which we already
could seek an intimation in the demand for freedom of mind in enjoyment, since this demand indirectly posits enjoyment as
subordinate to the life of the mind, we see happen with the later representatives of the Cyrenaic school, on whom Plato's and
Aristotle's polemic against the principle of pleasure has no doubt exerted its influence.

\emph{Theodoros} (Atheos), with the consciousness of the impossibility of bringing insight into a positive relation to the
particular pleasure, posits pleasure and toil (\gk{ἡδονὴ} and \gk{πόνος}) as indifferent
(\gk{μέσα}): significance for him belongs only to \gk{χαρὰ} [joy] and \gk{λύπη} [grief], which
% printed "χαρὰ" with a grave, p. 96 l. 35 (defect logged in the transcription)
\opage{97}arise from \gk{φρόνησις} and \gk{ἀφροσύνη} [folly], and \gk{χαρὰ}, which designates the
wise man's goal, designates an abiding state, not something momentary. But the content of \gk{φρόνησις} can be nothing other than
the consciousness of the indifference of pleasure and unpleasure, and the content of \gk{χαρὰ} becomes the wise man's empty, but precisely therefore
% printed "χαρὰ" with a grave twice more, p. 97 ll. 1 and 5 (defect logged in the transcription)
unlimited, self-feeling in the independence this consciousness gives. The wise man is \gk{αὐτάρκης} [self-sufficient]; all objective relations lose
their significance for him, and Theodoros's statements about the independent wise man sometimes meet with those of Cynicism. In \emph{Hegesias} the conviction of the contingency and relativity of all pleasure, and of the impossibility of realizing Aristippus's principle
even in the form Theodoros gave it (because the soul suffers with
the sufferings of the body, is constantly in motion and exposed to chance, and because the calculation of insight is unreliable), leads to
positing painlessness as the goal and determining the negative:
escaping evils, as the only thing in which the wise man has the
advantage; and despair of giving life a content of enjoyment
leads him to posit life as the indifferent, death as
something alluring. Through indifference toward life itself
the wise man becomes \gk{αὐτάρκης}.

It was only apparently a return to the original
form when \emph{Annikeris} made the positive and
particular pleasure the goal; for this pleasure he tied to the quality of an
action or of a disposition, let it spring
from it, and he taught that one should sacrifice the greater pleasure when
it came into collision with goodwill, gratitude, friendship,
or love of the fatherland. Here in reality,
% printed "Kjærllghed" with two undotted l's (defect logged in the transcription)
even though the sympathetic pleasure is once more sought to be led back to
the idiopathic, the content of pleasure is sought in a foreign domain, that
of the moral, and the original principle is abandoned. Annikeris
concludes the Cyrenaic school in time; its ideal conclusion
had already been given in Hegesias.

This Cyrenaic doctrine of pleasure proves, in its appearance
and development, to be a peculiarly Greek phenomenon. It is
the \emph{morality of beauty} that Aristippus proclaims, in that he, from the standpoint of
\opage{98}sensuous beauty, sees the life of the individual as the sum of the harmonious moments, in each of which sensuousness is harmoniously mirrored in the ideality of consciousness. The
aesthetic, which in Plato and still in Aristotle is a moment in ethical action, is here everything, and the world of reality
is by an illusion transformed into the timeless world
of beauty, which becomes the individual's stage until the oppositions of time and reality cause the illusion to burst. And
as a norm for life is now sought, the development moves hastily
away from that starting point in pleasure, and ends in its sharpest
contrast: in death. Plato too demanded a dying-away, but
through the Idea he rescued thought out of the world of the changeable
and let the soul die away from this only in order to find its life
again in the eternal existent. The Cyrenaic, who would keep himself within
the becoming of the world of sense, is forced by its dialectic out into
death, and into the death that is not the beginning of the Idea, but the
cessation of what becomes. The doctrine whose principle is enjoyment ends
by speaking so seductively of what stands in mysterious connection with enjoyment, of death, that those who heard
it chose it voluntarily, and Hegesias, the Cyrenaic school's
proper concluder, received the name Peisithanatos.

\divrule

From the brief outline we have given of the lesser Socratic schools it has become evident that the philosophical element in
them was weak, and quickly exhausted. None of them has been able
to create a fruitful development. Their philosophical significance lies
most nearly in this, that by carrying through with vigorous one-sidedness
the individual sides of the Socratic, they place the problems that lie in
it in a sharper light; that they preserve the forms of thought of an older time for the consciousness of a later time; and that they prepare
such phenomena as we, after the culmination of Greek philosophy
in Aristotle, see appear in Stoicism, Epicureanism, and Skepticism. To seize the thought that lay
hidden in Socrates's seeking, and that was the central thought of Greek \opage{99}consciousness, was granted among Socrates's hearers only to
\emph{Plato}.

\divrule

\begin{center}
\textbf{3.\enspace Plato.}
\end{center}

a. \emph{Plato's Life and Writings}. --- \emph{Plato} was, according to the
most reliable reports, born in Athens in 428 or 427 before our era. His
father Ariston was descended from Kodros; his mother Periktione was related
to Solon. As a boy and youth he was instructed in the knowledge and skills
that belonged to the education of a freeborn Athenian: in grammar,
gymnastics, and music; he took part in the Isthmian Games and, after his
eighteenth year, perhaps in some minor campaigns. At the same time he
occupied himself with poetry; but he abandoned this occupation when,
through his association with Socrates, he had come to the conviction that
philosophy was the most excellent of the arts of the Muses. But the poetic
gift in him, which he consciously subordinated to what was for him the
higher, has left its stamp on his writings; beauty became the form of his
thinking, his prose became a work of art. Already as a young man Plato
became acquainted through Kratylos with Heraclitus's doctrine, and perhaps
through private study with others of the older philosophers; but when, at
twenty years of age, he joined Socrates's circle, Socrates became the one
who, by his dialectical art and his ethical view of life, gave him the
impulse decisive for his whole life as a thinker; and Plato in turn became
the one who not only attached himself faithfully to his master with
personal devotion, but who, alone of all the Socratics, was able to find in
what Socrates had intimated the fruitful thought that was to transform
Greek philosophy; and, always mindful of the impulse from the master who by
his death was lifted up into a sphere of ideality, he carried back in his
writings his whole rich and distinctive life of thought to Socrates, whose
image he draws for us with poetic truth.

After Socrates's execution he stayed for a time, together with several
Socratics, with Euclid in Megara. From Megara \opage{100}Plato is said to
have begun his first great journey, which is said to have taken him to
Cyrene, to the mathematician Theodoros, and from there to Egypt. Having
finished it, he is said to have lived again for a time in Athens until,
about 389, he undertook his journey to Magna Graecia and to Sicily. In
Magna Graecia he presumably sought chiefly, through personal contact with
the school's most important representatives --- Archytas, Timæos, and
others are named --- to come to know Pythagoreanism from its theoretical and
practical side. In Sicily he formed a friendship with the young Dion; but he
tried in vain to influence the tyrant Dionysios the Elder, with whom the
relationship at last became so hostile that Dionysios had him treated as a
prisoner of war, handed over to the Spartan envoy, and sold as a slave on
Aegina, where one Annikeris of Cyrene ransomed him. When he had returned to
Athens (c.\ 387), he founded his philosophical school, which had its seat
first in the grove by the Academy, later in a garden close by the hill of
Colonus. Here he taught, surrounded by numerous disciples; but he twice more
left his native city to visit Sicily: the first time, c.\ 367, in the
illusory hope of winning the younger Dionysios for philosophy, for a moral
life, and for a lawfully ordered sole rule; the second time, c.\ 361, to
bring about a reconciliation between Dionysios the Younger and Dion, which
failed so completely that the attempt even put Plato's life in danger.
After his return home from this last journey he taught undisturbed until his
death, which is said to have overtaken him at a feast in his eighty-first
year, 348---347. From the political life of Athens, which agreed so little
with his conception of the true life of the state, he kept aloof, like
Socrates. And like his master he placed his highest political task in this:
himself to find, and to impart to his disciples, the true philosophical
knowledge from which alone, in his conviction, a true life of the state
could proceed.

As soon as Plato was dead, legend twined itself luxuriantly about his
figure, and the admiration, whose object none of the thinkers of antiquity
has been \opage{101}\setcounter{footnote}{0}in such a degree as he, gave the
legend growth through the centuries. Already in the first generation after
Plato, legend made Apollon, the knowing god, his father. His countrymen and
strangers honored him after his death with statues and monuments; the most
imperishable monument he has raised to himself in his writings, and time,
which has ravaged so mercilessly among the most glorious written works of
Greek antiquity, has reverently spared what Plato wrote.

Together with Plato's genuine writings, various others have been handed
down to us under his name, which were partly designated as spurious already
in antiquity, partly have been separated out as such by modern criticism.
They are all minor and less significant writings. About the genuineness of
the writings on which our acquaintance with the Platonic doctrine
essentially rests there prevails no doubt; only a smaller group of
importance, the dialogues Parmenides, the Sophist, and the Statesman, has
been attacked in the most recent times, yet not on sufficiently convincing
grounds.

Plato's writings are all, with the exception of the Apology, composed in
dialogue form. They mark a progression in his development from the Socratic
to the distinctively Platonic\footnote{In the foregoing the relation has been
elucidated between the Platonic speculative knowledge and the Socratic
ignorance, between Plato's doctrine of Ideas and Socrates's attempted
formation of concepts. As distinctive of Plato we further emphasize here,
provisionally, the transformation of the Socratic \gk{φιλία} [friendship]
into the concept \gk{ἔρως} [love]; the conception of recollection; the whole
view of nature conditioned by the doctrine of Ideas; and, in the domain of
the ethical, the articulation of the one virtue into the four cardinal
virtues.}. In one group of dialogues (Laches, Charmides, Euthyphron, and
Protagoras) Plato does not yet go beyond the Socratic, of which he presents
the negative-dialectical side and the general ethical principle. In these
dialogues the conclusion is either negative, that is, not merely without a
positive result, but such that one is cut off, or one arrives only at the
abstract formal determination of virtue as \opage{102}knowledge and as one.
In the dialogue Lysis (on friendship), on the other hand, traces, if faint,
of the distinctively Platonic come forth. This now becomes clear through the
remaining dialogues: partly through the development of the general
fundamental presuppositions of philosophy, as in Phædros; --- partly through
a thoroughgoing grounding of the distinctive standpoint of the Platonic
philosophy, either dialectically searching, as in Parmenides, or
polemical-critical, as in Theætetos, the Sophist, and the Statesman; ---
partly through the treatment of a single fundamental philosophical problem,
as in Menon, and in part in the Symposion and Phædon; --- partly through the
development of the fundamental ethical principle, as in Gorgias, or of the
metaphysical foundation of the ethical, as in Philebos; --- partly through a
coherent constructive presentation of the application of the doctrine of
Ideas to the state and to the philosophy of nature, as in the unfinished
trilogy: the Republic, Timæos, and Kritias. In the Laws, which were
published only after Plato's death and were no doubt heavily revised by one
of his disciples, Plato has depicted a form of state, the nearest to the
ideal one, but adapted to the given conditions. Of the remaining writings
not yet named, the little dialogue Kriton belongs to the purely Socratic
dialogues and attaches itself by its content to the Apology of Socrates;
Hippias major and Kratylos treat more isolated problems (the beautiful and
language); Euthydemos gives a lively picture of sophistry in its
degeneration into empty eristic and dialectic of word-play, and illuminates
it by the contrast with philosophy, the true educator of youth. Menexenos,
whose genuineness is very doubtful, is an occasional piece: a speech over
fallen Athenians.

With regard to the time at which the various writings came into being, one
has very few reliable points of support to hold to. In the ancient authors
only little information is found, and the allusions to definite historical
events which some have sought to use to determine the date of composition of
individual dialogues (e.g.\ the Symposion's) are not certain, since Plato is
said to have constantly worked over and altered his \opage{103}writings. The
inner development gives guidance only in broad outline. Only a few of the
dialogues have been placed by Plato himself in such a relation to one
another that their relative chronological order is clear. Moreover,
provided only that the writings in which Plato is still dependent are
separated from those in which his distinctive standpoint appears, the
question does not have the importance that the enormous literature heaped up
over it, especially in Germany, might lead one to suppose.

In antiquity there were still preserved notes of Plato's oral lectures on
the Ideas and the Good, i.e.\ the Good in and for itself, the Idea of the
Good. Such notes had been made by Aristotle, Herakleides Pontikos, and
Hestiæos. They contained a completion of Plato's doctrine of Ideas, in which
he had taken up strong elements from the Pythagorean doctrine of number.
These notes are lost except for fragments, which are preserved in Joh.
Philoponos and Simplikios, and have been collected and published by Brandis.
But these \gk{ἄγραφα δόγματα} [unwritten doctrines] signify not a higher stage
of development of Plato's doctrine, but a decline of his dialectical power,
for which a substitute is then sought in something inferior. These fragments
are therefore without essential significance for the presentation of
Plato's doctrine.

The form of the dialogue is not an accident in Plato's writings; it is
closely conditioned by Plato's conception of knowledge. In Plato's
doctrine of the soul, of its ideal pre-existence, and in his concept of
recollection, which makes all teaching merely a calling forth of what is in
us, it lay as a consequence that the Socratic conception, which made the
individual essentially his own teacher, must also become Plato's. For him
too all communication had to become an awakening of independent activity of
thought through an awakening of recollection, and the learner had to be
conceived as essentially learning from himself. When, therefore, Plato
wished to communicate himself through writing (which he regards as the more
imperfect means in comparison with the living and ensouled speech of the
one who knows, of which writing is only a faint copy, unable to defend
itself and liable \opage{104}to produce a sham knowledge), he had to seek a
form by which the recipient becomes essentially self-active and, as
recipient, must constantly be not merely reproducing but independently
judging. It is this that Plato means to achieve by the dialogue form and by
the artistic presentation of thought. Already in the dialogue form there
lies an annulment of directly instructive communication. Since the reader is
compelled by it to follow attack and defense, often through many and subtle
windings and through apparent digressions, self-activity is awakened; and
this holds especially of the form of the dialogue found in Plato's earlier
writings; for in some of the later ones (the constructive ones) the dialogue
assumes a different character and then essentially presupposes the
independent development already. But in Plato it is not merely the
conversational form that awakens independent productivity; it is just as
much, or even more, the artistic manner in which the result is frequently
hidden, so that now only single hints of it remain, now the whole
investigation ends in contradictions, to whose solution, however, guidance
is again found in the dialogue itself, but not directly. In order that the
reader may be able to understand the dialogue, it is therefore demanded of
him that he independently gather all these scattered hints, that he grasp
the real center of the investigation, set everything episodic in the right
relation to it, and then judge independently. Nothing is given immediately;
everything must be won by self-activity.

\divrule

b. \emph{The Genesis of the Platonic Philosophy and Its Relation to Earlier
Thought}. --- As those among Plato's predecessors who exercised the
strongest influence on him while his distinctive thoughts were taking form,
we must name, besides Socrates, Parmenides and Heraclitus. Early on Plato
had become familiar with Heraclitus's doctrine of the flux of all things. To
this doctrine he grants validity for the \opage{105}domain which it
originally concerned: the domain of the world of sense. Plato too makes
everything sensible subject to becoming, and he therefore determines it as
an intermediate between being and non-being, and as that of which there can
be no knowledge, but only an uncertain opinion. But instead of stopping
here, Plato, being led by Socrates to seek something abiding and
unchangeable in the universal essence of things expressed in the concept,
points to this as something other than the sensible, as something that is
higher than the world of sense and not subject to the power of becoming, and
which can therefore be the object of a true and real knowledge. This new
world, the world of the eternal and supersensible, is the world of the Idea.

It was already shown above (p. 81) how Plato in the concept of the Idea
grasped what constantly fled from Socrates's thought, and how he made the
content of true, \emph{real knowledge}, of conceptual knowledge, become the
\emph{objectively real object}, that which is in and for itself, so that this
in-and-for-itself existent, the supersensible, the Idea, is at once that
which alone possesses true being and alone can be the object of true
knowledge.

By this concept of something in and for itself existent, true and abiding,
independent of the change of the world of sense, Plato also touches on
Parmenides's doctrine, and Plato acknowledges the great Eleatic as his
predecessor. He meets the doctrine of the Eleatics in the determination of
abiding, unchangeable being, apprehended by thought, as the true reality.
But whereas the Eleatics, although so many of their determinations negated
space, could not raise the determination of the existent above the
spatially existent, Plato gives the existent purely ideal determinations.
Whereas the Eleatics held fast to the unity of being so unconditionally that
being became devoid of determination, Plato shows that, so conceived, it
would become equal to nothing and no longer be an object of knowledge; and
he therefore lets the unity of the Idea unfold itself into a circle of
Ideas, of which \opage{106}each has the determination of the true existent,
but at the same time a concrete content, by which it becomes a member of an
organized unity. And whereas, finally, the being of the Eleatics became
immovable, Plato seeks to vindicate for the Idea activity, motion, life,
and reason.

But if Plato received his strongest impulses from the thinkers named, he
has nevertheless, while standing in living interaction with the
contemporary Socratics, taken up and assimilated essential elements from
\emph{all} the earlier philosophical attempts. This holds first and foremost
of the doctrine of the Pythagoreans, but it holds also of that of
Anaxagoras, and, if in a lesser degree, of those of Empedocles, the
Atomists, and the older Ionians. Plato is as it were the personified memory
of Greek philosophy, and as he expresses in the Idea the thought of Greece,
he sums up, in the doctrine he builds on this thought, the principles that
governed the thinking of the past, and the essential content of Greek
philosophy.

\divrule

c. \emph{Polemical-dialectical grounding of the doctrine of Ideas}. --- Polemically, Plato grounds his standpoint through a demonstration of those defects of ordinary consciousness
and of sophistry by which these standpoints of consciousness
prove insufficient to lead to true knowledge and
true action.

With respect to the standpoint of ordinary consciousness
in its theoretical aspect, Plato shows that, as confined
to the domain of sensation and opinion, it does not attain to a
true and actual knowledge. It lies in the nature of sensation that
it presents one and the same thing to us in the most contradictory manner, and the sense-image in general expresses only
the form under which things appear to the subject (thus
something subjective, a \gk{φαντασία}), not the matter in its objective validity. And opinion (\gk{δόξα}), in which the soul as it were busies itself within
itself with what is (\gk{αὐτὴ καθ' αὐτὴν πραγματεύεται}
\opage{107}\gk{περὶ τὰ ὄντα}), and is to that extent independent of the external and more independent than in immediate sensation, is not, either,
even if it is right opinion (\gk{ὀρθὴ δόξα}), the goal of
knowledge. Since there is a false opinion (\gk{ψευδὴς δόξα}), it becomes evident
that opinion does not by its nature exclude the untrue,
which knowledge, on the other hand, does. What right opinion still
lacks is insight into the necessity of the matter. It states what is
right, but does not comprehend its ground. It therefore does not exclude
the possibility of being altered by persuasive specious reasons;
it is wavering and uncertain, an unreliable possession. The goal of knowledge must be a knowledge unshakable by appearance, secured by the knowledge of the ground of the matter (\gk{αἰτίας λογισμῷ}).

To the fundamental character of ordinary consciousness in theoretical
respect corresponds its peculiarity in practical respect. It does
not reach beyond the contingent and contradictory determinations
of the good given by opinion; it does not comprehend the good in its unity and universal validity. For, in accordance with its peculiar character, opinion conceives virtue as a multiplicity of single activities not brought back to one
concept, which in part come into conflict
with one another; and it mixes the good with the evil when
it teaches that one should do good to one's friends and evil to one's enemies,
and when it makes the good an aim not for its
own sake, but for the sake of pleasure and advantage, and thus
lets the good be done from base motives, be loved for the sake of its
opposite. The virtue that is attained at this
standpoint comes into being through habit, as it were instinctively; it
therefore, since it does not rest on a real insight, has
no security in itself; it can cease through chance and
circumstances, and all who content themselves with it stand on the
same level as soothsayers and poets. Only through true
insight is virtue grasped according to its truth, and only through
it is it freed from contingency. For the right knowledge of the good and the true moral will are inseparable; no one is voluntarily evil; no one does with consciousness and
intent what he has recognized to be harmful to himself.
\opage{108}For the defects and contradictions of ordinary consciousness
sophistry now has its eyes open, and it uses its dialectic
to confuse the former; but what it puts in the place of what it
overthrows is only the contingently subjective, which in the theoretical field finds its consistent expression in Protagoras's
proposition: Man is the measure of all things, and in the practical field in the
doctrine that posits pleasure as one with the good.

With respect to that theoretical proposition, whose significance
was developed above (p. 55 f.), Plato shows that, since it
makes everything, even the opposite, true, it destroys itself by
giving its own opposite the same validity as itself, and
that it is likewise cancelled when its presupposition, restless
becoming, is carried through, since then all sensation and every statement about it would become impossible. Under this presupposition,
namely, the sense-impression would have to change in every smallest part of time,
and one would have to deny and affirm simultaneously every determination of the thing and thereby destroy every statement about
it. Sophistry, taken as a whole, Plato leads back in theoretical respect to the very same basis in sensation
on which the ordinary consciousness that it confuses
rests, and he therefore posits knowledge of the essence of things
(\gk{οὐσία}), true knowledge of what in truth is, as
inaccessible to it.

In the ethics of the Sophists Plato sees only a conscious expression
of a principle that is the tacit
presupposition of the ordinary way of acting. Pleasure is posited by the Sophists as the
end of action, thus as the good. Against this identification
Plato directs his attack. Pleasure and the good relate in unlike
ways to their opposites: the good excludes its opposite, evil; whereas pleasure and its opposite,
pain, presuppose each other and pass over into each other; pleasure
and pain belong equally to the good man and the bad, but good
and evil do not. --- The adherents of the doctrine of pleasure themselves admit
that there are also bad and shameful pleasures; but thereby they themselves deny the unity of pleasure with the good. --- Finally,
\opage{109}pleasure is not, as the good must be, something existent in and for itself and
essential, but belongs to the field of becoming and falls
under the determination of \gk{τὸ ἄπειρον} [the unlimited]. If pleasure is posited as the
end of action, this leads, since pleasure is something merely subjective, individual, to the sublation of the morally universally valid, in the same way as the subjectivity of opinion led
to the sublation of a universally valid knowledge. --- Sophistry thus
becomes, in Plato's conception, both in the theoretical and
in the practical field a displacement of essence by appearance;
it replaces true knowledge by a sham knowledge, moral action by an untrue, egoistic action issuing from an untrue notion of the good. It is not philosophy, but an art of appearance, or rather, since even the name of
art is too noble, a juggling knack that flatters and
deceives the ignorant.

We can, with conceptual determinations from the older philosophy, but which Plato takes up, briefly characterize the defect of both standpoints, that of ordinary consciousness and that of sophistry,
by saying that they mix not-being and being and make the mixture into being proper. From the theoretical side
both forms of consciousness rest on sensation, and it is the \gk{μὴ ὄν}
contained in the becoming and plurality of the sensible that makes them
uncertain. For the uncertainty of opinion, too, is grounded
in its relation to the sensible, from which it draws its elements. In the practical field the good, as the determination and end of action, since it is conceived as a plurality
and is not kept apart from its opposite, and since it is identified with pleasure, is constantly conceived as a being mixed with not-being.
% errata, p. 109: printed "identificerer", corrected "identificeres"

The lack, in ordinary and in sophistic consciousness,
of a really reliable and secured knowledge now calls forth,
once the insufficiency of these standpoints has been shown,
the demand for a higher kind of knowledge, which must then be sought in
\emph{philosophy}, whose task, in contrast to those standpoints that
hold only to the subjective, to appearance, has already with Socrates
\opage{110}been indicated as this: to find the objective essence of the matter and to
find it through the concept.

As the means for the solution of such a task the
\emph{dialectical method} is determined. Its goal is precisely the concept
in its purity, separated from the sensuously individual phenomenon,
freed from all sensuous shape and presupposition, and as the
truly existent. The dialectical method becomes, in other
words, the organ for what lies beyond the sphere of the changeable
sensible (dialectic is \gk{ἡ περὶ τὸ ὄν καὶ τὸ
ὄντως καὶ τὸ περὶ ταὐτὸν ἀεὶ πεφυκὸς ἐπιστήμη} [the science concerned with being, with what really is, and with what is by nature always the same]). More precisely,
the functions of the method are to be determined thus: on
the one hand it is to gather the concept out of the manifoldness of the existing
(\gk{συναγωγή}), on the other hand, through
organic division, to lead the concept through the whole graded series of
its subordinate species right down to the limit of the individual
(\gk{διαίρεσις}). The first function is thus that of concept-formation,
the second that of conceptual division. Through \emph{concept-formation}
the What of the thing (\gk{αὐτὸ ὃ ἔστιν ἕκαστον}) is to be determined, the
distinguishing marks stated by which it differs from others;
the essential is to be brought out, for only with this does science
have to do. But this essential is the universal, what is common to the
individuals akin in species, that which is the presupposition of these individuals, without which they cannot be understood. The correctness
of the conceptual determination is to be tested not by single, arbitrarily chosen cases and examples, but by testing every
determination put forward, positively and negatively, in all its
consequences, inquiring what follows both from
the assumption of this determination and from that of its opposite;
and only after such a testing is what was provisionally put forward to be held fast, when it does not lead to contradictions and when
it agrees with and is required by what is given and by what
is already acknowledged as true. This hypothetical
dialectical procedure is to lead to an all-sided determination of the concepts, in order then, when the path of analysis has been traversed,
to give place to the direct development of thought and the syn\opage{111}thetic procedure. Here the second main function of the dialectical method,
\emph{division}, now gains its significance. It
is to state the differences through which a genus-concept divides
itself into species, its natural and complete articulation, and the division is to be carried through from the highest to the lowest concepts, to the point where the regular articulation of the concept
ceases and the indeterminate manifoldness of the phenomenon begins.
Division thus serves to determine the mutual relations of the concepts to one another, their compatibility or incompatibility. (In that the ground of division is determined, the articulation of what had been combined through induction
into a unity is \emph{grounded}, and the passage from the universal to the individual becomes
more than a merely repeated retracing, in the opposite direction,
of the path from the individual to the universal.)

Inasmuch as the dialectical (philosophical) activity of mind is seen
as leading to a real knowledge, it must in the first place be determined as the highest theoretical activity of mind, for
which the subordinate stages of development in knowledge, including
mathematical knowledge, become preparatory and serving.
Even mathematical knowledge has an object that
is not yet detached from the sensible, and it still proceeds
from presuppositions for which it gives no account;
only philosophical knowledge obtains an object freed from the
sensible, only it goes back to an \gk{ἀρχὴ ἀνυπόθετος} [unhypothetical first principle] (\gk{τὸ ἱκανόν} [what is sufficient]), and as it were rests in itself. But as
leading to a \emph{knowledge} of the true, dialectic cannot have merely
theoretical significance; for since action is determined
by knowledge, true knowledge must gain significance
also for human action, thus for human life
as a whole. Philosophy must thus comprise all ideal
human existence; it must, through true knowledge,
be the consummation of the whole of human life, the ``royal art'', which unites producing and insight into the use of
what is produced.
\opage{112}This higher standpoint of knowledge, this knowledge of
what in truth is, nevertheless stands here still only as a
demand. The task in its generality is determined, and with
it the object of the higher kind of knowledge in general.
But for the \emph{reality} of this object, by which for
Plato the \emph{reality} of the higher knowledge is immediately
secured, the definitive proof has not yet been given. This
takes place in Plato through a consideration of the nature of what
is, while the reflection on the relation of knowledge to sensation
makes precise the determinations that must be ascribed to the higher
existent in contrast to the sensible.

We consider first the consequences of the concept of knowledge. If
there is to be a knowledge different from insufficient sensation and opinion, its object must be \emph{another} than
that of sensation, and in such a way that the object of knowledge, of
true and actual cognition, must be \emph{what is in truth real}. That is: it must be \emph{an unmixed
existent}, and not, like the object of sensation and opinion,
stand midway between being and not-being. It must next
be an \emph{unchangeable existent in and for itself}, not, like
the object of sensation, something determined merely by our contingent and changing individual apprehension. And as the non-sensible, eternal and unchangeable essentiality it must finally be
something accessible \emph{only} to the beholding of the mind, something \emph{apprehensible only through
reason}. All these determinations can
be summed up in the determination: \emph{the Idea}, the objective concept existing in and for itself.

The reality of this Idea, more closely determined through knowledge,
is definitively established through the consideration of \emph{the nature of what is}. Everything sensible, says Plato, is something becoming,
but the end of all becoming is being. --- Everything sensible is a separated
plurality, whose essence must be what is common to the manifold, which makes the manifold things what they are, and this
essential must, as what is common to all the manifold individual things, be separated as archetype from these as copies. ---
\opage{113}No individual thing presents its essence purely (therefore
the notion of the common cannot be derived from the individual things either), but always only as mixed with its opposite; it therefore stands midway between being and not-being;
but imperfect being can be thought only on the presupposition of a perfect being and reality, of a being equal to
itself, freed from all opposition. From the becoming, divided and imperfect existence one is thus pointed
to what in truth is and is real, to the Idea in its
non-sensible, unchangeable unity, as something demanded with necessity. And in order also directly, by definite facts,
to establish the reality of such a non-sensible, Plato appeals to mental states and activities
that cannot be conceived as anything corporeal, e.g.\ reason
(\gk{φρόνησις}) and justice, and to the nature of the soul.

The doctrine of Ideas thus finds its final security in the nature of what is, but through the comprehending thought that knows
what is. Only in the reality of the Idea does Plato see the possibility of maintaining a true being and a true knowledge. To the
changeable, contradiction-filled sense-world, the world of ignorance, true reality could not belong; but by the very
imperfect being of the sense-world one was pointed to a true
and eternal being, to an object for the actual knowledge demanded by the human mind
and striven after through the philosophical \emph{Eros},
for the perfect knowledge to which the imperfect pointed by its defects; and this higher being,
the object of knowledge, was the Idea.
% printed "Ideen" (defect logged in the transcription)

\divrule

d. \emph{The Idea and the world of Ideas}. --- Through the foregoing developments the essential moments are given that
are enclosed in the concept of the Idea. \emph{The Idea} (\gk{ἰδέα}, \gk{εἶδος}) is in the first place
determined as \emph{the true existent}, which is \emph{the intelligible
object of the concept}. But the object of the concept, the
essential \emph{universal}, inasmuch as it is separated out from the many
\opage{114}individual things that have in it the pure (ideal) expression of their being, is \emph{intuited} by the mind as \emph{archetype} (\gk{παράδειγμα}) and
thereby acquires the typical individual existence of the ideal. In relation to
the archetype the sense-things then become copies, which participate in it
or imitate it, while it has a community (\gk{κοινωνία})
with them. As archetype the Idea separates itself, as substance existing for itself,
from the world of the phenomenon in \emph{essence}, as the eternal,
non-sensible, unchangeable, always equal to itself, unmixed with
oppositions, and it separates itself from that world in \emph{being}, in that it is \gk{χωρὶς} [apart], in the intelligible space (\gk{τόπος νοητός}).
The Idea is thus the expression of the common essence of individual things of like kind,
in its freedom from the limits of space and time,
from the obscuring and admixture of matter, as lifted up into a
sphere of eternity, and as set apart as an independent,
in-and-for-itself existent, unchangeable and perfect unity, which is
the model of the imperfect being that is its copy.

It has already been shown above (p. 105) how Plato,
in contrast to the abstract unity of the Eleatics, conceives the unity of the Idea
as a concrete unity, full of content, as a unity,
that is, in which plurality gets its place as bound by
unity. Therefore the Idea determines itself for him into a plurality
of Ideas, enclosed by the unity of a world of Ideas, and the single
Idea, as concrete, in turn unites in itself the determination of unity and of plurality.

As regards the first point, we find that Plato
almost never speaks of the Idea in the singular, but in by far the most
passages of the Ideas. And we then see him apply the designation Idea
not only to what we call something ideal, such as
the good, the beautiful, justice, life, and to the organic beings of nature, such as man and animal; but to all
kinds of natural objects, e.g.\ the Idea of fire, the Idea of water, indeed
even the Idea of dirt; to artificial products, e.g.\ the Idea of the bed, the Idea of the table; to properties and relations, activities and states; to mathematical figures and grammatical forms; indeed Plato speaks even of Ideas of the negative:
\opage{115} \setcounter{footnote}{0}
the Idea of what is not, the Idea of badness. Since the
Ideas were the hypostatized universal concepts, he had to posit as many
Ideas as there are concepts of genus and species, and since the Idea is
the real by which everything is what it is, nothing
can be or be conceived of which there is no Idea. As far
as something universal can be traced, the domain of the Idea must reach;
only where the trace of the universal is lost in the conceptless
manifoldness and unrest of becoming is the limit of the world of Ideas, and there the domain of the sensible begins. This whole manifoldness of
Ideas Plato then strives to conceive as forming a whole,
a comprehensive unity, a \emph{world of Ideas}. He seeks to conceive
the Ideas as a graded series which, from a highest Idea that is
the source of the being of all the subordinate ones, descends successively
to the lowest, which no longer have Ideas but only sensible
individual things beneath them, and he sets it as the task of science:
by an exhaustive enumeration, comparison and
ordering of all concepts to comprehend the whole world of Ideas.
As that highest, mightiest Idea Plato posits \emph{the Idea of
the good} (the Idea of the highest \emph{end}), which is determined as the source
of being and knowledge, and as standing above both; and with it
the concept of the highest divine merges for him.
% printed "Begreb sammen," (defect logged in the transcription)
But to show the branching of the Ideas out from this Idea,
Plato has done but little, and, by the nature of the matter, could only
do little. Since in the Idea the content of the concept is intuited
as archetype, it is hypostatized: the Idea becomes a resting
shape, the manifoldness of the Ideas a manifoldness of essentialities subsisting
for themselves, and a dialectical development proceeding with inner necessity out of the highest Idea, an unfolding of concepts, becomes impossible\footnote{Cf. Plato's statement about the Idea: \gk{οὔτε εἰς ἑαυτὸ εἰςδεχόμενον ἄλλο ἄλλοθεν} [neither receiving into itself anything else from elsewhere], \gk{οὔτε αὐτὸ εἰς ἄλλο ποι ἰὸν} [nor itself going anywhere into anything else].}. The unity of the concepts takes only
the form of community, participation; there is no question of a
proper development of concepts, but only of an ordering of the concepts by means of induction, and of a dialectic that proceeds
\opage{116}from given concepts as presuppositions. Perhaps it is
the feeling of this defect that brought Plato, in his advanced age, to attempt to lead the Ideas back to numbers.

With regard to the single Ideas, the union of
unity and plurality in them shows itself when one reflects on the mutual
relation of the Ideas. A connection of the different Ideas is
necessary if what is is to be capable at all of being determined and known, and not merely intuited in abstract, contentless
equality with itself. Every statement about being, even the
simplest, already includes subject and predicate, and therein
connection and plurality. And on closer consideration of
the mutual relation of the Ideas it also turns out that there are indeed
concepts that as it were repel one another, but that
on the other hand there are such as agree and even presuppose one another. With the concept of being, e.g.\ all those concepts agree that express a determinateness of being, even if they
mutually exclude each other, and as determinations of being
they presuppose the concept of being. Insofar now as concepts agree,
and thus allow themselves to be connected, they are one, i.e.\ the being of the one is
also the being of the other. Insofar as they repel each other, and
thus do not allow themselves to be connected, they are different, i.e.\ the being of the one
is the not-being of the other. But since each of the derived concepts can be connected with many others and excludes connection with many, being belongs to each
concept in many respects, but in many respects
also not-being, i.e.\ \emph{relative} not-being, the other-being of \emph{difference}. What truly is, then, is a
manifoldly determined concrete being, and each of the Ideas contains in itself, despite its unity, a plurality, governed by the unity, of positive and negative determinations and relations, and
stands in the most manifold relations of identity and difference, and
only on this presupposition can the explanation of the
manifoldness of existence be sought in the world of Ideas. ---

Through the determination that Plato gives of the highest Idea
as \emph{cause} of all being and knowing, a new moment
is \opage{117}brought out in the Idea, which supplements those indicated earlier. Just as the highest Idea is the cause of all being, so, according to
Plato, the Idea in general, in its relation to sense-things,
is to be conceived as the \emph{active cause} (\gk{αἰτία}) of their
being, and as \emph{power} (\gk{δύναμις}). And just as \emph{reason} must belong to the highest cause, because it is the cause of
knowledge and of the purposive arrangement of the world, so reason belongs also to the Idea, the perfect existent, in general, and reason presupposes \emph{life, soul and motion}.
But Plato's striving to claim these determinations for the Idea
remains only a striving, which is hampered by the nature of the Idea as
the plastic, resting archetype, as the self-enclosed
\emph{objective} existent. Over against the determination of the Idea as
the unchangeably existing archetype the determination of
activity must recede, and its lack in the Idea is as
a rule replaced by a \emph{mythical} or \emph{symbolic} element. And since
in the Idea the determination of objective being is essentially accentuated, the determinations \gk{νοῦς} and \gk{φρόνησις} can be carried through just as little
as the determination of activity. The Idea becomes
the object of thought, of the nature of thought, the intelligible archetype;
the determination of it as thinking, intelligent, recedes into the shadow.

\divrule

e. \emph{Natural existence as the world of sense and as cosmos}. --- Since the Idea is
determined as the archetype by participation in which --- a participation that gives it,
as it were, a presence (a \gk{παρουσία}) --- the sensible things, the Idea's copies, have
the being they possess, the Idea is thereby placed in a causal relation to the existence
that is not the Idea, but in its more imperfect being has a relation to the Idea. This is
\emph{the world of the sensible}. Whereas the world of Ideas was something eternally like
itself, something that has not come to be and is imperishable, not to be apprehended by
sight or the other senses, only an object for the knowledge of reason, the sensible, which
is named after the Idea and has a likeness to it, is something perceptible by
\opage{118}the senses, something that has come to be,
% errata, p. 118: printed "Tilblevet", corrected "et Tilblevet"
something that becomes and again passes away in space, an object for opinion, which rests
on sensation. The sensible is an \emph{intermediate between being and non-being}; its
fundamental determination is \emph{becoming}. It does not present the Ideas in their
purity, but always only as joined with their opposites. It does not present them in a
unity, but in a dispersed plurality of individual beings. To the determination of becoming
are thus added the determinations: \emph{relative being and dividedness}. But these
determinations of the sensible, since by them it is designated as an imperfect
\emph{being}, do indeed point to the Idea, the existent; but the \emph{imperfection} of
being cannot be explained by the perfect being of the Idea; it points to another factor,
which, in order to explain what in the sensible conflicts with the Idea, must be something
essentially different from the Idea. In order that absolute being may become relative,
likeness with itself become becoming and change, unity become dividedness, there is
required --- so Plato concludes --- a principle that has as its determinations non-being,
absolute changeability and multiplicity without unity; and these determinations
characterize the formless \gk{ἄπειρον} [the unlimited], which Plato describes as
\gk{ἡ δεξαμένη} [the receptacle], as the place of all coming-to-be, as the common element
that underlies all bodily elements and determinate formations of matter, and in their
unceasing becoming moves through them all as the common substrate, as the mass out of
which all sensible things are formed, but which itself, precisely because it is the
possibility of all forms, must be without any determinate form and property, an invisible
and formless nature (\gk{φύσις}, \gk{εἶδος}), and which he finally posits as equivalent to
space. As the formless, this \gk{ἄπειρον}, in its complete lack of being, can be an object
neither of sensation nor of thinking; it is, in Plato's expression, apprehensible only by
a kind of ``spurious thinking,'' and hard to arrive at clarity about
(\gk{μεθ' ἀναισθησίας ἁπτὸν}
\gk{λογισμῷ τινι νόθῳ}, \gk{μὸγις πιστόν}).
\opage{119}In the world of sense, which by its essence points to a twofold principle, two
opposites thus meet, and as opposites they both assert themselves; not only is the Idea
active, but matter as such becomes, despite the determination of non-being, active in
opposition to the activity of the Idea; that is to say: as limiting and hindering the Idea.
But while this hindering activity is brought forward where sensible existence is viewed
from its sensuous side, that is, in the parceling out of the sensible, it recedes when the
world of sense is viewed as a totality. The opposition then, for intuition, passes under
the harmonious unity; the hindering activity of matter becomes as it were hidden under the
reflected radiance of the Idea, and the world of sense is seen in another light, as the
living whole formed and ensouled by the Idea, as the divine world-work of art, as cosmos.
Seen as cosmos, the world of sense is apprehended as that in which the truly existent, the
Idea, which bears within itself as form measure and limit, is the ruling element; and
through the Idea, which through matter reveals \emph{itself}, it is something divine,
imperishable, harmoniously complete in itself, living, ensouled, endowed with reason.
Cosmos is, according to Plato, the perfect living being, formed in the likeness of the
Idea of the living (\gk{αὐτοζῶον}), after the image of the Godhead; it encloses in its body
the whole of the corporeal; it is by its soul in possession of a life of its own, i.e.\
belonging to itself, and of divine reason; it is never touched by age or destruction; it
is not only the perfect copy of the eternal and invisible Godhead, but itself a blessed
god, the only one of its kind (\gk{μονογενὴς}), sufficient to itself and needing nothing
else.

To the fact that cosmos exists through the Idea Plato gives, in his Timaeus, a
\emph{mythical} expression, depicting the formation of cosmos by the world-builder (the
Demiurge). Being good, the world-builder wished to make the totality of the visible as
perfect as possible, and therefore he wished to form, after the Idea of the living being
(\gk{αὐτοζῶον}), a copy that had all the \opage{120}perfections of the archetype,
% printed "Fuldkommenheder", ink blot over the first "n" (defect logged in the transcription)
so far as was possible. First, therefore, he formed the world-soul, the condition of the
reason of the universe, out of a mixture of the substance of the Idea and of the
non-existent, to which were added \gk{ταὐτὸν} [the same] and the \gk{θἄτερον} [the other],
which is hard to unite with it (the condition of uniform and of changing motion, and of the
knowledge of the identical and of the different); and this mixture was divided, formed and
spread out according to the harmonic numerical ratios, and into its frame, whose form the
equator and the ecliptic indicate, was fitted the chaotically surging material, after it
had been brought by the ordering Godhead into the mathematically determined spatial
boundaries of the physical elements (earth corresponds to the cube, fire to the pyramid,
water to the icosahedron, air to the octahedron; but all these stereometric figures are
thought of as formed from triangles). Thus arose the system of star-bearing celestial
spheres that enclose the earth as the resting center. In the various regions of heaven the
Demiurge placed the living and ensouled heavenly bodies, the immortal ``gods that have come
to be.'' Lastly he let all kinds of living beings come to be, so that the world might
enclose the same wealth of life as the Idea of the living. Of these, however, the
world-builder himself forms only the immortal and divine part of human souls, but leaves it
to the gods that have come to be to form the rest.

In the more fully executed depiction of the ordering of the universe, which, according to
Plato, owing to the sensible nature of the object, belongs not to the domain of certain
knowledge but to that of probable opinions, the conceptually determined opposition between
the activity of the Idea and that of matter everywhere comes forward. In the coming-to-be
of the universe as in its being, the ideal causes, the final causes, which everywhere
strive to realize the Idea of the good, become the proper causes. The material, mechanical
causes, which are traced back to natural necessity (\gk{Ἀνάγκη}), become merely contributory
causes (\gk{ξυναίτια}), conditions for the activity of the proper causes; they can indeed,
as we saw above, owing to their origin, limit and hinder in particulars the complete
activity of rea\opage{121}son; but their hindering activity has vanished in the universe as
a totality, so that cosmos stands as the completed work of art, resting in its divine life.

But with this Platonic determination of the world of sense, difficulties arise that Plato
cannot resolve, and that mark the rock on which Greek thought as a whole must strike. The
world of sense was not deduced from an ideal element essentially determined as
\emph{active} alone; rather, the Idea and matter were posited as the factors of the world
of sense. And these factors were to be opposites fixed \emph{absolutely} for intuition: the
Idea the purely existent, matter the non-existent. The character of this opposition is not
cancelled by what was developed above (p. 116), that Plato posits a non-being in the Ideas
too; for the non-being of the Ideas was only the relative non-being of difference, more
precisely: the non-being posited only by a comparing reflection, which does not get its
place as a determination of essence in what merely positively is, which excludes all inner
negativity, and becomes the different that is indifferent to this being-other. This
non-being is thus itself being; it is precisely the non-being of an other as the being of
this determinate thing. The non-being of matter, by contrast, is the absolute non-being
foreign to the Idea, a non-being in and for itself. But such a non-being, which cannot be
explained from the Idea, nor derived from its being, cannot be thought of as of itself
coming together into a unity with the being of the Idea; and yet the world of sense was
supposed precisely to express the unity. Plato must, since the opposites stand for him as
given presuppositions, seek to explain the being of the world of sense by finding a third
that joins and connects them; but this attempt proves, by the nature of the case,
impossible to carry through. Plato seeks the joining element partly as an active
connecting agent, which finds its expression in the world-artist (the Demiurge), partly as
a conceptual mediating element, which is represented by the world-soul. The notion of the
\emph{Demiurge} is, however, only the mythical expression for the determination of motion
and activity that was required by \opage{122}Plato in the Idea but not conceptually
secured, and which is therefore intuited outside the Idea (after which the world-artist
forms) and personified, but without the essence and being of the world-forming Godhead
being able, after all, to be any other than precisely those of the Idea; and that
personification therefore cannot explain what was not explicable by the concept of the
Idea. The explanation of the connection of the Idea and matter would then have to be sought
in the \emph{world-soul}. The world-soul, whose formation has been described above,
becomes, according to Plato, that which, as enclosing all the harmonic numerical ratios, is
the cause of the order prevailing in the universe; that which, as moving itself, moves the
world; and that which, as possessing life, consciousness and thought, becomes the principle
of life, of consciousness and of knowledge. All these determinations point to the Idea as
their source, but at the same time they point to the active pervasion of the sensible,
whose nature is a moment in the essence of the world-soul. The world-soul is designated as
an intermediate being in that it is determined as enclosing the numerical ratios. For the
mathematical element, for Plato, forms an intermediate sphere between the Idea that is for itself
and the corporeal, but nevertheless stands nearer the former than the latter, since,
although existing in the sensible and attached to it and to its multiplicity and its being
in space, it nevertheless, like the Idea, expresses something eternal and unchangeable.

But the middle term which was to explain the connection between the opposites in the
universe itself comes to be, according to Plato, only through the connection it was to
explain: through the mixture of the Idea and matter; and just as the opposites are
hypostatized each by itself, so it too is hypostatized and comes to stand as something
separate \emph{beside} the Idea and matter. Whereas it encloses in itself only something of
the whole of the Idea and something of the whole of matter, and therefore becomes only the
frame of the world of sense, it becomes a distinct member, which in its essence contains,
as it were, added-together elements of the essence of both opposites, but falls outside
them both. Yet by \opage{123}its fundamental determinations it stands so near the Idea that
it becomes difficult to maintain its separation from the latter; it becomes in reality an
expression for the Idea (for that which has measure and form in itself), seen from the side
on which the Idea stands in relation to the phenomenon and, as determining the phenomenon
according to the mathematical numerical ratios, appears as the abiding measure, as the
forming and ordering element in corporeal existence. What was to distinguish the
world-soul from the Idea would have to be, most nearly, that the former is represented as
spread out through space, while the latter is without relation to space. But the extension
of the world-soul in space becomes really only an image. Since it invisibly, as a unity,
pervades the world, its relation of dependence on space is essentially cancelled, and this
becomes still more explicit when the world-soul is determined as incorporeal and thinking.
The world-soul is the harmony of the existent as hypostatized and separated from what
harmoniously exists, but as harmony it is independent of space.

The connecting links, then, that in Plato were to connect Idea and matter prove
insufficient. The mythical figure of the world-builder could gain power and significance
only for the imagination. The activity that remained unconceived in the Idea hypostatized
into an archetype remains unconceived in him; and unconceived, even if the activity is
presupposed, remains its power to connect what is separated by an abstract opposition.
Here the conceptual middle term, the world-soul, was to fill the gap. But this middle term
did not explain the connection, because it was itself only explained by the connection;
and the incompatibility of the opposites showed itself in this, that its postulated
determinations could not after all keep it apart from one of the opposites: the world-soul
became the universal \emph{form} of the universe. What is absolutely different in essence
can have no true middle term: every attempt to maintain such a one must lead over to one of
the sides.

When, then, the problem of the possibility of the world of sense cannot be solved along the
path attempted, a new \opage{124}difficulty comes forward, in that the world of sense, which
forces itself upon us as something given, which Plato neither can nor will deny, comes, as
the relatively existent, to claim a place \emph{beside} the absolutely existent: the Idea.
If the concept of the Idea were held strictly in the form of thought, the world of Ideas
and the world of sense could not be posited as \emph{two} real worlds standing beside each
other. All reality belongs to the Idea; the sensible, therefore, insofar as being is
ascribed to it, could only be posited as existing in unity with the Idea. From this
standpoint spatial relations would not become separating (cf.\ Parmenides), the Idea and
the phenomenon would not be placed side by side, the being-for-itself of the archetype would
not become different from its being as indwelling in the copy. In cosmos we would then have
this unity; in it all existence would be intuited as an all-embracing living work of art,
existing through the Idea and ensouled by the indwelling Idea. And this unity would, by
virtue of the necessary relation between the factors of the work of art, be an eternal
unity. But since, in the conception of the Idea, intuition determines the form of thought
and transforms the universal essence with its \gk{παρουσία} into the hypostatized archetype,
the world of Ideas becomes a \gk{χωριστόν} [something separate], and as it separates itself
from the world of the phenomenon, it comes, as it were, to share reality with the latter,
which, as the product of the originally juxtaposed opposites, is represented as having come
to be. Despite the fact that Plato has, as it were, gathered in the concept of cosmos all
the determinations that express the unity of existence, he cannot after all maintain this
unity, and what prevents him from doing so is the peculiarity of Greek thought, conditioned
by the character of the Greek fundamental intuition.

In the fundamental tendency of the Greek spirit, to apprehend existence in the form of
beauty, lies an impulse that must constantly draw the opposites toward each other in order
to connect them in the intuition of cosmos, and that must drive thought to give the
intuited unity its form, to express it through the concept. But in order to apprehend the
unity as a unity of concept, thought must first give the moments of the unity the form of
the concept, \opage{125}which it can do only by loosening them from their immediate
connection, given for intuition. But in virtue of the moment of \emph{artistic intuition}
which the thought of a spirit resting in beauty, in the immediacy of harmony, must
constantly bear within itself, and by which its form is determined, the concepts of the
separated moments must fix themselves in the separation as resting opposites in the form of
immediacy, which thought is not able to set in motion and to lead back to unity, however
much it tries. The relation must become double and ambiguous, the activity of thought
divided between unification and separation, but in such a way that separation becomes
dominant, that thought, whose motion is bound in the resting immediacy of artistic
intuition, is not able to realize its striving for unity in the form of the concept, but
ends in dualism.

It is this peculiarity of Greek thought that in the Idea pushes back the determination of
activity and gives the Idea as archetype the form of resting being, just as it was this
peculiarity that prevented the world of Ideas in its ordered multiplicity from being seen
as proceeding from the highest Idea, and that made the world of Ideas a world of resting
Ideas whose ordering remained unconceived.

It is the same peculiarity that, when the opposite of the Idea comes forward with the
determination of non-being, brings it about that this non-being is hypostatized and, as the
opposite of the intuited existent, becomes the \emph{intuited} non-existent, which,
precisely as such, is fixed in \emph{absolute} opposition to the Idea, and stands as an
underived, ever-present existence which, as it were by force, asserts itself to intuiting
thought as an objectively existing and active non-existent, whose connection with the Idea
becomes incomprehensible.

It is, finally, this peculiarity that, even at the point where, in the intuition of cosmos,
the unity in existence found its fullest, as it were definitive, irrevocable expression,
where thought sought rest in the intuition of the \opage{126}harmonious unity of the whole
of existence, nevertheless keeps the Idea separated from this its harmonious revelation,
and removes it into a higher sphere of existence, into an intelligible space filled with
non-sensuous archetypes, whereby the reality of the world again becomes incomprehensible
for thought, since all real being is enclosed in the world of Ideas, which is posited as a
\gk{χωριστόν}. And through this involuntary separation the dualism in principle of the
fundamental opposites again betrays itself.
% printed "principielle", final "e" over-inked (defect logged in the transcription)

\divrule

f. \emph{The Human}. --- \gk{α}. \emph{Anthropology}. ---
The cosmic principles of the Platonic physics are mirrored in anthropology, which is the conclusion of the special physics and the foundation of ethics.

The human soul is, according to its eternal essence, of one essence with
the world-soul. Its immortal, divine part is namely formed
(according to the Timaeus) of the same constituents as the latter,
though in lesser purity. It is, without itself being an Idea,
so closely bound up with the Idea, in particular with the Idea of life,
that it cannot be thought without it. For its determinations it
therefore has self-motion, eternal life, simplicity, likeness with
itself. \gk{Νοῦς} [mind] is implanted in it. But to the divine and immortal part of the soul there has been added, through the union with
the body --- a union which Plato now determines as necessary, hence as essential, now as conditioned by a fall
in pre-existence --- a mortal part, and through this
union the purity of the immortal soul and its intuition of
the Idea are obscured and disturbed, and can only be regained successively through the contemplation of the harmony of the universe, and through thinking, as the mind, awakened to Eros by sensuous beauty, the imprint of the Idea, \emph{recollects} the Idea beheld in the soul's pre-existence, and through philosophy is led back to
the life in it.

The obscuring of the soul in its union with the body is caused by sensuousness and passion, by the mortal soul, bound to the \opage{127}body
and perishable with the body, which,
through the union of the rational and immortal soul with the body,
as it were grows fast to it and affects it. The mortal
soul in turn has two parts in it: a nobler one, \gk{θυμὸς} (\gk{τὸ θυμοειδές}),
which, though lacking rational insight, strives as if by instinct toward the true and the good, and, when it
has not been made refractory by bad habit, assists the rational
soul, --- and a less noble one, \gk{τὸ ἐπιθυμητικὸν} (\gk{τὸ φιλοχρήματον} [the money-loving]).
The first, ``the spirited,'' is the expression of the more vigorous,
nobler passions (such as anger, indignation), the second, ``the
desiring,'' of the lower sensuous drives. They each have
their own particular seat: \gk{θυμὸς} in the breast, \gk{ἐπιθυμία} in the abdomen;
while the thinking, cognizing soul (\gk{τὸ λογιστικόν}) has
its seat in the head. These parts of the soul, of which the
lowest is also found in the plant and the lower animals, the middle one in the nobler animals, Plato determines as different
stages of life, of which the higher, in order to \emph{exist}, always presupposes the lower; but under this conception, which suggests an
inner unity, he nevertheless retains them, in agreement with
the notion of their origin and their seats in the body,
as real \emph{parts}. The question of the \emph{unity of the soul's life}
within the tripartition Plato has not answered. The concept of personality
did not yet exist at all for the Greeks in its sharpness:
they viewed the life of the soul more as spread out through the plastically
formed body than as concentrated in ideal unity. Connected
with this, too, is the fact that the problem of the \emph{freedom of the will} comes forward only late and is answered only unsatisfactorily.
In Plato the problem is in reality only pushed back to
the state of pre-existence, but not solved.

\emph{The immortality of the soul} is for Plato an immediate
consequence of its relation to the Idea of life. In an earlier writing
(the Phaedrus) he lets the whole soul have pre-existence and immortality, but in a later one (the Timaeus) only the thinking soul,
which is more consistent and agrees with his conception
of the relation of the ideal to the material. In depicting
\opage{128}the soul's life after death, its wanderings and purifications, its
punishments and its blessed life, Plato, in the consciousness
that the domain of clear knowledge is here overstepped, makes use of the poetic form of \emph{myth}, which for him becomes an indispensable supplement to the properly philosophical exposition.

\gk{β}. \emph{Ethics and Politics}. --- The fundamental character of the Platonic ethics
is immediately determined by his metaphysical
fundamental view; its more detailed working-out is closely conditioned by
his anthropology, inasmuch as ethical action becomes the expression of the human being, determined according to its true peculiarity, in its activity.

The fundamental concept of ethics is the concept of the highest goal of human
striving: \emph{the highest good}, which is synonymous with \emph{happiness}. This highest goal becomes for
Plato the realization of the Idea in man, either --- which
corresponds to the conception of the Idea as what alone is
real in its separateness --- through complete liberation from the sensuous; or --- which corresponds to the conception of the Idea as
principle of a harmonious being, i.e.\ of an ideal reality,
in the sensuous --- through the harmoniously beautiful development
of man's whole sensuous-mental existence, in which
the nobler pleasure too finds its place. To realize the Idea
is in Plato synonymous with becoming like the Godhead (\gk{ὁμοιωθῆναι} \gk{θεῷ}).

The means by which alone the highest good can be attained is
virtue. \emph{Virtue} Plato conceives as the normal harmonious
activity of the parts of the soul, appearing in a plurality of
essentially connected forms. He is thereby led to set up
4 cardinal virtues. The virtue of the rational part of the soul is wisdom
(\gk{σοφία}), the clear knowledge of the true and the good
that guides the life of the soul. The virtue of the passionate part of the soul is courage
(\gk{ἀνδρία}), whose essence is, in the face of the promptings of pleasure and pain, to uphold the pronouncement of reason about what ought to be
feared and what not. The virtue of the desiring part of the soul is moderation (\gk{σωφροσύνη}), which consists in the subordination of the
% printed "bestaaer" (defect logged in the transcription)
\opage{129}lower in the soul to the higher and more rational. But to
the soul as a whole there also corresponds a virtue: justice
(\gk{δικαιοσύνη}), which expresses the normal state of the total life of the soul through the right mutual relation of the individual parts of the soul, and
thus becomes, as it were, the universal virtue. These virtues, which
are thus distinguished, work together in perfected virtue, which
through true knowledge develops out of ordinary ethical life, resting on
practice and tradition; and in all virtue
\emph{knowledge}, which, as pointing toward the Idea, expresses man's
essential determination, becomes the proper foundation and the properly
hegemonic element.

The virtue of the single individual presupposes for its realization the \emph{state}, which, in itself realizing the ethical on
a larger scale, educates individuals to it and secures
its dominion. The essential aim of the state is not external
goods and external security, but to make its citizens
good and thereby happy; and this it can do only when it
itself expresses and realizes virtue, and in particular the all-embracing virtue, justice, whose distinguishing mark is: \gk{τὰ ἑαυτοῦ}
\gk{πράττειν} [to do one's own]. Now since Plato, in his peculiar manner, views the fundamental activities of society as particular \emph{parts} of
the state, corresponding to the parts of the soul, this leads him to let
the \emph{estates} of the state (\gk{τὰ γένη}), which represent those fundamental activities, sharply separate themselves from one another, each with its particular work, and each represent a single virtue. While
the acquiring estate (artisans, farmers and traders), which corresponds to \gk{τὸ ἐπιθυμητικὸν} (\gk{φιλοχρήματον}), the
desiring part of the soul, has moderation for its virtue, the estate of the warriors, which corresponds to \gk{τὸ θυμοειδές}, has courage for its virtue, and the estate of the rulers, which corresponds to \gk{τὸ λογιστικόν},
has wisdom for its own. And the ethical life of the \emph{whole} state, its \emph{justice} and its happiness, is posited as conditioned by the
distinct estates, each attending to the work that is determined by its concept, standing in the right harmonious relation
to one another.
% printed "betinget" (defect logged in the transcription)
The rulers, who represent reason,
\opage{130}the Idea, must for this reason have power exclusively, power of which they must make themselves worthy through the possession of the highest
insight and cultivation.
% printed specks on p.130, none transcribed (defect logged in the transcription)
Only the dominion of true philosophy
can make the state answer to its task, and free the human race from unhappiness. The estate of the warriors, which is comprehended
together with that of the rulers under the common designation of guardians of the state
(\gk{φύλακες}), and comes to stand nearer to the rulers than
\gk{θυμὸς} stood to \gk{τὸ λογιστικὸν} in the anthropology, is to live solely for
the defense of the state at home and abroad. The third estate is to
live solely for acquisition, for the provision of all necessities of life, in strict and compulsory subordination to the others.
In the estate that rules without restriction, representing the dominion of the Idea, the idea of the state has a \emph{separate} existence, and
it is therefore consistent that the individual can be a member of this estate only by giving up individual interests
for the universal. For the ruling estate private property and family life cease, and the same law holds for its
``helpers'': the warriors. The authority of the state regulates sexual unions, fixes the number of children that may live, and selects that number from among those born, takes complete possession of the children
from birth onward, determines, according to the aptitudes they show, their
station in life, and directs the strictly
ordered upbringing and education of those selected as guardians of the state. (The first two to three years of life are appointed for the children's bodily care; from the 3rd to
the 6th they are educated through the telling of myths; from the 7th
to the 10th they are exercised in gymnastics; from the 10th to the 13th
in reading and writing; from the 14th to the 16th in acquaintance with
poetry and in music; the 16th to the 18th is devoted to
the mathematical sciences; the 18th to the 20th to military exercises. With this the instruction of those destined to be warriors is concluded. Those who show scientific aptitude, on the other hand, receive
from the 20th to the 30th year a more rigorous scientific instruction. The most excellent among them are then singled
out and occupy themselves, while the others pass over to
a practical activity, from the 30th to the 35th year with
\opage{131}the study of dialectic. From the 35th to the 50th year those so educated take over posts of command and office.
After the 50th year they attain to the contemplation of the highest
object of knowledge: the Idea of the good; they are received among the rulers, and hold in rotation the highest offices of state,
with which the supervision of the whole state's government is connected;
yet most of their time is devoted to philosophical contemplation, in comparison with which life in the state is the
subordinate thing.) The women of the two estates are brought up in the
same manner as the men and take the same part in affairs of state and in war. The difference of sex acquires no significance
over against the universally human. The material maintenance of the two higher estates
is provided by the acquiring estate.
Of this estate Plato speaks only little, and from a strictly aristocratic standpoint. The educative influence of the state upon it
seems to have to be limited to this: that the example and power of the higher estates
hold back the sensuous desires of the third estate, and that the state's care to remove an effeminate music and an unethical poetry from its domain must also acquire
significance for this estate.

Plato has formed his peculiar doctrine of the ideal state,
although he takes up motives from Doric political conditions and from
the Pythagorean league, in an independent manner
\emph{from the fundamental concept of his philosophy}, and in it,
through what he posits as unconditionally ruling over the striving of all
individuals, and through the ideal goal to which the whole life of the state
is subordinated, he has idealized the principle of Greek substantial ethical life;
and he has intimated much that only in a distant future, from other points of departure, has come to reality, and
much that stands as the task of a still more distant future.

\divrule

g. \emph{The Significance and Defects of the Platonic
Philosophy}. --- The abiding and imperishable
significance of Plato's philosophy lies in this: that through dialectic he discovers
the \emph{Idea}, the fundamental thought of Greece; that, by giving
\opage{132}thought, in the Idea, a true and eternal, ideally existent object, he gives rebirth to the dissolved philosophy as a new universal science, which, as it were through a recollection, preserves within itself the essential
content of the past, and which through the ideality of its principle receives throughout
the stamp of ideality; and that he lets the Idea, revealed to knowledge, rule the whole of human existence and shape
it into artistic unity. By being the first to proclaim, energetically and with enthusiasm,
the reality of the Idea and its accessibility to thought,
and by determining it as what rules in all existence
and is essential in it, Plato has become a representative of all ideal
striving in life and science, and has given the human mind an
impulse that has carried it far beyond the
bounds of his system, and has worked to awaken and inspire in all ages.

But beside the abiding
and comprehensive significance of the Platonic philosophy, there appears in it an essential defect, already
indicated above, which is connected with
the peculiarity of Greek thought, \emph{typically} pronounced in Plato. We saw above how the Idea and its opposite, as both were hypostatized for intuition, became fixed
as abstract opposites, whose fundamental determinations placed
them in a merely negative relation to each other, so that the Idea
became the purely existent, matter the purely non-existent. And
as the determination of non-being, of negation, was thus excluded from the Idea, the Idea had to lose the determination of activity, which could only have been maintained had negation, the incitement of motion, been posited in the essence, in the Idea;
it had to become the resting, substantial existent, the
tranquil archetype. From this, in turn, essential consequences spring. When the incitement of motion, of dialectical unfolding,
is lacking in the Idea, then the concrete content of the world of Ideas
cannot be reached by a deduction from the highest Idea,
but can only be won through an abstraction from the sensible world; and this sensible world, given for intuition, can
be comprehended neither according to its concrete content nor according to its existence
from the self-subsistent, resting, inactive,
\opage{133}negationless Idea. To explain the sensible world, matter was called in
to help; but since this new factor was the Idea's
abstract opposite and as such what is incompatible with the Idea, the explanation of the sensible world nevertheless became impossible, and
the sensible world remained standing as a riddle, with a contradictory existence, as a relative being beside the
absolute, which is a \gk{χωριστὸν} [something separate], and as a combination of that
whose essence excludes combination. And as contradictory as its
existence becomes that of the Idea's opposite, whose validity constantly
asserts itself anew through the fact of the sensible world:
it becomes an existent and active non-existent and stands
as such over against the absolute being of the Idea.

Through the defects and contradictions that remain in the
Platonic philosophy, the task is set for philosophy:
to find a new formula for the relation of the fundamental opposites, one that
sets these in a positive relation, in a relation of essential unity, to
each other, and to comprehend the ideal, form, as the active force, and as that which, indwelling in matter, forms
matter. It is by taking up this task that \emph{Aristotle}
carries Greek philosophy forward, while \emph{Plato's school}
loses itself in fruitless speculation.
% printed "der" (defect logged in the transcription; p.133, one site, two occurrences)

\divrule

\begin{center}
\textbf{4.\enspace The Older Academy.}
\end{center}

The school that attached itself to Plato followed in the course of
time essentially different directions, which have led to a
distinction between what has been called the older, the middle
and the newer Academy. The last two stages of the school's development
belong to the 3rd period; the first we discuss here in brief.

It is represented essentially by \emph{Speusippus} (Plato's
sister's son and head of the Academy after him), \emph{Xenocrates},
\emph{Heraclides Ponticus}, \emph{Philip} of Opus, \emph{Polemo},
\emph{Crantor} and \emph{Crates}. Of these, after Speusippus, Xenocrates, Polemo and Crates were scholarchs.

The philosophical activity of \emph{the older Academy} is of
\opage{134}subordinate value. The Pythagoreanizing form which Plato's doctrine of Ideas had assumed in his old age is held fast here,
and in its first generations the school loses itself in obscure speculations on the relation of the Ideas to the numbers, on the
elements of the numbers and on the origin of the numbers from the elements and
of things from the numbers. The Platonic doctrine of Ideas in its original
and pure shape disappears, and Pythagorean mathematical-religious mysticism takes its place; dogmatism replaces dialectic.
Weight is laid on encyclopedic learning, and as
interest in the abstruse speculations gradually fades, the
main interest turns from the theoretical to the practical, in
which domain there then appears a shyness of exaggeration and
one-sidedness, a striving to give external goods a place
beside the essential ones: the mental. This turn
toward an ethics that preserves a popular form is represented by
\emph{Polemo} and \emph{Crantor}. Polemo sets the practical above
the theoretical, teaching that one should cultivate oneself by
action, not by dialectic. --- After Crates the
Platonic school assumes another shape. The cast of thought of the middle Academy
is essentially skeptical.

\divrule

\begin{center}
\textbf{5.\enspace Aristotle.}
\end{center}

a. \emph{Aristotle's Life and Writings}. --- \emph{Aristotle},
son of Nicomachus, was born at Stagira in Thrace in 384 B.C.
His father belonged to a family of physicians who derived their
descent from Machaon, son of Asclepius. At eighteen years of age
Aristotle came to Athens to hear Plato, and for 20 years
he was his pupil, though at the same time he appears to
have come forward as a teacher of eloquence in rivalry with
Isocrates. Of Aristotle's personal relation to Plato only
unreliable reports are available. When Plato had died, he left
Athens and betook himself, together with his former fellow-disciple
under Plato, Xenocrates, to the tyrant Hermias in Atarneus
in Mysia, with whom he stayed until 344. After Hermias
\opage{135}had been captured through treachery and executed by the Persians,
Aristotle stayed for a short time on Mitylene, and went from there to
Macedonia, where King Philip made him tutor to his son Alexander. When Alexander, after his father's death, was preparing
his expedition to Asia, Aristotle returned to Athens, where
he founded a philosophical school in the Lyceum, which he directed
for 12---13 years.
% printed stroke above an "o", p.135 l.6 (defect logged in the transcription)
With Alexander he maintained a friendly connection, which cooled only in the king's last years,
and Alexander is said to have richly supported him with resources for his studies in natural science and for the collecting
of a library. Alexander's death gave the party in Athens hostile to him
the courage to raise an accusation against Aristotle for
\gk{ἀσέβεια} [impiety]. He then left Athens, ``so as not to let the
Athenians sin a second time against philosophy'' and died soon
after at Chalcis on Euboea, in the year 322 B.C.

Of Aristotle's writings, the popular ones, composed in dialogue form,
are lost; the more strictly scientific (acroamatic) ones
are very incompletely preserved, in part handed down in an
imperfect form and mixed with spurious components. The
genuine Aristotelian writings do not give us a picture of the development of his peculiar philosophy as the Platonic ones do, but
all show the same already fully clarified fundamental view, and were
probably all composed during Aristotle's second stay
in Athens. We do not have in these writings the Aristotelian philosophy lying before us as a closed system, equally
worked out and carried through in all its parts, whose individual sciences are also joined in an external manner into an
articulated unity. The writings form, according to the diversity of their content, a plurality of more or less complete groups,
which indeed show themselves as belonging together both through the unity of the pervading thought and through mutual references, but
which are nevertheless not joined into the architectonic
unity of a system. There are 4 main groups: a logical, a metaphysical,
a physical and an ethical-political. The \emph{logical}, which seems to
have been the earliest composed, includes the writings: On the Cate\opage{136}gories (the fundamental forms of statements about what is); \gk{περὶ}
\gk{ἑρμηνείας} (which treats of the proposition and the judgment); the two Analytics (on inference, proof, definition, division and the knowledge of principles); the Topics (on dialectical inferences
from the probable), and the writing on sophistical fallacies (\gk{περι σοφιστικῶν ἐλέγχων}). The collection of the logical
writings was later designated as a whole as \gk{Ὄργανον} [instrument]. ---
The \emph{metaphysical} group is represented in our collection of
Aristotle's works by one writing, which bears the later
name: the Metaphysics (\gk{τὰ μετὰ τὰ φυσικὰ}); but this writing
contains a plurality of independent drafts, more and less worked out,
which are put together in a not felicitous manner and
in part mixed with spurious components. The treatises that
belong to the fundamental science are probably the last
composed of Aristotle's writings, and they are the ones that lie before us in
the least completed and finished form. --- The \emph{physical}
group comprises the main physical work, \gk{φυσικὴ ἀκρόασις} [Physics] in
8 books; the writings \gk{περὶ οὐρανοῦ} [On the Heavens] and \gk{περὶ γενέσεως καὶ}
\gk{φθορᾶς} [On Coming-to-be and Passing Away]; the Meteorology; the extensive natural history of animals
(\gk{περὶ τὰ ζῶα ἱστορίαι}): a comparative physiology;
the writing on the parts of animals; the writing, important in philosophical respect, on the generation of animals; and some lesser physical treatises; in addition, the main psychological work \gk{περὶ ψυχῆς} [On the Soul]
and a number of lesser psychological writings that attach themselves to
it (\gk{περὶ μνήμης καὶ ἀναμνήσεως} [On Memory and Recollection], \gk{περὶ αἰσθήσεως και αἰσθητῶν} [On Sense and Sensibles], \gk{περὶ ὕπνου καὶ ἐγρηγόρσεως} [On Sleep and Waking], \gk{περὶ ἐνυπνίων} [On Dreams], \gk{περὶ}
\gk{μαντικῆς τῆς ἐν τοῖς ὕπνοις} [On Divination in Sleep], \gk{περὶ μακροβιότητος καὶ βραχυβιότητος} [On Length and Shortness of Life], \gk{περὶ ζωῆς καὶ θανάτου} [On Life and Death]). --- To the \emph{ethical-political} group belong the main ethical work, the Nicomachean Ethics in 10 books, the writing on the state, the Politics (in 8
books), and perhaps the 1st book of the Economics (the 2nd is spurious).
--- The Aristotelian \emph{Rhetoric} occupies a somewhat indeterminate
position, since from one point of view it belongs to the logical, from
another to the ethical-political group. --- The unfinished \emph{Poetics}
(the theory of poetry) stands isolated, since it alone, among the
\opage{137}writings preserved to us, belongs to the domain that Aristotle, in distinction from
that of the theoretical and that of the practical, designates as that of the poietic
(see below: the division of Aristotle's philosophy). Among
the lost writings of Aristotle, his writing on
animal anatomy (\gk{ἀνατομαὶ}) and his \gk{πολιτεῖαι}, a description
of the constitutions of over 150 states and cities, have surely been the most important.

\divrule

b. \emph{The historical position of the Aristotelian philosophy,
and its distinctive character}. --- In the survey of the second period
the historical position of the Aristotelian philosophy has been briefly
indicated; we here determine it more closely, partly in relation to the
Platonic, partly in relation to Greek philosophy in general.

With Plato Aristotle shares the conception of form, the ideal, as what
truly is and as the source of all actuality; but since in Plato's
conception of this ideal as the archetype, separated from things and
resting in itself, he could find no explanation of the being of things
and of their coming-to-be, and in the conception of matter as the
abstract \gk{μὴ ὄν} [what is not] saw a source of contradictions, he
determines the formula for the relation of the opposites differently
from Plato. Form becomes for him, as already indicated earlier, the
immanent form, the forming principle active in matter, and form's
opposite, matter, is brought under the same determination as form:
under that determination of being which in Plato involuntarily
comprised what is not as well. Aristotle thus wants to find the ideal
in concrete real existence itself, not through a flight from it or an
annihilation of it, and he seeks a unity of the opposites in that he
sees the ideal as realizing itself in matter, which is akin to it in
essence, differs from it only in the mode of its existence, and has its
goal in form.

Under the opposition to Plato, under the polemic against the particular
form in which the ideal principle was conceived by him, the fundamental
character of the Aristotelian philosophy is therefore the same as that
of the Platonic. The difference is only this, that whereas
\opage{138}the power of the ideal principle over thought shows itself
in Plato in the striving to let the being of the finite be dissolved
through dialectic so that the world of the Idea might raise itself up
above it and reveal itself in its purity and independence, in Aristotle
it shows itself in his striving to see the Idea, the form, in the
sensible, in concrete natural existence,
% printed "i" twice in "i det Sandselige, i den concrete" (p.138 l.6) (defect logged in the transcription)
as the immanent cause of its being, as the source of its life, as
actively realizing itself in the material. There is in the two,
Greece's greatest thinkers, as it were a double current of thought:
back to the Idea and out from the Idea, and in this circulation pulses
Idea-animated Greek thought. Aristotle is Plato's supplement: together
they comprise the whole essential content of Greek philosophy. As
representatives of Greek thought in its culmination they both are, and
both must be, idealists, and Aristotle is so even in a higher sense
than Plato, insofar as he has the assurance of that power of the ideal
which preserves its being even in the midst of the world of phenomena,
in the connection with its opposite.

In the determination of Aristotle's relation to Plato lies the
determination of his relation to Greek philosophy in general.
\emph{Aristotle is the ideal conclusion of Greek philosophy}. What Greek
thought strove for as its highest goal, to express in the form of the
\emph{concept} that \emph{unity of existence} which is immediately
present to intuition in the work of art that is the world, Aristotle
has attempted to realize with a greater energy of thought than any
other Greek thinker; and through the determination of form as immanent
and of matter as what is according to potentiality, the principles are
set in such a relation to one another that that striving has a real
foundation. And while Aristotle energetically strives for the unity of
principle, his philosophy embraces and appropriates---precisely as a
consequence of his conception of the relation of the fundamental
opposites---real actuality in as rich a measure as any Greek system,
while at the same time it \opage{139}most surely and most definitely
conceives the task of science in relation to the various domains of
existence, and most sharply and most completely clarifies the forms of
thought. Precisely as the completion of Greek thought, the Aristotelian
philosophy becomes the form through which the latter exerts its
influence through the Middle Ages.

But since Aristotle's philosophical system stands as the ideal
conclusion of an entire historical main form of philosophy, it lies in
the relation between the principle, and the distinctiveness determined
by that principle, of the individual philosophical main form and the
universal essence of philosophy working through this distinctiveness,
that this system, which as the ideally concluding one fully reveals the
character of the main form to which it belongs, must point beyond the
limits of that form and intimate concepts from the forms of thinking
which develop out of it with inner necessity. For this main form of
philosophy has its distinctive life only as a member in an organic
development, and must therefore in its completion reveal its inner
connection with the supplementary forms. Intimations of the thoughts of
the future appear, too, as will be developed more closely below, both
in the organic conception of existence conditioned by Aristotle's
determination of the fundamental relation of the principles, and in his
conception of the fundamental forms of knowledge. But wherever what is
thus intimated is to be fully realized, we see it hampered and pushed
back by the general fundamental peculiarity of Greek thought, which as
an insurmountable barrier halts the attempt to reach beyond the limits
of the Greek spirit.

The distinctive movement of thinking in Aristotle is conditioned by the
determination of the relation of the principles. According to
Aristotle's view, it is in the phenomenon that we are to seek the ideal,
% printed "i" in "det i Phænomenet" over-inked (defect logged in the transcription)
whose first form of appearance for us the phenomenon is; only from the
object of experience, what is ``more certain for us,'' is the mind to
be able to attain consciousness of the form, the concept, the
universal, what is ``more certain in and for itself''; and the task of
thought, when \opage{140}it has found the universal, is to be precisely
this: starting from the universal, to descend knowingly to the
individual, the concrete, and to explain it. Therefore Aristotle
attributes to \emph{experience}, the form of knowledge that relates to
the individual real, a different significance than Plato does; for him
it becomes an essential and indispensable form of knowledge, without
which a conception of the natural in its actuality becomes impossible.
And since natural existence has form immanent in it, it gains value in
its \emph{whole} extent for knowledge, which even in the most
insignificant and lowliest finds a task and a yield. The limits of
philosophy therefore become far wider in Aristotle than in Plato. In
Plato philosophy properly had as its object only the ideal in its
eternity and immutability; the knowledge of nature, whose object is
something becoming, fell outside the stricter philosophical knowledge
and became merely a matter of fitting recreation. Aristotle does indeed
let the knowledge which, through its relation to the eternal, to what
truly is, has the form of necessity and universality, be the
\gk{τέλος} [end] of knowledge and the form for the highest science of
philosophy: the science of the ultimate grounds of things; but his
interest in the actuality in which the ideal reveals itself in its
activity leads him to give the concept of knowledge and of philosophy
such an extension that it comprises not merely the strictly necessary
and universally valid, but also the probable and the customary, and
thus to draw under science that material, too, which can only become
the object of an incomplete knowing.

\divrule

c. \emph{The division of philosophy in Aristotle}. --- Aristotle gives,
in various places in his writings, various more or less elaborated
divisions of the domain of philosophy. The one that affords the easiest
survey of his system in its distinctive articulation is the division of
philosophy into \emph{theoretical} philosophy, with the
subdivi\opage{141}sions: first philosophy (metaphysics), mathematics,
and physics; \emph{practical} philosophy, with the subdivisions: ethics
and politics; and \emph{poietic} philosophy. While all philosophy as
such aims at knowledge, for theoretical philosophy the final end is the
knowledge of what is, for practical philosophy action, for poietic
philosophy what is technically or artistically produced. The
theoretical sciences have the highest rank; among them, again,
metaphysics, whose object is what is as such, and preeminently what is
in the highest sense: form, which in its purity is the matter-free
divine form. After metaphysics there follows in rank mathematics, whose
object is the unmoved in the corporeal, the determinations of magnitude
and number, which by an abstracting thinking are conceived as separated
from physical bodies; and after mathematics, again, physics, whose
object is what is moved and corporeal. When Aristotle in the
subdivision omits mathematics, which he has treated only little,
physics as \gk{φιλοσοφία δευτέρα} [second philosophy] receives its
place immediately after metaphysics as \gk{πρώτη φιλοσοφία}. In
practical philosophy the name of politics is used partly in a wider
sense, as comprising the ethical and the political, partly in a
narrower sense, as designating a separate science beside ethics.
Economics is conceived as an auxiliary science; rhetoric occupies, as
already indicated above, partly a similar position, partly it attaches
itself to the logical writings. Of poietic philosophy, of which
Aristotle has only begun the execution in the doctrine of poetry, he
has given no subdivision.
% errata, p. 141: printed "Underafdeling", corrected "Underinddeling"

In the division cited no place is assigned to what we sum up in one
word under the designation \emph{logic}, a designation that comprises
what Aristotle distinguishes as dialectic and analytic. The reason is
that he regards these disciplines as \emph{preparatory} in relation to
the study of the system of philosophy, as a \emph{propaedeutic doctrine
of method}, with which \opage{142}he who wishes to occupy himself with
the study of philosophy must be familiar \emph{before} he proceeds to
the philosophical sciences proper.

\divrule

d. \emph{Dialectic and analytic}. --- The concept of \emph{dialectic}
Aristotle applies in a far more limited meaning than Plato, for whom it
was synonymous with the concept of philosophy in general. Aristotle
conceives dialectic as a preparatory exercise in the development of
thought and as an aid toward clarifying, on the basis of experience and
the probable (\gk{τὰ} \gk{ἔνδοξα}, i.e.\ what seems true to the
multitude, or to the wise, or at any rate to some eminent individuals),
concepts and scientific principles. For this purpose it is to take its
starting point from the current views, which are conceived as the
expression of a comprehensive experience, from the results of the
earlier philosophical attempts, and from linguistic expression and
linguistic designations; and as the various points of view that have
asserted themselves through these are now set against one another in
order to be tested, it brings forth aporias, through whose provisional
solution by reasoning from probability it seeks a foundation for the
scientific development. Dialectic teaches \gk{ἐξετάζειν καὶ ὑπέχειν λόγον} [to examine and to submit to examination]; by its character it
becomes \gk{πειραστική} [tentative], whereas philosophy is
\gk{γνωστική} [cognitive]. That dialectic can lead by the path
indicated to science proper presupposes the self-activity of the mind,
which uses the given only as starting point and occasion, and has in it
only a means by which it is led to the independent production of the
concept, which is present in the mind as a disposition.

From dialectic \emph{analytic} is distinguished as the doctrine of the
method of \emph{science}. Through his analytic Aristotle founded the
science that later received the name of \emph{logic} and gradually took
on the character of a \emph{formal} logic. He has not yet given a
complete and evenly balanced development of the whole activity of
thinking, but, with particular regard \opage{143}to the demands of
scientific grounding and exposition, has elaborated in detail the
doctrine of inference and proof, while the doctrine of judgment and of
the concept is treated by him only briefly, and in his chief logical
work only in dependence on the doctrine of inference. Since the
doctrine of grounding through inference and proof points back to the
ultimate grounds, it leads to the point where analytic meets the
fundamental philosophical science, metaphysics, and to this it stands
in relation throughout, since Aristotle, even though he separates the
logical investigations from the metaphysical ones and gives the former
a separate place as preparatory to the other philosophical sciences,
nevertheless does not, like the later formal logic, conceive the forms
and laws of thought in abstraction from actuality, but develops them in
relation to real actuality, whose principles and fundamental relations
metaphysics determines.

Aristotle's chief logical work is the two Analytics, which bear this
name because they give an exposition of the elementary forms of
scientific thinking (\gk{ἀναλύειν} = to lead something given back to
its constituents), and to these are joined the writings on the
categories and \gk{περὶ ἑρμηνειας} (on propositions as the expression
of judgment), yet in such a way that the writing on the categories
stands at least as near to metaphysics as to logic.

By \emph{categories} (\gk{κατηγορίαι τοῦ ὄντος, γένη τῶν κατηγοριῶν}
[predications of what is, kinds of predications]) Aristotle understands
the statements about what is that express the main kinds of its formal
determinations. In the writing on the categories he enumerates 10:
\gk{οὐσία, ποσὸν, ποιὸν, πρός} \gk{τι, πού, ποτέ, κεῖσθαι, ἔχειν, ποιεῖν} and \gk{πάσχειν} [substance, quantity, quality, relation,
where, when, position, having, acting, and being acted upon]; in later
writings \gk{κεῖσθαι} and \gk{ἔχειν} drop out, because they can be
subsumed under others of the categories. \gk{Οὐσία}, as the expression
of the independent and substantial, forms an opposite to the other
categories, which express determinations of \gk{οὐσία}
(\gk{συμβεβηκότα} [attributes]). Between the categories and the
grammatical parts of speech there is a certain correspondence, which,
however, is not definitely brought out by Aristotle and is not
consistent throughout either.

\opage{144}Of the \emph{fundamental logical forms} the \emph{concept}
(\gk{λόγος}) is determined as the expression of the \emph{real essence
of the object} (it is \gk{λόγος τῆς οὐσίας} [account of the essence]).
To the \gk{οὐσία} of things there attaches \gk{τὸ} \gk{συμβεβηκός},
which is either a necessary consequence of the \gk{οὐσία} (as the
properties of the triangle are of the essence of the triangle) and is
therefore called \gk{συμβεβηκὸς καθ' αὑτό} [attribute in itself], or
something purely accidental that cannot be deduced from the essence.
The definition (\gk{ὁρισμὸς}) states the essence of things by stating
the proximate generic concept as determined by the distinctive specific
difference (\gk{εἰδοποιὸς διαφορά}).

\emph{Judgment} arises through a combination of notions; its expression
is the declarative sentence (\gk{ἀπόφανσις}), which may be affirmative
(\gk{κατάφασις}) or negative (\gk{ἀπόφασις}). Only with the statement
does there arise any question of truth or untruth, not with the still
uncombined elements. The truth or untruth of the judgment depends on
whether or not it combines what is combined in reality and separates
what is separated in reality (\gk{τῷ γὰρ εἶναι τὸ πρᾶγμα ἤ μὴ ἀληθὴς ὁ λόγος ἤ} \gk{ψευδὴς λέγεται} [for it is by the thing's being or not
being that the statement is said to be true or false]).

\emph{Inference} is the derivation of one statement (one judgment) from
others. It proceeds either from the universal down to the individual,
and is then \emph{syllogism} (in the narrower sense), or from the
individual by summation up to the more universal, and is then
\emph{induction}. The syllogism Aristotle treats with great
thoroughness, and he is the first to have given its theory. He treats
the \emph{categorical} inference with its 3 figures (\gk{σχήματα}),
which are distinguished by the place of the middle term in the
premises, and states the conditions for correct inference in the
various figures. Inference is placed in relation to the real in that it
is required that the middle term express the cause (\gk{τὸ γὰρ αἴτιον τὸ μέσον}).

\emph{Proof} is the scientific (knowledge-producing) inference
% errata, p. 144: printed "apodiktiske", corrected "videnskabelige (en Viden frembringende)"
from true and certain premises, either such as, being the highest
(uppermost), are certain in themselves,
% errata, p. 144: printed "høieste", corrected "høieste (øverste)"
or such as follow with necessity from these. Proof thus stands above
the dialectical inference, which rests \opage{145}on \gk{τὰ ἔνδοξα},
and above the rhetorical, which proceeds from \gk{εἰκότα ἤ σημεῖα}
[likelihoods or signs] and is only meant to persuade.

\emph{Induction} (\gk{ἐπαγωγὴ} or \gk{ὁ ἐξ ἐπαγωγῆς συλλογισμὸς} [the
syllogism from induction]) concludes from the fact that a predicate
belongs to all (or many) individuals that it belongs to the specific
concept of these individuals. Only complete induction does Aristotle
acknowledge as strictly scientific. The inductive inference, which
proceeds from the individual that is ``more certain for us,'' has ``for
us'' a greater distinctness and clarity than the syllogism; but the
latter, because it proceeds from the universal that is more certain in
and for itself, is in and for itself the more excellent and the more
convincing kind of inference.

In both kinds of inference we have as a limit something that has
certainty \emph{immediately}. In induction it is the sensible
individual, which appears with immediate distinctness; in the
syllogism it is the most universal: the ``principles'' (\gk{ἀρχαὶ}),
which give the derived premises their certainty but cannot themselves
be proved, since there is no more universal ground of proof (\gk{τὰ πρῶτα ὁρισμοὶ ἔσονται ἀναπόδεικτοι} [the first definitions will be
indemonstrable]). Without something thus immediately certain we would
be led into an infinity, and all proof would be abolished. The
immediately certain \gk{ἀρχαὶ} are the object of \gk{νοῦς} [mind], the
highest form of knowledge.

\emph{As the highest logical} (and at the same time metaphysical)
\emph{principles}, which condition all certain knowledge and all
demonstration, Aristotle posits those which formal logic designates as
the \emph{principle of contradiction} and the \emph{principium exclusi
tertii}. The principle of contradiction as a logical principle states
that of two statements, of which the one affirms and the other denies
the same thing, the one must always be false, the other true; the
principium exclusi tertii (medii) expresses the same thing thus, that
nothing lies midway between the two contradictory statements, but that
the one of them must necessarily be affirmed, the other denied, of
every subject. Since these logical principles are the highest, no
direct proof can be given of them, but only
\opage{146}\setcounter{footnote}{0}the indirect one, which lies in the
fact that no one who thinks can deny them.

\divrule

e. \emph{Metaphysics}. --- What, with a designation accidental in its origin, we call metaphysics (\gk{τὰ μετὰ τὰ}
\gk{φυσικά} [the things after the physics]), is called by Aristotle, as indicated above, \gk{πρώτη}
\gk{φιλοσοφία} [first philosophy], or else simply \gk{φιλοσοφία}, or \gk{σοφία} [wisdom], or
\gk{θεολογία} [theology]. It is the fundamental philosophical science which,
while the other sciences each treat their own, peculiarly determined, domain of what is, considers what is
as what is (\gk{τὸ ὄν ᾗ ὄν}), and, since what is is comprehended
through its grounds, goes back to the highest and ultimate
grounds, since only these can be the grounds of all being. In these,
then, metaphysics has its essential object. But the unconditionally highest ground of being it finds in the concept of the
pure, matter-free form, the divine \gk{Νοῦς} [Mind], which is the supreme source of all actuality
and is itself the highest actuality. As \gk{φιλοσοφία}
it thus becomes at the same time \gk{θεολογία}. And while it thus
considers the highest and most universal grounds of being, the principles,
it at the same time demonstrates, by the indirect way indicated above, the validity of the most universal scientific axioms:
the principle of contradiction\footnote{Its \emph{metaphysical} form is: \gk{τὸ αὐτὸ ἅμα ὑπάρχειν τε καὶ μὴ ὑπάρχειν ἀδύνατον τῷ αὐτῷ καὶ κατὰ τὸ αὐτό} [the same cannot at the same time both belong and not belong to the same thing and in the same respect].} and the principle of the excluded middle, which have logical and metaphysical validity for all
that is.

The solution of the fundamental metaphysical problem, \emph{to determine the concept of what is, the concept \gk{οὐσία}}, Aristotle prepares through a critique of his predecessors'
determinations of what truly is, and in particular through
his critique of the Platonic doctrine of Ideas, which was the first to put forward
the principle of form, even if it did not determine it in the
right way. He appropriates the Platonic determination of
\opage{147}what truly is as something eternal and essential, as well as
the conception of the universal essence of things as the condition of knowledge
and its adequate object; but the critique is directed against the separation by which Plato places the universal, hypostatized as
Idea, outside the world of phenomena, and does so in such a way that precisely
this separation is to become the condition of its true actuality.

If we sum up the objections that Aristotle makes against the doctrine of Ideas in various places in his writings, so that only
the essential points of his critique are brought out, they are
the following:

1) The assumption of Ideas --- an assumption which rests altogether
on an unjustified substitution of a \gk{ἓν κατὰ}
\gk{πολλῶν} [one over many] for a \gk{ἓν παρὰ τὰ πολλά} [one beside the many] --- is of no use for
the knowledge of concrete actuality, since the content of the Ideas
is nothing but a superfluous repetition of the content of that actuality. The
content of sensible things is carried over uncomprehended into the Idea, merely
with the added determination of eternity, of being-in-and-for-itself, so
that the Ideas of things, e.g.\ the Idea of man, the Idea of the horse, become
eternal sensible things (\gk{αἰσθητὰ ἀΐδια}) and come to resemble the anthropomorphic gods of the popular religion, which are nothing but
\gk{ἄνθρωποι ἀΐδιοι} [eternal men]. Instead of a knowledge of the essence of things
one learns only what explains nothing, that beside
things there exist still other substances. And while
no relief is gained for knowledge by this doubling of the content, this content moreover acquires in
the Idea, by taking on the form of an archetype existing for itself,
a determination of singularity which would have to withdraw it from
knowledge.

2) In placing the universal altogether, in its separateness from the individual, as Idea, i.e.\ as something substantial, Plato entangles himself in a multitude of absurdities. For since,
to be consistent, he must posit Ideas for everything in which something common
can be shown, he must in the first place obtain Ideas also of the
accidental, of mere concepts of relation, indeed of negations.

\opage{148}He must next obtain Ideas of subordinate Ideas, which
latter would thus be at once archetypes and copies.
He will have to posit several Ideas for a single being, inasmuch as
there are several common concepts in the determination of this being, e.g.\
in the concept of man the concept of the living being, of the
two-footed, etc., and must therefore compose the Idea of this being
out of several Ideas. Finally, since the Ideas are posited as
\gk{χωριστά} [separate], so that the archetype and the copy come to
stand side by side, he will again have to seek a higher common concept for the Ideas and the things, for the Idea of man, e.g.,
and the sensible man a \gk{τρίτος ἄνθρωπος} [third man], and so
on to infinity.

3) The doctrine of Ideas does not explain the existence of things; for when
the Idea is posited as separate from things, it does not become a constituent of their essence, and so does not contribute to their being.
In that Platonic separation it has been overlooked that it is a
contradiction not to place the Idea, i.e.\ the substance, in that whose
substance it is, not to place the concept of the thing in unity with
that whose essential determination it expresses; and it has likewise been overlooked that in natural objects matter and becoming
belong to their essence, so that an incorporeal essentiality subsisting for itself,
such as the Idea, cannot express this essence. Once
the separation has been made, the relation of things to the Ideas can only
be determined by empty, figurative expressions such as imitation,
participation; and the latter determination in turn leads to
the contradiction that the Idea, as being in something other, different from
it, namely in that which is posited as participating in it,
would become something accidental.

4) As little as the being of things, is their coming-to-be (first coming-to-be and subsequent development) explained by the Ideas,
since the principle of motion, of activity, is sought in these, which
are only resting archetypes and would rather have to be a principle
of standstill and rest than of activity. Here too
Plato has been able to give only figurative, mythical determinations of
the relation, but has not made it truly comprehensible.

\opage{149}When account is taken, not of what Plato aimed at
in determining the concept of the Idea, but of what he was actually able to carry through by way
of thought, and when the involuntary influence of the moment of intuition, which turns the conceptual separation of the Ideas and
things into an actual separation of externality,
is taken into consideration, these objections against Plato's
doctrine of Ideas strike its essential defects, as these were
brought out above in the critique of his philosophy. Aristotle himself lays the greatest stress on the two objections:
that Plato has separated the essence (the substance) from that whose essence it is, and that he has not been able to give the Ideas the determination of activity. Since, therefore, he is himself to develop the fundamental concepts, his endeavor must first and foremost be directed toward
avoiding these defects.

This shows itself also in the conception of the concept which metaphysics, as the science of what is as such, has
as its immediate task to determine: the concept \gk{οὐσία}. In opposition to the Platonic separation of the universal from the
individual, and to the conception of this separated universal as what
truly is, Aristotle determines the individual, in which
the universal is inherent (\gk{ἐνυπάρχον}), as that which is something that is, an \gk{οὐσία} in the original sense, the substantial, which
bears all other determinations of being. The individual, the
concrete single being, is \gk{πρώτη οὐσία} [first substance]; the species and genera of these single beings
are \gk{δευτέρα οὐσία}, i.e.\ substance in a derived sense; the remaining universal determinations are not substantial.
The individual essentiality is constituted by the union of form and matter (it is a \gk{σύνολον} of \gk{εἶδος} and \gk{τὸ ὑποκείμενον}),
in such a way that in this union the primacy belongs to the form,
which, in its determinateness as the generic concept stamped by the specific difference (\gk{εἰδοποιὸς}
\gk{διαφορά}), is what is properly constitutive;
for which reason the definition too, which states the generic concept so
determined, expresses the essence of the thing.

The determination of \gk{οὐσία} as the unity of form and matter leads
us immediately to \emph{the prin\opage{150}\setcounter{footnote}{0}ciples of Aristotelian metaphysics}. As the fundamental causes (\gk{ἀρχαὶ}) of all that is,
Aristotle names four: 1) matter, \emph{\gk{ὕλη}} (\gk{αἴτιον ἐξ οὗ, τὸ ὑποκείμενον} [the cause out of which, the substrate]);
2) form, \emph{\gk{εἶδος}} (\gk{μορφή, τὸ τί ἐστι, τὸ τί ἦν εἶναι} [shape, the what-it-is, the what-it-was-to-be], i.e.\ the essence of the thing
expressed by the pure concept of form, in abstraction from
matter\footnote{Cf.\ Metaph. VII, 7: \gk{λέγω δὲ οὐσίαν ἄνευ ὕλης τὸ τί ἦν εἶναι} [by substance without matter I mean the what-it-was-to-be]. The abstract concept states \gk{τὸ τί ἦν εἶναι}; the definition gives the content of the concept.}, \gk{λόγος τοῦ τί ἦν εἶναι, λὸγος τῆς οὐσίας, οὐσία} [the account of the what-it-was-to-be, the account of the substance, substance]);
3) the principle of motion, \emph{\gk{ἀρχὴ τῆς κινήσεως}} (\gk{αἴτιον ὑφ' οὗ, τὸ}
\gk{κινοῦν, ἀρχὴ τῆς μεταβολῆς, τὸ ὅθεν ἡ ἀρχὴ τῆς κινήσεως} [the cause by which, the mover, the principle of change, that from which the beginning of motion comes]);
4) the end, \emph{\gk{τὸ οὗ ἕνεκα}} (\gk{τὸ τέλος}). But these 4 principles
are reduced, even if the distinction acquires its relative validity in
certain relations and within certain spheres, e.g.\ in mechanical production, essentially to the two mentioned above, since
the principle of form in the domain of the organic --- and the organic
is for Aristotle the standpoint for natural existence as a whole --- includes in its concept the principle of motion and the end, since the form is what is active and what has as the goal of
its activity its own realization. The fundamental difference
among the principles thus becomes the difference between \emph{form} and
\emph{matter}, a difference which, as already indicated earlier, falls
within being as the difference between \gk{τὸ ἐνεργείᾳ ὄν} and
\gk{τὸ δυνάμει ὄν}, between the developed, fully actualized, active being and the mere disposition, the potentiality, between the
actual and the potential.

More precisely determined, \emph{form} for Aristotle becomes equal to
the conceptual, ideal essence of things, which is the object of knowledge, and which, as thought, is ideally separated (by a conceptual separation,
not, like the Platonic Idea, by a separation of intuition) from
the material element that conditions sensible individual existence.
It expresses the objective unity in the manifold things of like
kind. As the ideal essential it is eternal and free of becoming, so that even where, as forming principle, it determines the development of something that becomes, it is not touched by becoming. It is
\opage{151}
akin to thought, whose name it takes where it comes forth in the Godhead in its highest purity, as the matter-free form.

The more precise determination of the concept of \emph{matter}, of material stuff, Aristotle gains through an analysis of the conditions of becoming. The possibility of becoming cannot be explained by what is
alone or by what is not alone, but only by something that unites being and not-being in itself, which is thus in a
certain sense existent, in a certain sense non-existent.
But such a thing is what is according to potentiality, i.e.\ that whose
being is a being not yet existent; but this determination
is precisely that of matter. Every becoming presupposes a substrate
in which a transition takes place from one state to another,
from one property to another, and this basis, which can
receive the different determinations, is, as basis, precisely
only the potentiality for both, and has neither of them actualized in
itself. And what holds with regard to becoming in the
individual case holds with regard to becoming in general: it
requires a substrate which is nothing, but can become everything,
which is everything according to potentiality and nothing according to actuality. This
is \emph{the first matter}, \gk{πρώτη ὕλη}, which, as the condition of
all becoming, cannot itself have come to be, and which, as the opposite of form, from which all determinations come, is \emph{the
absolutely indeterminate}: \gk{τὸ ἄπειρον} in the qualitative sense. As this \gk{ἄπειρον} it cannot at any point of time
exist by itself alone; for in such detachment from the relation to form it would become the absolutely non-actual. As
the indeterminate separated from all form, matter could
not become an object of knowledge either (\gk{ἡ ὕλη ἄγνωστος
καθ' αὑτὴν}); only in connection with form is it apprehended
through sensation; its universal essence is thought
only through an inference by analogy from the individual material thing.

It is thus by \emph{the relation to form} that matter
has the determination that goes beyond the merely negative,
namely the determination of \emph{potentiality}, and by this it comes,
in Aristotle, under the concept of being, becomes in a certain \opage{152}sense \gk{οὐσία} (\gk{οὐσία πως}, whereas \gk{στέρησις}, the negation of formal determinateness,
is \gk{οὐδαμῶς} [in no way] such), so that it becomes a \emph{non-existent} only in a \emph{relative} sense, whereas in Plato it was
so in an \emph{absolute} sense.
But it is not only by this conception of matter as what is according to potentiality, by which
it is drawn under the determination of being that holds also for form,
that Aristotle sublates the abstract separation
between matter and form; he moves the principles still
closer to each other by giving the concept of matter, in its application,
an elasticity by which it can comprise everything
that stands in a relation to something other analogous to the fundamental relation between matter as the passive, potential, undeveloped,
and form as the active, actual, completed. Thus,
e.g., the female in relation to the male,
the lower spheres of the cosmic structure in relation to the higher,
the soul in relation to \gk{Νοῦς}, the more abstract generic concept in relation to the more concrete specific determinateness, are conceived as matter in relation to
form, and one and the same thing is, in its different relations, seen now from the standpoint of matter, now from that of form. Thus the soul becomes matter in relation to \gk{Νοῦς}, form
in relation to the body, and in the successive formation of matter each
single stage is form in relation to a preceding one, matter in relation
to a succeeding one. By the inner mobility thus
given to the concept of matter, it can, since the whole of existence is seen
as a graded series of formations of matter, be applied to everything
determined by form, right up to the point where the highest matter-free
form reveals itself in its purity as divine \gk{Νοῦς}
and as active \gk{Νοῦς} in man. But in all this approximation of matter to form, which leads Aristotle to
speak not only of a \gk{ὕλη αἰσθητή} [sensible matter] but of a \gk{ὕλη νοητή} [intelligible matter], matter nevertheless retains in its essence, for him, the original
difference from form which it has as a coordinate \gk{ἀρχὴ}. Matter
becomes, as indifferent toward the determinate form, the principle of contingency, mutability and imperfection; it
becomes, as the external condition of the realization of form,
\opage{153}the principle of natural necessity; it becomes, as equally original with form, \gk{συναιτία τῇ μορφῇ τῶν γιγνομένων} [joint cause, with the form, of the things that come to be],
a co-determining factor in everything natural, where the differences of individuality derive from it.

Since form and matter are determined by actuality and potentiality, they stand through this determination in a \emph{necessary}
relation, and point to each other. Matter, as the potential,
strives toward form, i.e.\ toward that for which it is potentiality, and which is itself the actuality of the potential; but this
striving toward form is awakened only by form and obtains
actualization only through it, which in this relation determines
itself as the principle of activity, as \emph{entelechy}, i.e.\ as the
inherent, actively forming, which has the goal of its activity in
itself. As entelechy, form shows itself from the side of being
that which leads matter from potentiality to actuality. This
transition from potentiality to actuality is \emph{motion}: the activity (the entelechy) of the potential \emph{as} potential.

Motion, which in its universal concept comprises all change and becoming, thus designates that transition: the intermediate stage between potentiality, as relative or absolute, and actuality; it is the combination of being and not-being which
strives toward being, an \gk{ἐνέργεια ἀτελής} [incomplete actuality], whose limit
is complete actualization. It presupposes something moving and something movable, an actual and a potential being;
it cannot be explained merely by the potential or merely by the
actual, but only by an action of the latter upon the former.
The actual must, as moving cause, have a priority
in relation to motion, even in the case where motion
seems to proceed from the potential, e.g.\ where an organic individual
comes to be from the seed. Wherever an actual being meets
with a potential one and no external hindrances act to
impede, motion arises with necessity in the potential as a joint effect of the \gk{ποίησις} [acting] of the mover and
the \gk{πάθησις} [being acted upon] of the moved. The action of the former upon the latter
Aristotle posits as conditioned by a contact (\gk{θτξις}), which \opage{154}he even lets take place in the case of the purely immaterial, though
in such a way that the immaterial unmoved, which moves the world,
is indeed to touch what is moved, but not to be touched by it.

In motion the original relation of form and matter is
expressed; motion must therefore be \emph{eternal}, like the terms of the relation. Every motion also requires, as the condition of
its coming-to-be and its ceasing, another motion, whereby motion in its totality shows itself to be necessarily without beginning and end. But in order that the \emph{actuality} of motion can be thought at all, the series of moving causes must
be limited; one must not go on to infinity, but
\emph{a first eternal actual starting point} must be posited,
\emph{a first mover}, which itself, as the presupposition of all motion, must be unmoved, not only in the sense that
it is not moved by another, but also in the sense that it does not move itself. While matter is something that is only moved without itself moving, something merely potential, and while
nature is something that is both moved and moves, that
unites potentiality and actuality, the first mover, the Godhead, is thus something that only moves but is not moved, and it is pure actuality.

In the concept of \emph{the Godhead}, as the first unmoved originator of motion, Aristotelian metaphysics is concluded, and in it
the principle of form is united with the principle of motion and the end.
And this concept of the Godhead receives in Aristotle more precise determinations, which reveal the proper essence of form,
its kinship with mind, with absolute subjectivity.
From the fact that motion, as eternal, must be continuous (\gk{συνεχὴς})
and as such one, it follows, since the unity of motion presupposes the unity both of the mover and of the moved, that
the first eternal mover must be \emph{one}. As the source of the
continuous and single motion it must next be unchangeable, being with absolute necessity what it is. As unchangeable it must be absolutely immaterial, indivisible and free of space, and as immaterial, again, without passivity, absolute actual\opage{155}ity, matter-free form, pure energy, \gk{τὸ τι ἦν εἶναι τὸ πρῶτον} [the first what-it-was-to-be].
But this highest energetic form is \gk{Νοῦς}; for only thought
in its purity is absolutely free of the material, absolute energy;
only it is something divine, because it does not have its \gk{τέλος} [end]
outside itself. The divine thinking is an \emph{uninterrupted,
unchangeable activity of thinking}, and precisely as determined by this absolute energy of thinking is the Godhead
the absolutely living (that which has an \emph{immortal life}: \gk{ἡ γὰρ
νοῦ ἐνέργεια ζωὴ} [for the actuality of mind is life]), and is the primal source of all life. The object of the divine thinking can only be the highest; but
this highest is the Godhead itself, the divine thought.
\gk{Νοῦς} is thus eternally one with its object; it is an \emph{eternal}
\gk{νόησις τῆς νοήσεως} [thinking of thinking]. This self-enclosed energy of thought
is the undisturbed \emph{blessedness} of the Godhead, and only the blessedness of theory is divine blessedness.

Since the Godhead, whose unity conditions the unity of the universe
and its harmonious coherence, is conceived as this self-enclosed eternal thought, whose object is only thought itself, and since
motion and all outward-going activity (all practical and poietic activity, both of which presuppose a need for something, an
end outside the agent) are excluded from its concept, its
moving influence on the world can only be conceived in such a way
that it \emph{awakens the desire of matter}, and through this
calls forth its motion. It moves, therefore, in the same
way as the good and the beautiful, which are desired, as the
thinkable, which is thought. The highest form, which has shown itself
to be the highest moving cause, shows itself by this determination also to be the highest \gk{τέλος}, toward which there is striving.
What the Godhead immediately calls forth in matter is motion in space: the original form of motion, and more precisely a circular motion, the only absolutely continuous and
uniform motion in space; and this it calls forth by,
touching the universe, setting it in motion from its circumference. This calling forth of a uniform, unchangeable circular motion corresponds to the uniform,
\opage{156} \setcounter{footnote}{0}
unchangeable relation of the divine thought to itself; but it is not sufficient
to explain the multiplicity of determinate motions in space in
the universe, still less to explain the manifold concrete formations of matter. This defect Aristotle, in the
writings that have been preserved to us, has sought to make good only insofar
as, in order to explain a plurality of circular motions with
different directions, and a general mutability conditioned thereby in the sphere of the elements, he posits separate eternal principles of motion for the celestial spheres beside the first
mover; but he has not clearly determined the relation of these principles
to the first mover, and just as little the relation between the highest
form, the divine \gk{Νοῦς}, and the separate forms, the
essentialities of individual things, by which alone a becoming of the concrete could be made comprehensible. Aristotle does indeed let the
\emph{transcendent} \gk{Νοῦς} act as \gk{φύσις} [nature], as \emph{immanent} in
existence\footnote{The good (\gk{τὸ ἀγαθὸν καὶ τὸ ἄριστον}) is, on Aristotle's conception, present in the universe as its order and formal determinateness, but
exists still more perfectly outside it in \gk{Νοῦς}, through whom the order is.}, as its entelechy (see below); but thereby
the question of the grounding of the many forms in the
one form is not answered.

\divrule

What at once shows itself as the essential thing in this brief
presentation of the fundamental concepts of the Aristotelian metaphysics
is the endeavor, while the primacy of the ideal is maintained, to see the
ideal, the form, in essential unity with matter, and to explain the life
and development of existence, the becoming of the concrete multiplicity,
from the essence and relation of these principles. That essential unity
is sought by bringing both form and matter under the concept of being,
and by that conception of the relative, successive stages of development
in existence which, in seeing them as stages of the
\opage{157}self-actualization of the form, the ideal principle of
activity, through the potentiality that itself gains actuality through
the realization of the form, lets them appear as a unity of potentiality
and actuality. Through the concept of \emph{development}, which Aristotle
so strongly accentuates, the unity of form and matter is most
definitely pointed to. The principle of actuality is, at the beginning of
the development, present in the disposition: nothing becomes actual that
was not in this, and no actualization can take place without the ideal
pre-existence of what is to have actuality. But what is the goal of the
actualization, its \gk{τέλος}, exists in the disposition only under the
form of potentiality, in unity with the potentiality, which, as the
potentiality of the form, of actuality, becomes living potentiality, a
germ of unfolding. And at every stage of the unfolding that lies prior to
the goal, what is to actualize itself is only relatively actualized; it
still has a remnant of potentiality in it, or rather: it is still
throughout potentiality for a higher form of actuality. Only where the
goal is reached is the potentiality exhausted, but only because it is
unfolded into unity with actuality. --- As regards the explanation of the
becoming of the concrete and of the living development of existence,
Aristotle, who, conscious of the decisive significance of the concept of
becoming, energetically maintains it against those who denied it, and
with it the concept of potentiality and of development, has, as it were,
given all the premises for it in his determination of the concept and
the conditions of motion. But both his striving to hold fast the
essential unity of the fundamental opposites and his striving to
comprehend the becoming that becomes possible only through that
essential unity are checked, as if by force, at the decisive point by
the barriers that bound Greek thought in general, and the determinations
of the Aristotelian metaphysics are definitively forced back to the point
where Plato stopped.

This first shows itself when we consider more closely the explanation
Aristotle gives of \emph{becoming}. His explanation can be summed up in
the words: \emph{Through} the form (the active) there is motion, becoming
in the material (the passive), and what becomes \opage{158}is the
composite of matter and form, which in the motion has the form in itself
as entelechy. But now the further determinations are added that, while
the form calls forth motion, the motion is to be \emph{only} in the
passive, not in the active, and that, when the development has reached
its goal and the form is realized, it shall be possible to say of the
form only \emph{that it is, not that it has become}. What is the
principle of becoming thus falls outside becoming, as untouched by it;
but by the same token \emph{becoming itself is left unexplained}, when
what is to bring it about has no place for it in itself. Aristotle has
not explained, and cannot explain, how the incomplete realization of the
form during the development, which is as it were a partialness of the
form, a partial negation, is to be thought, when the form itself is held
outside becoming, and how the form can then become entelechy. But
precisely the process of the \emph{inner} principle of motion must first
be explained before the determinations Aristotle gives about the
realization of motion in the one existent through \emph{another}
actually existent can gain significance. And just as little as becoming
in the individual is explained by the individual form is the universal
ground of becoming explained by the universal principles: matter or the
highest \gk{εἶδος} [form]. Matter is the determinationless, the passive,
in which no activity whatsoever, not even such as lies in a striving, a
``desire,'' can be thought. And if we seek the origin of motion in the
divine \gk{Νοῦς} [mind], we come up against insoluble difficulties. We
are to think of \gk{Νοῦς} as the original source of motion and becoming,
although it is a \gk{χωριστόν} [something separate] that relates only to
itself; we are to think of it as touching the world although it is
absolutely immaterial, and as calling forth motion from the circumference
of the world although it has no relation to space;
% errata, p. 158: printed "har Rummet", corrected "har til Rummet"
we are to think of it as effecting motion by awakening the desire of the
absolutely determinationless; and we are finally to think of it as the
first and, as such, universal principle of motion, although several
particular eternal principles of motion stand underived beside it. On no
point, \opage{159}then, do we find in Aristotle a real solution of the
problem of becoming, and such a solution is cut off outright by the
statements that exclude motion from the active form and let the form in
its completion be not something that has become, but only something that
is. For in this way we again get in the form something that is, something
actual, outside which negation falls, and which therefore remains without
inner mobility; and by the same token potentiality is transformed into an
\emph{empty} potentiality, a \emph{potentiality for nothing}, a
\emph{non-being}. The Platonic, conscious, absolute opposition between
being and non-being comes back involuntarily in Aristotle. Matter, as
this potentiality whose principle of development has been annihilated,
becomes not the being that is not yet, but what is not, whose non-being
is abiding.

Since the task of explaining coming-to-be and becoming miscarried for
Aristotle, and miscarried although he already stated the conceptual
determinations by which it was to be solved, the lack of its solution
implies that the task of explaining \emph{the concretion of real
existence} must also miscarry for him; and it points back to the fact
that \emph{dualism}, which he sought to remove by his determination of
the fundamental moments, is not in reality removed, and must come to
view again in the development of the system.

As regards the explanation of the concretion of real existence, it was
demanded by Aristotle's whole striving to embrace with thought all the
wealth of existence and to take it up into knowledge, and in the concept
of development and of the principle of development he seemed to possess
the condition for it. But that his striving must miscarry on this point
too, we shall see by a consideration of the fundamental opposites.

The concrete multiplicity of existence can be explained from matter
neither by a direct positive working nor by an indirect working through a
passive resistance. A positive working would conflict with the general
determination of matter as the passive, and with Aristotle's statement
that the becoming by which the determinate, concrete thing comes to be
always presup\opage{160}poses a determinate opposite as an already
existing property or state in matter, hence an already form-determined
matter. A negative working of matter through a passive resistance, which,
as the realization of the form in existence is checked by matter, lets
the form in its unity divide itself, in the successive overcoming of this
check, into a graded series of forms, is indeed assumed by Aristotle, but
without his being able to make it comprehensible. For the resistance can
be comprehended only when a determinateness is already presupposed, which
could be due only to a determinate form, not to the absolutely
determinationless matter. As the absolutely determinationless
potentiality, matter would have to be absorbed without resistance into
the pure form, without being able to call forth a multiplicity of forms
of existence. Of this, too, Aristotle has a feeling, and he is thereby
led to let matter, as it were, appear from the very outset with a
determinateness; but this determinateness, the determinateness through
the first elementary opposites, is already a determinateness of form,
which stands unexplained.

Since, then, neither a positive nor a negative working of matter could
explain the concretion of existence, this concretion could be explained
only from the form. As a reality, something actually existent, it would
indeed have to have its ground in what is being in the proper sense; and
Aristotle's statements about the conditions for the existence of the
qualitatively determinate also point back to the form. The explanation
through the form would then more precisely be an explanation through the
\emph{highest} form as unity; for the multiplicity of existence would, if
the multiplicity of forms that conditions it were not led back to unity,
stand as an uncomprehended given. But what makes the explanation
impossible here too is that the highest form, since, as shown above,
negation falls outside form in general, cannot become dialectical, cannot
get the possibility of inner unfolding in itself. Aristotle's striving
does indeed aim at conceiving it as activity; but this striving is
checked not only by that conception of the form's relation to
\opage{161}becoming, but already by the fundamental peculiarity of Greek
thought, for which form could not become mobile for the reason that it is
originally the form of the work of art, and as such resting form: a
relation that also shows through in Aristotle when \gk{εἶδος}, which as
the determination of the Godhead became \gk{Νοῦς} and \gk{νόησις τῆς νοήσεως}
[thinking of thinking], so that the determinations of mind and of
subjectivity are stated, gets, as the \gk{εἶδος} of nature, the
determination \gk{μορφή} [shape] as its synonym. And even when Aristotle
conceives the pure \gk{Νοῦς} as pure, uninterrupted energy, it is
nonetheless resting insofar as in its relation to itself it lacks
otherness, opposition, and therefore eternally, unchangeably thinks only
itself, produces only its own undifferentiated reflection. A possibility
of solving the problem of the unfolding of the one form into the many
forms might seem to show itself if we presupposed --- what we are
nonetheless not entitled to assert on the historical evidence, even if
certain hints and analogies point toward it --- that the view had hovered
before Aristotle that the divine thought's, the pure form's, thinking of
itself was to be a thinking that embraced all the intelligible as a
unity in which the whole teleologically coherent content of the world's
reality was concentrated. The divine thought's thinking of itself would
then become a full unity, in which everything that appears in the world
separated in multiplicity was joined inseparably; and the world would
therefore, since the individual forms realized in it would become
determinations in the one divine form, not become something detached
from the Godhead, nor God's thinking of it a thinking of an other. But
this possibility would nevertheless vanish again, since Aristotle would
not be able to remove the difficulties which, in consequence of his
conception of the Godhead's separation from the world and of the form's
relation to negation, as well as of his paralleling of the relation of
universal concepts to particular ones with the relation between
\gk{δύναμις} [potentiality] and \gk{ἐνέργεια} [actuality], must stand in
the way of that develop\opage{162}ment and carrying-through of that
thought which would be necessary for a solution of the problem.

In Aristotle's conception of the concept that expresses absolute
actuality there thus appears most clearly the fundamental defect of the
Aristotelian metaphysics, and of Greek thought in general: the lack of
power to make what is living and moved through negation. By letting the
concept of form become that of the divine thought, by positing this
thought as pure energy and as the \gk{νόησις τῆς νοήσεως} independent of
external conditions, by determining the human \gk{Νοῦς}, identical with
the divine, as \gk{εἶδος εἰδῶν} [form of forms], as embracing in its
unity the lower forms of soul, and as thinking itself in its object,
Aristotle had come as far toward a conception that could give the
principle of form inner motion, and a fullness proceeding from it, and
lead to seeing the concretion of existence develop out of the principle,
as was possible for Greek thought. But at this decisive point the
striving of his thought is checked as by an invisible power: absolute
being, which separates negation from itself, becomes merely positive
being; the eternal \gk{νόησις τῆς νοήσεως} becomes only as it were a
staring of thought at itself in its uniformity; reflection sinks back
into the immediacy of rest, and the wealth of the world's reality stands
unexplained beside the principle.

When the failure of the solution of the problem of becoming for
Aristotle inevitably entailed that the solution of the problem of the
concretion of existence must fail, this points back to the fact that
\emph{dualism} has not been removed; for it was precisely the involuntary
abstract separation of the fundamental moments that made becoming
incomprehensible, since becoming is the first abstract expression of a
connection of the determination of non-being, of negation, with being,
and thereby represents in general the unity of opposites. Aristotle had,
by the determination of the conceptual relation between form and matter,
and especially by the conception of form as principle of development,
decisively pointed \opage{163}toward the unity in principle; but dualism
nevertheless breaks forth again and separates what was apparently united.
The separation shows itself on the side of the form chiefly in this, that
in its highest completion it becomes the form separated from all matter,
and that in general, even where it is seen in relation to development, it
separates itself from the becoming of the material; --- on the side of
matter chiefly in this, that a power of resistance is ascribed to it that
can occasion what deviates from the form, hence an independent working,
through which a fundamental opposition to the form betrays itself. In
this fundamental opposition matter and form stand as equally original;
their connection for intuition becomes the immediate connection of what
is different, their unity for thought only a postulate. This also lies in
Aristotle's statement that the two principles, form and matter, can be
led back neither to each other nor to the same generic concept (for the
determination of being, which becomes common to them, Aristotle does not
conceive as a generic concept, any more than the concept of unity).

In a conspicuous manner the dualism Aristotle combated breaks forth in
his doubling of existence into \gk{Νοῦς} and the universe. Although his
striving tends toward a penetration of matter by form which, consistently
carried out, would lead to seeing them as inseparable, so that the
highest form would be the one that most completely penetrated and
mastered matter; to seeing the eternal as eternal in what is determined
by time and motion, and what is independent of space as being in space in
a spaceless manner, as unity; yet this striving is checked by the
original separateness of the principles. For Aristotle the whole of the
world's reality takes shape for intuition as the cosmos, the
all-embracing, living work of art ensouled by the form. This formed
universe, whose matter is the potentiality of what forms it, and which is
eternally moved through its relation to its divine mover, includes in
itself a complete graded series of forms, which rises from the lowest to
the active \gk{Νοῦς}, identical with the divine thought, \opage{164}in
man, and in this completeness expresses the essence of form wholly.
Aristotle here comes as close as possible to letting form, as divine
\gk{Νοῦς}, have its place \emph{within} the world's reality. Nature in
\emph{its wholeness} is seen as ensouled, form is posited as the world's
inner moving entelechy, and this form, which unfolds itself
teleologically in the series of particular forms, shows, through the
culmination of the series in man's \gk{Νοῦς}, its unity with the divine.
In man, therefore, the divine \gk{Νοῦς} as such has a being in the world;
and since, by reason of the relation of its immaterial essence to spatial
determinateness, its spatial separation from the universe cannot be held
fast by thought, there seems to be no obstacle to thinking of the divine
\gk{Νοῦς} in general as \emph{immanent} in the world (cf.\ also the
development of the doctrine in the Peripatetic school), and all the less
so as Aristotle indeed also thinks of the eternal movers of the heavenly
spheres as unmoved in the moved spheres, the soul as unmoved in the moved
body, form as unmoved and eternal in changeable and moved matter, and the
active \gk{Νοῦς} in mankind as without partition and becoming.

But Aristotle nevertheless does not carry the unity through to this last
point. The cosmos does not become the whole of existence, nor does the
divine \gk{Νοῦς}, which yet in its essence includes all actuality, become
the sole existent; but actuality is doubled into the pure, matterless
\gk{Νοῦς} and the world's reality, in which matter, which could not be
pushed aside by thought, gets its place: the world in its fullness comes
to stand as an actuality beside the actuality that according to its
concept is all-embracing; the cosmos, in which actuality is eternally
reproduced from potentiality, asserts itself beside the pure energy in
which all potentiality has vanished.

This doubling of actuality is an unavoidable consequence of the
fundamental determinations. At the same moment as dualism seems to have
vanished in the world's reality, since this reality is apprehended by
intuition as the harmonious unity of the cosmos, it must \opage{165}break
forth in the doubling, since thought again separates out the form in its
purity, and intuition fixes it in its separation from the cosmos by
involuntarily making this a spatial separation. Dualism must do so
because the moments immediately united in the intuition of the cosmos,
which were what was originally different, at once betray their
dualistic separation as soon as thought determines them by its forms.

\divrule

f. \emph{Physics.} --- We have in the foregoing stressed the great
difference there is between Plato and Aristotle in their relation to
the physical. Plato hastens away from natural reality in order to lose
himself in the eternal Ideas, the only adequate object of philosophy;
Aristotle is constantly drawn toward the reality through which form
reveals itself, and which precisely for that reason everywhere has
something divine and admirable about it and holds within itself an
inexhaustible wealth for knowledge; and he does not deny this
natural-scientific knowledge the name of philosophy, even if it leads
only to the probable, the possible and the relatively universal, and
even if it must stop before it has completely mastered its object.
Aristotle is at once philosopher and natural scientist: he is a master
in the observation of facts; by his striving to explain them from their
peculiar physical grounds he stands highest in his period; and he has
not merely appropriated all the knowledge of his age concerning all
the domains of nature, but has himself in many directions given it
greater scope and depth.

Aristotle's physical writings have been named earlier. They treat
partly the most general principles and conditions of natural
existence, the order of the structure of the world, the heavenly
bodies, the elements and their relations, and the meteorological
phenomena, partly the organic; and the writings that have the organic
for their object are either descriptive, like the
Zoo\opage{166}logy, or investigative, like the books on the soul and
various treatises of physiological and psychological content.

In this presentation we separate psychology, which with Aristotle
forms the concluding portion of physics, from the other parts of
physics, and indicate briefly the fundamental features of the
Aristotelian conception of nature.

The fundamental philosophical science, metaphysics, had for its object
what is as such, and preeminently that which was in the highest sense
something that is: the unmoved and incorporeal form. Physics has for
its object the existent that is moved and corporeal. Everything that
is to be called a natural being (\gk{φυσική οὐσία}) must either be a body, or have a
body (e.g.\ man), or be the principle of something that has a body
(e.g.\ the soul), so that form is considered by physics only in this
relation. It must, as distinct from the object of mathematics, which
is the unmoved in the corporeal, be something moved. It must finally,
as distinct from the products of art, which have the ground of their
motion outside themselves, have it within itself. Nature (\gk{φύσις})
Aristotle therefore defines as the cause of motion in that to which
motion belongs originally and not merely \gk{κατὰ συμβεβηκός} [accidentally], and the
natural thing (\gk{τὸ φύσιν
ἔχον}) as that which has such a principle of motion
within itself.

Aristotle's definition of nature implies that, despite the significance
that matter and motion acquire for the natural, yet, in consequence of
the relation of both to form, form, as the principle of motion and of
actuality, must become the proper essence of the natural; and when in
the natural a distinction is made between two kinds of causes, final
causes and material (mechanical) causes, the former, which express the
activity of form, must be posited as the proper causes, the latter only
as conditions, albeit as indispensable conditions, for the realization
of the end. Aristotle therefore also designates nature as \gk{οὐσία}
[substance], he posits \gk{μορφὴ} [shape] as equal to \gk{φύσις}, while \gk{γένεσις}
[coming-to-be] becomes a \gk{ὁδὸς εἰς φύσιν} [way to nature]. The relation between the
active \opage{167}form that is immanent in nature, unfolded in a graded
series and constituting the essence of nature, and the completed
actuality of form in the divine \gk{Νοῦς} [mind], we can determine thus:
that natural form, which is the proper \gk{φύσις}, relates to the Godhead as
the immanent entelechy to the fully realized energy, or as in the human
the soul relates to \gk{Νοῦς}. God and nature can be used as synonyms when
the talk is of the activity of nature.

This conception of the relation of nature to form makes Aristotle's
conception of nature dynamic. Nature is not a composition of
unchangeable, merely quantitatively different atoms; its motion is not
the contingently-necessary mechanical motion of such atoms; but nature
has qualitative differences and qualitative change (transition) within
itself; it is, as moved, animated in its totality by a forming force,
and motion is the immortal life of nature.

As determined and governed by form, natural existence has its \gk{τέλος}
[end] within itself, and as \gk{τέλειος} [complete] it is bounded, artistically
rounded off in itself. The unlimited, the quantitatively infinite, would
be the formless and as such the imperfect. Infinity is present in the
universe only potentially: in space as the possibility of infinite
division, in motion and in time (``the number of motion with respect
to the earlier and later'') as the possibility of infinite
continuability, in number as the possibility of infinite
increasability. Infinity is thus in nature always only becoming.

The unity of form and end entails that the form active in the
individual natural thing and in nature as a whole is seen as acting
teleologically. The teleology of nature is, as grounded in the immanent
principle, an immanent teleology, an inner purposive activity, and more
precisely it is to be determined as an unconscious activity, in analogy
with unreflective artistic production. Nature is in its teleological
activity \gk{δαιμονία} [daemonic], not \gk{θεία} [divine].

The purposive activity of form cannot, however, since it is conditioned
by matter, everywhere in the particular reach its goal
\opage{168}unhindered. Because matter, which is the source of
contingency and of blind natural necessity, is a factor in natural
existence, nature frequently produces something beside the goal,
something for which no end can be stated, and frequently it realizes
the end only incompletely; and the disturbance that is thereby brought
into the purposive working of nature is precisely what prevents
natural science from reaching strict universality. In accordance with
its peculiar being, matter can exert a resistance against form, and
this resistance Aristotle conceives not merely as the cause of the
abnormal phenomena, which are as it were a distortion of the end, and
have their ground in this, that form does not have a matter of the
right quantitative and qualitative determinateness to form, so that the
form-moment of nature (\gk{ἡ κατὰ τὸ εἶδος φύσις}) does not gain mastery over its material
side (\gk{τὴν κατὰ τὴν ὕλην φύσιν}), whereby the abnormal must come to be; but he also points
to it as the cause of the lower stages of nature in general, of the
successiveness in the realization of form in the series of natural
things. But in this graded series of existence the resistance of matter
is precisely overcome successively by form, which as it were gains
strength by exercise and reveals itself ever more clearly. From the
inorganic, which nevertheless already has a degree of life, albeit the
lowest, there is an advance through the species of the organic, which
in an unbroken sequence rise from the less perfect to the more perfect,
until at last form attains its pure revelation in the human mind.

The totality of the natural, the universe (\gk{Κόσμος} or \gk{Οὐρανός}), whose unity
was for Aristotle given by the unity of the highest principle, and whose
eternity follows from the eternity of motion, he conceives as
spherical, as resting as a whole, in motion only with respect to its
parts, and as comprising all matter, the whole mass of the elements,
which is ordered according to the naturally determined place of the
elements in such a way that the four elements that are changeable and
opposed to one another by their fundamental qualities and motion,
earth, water, air and fire, have their spheres nearest the center,
\opage{169}while the higher spheres are filled solely by the ether,
which is without opposition and unchangeable, the divine in the
material, and which properly, even if it is named \gk{πρῶτον στοιχεῖον} [first
element], is not in the strict sense an element, but ``another kind of
body, and more divine than the so-called elements.'' The earth, as a
resting, spherical body, occupies the center of the universe. Water,
whose natural place is nearest the earth, fills the depressions of the
earth's surface, so that the surface of the sea also lies in the
circumference of a sphere. Air and fire surround the two first-named
elements in concentric circles. The region of the ether, which encloses
the sphere of fire, is according to Aristotle divided into a larger
number of concentric spheres (56 in all), of which some carry the
spherical heavenly bodies that are fixed in them, while the others
serve only, by their peculiar motion, to produce the path the planets
describe. The higher circles relate to the lower as form to matter, as
the active to the passive, as the mover to the moved. Each individual
ether-sphere has its own peculiar principle of motion, and these
principles of motion, which are the origin of an eternal motion, must
themselves be eternal and unmoved.

In the universe there appears a distinctly marked opposition between
the sphere of the four elements on the one side and that of the ether
on the other, between the terrestrial and the celestial region. The
region of the elements is the domain of the various mutually opposed
motions and of changeability; that of the ether is the domain of the
eternal, opposition-free circular motion and of unchangeability. Within
the region of the ether a distinction again appears between the
uppermost sphere, \gk{πρῶτος οὐρανός}, the heaven of the fixed stars, and the other
spheres, the spheres in the region of the planets. The former is
nearest the first mover, is not in space (which Aristotle conceives as
the limit of the enclosing body against the enclosed), but encloses
everything without being enclosed, communicates its motion to all the
lower spheres, and contains a multitude of heavenly bodies, which it
carries, in an undisturbed, uniform \opage{170}circular line, in the
most perfect direction, from right to left, i.e.\ from east to west.
The other ether-spheres contain, as stated above, either none or only
one heavenly body, and this heavenly body, the peculiar motion of whose
sphere goes in the direction opposite to that of the heaven of the
fixed stars, from left to right, has, on account of the influence of
the higher spheres on the lower, a compound motion deviating from the
pure circular line. But the planetary spheres are, despite their
difference from the uppermost heaven, like it free from change and from
being acted upon in the strict sense: they are animate and, by their
divine principles of motion, of a far higher nature than man. The
deviation of the planetary orbits from the uniformity of the orbits of
the fixed stars, and in particular the inclination of the sun's orbit,
acquires, according to Aristotle, the greatest significance for the
sublunary region. While the motion of the ether-spheres, in particular
that of the sun's sphere, calls forth an alternation of heat and cold
in the atmosphere, and thereby the transition between the elements and
their uninterrupted mixture, the direction of the orbits, in particular
the inclination of the sun's orbit, calls forth a law-governed
alternation of coming-to-be and passing away in the region below the
moon, and thereby \emph{the endless cycle of becoming.} Endless
becoming is for the perishable a substitute for the infinity of
imperishable being; through it the perishable takes part in its own
way in the perfection of what eternally is, and through it the
completeness of the universe is preserved. ---

Both by his whole view of nature and by the demands he makes on the
knowledge of nature, Aristotle points beyond what is peculiarly Greek;
but here too, as with the fundamental metaphysical determinations, his
striving is not carried through, but is forced back, as it were by
violence, within the bounds of the Greek spirit.

It has already been shown earlier how Aristotle's conception of form as
the immanent form, which gives itself actuality in concrete, real
existence, had to lead to an essential emphasis being laid on this
existence. The sub\opage{171}stantial existent, \gk{οὐσία}, in the proper
sense, was determined by Aristotle as a \gk{τόδε τι}, as the individual
subject, natural individuality in its concretion; and this concrete
individual had, as essential and substantial, to become an essential
object of knowledge, just as the form of knowledge that takes up this
individual, \emph{experience}, had to claim an essential significance.
And since the immanent form was seen as a graded series of forms
extending through the whole of natural existence, this existence
acquired, in its whole extent, value for knowledge, and it acquired
this value \emph{as} natural existence, i.e.\ as the existence for whose
completeness matter, in which form is immanent, as well as motion, are
necessary, and which is therefore known in its actuality only when it
is known as the unity of form and matter. Despite all the weight,
therefore, that Aristotle, in accordance with his idealistic
fundamental view, must lay on the formal principle of natural things,
his striving to see the ideal as essentially one with the material and
as immanent in it nevertheless leads him to lay an accent on the
sensible phenomenon, from which alone knowledge of nature can be
reached by the way of observation and experience, and to demand for
this knowledge everywhere an investigation of the peculiar physical
grounds and concrete conditions of natural things (\gk{φυσικῶς σκοπεῖν}), in
opposition to an abstract-universal conceptual reasoning (\gk{λογικῶς σκοπεῖν}).
Aristotle's determinations intimate a striving toward the knowledge of
nature which only a more recent age carries through, toward the
knowledge of nature through experience demanded in principle. But the
realization of this striving is hampered already by Aristotle's
metaphysical determinations, and in the last instance by the general
character of the Greek conception of nature.

When Aristotle in his Metaphysics determines the individual as \gk{πρώτη οὐσία}
[first substance], and lets the individual being be constituted by the
union of form and matter into a \gk{σύνολον} [composite whole], he posits form
in this union as the essential and properly constitutive element. But
in the \gk{εἶδος} [form] which thus constitutes \opage{172}the essence, the
element of universality asserts itself (cf.\ the Metaphysics), and
since only the universal becomes the proper object of knowledge, while
sensible things as such are not objects of knowledge, one can indeed
with knowledge descend to the most concrete concept of species, to the
lowest concept, to use the expression of formal logic, but cannot go
further without coming outside the limit of knowledge; and this will
mean, in accordance with the relation between knowledge and being that
Aristotle too maintains, that the universal after all claims for itself
the proper essentiality. Aristotle thereby comes to designate form as
\gk{πρότερον} [prior] and \gk{μᾶλλον ὄν} [more truly being] than the individual thing in
which it is immanent; indeed, he even uses outright of form the
designation \gk{πρώτη οὐσία}.

The fundamental determinations of his metaphysics made it impossible
for him to carry through the precedence of the individual being to the
determination \gk{οὐσία}. For this to be able to happen, a positive more
would have had to be shown in individuality, which could distinguish it
from the species; but such a thing Aristotle cannot claim for it. Since
form is posited as the universal, that plus would have had to be sought
in matter, which Aristotle conceives as the ground of individuality.
But since matter is the absolutely indeterminate and what is in itself
absolutely devoid of actuality, it cannot be thought to determine form
in such a way that the individual could acquire a content of actuality
peculiar to it, and the individual substantiality of the individual
being thus becomes something inexplicable. Aristotle must maintain it
according to the peculiar direction of his thinking; but by the kind of
thinking that bears and encompasses his individual thinking, and which
lets form, the ideal universal, \emph{immediately} master its opposite,
so that this opposite vanishes, as mere potentiality, in form, and form
cannot gain concretion through its relation to it, he must be led to
give determinations concerning the factors constituting individual
existence that cancel what was maintained and make its possibility
incomprehensible. And as the individual being again recedes into the
shadow of the universal, experiential \opage{173}knowledge must step
back before a speculative conception of nature through abstract
concepts.

What has shown itself as immediately connected with the determinations
of the Aristotelian metaphysics will now also show itself as more
deeply grounded in the general character of the Greek world-view and of
the thinking determined by it. Since the Greeks intuited existence as
an immediate unity of spirit and nature under the primacy of the
spiritual, the natural had, by reason of the immediacy of the unity, to
be seen immediately as determined by the spiritual, as mirroring it.
Nature was seen immediately in the light of the human mind, was
anthropomorphized, acquired a human character, a human being. Therefore,
in the conception of nature, \emph{experience} does not come into its
own. Experience, in the pregnant sense in which a modern natural
science takes the concept, presupposes a deeper separation between
thought and its natural object than the Greeks knew; it presupposes
that the natural, the material, has as it were been liberated to a
relative independence, that it has been able to make its own laws
valid, and thereby has asserted itself as a peculiar source of
knowledge for the mind. But with the Greeks matter is only the
non-existent in relation to thought, and therefore it, and that in
which it is the characteristic determination, cannot win for itself
such an independence and such a significance. The Greek regards
nature, as a whole and in the particular, essentially not with the eye
of the observing natural scientist, but with the \emph{artist's} eye.
Just as the totality of nature was intuited as a cosmic work of art, so
the particular was intuited in analogy with the singularity of the
artistic ideal, which immediately reveals the idea and lets the
contingencies of finitude vanish; the singularity of concrete natural
reality, natural individuality, is transformed into the distinctiveness
of beautiful individuality, and intuition, in its immediate unity with
thought, becomes an intuition that relates itself to the distinctive,
not to the individually singular natural. Through \opage{174}intuition
so determined, thought determines nature: the conception of the
physical becomes, where thinking comes to clarity about the principles,
determined by these \emph{prior to experience}: \emph{the conception of
nature becomes metaphysical, not physical.} Greek philosophy therefore
did not know the art of questioning nature through experiment either;
for by such questioning nature would precisely be recognized in its
relative independence. Only late and to a small extent does the
experiment come forward among the Greeks, and then outside philosophy.

This kind of knowledge of nature we now recognize also in Aristotle,
despite his striving for a knowledge from experience, and despite his
extensive knowledge, to which his natural-scientific writings bear
witness. Of his knowledge of nature too it holds that it essentially
acquires a metaphysical, not a physical, character. The fundamental
metaphysical concepts, form and matter, are applied everywhere in
physics as metaphysical concepts, without being made more concrete by
the relation to the physical; the intuition of the universe as a whole
is determined immediately by them, not by a summing-up of the
conceptually grasped real; the nature of the fundamental forms of the
material, the elements, is determined not through observation but
through an a priori deduction; and thought is involuntarily drawn
toward those domains of nature where the fundamental metaphysical
relation finds the easiest attachment, toward the domain of the
organic, and the lower spheres of nature are seen as it were only in a
light reflected from the organic, while in the organic itself the
mechanical causes are indeed stressed, but without their claiming a
conceptually determined positive significance in relation to the
teleological. And that the Aristotelian knowledge of nature acquires
this character may in part be conditioned by his age's lack of the
numerous technical and scientific aids to observation and calculation
which the natural research of a more recent age commands, and by its
fragmentary acquaintance with the phenomena of nature; but it has, as
all the main determinations of the Aristotelian philosophy indirectly
show, its last and most essential ground \opage{175}in the general
character of the Greek conception of nature, which itself becomes a
co-condition of those deficiencies. ---

What holds of the Aristotelian \emph{knowledge of nature} in relation
to the general Greek knowledge of nature holds also of his
\emph{general view of nature}, namely, that while it strives beyond the
bounds of the Greek view of nature, it is nevertheless not able
actually to overstep them, so that its categorial determinateness
becomes the same as that of the Greek view. The Aristotelian view of
nature is the expression of the Greek view of nature in its highest
development and at the same time in its \emph{transition} to a new and
higher form. The \emph{beauty-view} of the universe, the intuition of
real existence as a cosmic work of art, as cosmos, is preserved in
Aristotle, and form, although in its ideality it has no spatial
determinateness, nevertheless becomes spatial form in nature. But on
the foundation of the beauty-view a richer, deeper, more living
conception is in the process of working its way forward: \emph{the idea
of life} is in the process of replacing \emph{the idea of beauty} as
the idea that determines everything for real existence. Nature, which
Aristotle determines as animate in its totality, and that not merely by
an involuntary anthropomorphizing but in accordance with the conceptual
relation to the immanent active form, is by this conception brought
under the point of view of the organic, and the concepts of form as the
principle of development and of the inner teleology in the particular
and in the whole of nature point toward the organic. The \emph{organic}
conception of nature, which, consistently carried through in the form
of thought, in accordance with the essential relation between the side
of ideality and the side of reality in the living, includes the
positive significance of the latter and therewith the significance of
knowledge of the real, is present in Aristotle in a living intuition,
and the conception of a modern age is touched upon. But the previously
demonstrated deficiencies of the Aristotelian philosophy make
themselves felt as soon as the development determined by form is to be
made clear to thought, and as soon as form, the merely ideal and
eternal, which as such lacks the determination of reality, is to be
seen as natural energy. Just as the determination of the
ener\opage{176}getic activity of form could not be held fast, but
involuntarily went back to the conception of form as the eternally
finished \emph{resting} form, untouched by becoming, which did not
essentially differ from the Platonic archetype with its postulated
\gk{δύναμις} [power], so the intuition of \emph{the living, organic universe,}
when the principle of development is annihilated by thought, must again
go back to the intuition of \emph{the artistically formed cosmos
resting in the immediacy of beauty.} And to this cosmos corresponds that
\emph{knowledge of nature,} which immediately lets matter be shone
through by form and be absorbed into the unity of beauty with it: the
\emph{aesthetic-metaphysical} knowledge of nature.

\divrule

g. \emph{Psychology.} --- The doctrine of the soul, which Aristotle was the first to treat separately and in scientific connection, forms with him the concluding part of the Physics and, at its own close, bends back into the Metaphysics through the concept of the active \gk{Νοῦς} [mind]. The domain of psychology is, on Aristotle's conception, the organic, the living, as a whole. Everything living has soul, and the soul is what distinguishes the living from the lifeless. The soul is the body's form, principle of motion and end, and the ensouled body, whose parts are disposed toward the end and serve as instruments (\gk{ὄργανα}) for a determinate activity, is precisely thereby an organic body. The soul can therefore be defined as the first entelechy of a natural organic body, i.e.\ the entelechy comprising the activities of all the organs. Its relation to the body, and more particularly to this determinate, peculiar body, has the same necessity as the relation of form to matter in general (for this reason Aristotle rejects the notion of a transmigration of souls).

The soul, which as form is an immaterial unity and has for its nearest substrate the vital heat (\gk{πνεῦμα}), an ether-like matter, effects as the body's entelechy the combination of its material constituents into a whole; it determines the body's
\opage{177} \setcounter{footnote}{0}
peculiar constitution, governs its nourishment and the development of its organs. And the vital activity of the soul shows itself, within the domain of the organic, in a constantly advancing unfolding to higher forms, and in such a way that, through the imperceptible transitions of quantitative differences, qualitative main distinctions assert themselves, which constitute essentially different stages of the life of the soul. Of such main stages Aristotle distinguishes three: that of the plant, that of the animal and that of man. The plant has only a \emph{nutritive} soul, i.e.\ a soul whose activity is limited to the nourishment of the individual and the propagation of the species. In the animal the soul becomes a \emph{sentient} soul, with feeling of pleasure and pain and of desire, and in most animals the soul also produces the capacity for free motion in space. In man, finally, the soul becomes \emph{rational and thinking.} These forms of the life of the soul stand in relation to one another as lower and
% printed a small thin speck above the full stop after "tænkende." (defect logged in the transcription)
higher stages of development, of which the lower can exist without the higher, but the latter not without the former, since they have them as their necessary foundation and include them within themselves as their moments\footnote{In being taken up into the higher, the lower is modified. Thus in man free motion becomes determined by purpose and deliberation, and sensation and notion become accompanied by a certainty (holding-for-true, conviction).}. Between all the domains of organic nature there is thus an inner connection: one life unfolds itself through them all to ever higher perfection, until it attains to unity with the divine itself. And the organic as a whole we have in turn seen as the goal (the end) of the inorganic, which, in striving toward the living, itself shows the first intimations of the universal life of nature. Motion, which is a determination of everything natural and has form as its source, becomes for Aristotle the immortal life of nature, and in a certain sense he would call everything ensouled; yet life and ensoulment, when ascribed to the inorganic, become only analogues of the life and ensoulment of the organic, i.e.\ of life and ensoulment in the proper and strict sense.
\opage{178}As the terminal point in the series of stages of the organic, \emph{man} is the most perfect living being in bodily and psychical respects. His bodily perfection Aristotle places in his having the highest degree of vital heat, the purest and, in relation to the size of the body, the most blood, the relatively largest brain, the most harmonious shape, an upright gait, the most perfect sense organs, organs for articulate speech, and in his hand an instrument of instruments. But these bodily advantages man has only because he has the highest and most perfect form of soul: the rational soul (\gk{Νοῦς}), which comprises within itself the lower forms of soul. The peculiar function of the rational soul is essentially thinking, in particular that thinking which immediately, without error, relates itself cognitively to the highest principles. But this activity of thinking is conditioned in the individual by a development in time, in which the mind, self-active according to its essence, unfolds the endowment of the human through the relation to objects. The human soul, as \gk{εἶδος εἰδῶν} [form of forms], carries within itself potentially the whole content of knowledge: the forms of the sensible and of the intelligible; it is in a way all things (\gk{ἡ ψυχὴ τὰ ὄντα
πώς ἐστι πάντα}); but its potential knowledge attains actuality only through a progressive cognitive activity proceeding from the relation to the sensible.

The first form of knowledge is \emph{sensation}, whose nearest and immediate object is the singular, but which through the singular relates itself mediately to the universal that is present in it. For what exists for sense-consciousness is not the singular being as such, but determinate properties (determinations of form) of it, which already relate to it as something universal to the singular, and which, although in being sensed they are always intuited only in the singular and individual, nevertheless by their essence retain the character of something universal (\gk{αἰσθάνεται τὸ καθ' ἕκαστον, ἡ δὲ αἴσθησις τοῦ
καθόλου ἐστίν}), so that the thought of the universal can develop from the sensuous observation of them. The first step
\opage{179}
thereto is that, in sensation itself, through the activity of the \emph{common sense}, the sensible properties are distinguished from one another in their relative generality; the next, that with the help of \emph{imagination and memory,} as what is uniform and constant in the sense-feelings is involuntarily or voluntarily retained, a \emph{general image} is formed, matter for an \emph{experience}; and from this stage one now advances further, as a plurality of experiences is combined, to an ever more and more general and comprehensive knowledge, in which \emph{thinking} relates itself to the universal forms and \emph{causes}, and in the highest instance to the \emph{most general principles.} Aristotle has thought itself accompanied by a notion-image, a \gk{φάντασμα}, which relates to the thought as the mathematical diagram does to what is proved. Throughout the whole development the soul is at no point merely suffering, receptive; even the sensation of the object is not a passive reception of it, but a self-activity which is only called forth by it, an actualizing of what was already potentially present in the subject; and still more decidedly thinking appears as energy, activity, since, as relating itself to the immaterial, the pure form, it has not an external object, but an object that is identical with itself.

With the concept of a development in time of what is peculiar to the human soul, of human knowledge and thinking, an opposition is posited within the concept \gk{Νοῦς}, which Aristotle expresses by the distinction between \gk{νοῦς παθητικὸς} [passive mind] and \gk{νοῦς ποιητικός} [productive mind]. The relation between the two forms of \gk{Νοῦς} becomes clear from what was developed above, under the Metaphysics, about form as the principle that step by step determines the development and yet at the close of the development stands untouched by becoming, and from what was likewise stated there and in the Physics regarding the relation of the divine \gk{Νοῦς} to the series of stages of forms. Just as Aristotle, in the advance of the life of the soul, constantly sees the lower forms as taken up into the unity of the higher, so too, \opage{180}in the case of the human soul, he lets the designation for what is its peculiar essence become the designation for all the lower forms as taken up into the unity of the rational soul. \gk{Νοῦς παθητικός} designates with Aristotle precisely the sum total of those stages of development of the human soul that have not yet reached the final goal, the completed realization, and it comprises not only all the lower forms of life and consciousness that become moments in the unity of the rational soul, but also the reflective, mediate thinking that rests on sensation and experience (\gk{διάνοια} as distinct from \gk{νοῦς}: the immediate intuition of the principles). Through all the lower stages of development there runs the unity of one essence: the form that at the close of the development reveals itself as \gk{Νοῦς}, as intuitive mind; they therefore all participate in its essence and fall under its determination; but since, as stages in the development of the not yet fully realized entelechy, they still have potentiality and becoming in them, and thereby come into contact with the material, they are \gk{Νοῦς} only as that which is in development and becoming, as \gk{παθητικός}, i.e.\ as that which is realized by something actually operative. This actually operative thing, which distinguishes itself from \gk{Νοῦς} in that determination, is then \gk{Νοῦς} as \gk{ποιητικός}, as pure energy. \gk{Νοῦς ποιητικὸς} is man's \gk{εἶδος} [form]; it is the goal toward which the whole development strives, and, as the goal, the operative, actualizing principle. It relates to what is comprised under the determination \gk{νοῦς παθητικὸς} as energy does to the not yet fully realized entelechy. \gk{Νοῦς ποιητικὸς} is thus the form independent of becoming; in accordance with its relation to the series of stages of the forms of nature, as the one that concludes it, it is in its essence one with the divine \gk{Νοῦς}, which it expresses in the form of human existence; it is simple and unmixed, in uninterrupted energy a thinking of itself in its thinking of the intelligible object; it is independent of the body, a peculiar, separable principle, immortal, but not in the form of individuality, whereas \gk{Νοῦς}
% the Danish has a blank line here, mid-sentence, before the page turn; not treated as a paragraph break (see odd-spots list)
\opage{181}
\gk{παθητικός}, which is conditioned by the body, is perishable with the bodily.

As the peculiar principle in the human being, \gk{Νοῦς} becomes, according to its concept, what is determining in the \emph{practical} as in the theoretical. Already in the animal soul, the sentient soul, the presupposition for \emph{desire} is given through the feeling of pleasure and pain connected with sense-feeling; desire is a seeking or fleeing, called forth by the notion of the pleasant (here the expression for the practically good), by which a motion is effected in the bodily organs. This desire, which is the presupposition of the life of the will, develops in the human soul into \emph{will}, inasmuch as the notion through which it is determined by \gk{Νοῦς} bears the stamp of the rational, and desire is thus directed toward what is truly good. Reason, which enters into relation to desire as determining it, as setting ends and choosing means, is in this function \emph{practical reason} (\gk{νοῦς πρακτικός}). Subordination to reason is the normal relation of desire: the relation that makes its operation ethical, and through this subordination desire itself becomes a participant in reason. Just as, however, the material had the possibility of withdrawing from the dominion of form, so desire can oppose itself to the determining of reason and operate in contradiction with it. Man has the possibility of acting rationally or irrationally, well or badly; how he acts is a voluntary matter; man is therefore himself responsible for the direction of his action. ---

In the Aristotelian psychology we see the striving for unity that characterizes his whole philosophy come forth in the most energetic manner. The soul, as the body's entelechy, is in its individuality indissolubly bound to a corresponding bodily thing, which, in the unity of essence with the ideal, is subject to its dominion. The various stages of life do not become, as with Plato, separate parts of the soul, but connected members in a teleological development, whose higher stage takes up the lower into its unity as a moment and determines it
\opage{182}
in this taking-up. Thus in man \gk{Νοῦς} becomes the expression of unity for all the psychical functions that the human being contains, and Aristotle, in designating \gk{Νοῦς} as that which we, as men, properly are, points to something central in the human being, intimates a unity of the personality. But yet in this domain too, where the richer content of intuition as it were compels the decisive determinations of thought, that striving for unity is not carried through. As \gk{Νοῦς} as \gk{ποιητικός}, as something eternal, untouched by becoming, distinguishes itself from the development and its stages, the unity is broken. \gk{Νοῦς}, which was to be the full expression of unity for the human being, separates itself dualistically from the bodily and from all the psychical forms conditioned by the bodily, which it was to comprise within itself, and stands as that which is untouched by development, the impersonal, the abstract. Through this separation the development itself and the connection of the stages remain unexplained, and the concept of the unity of the self, of the personality, escapes Aristotle. The unity in the forms of development of the life of the soul becomes at any rate the unity of harmony in the diverse, and that which was to express the true, living, inner unity dissolves into an indeterminate fusion with the divine, in which no explanatory ground for individuality lies.

With the lack of a concept of the self in its inner unity and concentration there is connected, in Aristotle, the lack of a sharper conception of the \emph{concept of will}. With respect to this concept too Aristotle strives to get further than his predecessors, and seeks in particular to secure for the will a relative independence in relation to knowledge by accentuating the element of desire and the peculiarity of the practical. But yet the concept of will remains only hovering and uncertain. When the will is conceived as a \emph{desire determined by reason}, then, since according to Aristotelian presuppositions all desire can come to be only as it is called forth by form, and a determinate limit therefore cannot be drawn between mere desire and will, it must remain unclear what scope the life of the will acquires within
\opage{183}
the human, and what in the activity of the will is the properly decisive thing: desire, the basis of the will, or the influence that reason exerts on desire through the notion-image, and by which desire is first to become will. Likewise it remains unclear how man's determinateness toward the rational or the irrational, the good or the bad, can, as Aristotle would have it, become a voluntary matter, and how the freely choosing subject is to be determined. Thus only a start has been made toward determining the concept of will in its peculiarity, and while the question of the freedom of the will in the choice between good and evil is not penetratingly treated, but the voluntariness of action is in reality only postulated in order to save the concepts of imputability and responsibility, the question of the freedom of the will in its relation to the divine, from which the determination of the will is excluded, and to the operation of nature, is passed over.

\divrule

h. \emph{Practical philosophy.} --- \emph{Practical}
philosophy, which has action as its ultimate aim and as its
domain that which, as dependent on the human \emph{will},
lacks the character of necessary being that the object of metaphysics had, acquires, in accordance with this nature of its object, a lesser degree of stringency, but nevertheless falls
under the concept of philosophy in its comprehensive meaning.
It divides into the doctrine of the moral activity of the individual:
ethics, and the doctrine of the state, the moral community: politics; but since ethics as it were issues into politics, the
name of the latter can become the common designation for the whole domain.

\gk{α}) \emph{Ethics.} Aristotle's ethics stands, like Plato's,
in the closest connection with his anthropology, and through
it with metaphysics. The distinctive character of its fundamental concepts is determined by this relation, its deficiencies conditioned
by it.

When the \emph{ultimate} end of human action, i.e.\ the end striven for for \emph{its}
\opage{184}\emph{own} sake (\gk{τὸ πάντων ἀκρότατον τῶν πρακτῶν ἀγαθῶν}), is to be determined, it can, according to what was developed in the metaphysics, only be the \emph{energy} in which the
distinctive character of the human being comes fully
to actuality, and in which man must therefore find the
highest satisfaction: \emph{happiness}, the highest \gk{πρακτὸν ἀγαθόν} [practicable good]. Now, the energy in which the distinctively
human asserts itself was the energy of \emph{reason},
which in turn can be either the pure energy of reason, the energy
in \gk{νοῦς} as such, the highest and most perfect activity of mind, that by which the human mind lives a divine
life, and which is therefore the most excellent, properly
superhuman constituent of happiness; or the energy of the will
determined by reason, the properly ethical activity, the activity of reason which man \emph{alone} possesses.

By this determination of energy the \emph{essential} content of happiness (of the
highest good) is stated; for the completion of the concept of happiness, however, further
determinations are required. The energy must comprise \emph{a whole and complete life}, i.e.\ one that has not been broken off at an
imperfect stage of development. As a \emph{condition}, happiness requires a certain measure of \emph{external goods} (health, beauty,
prosperity, friends, etc.), which become aids, as it were a
\gk{χορηγία} [furnishing of means], to the realization of moral activity, and which, owing to
the individual's place within the relations of existence, cannot be entirely
dispensed with, even though the truly moral man needs the external only in small measure
and can never become entirely unhappy.
Finally, to full happiness belongs \emph{pleasure}, not as
something that is an aim for its own sake, but as something
that, as a natural completion, supervenes upon the unhindered energy that
reaches its goal, is inseparable from normal activity as health and beauty are from the normal constitution of the body, and is determined in value by the activity, which
comes to intensive self-feeling in it.

When the concept of happiness is determined as \gk{ψυχῆς ἐνέρ\opage{185}
γεια κατ' ἀρετὴν τὴν ἀρίστην καὶ τελειοτάτην ἐν βίῳ τελείῳ} [activity of the soul in accordance with the best and most complete virtue in a complete life],
then in this definition the concept of energy, which constituted the essential and original content of happiness, is set in unity
with the concept of \emph{virtue}, and the development of the latter attaches itself
immediately to the development of the concept: the highest good.
Virtue is in general excellence, the undisturbed activity of a being and its self-assertion through this activity; but the
virtue that ethics is to investigate is \emph{human}
virtue: the \emph{normal energy} of the \emph{rational} soul, or, as
Aristotle expresses it: the activity by which man performs his \emph{distinctive} work \emph{well}. Virtue,
as the expression of the distinctively human in normal
operation, then becomes, according to what was touched on above, either \emph{dianoetic} virtue, i.e.\ the normal operation of the rational life of the soul as such (in and for itself or as determining the lower
psychical activities), or \emph{ethical} virtue: the operation of the appetitive part, which is in and for
itself irrational but, as a moment in the human soul, receptive to reason,
in its lasting harmonious determinateness by rational insight. \emph{Ethical} virtue,
the specifically human virtue in comparison with the
human-divine dianoetic virtue, becomes the proper central concept of the Aristotelian
ethics (\gk{ἡ ἠθικὴ ἐστι περὶ πάθη καὶ
πράξεις} [ethical virtue is concerned with passions and actions]).

According to Aristotle, it has as its presuppositions certain
natural conditions and natural dispositions, which are already present in the child, and to which analogues are found even in
animals, but which first become virtues in the proper sense when
rational insight comes to the natural drive and forms it. The boundary of the ethical is thus not
reached so long as we are within mere natural disposition, within the natural determinateness that is something given, which is not in
our power, not, like virtue, something self-produced. We do not even
reach it through such inclinations and moods as,
like compassion, abstemiousness, modesty and the like, already
touch on the moral. The moral is first reached when ra\opage{186}tional
insight, which represents the formal principle, comes to
that matter of the moral, consciously forming it and determining the content of the will. But ethical virtue as such
is nevertheless not identical with insight: its concept comprises
both moments: \gk{φύσις} [nature] and \gk{λόγος} [reason]. The Socratic definition
of virtue as knowledge overlooks, according to Aristotle, the pathological moment, the significance of desire for action, and the necessity of moral self-control (\gk{ἐγκράτεια}) for
the congruence of insight and action, just as Socrates's conception
% errata, p. 186: printed "hans", corrected "Sokrates's"
of knowledge as immediately determining action overlooks
that the fact that knowledge can be merely potential, or that the application of
universal knowledge to the single concrete
case can fail to occur, also makes it possible to act against knowledge.
More precisely, virtue must be determined, not as something momentary
or transient, but as an abiding disposition (\gk{ἕξις})
in the will that acts according to purpose (conscious principles),
as a habit acquired through free activity: morality
becomes actual through custom; virtue is not innate,
nor yet imparted through instruction, but acquired through
practice, which lets what was originally the affair of a single free decision become fixed as determinateness of character, as
an abiding cast of mind. Morality would be complete where
free virtuous action had become as it were a natural necessity for the individual. Practice also becomes the condition for knowledge of the moral; for in order to know it one must
\emph{have} acted in a moral way (in accordance with
the moral), since knowledge of it must proceed from \gk{τὸ ὅτι} [the that]
and become real in the same way as theoretical knowledge of the universal became real by means of the sensation of
the particular. The moments in the ethical are thus \gk{φύσις},
\gk{λόγος} and \gk{ἔθος}, and Aristotle's complete definition of
ethical virtue determines it as the abiding disposition
of the will acting according to purpose, which keeps to the mean that is right for us
(determined with respect to us), as this
is determined by reason, by the man of insight (\gk{ἕξις προαιρετικὴ
\opage{187}
ἐν μεσότητι οὖσα τῇ πρὸς ἡμᾶς, ὡριςμένῃ λόγῳ καὶ ὡς ἂν ὁ
φρόνιμος ὁρίσειεν}).

In the concept \gk{μεσότης}, which Aristotle employs in the definition of ethical virtue, and which some have wished to take as
a purely external and empirical determination, there appears a
reflection of the metaphysical determination of the relation of the principles and of the Greek idealistic fundamental view. For Aristotle
conceives the relation in the irrational but
reason-receptive part of the soul (the sentient soul) thus: that
the affect and desire bound up with the feeling of pleasure and pain, and the action called forth by them, must, owing to their
being affected by the sensuous, and in general owing to their
connection with the domain of the material, of what is in itself indeterminate,
lack the secure norm and the right determinateness until they are subjected to the conscious forming influence of reason, and that therefore it is reason that first brings measure and harmonious determinateness, i.e.\ a \gk{μεσότης}, into affect and desire,
and through them into action, and thus ethicizes them. In
this relation \gk{νοῦς} works as practical insight, as \gk{φρόνησις}.
Just as, according to the Aristotelian theory of knowledge, \gk{νοῦς} immediately grasps the highest theoretical principles, the presuppositions of all knowledge, so in the domain of the ethical
\gk{φρόνησις} is to grasp, by an immediate tact, a kind of \gk{αἴσθησις} [perception], the
good and right in this entirely concrete case, and that by an
operation of \gk{νοῦς}, which as the highest energy must be what is
operative in all functions of reason. Since action is determined
by the regulation of affect and desire by insight (\gk{ὄρεξις
διανοητική, ὄρεξις ὀρθή}), and itself falls under the determination \gk{μεσότης}, it gives, insofar as it appears in
the external, to \gk{μεσότης} the shape of a middle way, and therefore
the concept \gk{μεσότης} can be carried over to domains where it does not belong
according to its proper meaning. When the concept is applied in the determination of ethical virtue, what is thus
expressed by it is the harmonious determinateness of the operation of man's
sensuous nature by form, the rational, (its
\opage{188}lawfulness: \gk{ὁ γὰρ νόμος τὸ μέσον} [for the law is the mean]) --- not the immediate natural determinateness, the immediately given natural disposition, but
the immediate, as it were aesthetic, determinateness through
insight. \gk{Μεσότης} expresses, with this more precise determination of a relation of consciousness, the same thing that Plato in some
of his later dialogues called \gk{τὸ μέτριον} and \gk{τό μέσον}, i.e.\ that
which lies midway between ``the more and the less,'' the determinations of the material, and in Aristotle too these
concepts alternate with the concept \gk{μεσότης}. From this meaning of the concept it then also follows that in Aristotle it can have
validity only for the ethical virtues, not for the dianoetic; for
in the latter there is no relation to something material that is to be
limited; therefore in their case there can be no question of excess or deficiency. ---

Aristotle does not deduce the \emph{particular} ethical virtues,
as Plato had done, but enumerates them (in such a way that in
their order some regard is had to the value of the functions and drives
to which they relate), and determines them individually with clear, sharp, characteristic traits in their relation
to the extremes. They are: \gk{ἀνδρεια} [courage], \gk{σωφροσυνη} [temperance], \gk{ἐλευθεριοτης} [liberality],
\gk{μεγαλοπρεπεια} [magnificence], \gk{μεγαλοψυχια} [magnanimity], \gk{φιλοτιμια} [ambition], \gk{πραοτης} [gentleness], \gk{ἀληθεια} [truthfulness],
\gk{εὐτραπελεια} [wittiness], \gk{φιλια} [friendship] and \gk{δικαιοσυνη} [justice]. Justice receives in
Aristotle a more definite delimitation than in Plato. Even when it is taken in the widest sense, it is indeed \gk{τῆς} \emph{\gk{ὅλης}}
\gk{ἀρετῆς χρῆσις} [the exercise of whole virtue], but: \emph{\gk{πρὸς ἄλλον}} [toward another]. It is thus indeed \gk{ἀρετὴ}
\gk{τελεία} [complete virtue], but not \gk{ἁπλῶς} [without qualification], but \gk{πρὸς ἕτερον} [in relation to another]. In the narrower sense it relates to \gk{τὸ ἴσον καὶ τὸ ἄνισον} [the equal and the unequal] and is partly
the \emph{distributive} justice (according to a geometric proportion):
\gk{τὸ ἐν ταῖς διανομαῖς δίκαιον}, partly the \emph{equalizing} justice (according to an arithmetic
proportion): \gk{τὸ ἐν τοῖς συναλλάγμασι δίκαιον}. The \emph{equitable} (\gk{τὸ ἐπιεικὲς}) supplements the just where the latter, owing to the generality of the provisions of law, becomes insufficient in the face of the concrete case.

With respect to the distinctive character of virtue in different
individuals, it is characteristic of Aristotle that, while
\opage{189}he determines perfect virtue as one in its essence,
and with moral insight (\gk{φρόνησις}) lets all ethical
virtues be given, he nevertheless, because the natural
basis of virtue differs in different persons, maintains that the moral task must necessarily become different, since different
directions of virtue must develop with a certain one-sidedness in
different persons, and different stages in virtue become the limit of
the individual's striving. Only the man, and the man as freeborn, can fully realize ethical virtue.

The \emph{dianoetic} virtues Aristotle treats briefly,
and chiefly only in order to determine, through a distinction of concepts,
what \gk{φρόνησις} is. Distinguishing between two kinds of conscious activity of reason: that which relates to the
necessary and unchangeable, and that which relates to
what we can influence at will (\gk{τό ἐπιστημονικόν} and
\gk{τὸ λογιστικὸν}), he assigns, of the dianoetic virtues, to the former \gk{ἐπιστήμη}, \gk{νοῦς} and \gk{σοφία}, to the latter \gk{φρόνησις} and \gk{τέχνη}. They
all aim at truth, either at truth in and for itself,
or at the truth that comes forth in right execution.
\gk{Ἐπιστήμη} is the mediate knowledge of the necessary; \gk{νοῦς}
the immediate knowing of the highest and most universal
principles; \gk{σοφία} unites both kinds of knowledge, and thus
acquires in Aristotle a position with respect to the two preceding virtues similar to the one justice had in Plato with respect to the
other cardinal virtues. \gk{Τέχνη} relates to production,
the poietic activity, whose aim is the external work;
\gk{φρόνησις}, finally, relates to practical activity,
whose essential aim is the energy of the action itself, the
moral in the subject, and its operation is the above-stated,
immediate expression of a moral artistic tact, which with
the sureness of genius grasps the universal in the entirely concretely
determined case. ---

In order for virtue and morality to be realized in the
individual, an insight existing in accordance with energy, i.e.\ actual, moral
insight is required, --- just as in general for Aristotle energy was
\opage{190}the condition for the realization of possibility --- but this actual
insight exists in the \emph{state}, which through upbringing, instruction and laws calls forth insight and morality in the
individual. Ethics thus leads over into politics: the doctrine of
the moral community.

\gk{β}) \emph{Politics.} --- The state is, according to Aristotle, in its moral action greater and more divine than the individual. It
is the distinctive end of human life, since man by
his nature is a \gk{ζῶον πολιτικὸν}, i.e.\ a being disposed to a communal
life. It is, by its nature as the moral whole,
the conceptual presupposition of its parts, of the individual and the family,
which with respect to time precede it as its conditions.

The aim of the state is not merely, and not even chiefly,
to provide the material conditions of life and security for
its members, but essentially to make its citizens virtuous and
thereby to realize the highest human happiness. It is a community in moral activity for the realization of a perfect and independent human life.

The attainment of the aim of the state is conditioned partly by
the constitution, partly by the insight at work in statecraft. The \emph{constitution} expresses the form of the state as a concrete organic whole; it conditions the distinctive character of the laws, and while, precisely as the form of the living whole of the state,
it must constantly be adapted to the concrete circumstances and the given
conditions, so that one and the same ideal form of state cannot be
imposed on all states, it acquires a difference in value according to
the value of those into whose hands it places power. The true and
legitimate forms of state: monarchy, aristocracy and politeia
(i.e.\ the state in which all who are fit to bear arms have equal rights, and the well-to-do middle class has the preponderance), are distinguished from the perverted ones: democracy, oligarchy and tyranny, by the fact that the former all have the common good, the latter the
advantage of one individual or of a few individuals as their aim. The perverted
forms of state, the transgressions (parekbases), are the
\opage{191}caricatures of the true ones, and they are worse in the same proportion as what they
distort is more perfect. The absolutely best form of state is an
aristocracy of those most excellent in intellectual and moral respects, or, where the case arises that there is a single individual
who towers beyond all comparison over all others, his
sole rule. --- \emph{Statecraft} has as its most excellent
task the upbringing of the citizens. The state is, by laws on
marriage and on the begetting of children, to secure for itself the right
number of new individuals of normal constitution; it is then
to bear care without interruption, through childhood and youth, for their physical and mental development, and even
among the older it is to exercise supervision over morals. But
the state is nevertheless not to abolish family life or the individual's
right to individual property. In Plato's conception of these relations Aristotle sees a danger to morality and concord, and a fetter on free helpfulness; he wants
unity in the state brought about, not by an abolition of
individual interests and of family life, but by unity of
disposition among the citizens. The unity of the state is to be precisely a
unity that binds together what is different. The virtue to
which the state is to educate the citizens is primarily civic virtue;
but while in the imperfect states this may still be
incongruent with true ethical virtue, in the
perfect state it must coincide with it: the laws of the state are there
moral laws. Since theoretical virtue is the highest, it is
not training for warlike activity but for the right
use of peace that becomes the most important task. For Aristotle too the scientific life stands as the highest goal of life. ---

Where Aristotle in his social doctrine develops the \emph{relation
between master and slave}, it is of interest to see how the concept of \emph{humanity}, the recognition of man's
right and dignity \emph{as} man, begins to come forth,
in the distinction between what is incumbent on the master in relation to
the slave as slave, i.e.\ as living property, and to the slave as
\opage{192}man. But these intimations of the humane could nevertheless
not be maintained in the face of the fundamental view that sees slavery as something necessary, and as something legitimate in accordance with the original natural difference between peoples and individuals.
The contingency of natural determinateness becomes, precisely owing to the view
of the immediate unity of the spiritual and the natural, an
unremoved barrier for spirit, and makes freedom of spirit, man's essential determination, into something particular and contingent. ---

The Aristotelian ethics has, in comparison with the
Platonic, an advantage in the sharper distinction and more exact determination of concepts, in the endeavor to secure for the life of the will
an independence in its distinctive character, and in the greater regard
that, in accordance with the whole character of Aristotle's philosophy,
is paid to the natural moment in the human being.
But it nevertheless preserves the general character of Greek ethics,
and is unable to overstep its bounds, and the reason
for this lies in what was shown above in connection with the anthropology: that the concept of the unity of personality concentrated in its inner infinity
was lacking, and that the concept of will was only
uncertainly and vaguely determined. Ethical action, then, can
not be conceived as proceeding from the ethical decision of a personality that is in itself infinitely transparent, clarified in exhaustive reflection, fully conscious of \emph{ethical obligation},
but must be conceived as coming to be as it were suddenly, in an artistic-genial way, through a concretion of the universal law of reason, through
the immediate operation of \gk{Νοῦς}, which withdraws from conscious reflection and is only the object of an \gk{αἴσθησις}. The
moral life therefore becomes a series of such moments of immediacy: in each of them the moral individuality stands, in its immediate determinateness by the universal,
as an individual work of art, closed in itself; but continuity for consciousness is lacking; ethical development is transformed into the successive revelation, governed by beauty and appearing harmoniously in a series of
discrete moments, of undisturbed human nature. Ethical
\opage{193}action becomes an occurrence, an event, and the ethical
good must, in principle, be drawn toward the metaphysical.
The character of the Aristotelian ethics thus becomes, like
that of Greek ethics in general, \emph{aesthetic-metaphysical}.

\divrule

What gives the Aristotelian philosophy its historical
significance, and what its deficiencies are, has already become
clear through the exposition of its main parts. We have
everywhere seen it strive beyond the bounds of Greek thought,
but everywhere be forced back within them. What was
the fundamental deficiency of Greek thought: the lack of ability
to reconcile the dualism that hid beneath the intuited harmony of the work of art, was also the pervasive deficiency in
the Aristotelian philosophy, and since through the latter the dualism
is brought precisely and clearly before consciousness, the task and character of the
subsequent philosophy are thereby conditioned.

\divrule

\begin{center}
\textbf{6.\enspace The Peripatetic School.}
\end{center}

Aristotle's school has, like Plato's, only a subordinate significance for the further development of philosophy. Metaphysical investigations recede before empirical
knowledge and before a popular ethics. Only at certain points
are consequences drawn from the Aristotelian doctrines that
touch on essential problems and mainly go in the
naturalistic direction that post-Aristotelian philosophy
follows in its first form. The later Peripatetics earned
merit by editing and interpreting the Aristotelian writings.

In the first generations after Aristotle the school is represented
essentially by Theophrastus and Strato. \emph{Theophrastus} of Lesbos (c. 373---288) was Aristotle's successor as
head of the school; he is significant especially for his natural-historical writings on plants and on stones. In metaphysics he departs from Aristotle in his conception of
\opage{194}motion, which he seeks to bring into a closer relation
to energy and to make a determination of the immaterial as well. From the not entirely clear reports on this
point it seems as though the difficulty pointed out above in Aristotle
of conceiving form as principle of motion
and yet as untouched by becoming led Theophrastus to this
approximation of the concepts, which also brought him to maintain that at every point in the development of the possible
an energy is realized. In psychology there is in Theophrastus an endeavor, although \gk{νοῦς} is in a certain sense conceived
as \gk{χωριστός} [separable], to maintain it as \gk{σύμφυτος} [congenital] with the human, though without the relation really becoming clear. To the activity of thinking he applies the concept of motion, but not in
the sense of spatial motion. In ethics greater
weight is laid on the external aids of virtue than in Aristotle.

\emph{Strato} of Lampsacus was Theophrastus's successor
as head of the school. He cancels the doubling that occurs in
Aristotle of form as divine \gk{Νοῦς}
\gk{χωριστὸς} and as indwelling form (\gk{φύσις}), and carries through a
standpoint of immanence, from which all activity and all life in
the world is conceived as deriving from the force originally indwelling in matter and working with unconscious necessity: nature,
which is the divine. Correspondingly he denies
in man a \gk{νοῦς χωριστός}, maintains the unity of the life of the soul
as rational and sensuous, gives reason, the universal
fundamental power of the soul, whose particular expressions the activities of sense are,
a bodily organ, and lets the soul as a whole be dissolved with
the body. The activities of the soul are conceived as motions; \emph{all}
living beings are posited as partaking of reason.

Besides Theophrastus and Strato, among the Peripatetics of the first generations may be named \emph{Eudemus}, \emph{Aristoxenus} and \emph{Dicaearchus}, of whom the latter two essentially stand on the same standpoint
as Strato in their conception of the psychical.

\divrule

\opage{195} \setcounter{footnote}{0}

\begin{center}
{\Large\textbf{Chapter Four.}}
\end{center}

\divrule

\begin{center}
{\Large The Third Period of Greek Philosophy: The Dissolution\\
of the Doctrine of Ideas.}
\end{center}

\divrule

\begin{center}
\textbf{1.\enspace Survey of the Third Period}
\end{center}

\markboth{The Dissolution of the Doctrine of Ideas}{The Dissolution of the Doctrine of Ideas}
The peculiarity of the third period was earlier stated thus: that, with a
clearer or dimmer consciousness of dualism, it strives after a unity that
sacrifices one of the opposites and therefore cannot lead to their
reconciliation.

In this striving after unity in principle, post-Aristotelian philosophy
develops from a materialism, through a dualism expressed in principle as a
short-lived transitional link, to a spiritualistic doctrine of unity, which
strives powerlessly beyond the limits of the Greek spirit and ends in
fantastic scholasticism.

In the general relation between the fundamental concepts of Greek thought,
and in the direction toward the real that philosophy had taken in Aristotle,
lay the reason why that striving after unity had to begin from
\emph{materialism}. Materialism is expressed dogmatically, in essentially
different forms, in Stoicism and Epicureanism, and it shines through beneath
Skepticism. This materialism, in whose dogmatic representatives dualism had
proved unovercome: in the Stoics already in the very determination of the
principle, in the Epicureans in anthropology, where the opposition that could
be kept hidden in the superficially treated general part of the system broke
forth through the more concrete determinations, is replaced at the beginning
of the Christian era by \emph{the dualism of principle}, which through the
opposition between the transcendent spiritual God and matter\footnote{When the
evil world-soul is posited as the opposite of the divine, it is as it were the
personification of the hindering and disturbing power lying in matter, by
whose separation matter becomes something indifferent. But in other
conceptions matter itself becomes the properly negative. Its determination as
a \gk{μὴ ὄν} [non-being] then acquires a significance that goes beyond the
Platonic: matter becomes not merely the negative, but the negating.} prepares
\opage{196}the spiritualistic doctrine of unity. The dualism expressed in
principle, whose most essential representatives were the Neopythagoreans and
Philo, was, precisely because it was dualism, of only lesser duration. It had
in turn to give way to a monistic philosophy, which now, after the other
possibilities had been tried, had to be of a \emph{spiritualistic} kind, and
which was completed in Neoplatonism, in which Greek thought, influenced by the
spirit of the Orient, strove to derive all that is, and matter itself, from
the purely spiritual Godhead, transcendent over thinking and being, to whose
undifferentiated unity man sought to return in order to be deified in ecstatic
surrender. In the struggle against Christianity, which carries out more
consistently what Neoplatonism strove for, Neoplatonism concludes Greek
philosophy, now fully unfolded according to its whole disposition.

\divrule

\begin{center}
\textbf{2.\enspace The General Character of the Third Period.}
\end{center}

In the second period, the culminating age of Greek philosophy, the essential
interest of philosophy had been the knowledge of the ideal objectivity that is
true and actual in and for itself. Theory was posited as the highest; the
possibility of knowledge was not doubted; no domain of existence was excluded
from thought, but philosophy unfolded into a system embracing all reality. The
ethical received the stamp of ideality, and the thought of a universally human
morality worked its way forward; but there was no detachment of the individual
from nature and the state. The Greek spirit has here reached its richest and
most complete unfolding, and it strove with energy after the solution of the
task: to bring the oppositions of existence \opage{197}under a unity of
thought. But as this striving failed and dualism proved insuperable, the
process of dissolution in Greek philosophy begins, and the third period is the
span of time in which this process is consummated.

The endeavor, after the consciousness of dualism had been awakened, to restore
unity by bringing both opposites under the determination of one of them, then
becomes the \emph{fundamental tendency} in this period, the one that gives it
its peculiar stamp. But beside this striving, and in inner connection with it,
other characteristic peculiarities come forward, by which the third period as
a whole distinguishes itself from the preceding one.

Whereas the essential interest of the second period had been the very
knowledge of the true, objective existent in the sphere of nature and of
spirit, the essential interest of the third period becomes the effect of that
knowledge on the subject, so that knowledge of the objective definitively
becomes a means of procuring the subject's virtue and happiness. The center of
gravity in the systems is thus shifted from the object to the subject, and
more precisely to the subject conceived from the standpoint of the practical.

But as man is essentially conceived from the standpoint of the practical, and
the ethical problems, man's moral task and the conditions for its realization,
thus come forward, an essential difference from the earlier philosophy again
shows itself. What Aristotle had intimated, the breaking through of the
barriers that prevent the human from coming forward as such, is carried
further here. Man is now seen, in the infinity of his self-consciousness,
essentially \emph{as} man; the concept of humanity, in opposition to that of
citizenship and nationality, begins to become clear. The boundaries of human
society reach beyond those of the state to the boundaries of the cosmos. But
while the ethical acquires a stamp of universality, it acquires at the same
time a stamp of the abstract and negative. In the infinity of its
consciousness the subject knows itself independent of everything, but by the
same token everything external, including that in which \opage{198}the Greek
formerly found the substantial content of his life, family life and the life of
the state, becomes alien and indifferent to it. The subject becomes as it were
its own world: apathy and ataraxia in relation to the external world become
its task. It does not set itself the goal of realizing the moral ideas in the
world, but that of winning through philosophy autarky, independence of the
world. Subjectivity, to which everything returns, becomes the abstract
subjectivity that asserts itself, in its separation from nature and human
society, as the unconditionally highest, and in this autarky knows itself
perfectly free and satisfied.

As a phenomenon that accompanies the receding of the theoretical in comparison
with the practical, we see that all positive, dogmatic systems in the third
period attach themselves in a dependent manner to older philosophical forms,
% errata, p. 198: printed "knytte", corrected "knytte sig"
especially in the domain of physics and metaphysics, while in the domain of
the ethical there is greater independence. The productive power of thought
proves to be in the process of being exhausted, and its mistrust of itself
comes forward indirectly in the investigations concerning the criterion, while
skepsis directly declares man's lack of capacity for knowledge of the truth.

A last peculiarity of this period, which is connected partly with the
exhaustion of theory, partly with the preponderance of the practical, is the
greater influence of the positively religious on thought. Plato and Aristotle
had stood freely over against popular belief and had reshaped it according to
their philosophical fundamental thoughts, in which their minds found
undisturbed rest; post-Aristotelian philosophy is in many ways dependent on
the religious form of notion, either as defending it or as combating it: it
occupies itself to a great extent with the problems of natural theology and
ascribes essential significance to them for the moral; and in its later forms
it seeks in the religious revelations a source of truth higher than human
knowledge. Among \opage{199}the Neoplatonists philosophy finally becomes a
scholastic theology.

For this great change in the whole physiognomy of Greek philosophy the reason
must, in the nature of the case, be sought first of all in the development of
the earlier philosophy. In that development the greatest significance then
belongs to the dualism that comes forward in both Plato and Aristotle, by
which mind, as pure form existing for itself, distinguishes itself from the
sensible, and in this distinction becomes the essential and absolutely real.
By this dualism, as will be shown below, the peculiarities of the period can be
explained. But beside the influence from what we may call the main current of
the earlier philosophy, account must also be taken of the personal influence
that the founders of the first post-Aristotelian systems received from the
imperfect Socratics. Stoicism and Epicureanism stand in immediate connection
with Cynicism and Cyrenaicism, in which the practical was the essential
interest. But even if this connection can explain the direction toward the
practical in Zeno and Epicurus, and even if their schools' continuation of
this direction can be explained by their dependent attachment to the founders
of the schools, the reason why the \emph{whole} philosophy of the
post-Aristotelian age takes this direction is nevertheless to be sought in
something more general, and the connection with this we shall strive to
demonstrate in what follows.

By the external political conditions under which philosophy developed in the
third period, the change in its character was also bound to be favored. For
as the individual was no longer borne up by a vigorous, free common spirit,
and no longer had in a free political life a stage for its activity, the
classical ground was as it were drawn away from under it; the development of a
free and many-sided theoretical world-view, corresponding to that free and
full social life, was made impossible, and the subject was as it were thrown
back into itself, and confined to itself as the last firm point of support.
The classical Greek spirit be\opage{200}gins its period of dissolution from
the fourth century, and on its dissolution as a thinking spirit the mingling
of the Greek with foreign elements, Roman and Oriental, then has its effect.
In particular the contact with the Orient becomes significant for the last
phases of Greek philosophy, in which an attempt is made to carry through an
abstractly spiritualistic view of life.

When we have stressed above the dualism coming forward in the early philosophy
as the ground of explanation for the peculiar stamp of the third period, this
is to be substantiated more closely as follows. The most immediate consequence
of the consciousness of dualism becomes the striving, emphasized above, after
a unity in principle, which characterizes the period taken as a whole, and
which we have already come to know in Aristotle's school. But in this striving
it was the presupposition, even where mind was determined materialistically,
that the mental, which in both Plato and Aristotle distinguished itself, in
its being-for-itself, from the material as the only perfect reality, was the
highest. The next consequence of dualism was therefore that the rational
being, man, in his being-for-itself, drew the main interest to himself. And as
now, after the fundamental thought of the Greek spirit has been found and
carried through, the theoretical power of the spirit is exhausted, so that the
later systems must attach themselves dependently to older, theoretically
productive ones, the interest of the thinkers turns preferably to the
practical, and the human is conceived essentially from the side of the
practical. And as, further, the dualism present in concealment even in the
monistic systems --- which also has its effect on the investigations
concerning a criterion of truth --- makes the conception of the human
abstract, just as the thought separated from matter became abstract, man's
practical goal is conceived essentially as a resting in the infinity of an
abstract self-consciousness, whether this is now determined from the
standpoint of the single individual, or from that of the immanent universal
reason, or from that of the abstractly transcendent. The widening of the
ethical \opage{201}standpoint, the negative character of ethics, and the more
determinate emergence of subjectivity in the posing of the ethical problems
are a simple consequence of this, just as the greater significance that the
positively religious forces for itself is a consequence partly of the
exhaustion of theory, partly of the heightened interest that comes to rest on
subjectivity.

\divrule

\begin{center}
\textbf{3.\enspace Stoicism.}
\end{center}

a.\enspace \emph{History of the School.} --- \emph{Zeno} of Citium in Cyprus
(c. 350---c. 258) was the founder of the school which, after the place where
he taught, the \gk{στόα ποικίλη} [Painted Stoa], was called the Stoic. He
first heard the Cynic Crates, in whom he thought he had found the man who in
his own time most resembled Socrates; later he heard the Megarian Stilpo and
the Platonists Xenocrates and Polemo. About 310 he came forward as a teacher
and worked as such until his death, honored by the Athenian people for his
\gk{ἀρετὴ} [virtue] and \gk{σωφροσύνη} [temperance]. Of his writings, of which
the earlier stood closer to Cynicism, none are preserved. Among his disciples
the most significant were \emph{Aristo} of Chios, \emph{Herillus} of
Carthage, and \emph{Cleanthes}, who became head of the school after Zeno. Of
Cleanthes a hymn to Zeus is preserved. A greater name than these was won by
\emph{Chrysippus} (282---209), who led the school after Cleanthes. He first
gave the Stoic doctrine its systematic working-out and dialectical grounding,
so that he became as it were the second founder of the school. He displayed an
enormous activity as an author, but his numerous writings --- according to
ancient reports over 700 --- are lost. After Chrysippus, \emph{Zeno} of Tarsus
and \emph{Diogenes} of Babylon directed the school. With \emph{Panaetius} of
Rhodes (c. 180---111) eclecticism begins to make its way into the Stoa. His
disciple \emph{Posidonius}, who taught at Rhodes, shared this eclectic
tendency with his teacher. Among the Romans Stoic philosophy found
considerable diffusion in a popular and eclectic form; its most significant
representatives there are L. \emph{Annaeus Seneca} (3---65 p. Chr.), C.
\opage{202}\emph{Musonius Rufus}, the slave \emph{Epictetus} (from Phrygia)
and the Emperor \emph{Marcus Aurelius}. Of Seneca, Epictetus and Marcus
Aurelius we have writings preserved, of the last two in the Greek language.
For several centuries after its rise Stoic philosophy becomes the proper chief
leader of philosophical development, and the school maintains itself until
Neoplatonism absorbs all the earlier systems.

\divrule

b.\enspace \emph{The Stoics' Division of Philosophy.} --- The Stoics as a rule
divide philosophy into logic, physics and ethics, logic with them, as the
science of \gk{λόγοι}, i.e.\ of thought and speech, comprising theory of
knowledge and logic as well as rhetoric; physics, the doctrine of principles,
thus what in Aristotle had its place in metaphysics. In fact the two
first-named sciences are subordinated to ethics (since the whole of philosophy
becomes an \gk{ἄσκησις τῆς ἀρετῆς} [exercise of virtue], and since physics too
becomes a means to the knowledge of good and evil), although most Stoics, when
they pronounce on the mutual relation of the sciences, assign precedence to
physics because it treats of the grounds of things and of the divine. Logic
becomes the instrument for the other parts of philosophy.

\divrule

c.\enspace \emph{Logic.} --- As the Stoics divide logic into dialectic and
rhetoric, dialectic comes to comprise theory of knowledge and logic (in the
narrower sense), as well as grammar, with which the Stoics busied themselves a
great deal. The Stoic \emph{logic} rests on the Aristotelian, which it
supplements especially in the doctrine of hypothetical and disjunctive
syllogisms. The Aristotelian \emph{doctrine of categories} is transformed by
the Stoics, the 10 (or 8) Aristotelian categories being reduced to 4:
\gk{τὸ ὑποκειμενον} (matter), \gk{τὸ ποιὸν} (the essential form),
\gk{τὸ πὼς ἔχον} (the accidental property, which is predicated of the subject
without relation to another), and \gk{τὸ πρός τί πως ἔχον} (the accidental
\opage{203}determinations that express a relation). Of these, each succeeding
one is conceived as a determination of the preceding, so that \gk{τὸ ποιὸν}
fully expressed would read \gk{ὑποκείμενον ποιὸν}; \gk{τὸ πὼς ἔχον}:
\gk{ὑποκείμενον ποιὸν πὼς ἔχον}; \gk{τὸ
πρός τί πως ἔχον}: \gk{ὑποκειμενυν ποιὸν πρός τί πως ἔχον}.
The Stoics call these concepts \gk{τὰ γενικώτατα} (highest generic concepts
under \gk{τὸ τι} [something] or \gk{τὸ ὄν} [being]).

In the Stoics' \emph{theory of knowledge} the greatest significance attaches
to their doctrine of the significance of sensation for knowledge and of the
criterion of truth.

All our knowledge proceeds from sensation, inasmuch as the soul, which is
originally like an unwritten sheet, is filled with notions through the action
of the external objects, which as it were leave an impression in the soul, so
that the notion becomes a \gk{πάθος} [affection] of the soul (a
\gk{τύπωσις ἐν ψυχῇ} [impression in the soul] or an \gk{ἑτεροίωσις τῆς
ψυχῆς} [alteration of the soul]). Sensation leaves behind a memory; through
the combination of like memories an experience is formed (\gk{ἐμπειρία}
= \gk{τὸ τῶν ὁμοειδῶν μνημῶν πλῆθος} [the multitude of memories of like
kind]). From sensations, which also relate to our inner states, the general
concepts come to be, including those that lead beyond the immediately
observable and are reached by inferences from experiences, inferences that
rest partly on likeness, partly on analogy, partly on combination. Concepts
are partly those formed spontaneously, as by a working of nature,
\gk{κοιναὶ ἔννοιαι} [common notions] and \gk{προλήψεις} [preconceptions],
partly the concepts formed consciously and methodically (\gk{ἔννοιαι},
\gk{καταλήψεις} [apprehensions]). As general, concepts express only a
subjective abstraction, since only the individual has real existence. Through
the artful formation of concepts, for which dialectic gives the rules, science
arises, which is a system of certain and irrefutable concepts that stand in
harmony with those naturally formed.

The Stoics' theory of knowledge stands in close connection with their
materialistic and pantheistic doctrine of principles, which will be developed
under physics.
\opage{204}The question of the \emph{criterion}, i.e.\ of the distinguishing
mark of the truth of our notions, the Stoics answer thus: that notion is true
which by its clarity (\gk{ἐνάργεια}) shows itself as it were to seize (take
in) the object, the real (\gk{τὸ υπάρχον}). Such a notion is called
% errata, p. 204: printed "(opfatte) Objectet", corrected "(optage) Objectet, det Virkelige"
\gk{φαντασία καταληπτική} [apprehending presentation]. By what is clear and
evident in it, it is supposed as it were involuntarily to wring from us our
assent (\gk{συγκατάθεσις}). The mere images of imagination, which do not
derive from an actual object, are not supposed to have this clarity. Since,
however, it could not be denied that in certain cases the merely subjective
notion can have the same sensuous force as the one deriving from a real
object, later Stoics, in the controversy with Academics and Skeptics, were
compelled to add the determination that the true notion was only that
\gk{φαντασία καταληπτική} against which there was no counter-proof
(\gk{μηδὲν ἔνστημα ἔχουσα}). It is chiefly in the interest of the practical
that the Stoics fought for this doctrine, and they held fast that only the
possibility of a knowledge of the truth could secure for action a clear goal,
and for moral conviction a firm foundation.

\divrule

d.\enspace \emph{Physics.} --- The physics of the Stoics attaches itself
closely to Heraclitus's world-view, yet in such a way that acquaintance
with the concepts of Aristotelian physics shows itself everywhere. The
fundamental endeavor aims at leading everything back to one principle,
which is at once the element of things and their forming force. An
immaterial form-principle, as in Aristotle, the Stoics do not recognize:
everything that has reality (being) must in their conception be something
corporeal, because only such a thing can act and be acted upon (suffer),
which are the distinguishing marks of reality. What Aristotle determined
as immaterial form, the Stoics determine as an active corporeal thing,
which permeates the passive substrate and gives it its properties.

When it is said that the Stoics set up two principles (\gk{ἀρχαὶ}):
\opage{205}the active and the passive (\gk{τό ποιοῦν} and \gk{τὸ πάσχον}),
these are not to be conceived as originally separated principles; they
are posited not merely as inseparable, but as expressions of the same
fundamental unity. This fundamental unity (\gk{στοιχεῖόν}), which is at
once the primal force not yet unfolded into determinate forms and the
material germ of things, is the divine primal fire (\gk{πῦρ τεχνικόν}
[artistic fire]). From this everything in existence develops according
to an inner law, with the same necessity as the plant develops from the
seed. The primal fire transforms itself first into air (airy vapor), then
into water; of the water, then, one part becomes earth, one part remains
as water, one part evaporates into air, from which in turn the elemental
fire is kindled. From the mixture of the elements the world is formed.
The primal unity the Stoics conceive as transformed \gk{τὴν πᾶσαν οὐσίαν}
[in its whole substance] through air into water, but it is nevertheless
to remain behind in the moist, just as the \gk{σπέρμα} [seed] is enclosed
by the \gk{γονὴ} [generative fluid], as \gk{σπερματικὸς λόγος} [seminal
reason], and thereby to make matter fit for further forming. Only with
the coming-to-be of the elements does an opposition emerge between the
active and the passive principle: the world's soul and body, between the
Godhead (\gk{θεὸς}) and matter (\gk{ὕλη}). The active principle is the
primal fire as such, which remains as something inwardly active and
forming; the passive is the substrate for the forming activity, which is
stamped into the elements, among which fire and air stand nearer to the
active principle, water and earth to the passive.

The Godhead (\gk{θεὸς}), the immanent creative force, the forming
principle, which in the Stoics' conception cannot be thought of as
immaterial, is designated at once as \gk{πνεῦμα} [breath], ether,
\gk{πῦρ τεχνικὸν} (in contrast to the elemental fire), \gk{πῦρ διῆκον}
[pervading fire], and as soul, reason (\gk{νοῦς} and \gk{λόγος}), as
\gk{εἱμαρμένη} [fate] and \gk{πρόνοια} [providence] and as the universal
world-law. The more spiritual expressions too designate something
corporeal, the warm fluid that permeates everything. Heat is the ground
of life, but the vital heat of the universe must be mind and reason,
since the world is purposively arranged; and since there are already in
the parts of the world, notably in man, feeling and
rea\opage{206}son, the universe, which is more perfect than its parts,
must have an infinite reason for its forming force.
% errata, p. 206: printed "maa ogsaa", corrected "maa"
The excellence of this reason is grounded in the purity and mobility of
the fire-substance.

Since the world is the product of the unity of primal matter and primal
force, it is itself \emph{one}, but in its unity it comprises the richest
multiplicity of forms. It is perfect through its beauty and thoroughgoing
purposiveness, which point back to a rational originator. It is bounded,
spherical in form, with concentric circles. Around it spreads the
infinite, empty space.

All that happens in unfolded existence happens according to a necessary
and unalterable connection of causes and effects, according to the
nature and law of the whole: \gk{εἱμαρμένη}, which is the expression of
the primal being itself, the world-reason, and therefore equal to
\gk{πρόνοια} and \gk{λόγος}. Providence designates the primal being as
the interior of the world; it permeates the universe as the warm breath
spread everywhere (\gk{πνεῦμα}), as the forming, rational soul-principle,
as the \gk{λόγος σπερματικός} that in its unity encloses the
\gk{λόγοι σπερματικοί} of individual things, (= \gk{ποιότητες}
[qualities]), their rational germinal forms, their productive forces
(\gk{δυνάμεις γόνιμοι}). The inner form of things is something abiding
amid the change of the material parts, and by the active principle the
identity of the universe is preserved. The existence of the universe is
subject to a periodicity. The universe that has proceeded from the primal
fire according to a law-determined order returns into it, when a
world-period is concluded, by the same steps by which the unfolding took
place. There then takes place a world-conflagration (\gk{ἐκπύρωσις}),
and the divine primal fire then exists alone, but only in order to bring
forth anew, according to the eternal world-law, a new universe, in which
the forms and fates of the previous one repeat themselves. ---

In the Stoics' doctrine of principles, despite their striving after a
unity of principle, dualism nevertheless breaks through at once. When
the primal ground is separated into an active and a passive principle,
\opage{207}this becomes, despite all striving after unity, more than a
conceptual distinction: involuntarily the passive moment separates itself
as \gk{οὐσία} [substance] (\gk{ὑποκειμενον} [substrate]) from the active,
which permeates and determines it, and precisely in this relation,
despite the inseparability, stands over against it as an other. The
possibility of the real unity of the active and the passive is not
comprehended; neither of them is derived from the other. And when the
primal fire is conceived as developing itself into the determinate forms
of existence, in which it remains as active Logos, the same thing repeats
itself here; the coming-to-be of the passive through the active is not
explained, and there remains a duality whose consequences can be traced
in the Stoics' theodicy and in their ethics. That the opposites are
constantly posited as inseparably bound together does not prevent the
moments of the principle of unity from in reality becoming two principles
dualistically separated in their essence. ---

Against the Stoics' doctrine of the world's perfection as the effect of
the rational primal principle that embraces all being, physical and moral
evil, which seem to point to the activity of an opposite principle, stood
as an objection. The Stoics therefore occupied themselves very eagerly
with the problem of \emph{theodicy}. Physical evil they did not recognize
as a real evil, but only as a naturally necessary and unavoidable
consequence of purposive arrangements, and they pointed to the useful
effects of many physical evils. Moral evil caused them greater
difficulties. They could not in the last instance shift it onto man,
although they refer to man's intention and free will; but they had partly
to appeal to the fact that not even the Godhead was able to keep human
nature completely free from fault, since the individual in existence is
necessarily limited and imperfect; --- partly they had to stress that
evil, as the opposite of the good, was necessary for the good itself,
which can exist and come to consciousness and activity only in this
opposition; --- partly they had to urge that evil is guided by the
Godhead in its effects toward the good, so \opage{208}that its disharmony
is ordered into the harmony of the universe. Yet they sometimes make
concessions with regard to particulars in existence that cannot be
reconciled with their fundamental view.

In the treatment of this problem the closer attachment to the popular
religion's conception of the divine comes forward, which on the whole
becomes characteristic of the Stoics. This attachment had a motive partly
in respect for what is generally assumed according to
\gk{φυσικαὶ ἔννοιαι} [natural notions], partly in regard for the moral
usefulness of popular religion for the multitude. The principle of the
Stoic system had consistently to lead to a strict monotheism, and the
name of Godhead in the highest sense they also ascribed only to the one
primal being; but in a derived sense they used the name Godhead of
everything in which the power of the divine world-soul expressed itself
in a special way. Thus the stars, the elements, and everything that by
its usefulness particularly shows the indwelling divine power (e.g.\
corn and wine), as also the benefactors of the human race, the heroes and
the daemons, were designated by the name of god. But all these deities
are only forms of appearance of the highest god, Zeus, from whom they
have proceeded, and to whom they return in the world-conflagration. In
order to mediate between their philosophical views and the notions of
popular religion the Stoics employed an \emph{allegorical interpretation}
of the myths, and they distinguished between what is said in these
\gk{κατὰ δόξαν} [according to opinion] and \gk{κατ' ἀλήθειαν} [according
to truth]. But although they sought in this way to rationalize popular
religion, they did not avoid its having an obscuring effect on their
philosophical view. In the validity they ascribe, e.g.\ to divination,
there betrays itself, despite all striving to maintain a rational and
natural connection between the sign and what is foretold, the influence
of the irrational and superstitious in popular belief.

The \emph{psychology} of the Stoics forms essentially a part of their
physics, although some portions of it used to be treated under ethics.
The human soul is conceived as a part \opage{209}(\gk{μέρος, ἀπόσπασμα})
or an emanation (\gk{ἀπόῤῥοια}) of the Godhead, with which it is
constantly in reciprocal action. It is, like the Godhead, an ethereal
substance, a warm breath (\gk{πνεῦμα}), which is spread through the body
as the Godhead is through the world, and holds the body together. As
\gk{πνεῦμα} it is corporeal: something incorporeal could not suffer with
the body or act upon the body. The human soul (or, according to some
Stoics, only the soul of the wise man) survives the body, but returns at
the end of the world-period to the Godhead. Its \emph{parts} are the five
senses, the faculty of speech, the generative power and the governing
faculty: reason (\gk{τὸ ἡγεμονικόν}, \gk{τὸ διανοητικόν},
\gk{τὸ λογιστικόν}). The ruling faculty, whose seat most Stoics place in
the breast because speech, the expression of thought, comes from there,
is the source of our notions and thoughts, our desires and resolutions.
In the last instance it is the fundamental force from which all the
powers of the soul are derived, and it is the expression of the
distinctively human. By it we know things and are able to lead a
rational and moral life. The human soul cannot, in virtue of its relation
to the world-soul, withdraw itself from the necessity of the world. What
are regarded as arbitrary and accidental effects proceed with necessity
from hidden causes. Yet man's own activity does not become dispensable;
it too is taken up as an intermediate cause into the nexus of fate.
Imputability is maintained, inasmuch as fate is posited only as the
occasioning cause, man's own will and moral constitution as the
definitive cause. Freedom, however, is so determined that it becomes an
inner, necessary self-determination. The Stoics strongly emphasized the
community among men, because they are of one race and partake of the
same rational nature. Man, as a \gk{ζῷον λογικὸν} [rational animal], is
at the same time a \gk{ζῷον κοινωνικὸν} [social animal] and
\gk{φιλάλληλον} [loving one another].

\divrule

e.\enspace \emph{Ethics.} In the Stoic ethics the influence of Cynicism
is strongly prominent; yet with the Stoics, since \opage{210}they demand
as the condition of true morality a comprehensive scientific knowledge of
the world and of the self, the moral subject gains a wider horizon and a
truer liberation.

The highest principle of Stoic ethics is conditioned by the fundamental
view of existence. If existence forms a whole animated and governed by
the divine world-law, by which the individual existences are enclosed as
parts and members that are subordinated to the universal law, then the
ethical task for man can only be consciously and voluntarily to fulfill
this subordination under the law of reason, and to express this
subordination in his action. The task becomes
\gk{ὁμολογουμένως τῇ φύσει ζῆν} [to live in agreement with nature], and
its meaning remains essentially the same whether by \gk{φύσις} one thinks
of universal, reason-determined nature, or of human nature, which, as
rational, is of one essence with the divine principle of existence. The
fulfillment of this moral task also becomes one with the satisfaction of
the fundamental drive in us: the drive of self-preservation; for its aim
must essentially be what is essential in man: reason, which is a part of
the universal world-reason. To live according to the law of nature is
one with living according to that of reason.

The life that is lived in agreement with the law of reason in existence
and in ourselves is the virtuous life, a living \gk{κατ'
ἀρετήν}. Virtue is the highest good (\gk{τέλειον ἀγαθόν}), and that in
such a way that its concept covers the concept \gk{ἀγαθόν}. Virtue as
such is happiness, not merely an element, albeit the most important, in
happiness. \emph{Only} virtue is a good in the proper sense of the word;
only its opposite, vice, is an evil. Everything that does not fall under
one of these two concepts is indifferent with regard to happiness or
unhappiness (an \gk{ἀδιάφορον}). Such things as health, wealth, noble
birth and the like can just as well be misused as used in the right way;
sickness, poverty and lowliness can just as well become a spur to moral
strength as a hindrance. Yet the Stoics, because the drive of
self-preservation and \opage{211}the practical relation to external
things demanded it, made within the wide concept of the indifferent a
distinction between what was to be preferred (\gk{τὸ προηγμένον}), what
was to be rejected (\gk{τὸ ἀποπροηγμένον}), and what was unconditionally
indifferent. A \gk{προηγμένον}, something valuable, are certain mental
and bodily qualities (e.g.\ good endowments, health, strength), and the
so-called external goods, because they have a relative value
(\gk{ἀξία}) as aids (as \gk{εὔχρηστα}), and because they are more in
agreement with our nature (\gk{κατὰ φύσιν}), objects of natural drives.
An \gk{ἀποπροηγμένον} is the opposite of this, which has an
\gk{ἀπαξία} [disvalue] as rather disturbing and hindering (as
\gk{δύςχρηστον}) and as being \gk{παρὰ φύσιν} [contrary to nature].
Absolutely indifferent is that which exerts no influence whatever on our
mental and our bodily condition.

With regard to \emph{pleasure} the Stoics denied that it could be a good,
let alone the highest good. As a passive state of the soul
(\gk{πάθος}), pleasure is a disturbance of the equilibrium of rational
consciousness. Cleanthes therefore determines pleasure in general as
something not in accordance with nature; and when other Stoics make a
distinction and allow value to that pleasure which proceeds (as an
\gk{ἐπιγέννημα} [by-product]) from moral action in harmony with nature,
and which they distinguish as \gk{χαρά} [joy], \gk{εὐφροσύνη}
[cheerfulness] and \gk{ἀλυπία} [freedom from pain] from ordinary
\gk{ἡδονή}, they would nevertheless not make it the \gk{τέλος} [end] or a
part of the highest good.

Virtue the Stoics determine in the Socratic manner as an \gk{ἐπιστήμη}
[science]. But knowledge becomes an aim for them only as the
\emph{immediate} source of action: not theory, cognition as such, is the
goal, but the unity of insight and moral action, the obedient following
of the known law of nature, the strength of will that rests on right
insight. A life merely for theory the Stoics regarded only as a finer
kind of enjoyment. Virtue is in its essence a unity: it can appear in
several forms, but in each single virtuous activity all virtues are
present: he who has one virtue has them all; therefore \opage{212}all
virtues are equal. As the main forms of virtue are determined
\gk{φρόνησις, ἀνδρία, σωφροσύνη} [prudence, courage, temperance] and
\gk{δικαιοσύνη} [justice], which all become forms of \gk{ἐπιστήμη} (or of
\gk{φρόνησις}). Since virtuous action is the \emph{conscious} and
\emph{free} subordination under the world-law, and insight and
disposition become essential moments in its concept, only that action by
which the right, i.e.\ what is in agreement with the law of our nature,
is done in the right way, i.e.\ with the right disposition, is moral
action in the proper sense of the word: \gk{κατόρθωμα} (recte factum),
whose opposite is \gk{ἁμάρτημα} [error]. What in its content agrees with
the moral law, what therefore it behooves us to do as duty, is
\gk{τὸ καθῆκον}. This in turn can be something that by its nature ought
always to be done (\gk{καθήκοντα
τέλεια --- ἄνευ περιστάσεως} [perfect duties --- without regard to
circumstances]) e.g.\ \gk{τὸ φρονεῖν}, \gk{τὸ δικαιοπραγεῖν} [being
prudent, acting justly] (which is only the expression of virtuous,
rational action itself) --- or something that \emph{can} be duty in
certain cases (\gk{καθήκοντα μέσα --- περιστατικά} [intermediate duties
--- dependent on circumstances]) e.g.\ to marry, to take up an office.
% errata, p. 212: printed "(...). Den", corrected "(...) f. Ex. at gifte sig, at overtage et Embede. Den"
The action by which this latter kind of \gk{τὸ καθῆκον} is performed
without perfect insight and the disposition thereby conditioned,
% errata, p. 212: printed "dermed", corrected "derved"
the merely legal action, is a \gk{μέση πρᾶξις} [intermediate action],
neither \gk{κατόρθωμα} nor \gk{ἁμάρτημα}. It first becomes moral,
virtuous action by being carried out in the right way.

He who realizes the moral demand through insight and strength of will is
\emph{the wise man}. Only he can act morally in the proper sense of the
word. Since there is no intermediate between virtue and vice, any more
than between straight and crooked, the wise man, as the virtuous man,
comes to stand in abstract opposition to all those who have not yet
reached the right moral insight and the firm willing of the good. He too
who approaches insight and virtue is still outside virtue,
% errata, p. 212: printed "Selv den", corrected "Ogsaa den"
just as fully as the vicious man, and falls under the concept
\gk{φαῦλος} [base] (or \gk{μῶρος} [fool]). Every action of the wise man
is good, every action of the unwise man bad. The image of the wise man is
raised by the Stoics to a paradoxical ideality. The wise man is first and
foremost, as is in general the condition of virtue,
\opage{213} \setcounter{footnote}{0}
completely free from passions, affects (such as fear, desire etc.),
because these always proceed from a false judgment about good and evil,
and, as diseases of the soul, hinder the moral and are something contrary
to reason. Yet his \emph{apathy} is not insensibility, but the effect of
a moral force that prevents the passive state from calling forth an
emotion or influencing action. He alone therefore possesses freedom; he
alone is independent of every external fate. He knows in all
circumstances the right and the good and understands how to carry it
out; he unites in himself all perfections; he alone is the true poet,
orator, philosopher, priest, general, king, rich man, benefactor. He
alone is happy, and in inner dignity he does not stand behind Zeus. He
has the moral law within himself, and is therefore not bound by external
laws, but everything that he does according to his higher insight is
good. He has the right, when life becomes unbearable to him through
external circumstances independent of him, to assert his freedom by a
voluntary death (a \gk{εὔλογος ἐξαγωγὴ τοῦ βίου} [reasonable departure
from life])\footnote{Suicide becomes the highest confirmation of moral
freedom in relation to everything external. The individual negatively
asserts his independence by casting away life, which is an adiaphoron.
Here there comes forward the abstract relation of moral
self-consciousness to what gives moral action concrete determination.}.
Despite the wise man's independence he nevertheless stands in community
with all other rational beings, since man by his nature is ordered into
the community of reason that embraces gods and men. With the Stoics the
concept of \emph{cosmopolitanism} gains a content through the conception
of human society as an organized rational whole, in which the individual
becomes an integral member. Yet a wavering constantly enters here,
inasmuch as on the one hand the conception of the individual as part of a
whole leads to the demand for a union, but on the other hand the abstract
conception of the wise man makes him \gk{αὐτάρκης} [self-sufficient] and
lets the union lose its significance. The life of the state had, for the
Stoics, \opage{214}even if the earlier of them grant the wise man the
right to take part in it in the states advancing toward the ideal, to
lose its proper significance. The incongruity of the state with the
ideal had to lead to collisions between reason and the law of the state.
% errata, p. 214: printed "Incogruents", corrected "Incongruents"
And inasmuch as the relation of community was widened to the universe,
but cosmopolitan activity does not find its place within the single,
limited and imperfect state, cosmopolitanism leads to individualism.

Since the Stoics had forced the ideal of the wise man up to that
paradoxical height, they were compelled to concede that this ideal had
not been realized by anyone. The highest that had in fact been reached by
the most excellent, such as Socrates, Antisthenes and the like, was only
an approximation, a constant advance toward the ideal (a \gk{προκοπή}).
But if the sharp distinction between the wise man and the fool was
maintained, even the most excellent would still be fools, virtue would
not be realized among men, and philosophy would thereby lose its
significance. This consequence leads the later Stoics to concede
differences of degree in good and evil, and to move the state of
progress so near to real virtue that it can receive all its essential
determinations, --- and from there, in despair of human strength, which
was hampered by the dualism between the rational and the irrational in
the human being, to seek outside the human the source of the virtue that
man could not realize by himself.

\divrule

\begin{center}
\textbf{4.\enspace Epicureanism.}
\end{center}

a.\enspace \emph{History of the School.} --- The founder of the school,
\emph{Epicurus}, was born in 342 in Attica or on Samos, whither his
father went with an Athenian colony. As a young man he heard on Samos
Democritus's follower Nausicrates, and took over, with an insignificant
transformation, his atomic doctrine, which he combined with a hedonism
for which the Cyrenaic doctrine was the foundation, but which
nevertheless has a greater independence than his physics.
\opage{215}After having worked for some years as a teacher in other
places, he founded about 306 a school in Athens, of which he remained
the head until his death (270). Of his extraordinarily numerous but
carelessly composed writings, individual fragments, notably larger
portions of the work on nature, are preserved in the Herculanean
manuscripts, in which several works of the other Epicureans have also
been found, especially of \emph{Philodemus}: Diogenes Laërtius gives (in
the 10th book) a detailed account of Epicurus's doctrine.

Epicurus's school displays an extremely slight philosophical
productivity. The scientific interest recedes so strongly behind the
practical that they content themselves with repeating the master's
doctrine, which he himself had brought into compendious form, in which
it was learned by heart by his disciples. Precisely because of this
lack of independence the school held together more closely, while the
other schools gave eclecticism entry. It still flourished in the
Christian era and found numerous adherents in the Roman empire. The
members of the school lived, in a close league of friendship, a life
devoted to cheerful enjoyment, but to an enjoyment ennobled by humanity
and culture. Epicurus himself is praised for his moderation and his
pure morals, and for the strength of soul and the cheerful courage with
which in the last years of his life he bore bodily sufferings.

Among the Epicureans may be named Epicurus's favorite disciple
\emph{Metrodorus}, \emph{Zeno} of Sidon (c.\ 150 B.C.) and his disciple
\emph{Philodemus}, Cicero's contemporary \emph{Phaedrus}, and the
gifted Roman poet T. \emph{Lucretius Carus} (95---52 B.C.).

\divrule

b.\enspace \emph{The Division of Philosophy.} --- Epicurus divides
philosophy into Canonic (doctrine of the norms of knowledge), physics
and ethics. But since philosophy is conceived purely practically as an
activity that produces a happy life by means of conceptual knowledge and
deliberations, the first two sciences are taken entirely into the
service of ethics as auxiliary
sci\opage{216}ences. Every scientific occupation that has no practical
use is rejected. Therefore Epicurus censures occupation with
dialectic, and even physics is declared to be such as would be useless
if it did not free us from superstitious and false notions and from the
fear these instill in us.

\divrule

c.\enspace \emph{The Canonic,} which with Epicurus takes the place of
logic, treats essentially the question of \emph{the criterion of
truth.} This Epicurus seeks, as regards the theoretical, first of all
in sensation, next in the notion (\gk{πρόληψις}), and as regards the
practical, in the feeling of pleasure and pain (\gk{πάθος}), which shows
what is to be striven for and what avoided.

Every sensation is true; even our dreams and the sense-images of the
insane have truth; for they set the soul in motion, which what is not
cannot do. Not our senses, but our judgment about the relation of
sensation (of our being affected) to the object, deceives us: the senses
% errata, p. 216: printed "Sandsningens", corrected "Sandsningens (vor Afficeretheds)"
always speak the truth, and their statements have an immediate
% errata, p. 216: printed "og have", corrected "og deres Udsagn have"
certainty (\gk{ἐνάργεια}). Sensation is also by its nature irrefutable;
for it can be refuted neither by another sensation, since this has no
greater validity and does not refer to the same thing, nor by reason,
which itself develops out of sensation. If we did not put our trust in
sensation, we should fall into an unconditional doubt, which would be
self-contradictory (as a knowledge that we could know nothing),
% errata, p. 216: printed "Vished", corrected "Viden"
and would cancel the possibility of action. How the true images of the
object are to be distinguished from the false, Epicurus cannot decide.

From the sensations arise the notions, the concepts
(\gk{προλήψεις}), in that several similar sensations are gathered in
memory into a general image (\gk{καθολικὴ νόησις}). This image has its
expression in the name (\gk{ὄνομα}) and is called up for memory by it.
The notional images have truth as products of earlier sensations.
\opage{217}Opinion (\gk{δόξα}, \gk{ὑπόληψις}) forms itself through
after-effects of the impression of objects, with the cooperation of
reflection (\gk{συμβαλλομένου τι καὶ τοῦ λογισμοῦ}). It goes beyond
the given and relates itself to something future and to what cannot be
immediately observed. It can be true or untrue; it is true when it is
confirmed by sensation, or when no sensations testify against it. The
logical processes by which one passes from what is sensed to what lies
beyond sensation Epicurus treats carelessly, although his theory rests
on what cannot be sensed (the atoms). A detailed doctrine of the logical
functions he holds to be superfluous. Later Epicureans, notably
\emph{Philodemus}, do, however, occupy themselves with inference by
analogy and by induction as inferences that lead from the phenomena to
\gk{τὰ ἀφανῆ} (\gk{τὰ ἄδηλα}) [the non-apparent (the unclear)]. They
also reckoned among the criteria the intuitive apprehensions of the
understanding.

\divrule

d.\enspace \emph{Physics.} --- In accordance with the practical
conception of the aim of physics, the main task for Epicurus becomes to
set up a world-view that removes supernatural causes and leads
everything back to natural grounds. A thoroughgoing natural-scientific
knowledge of concrete phenomena is not his concern; it is enough for him
when, with respect to the particular, various natural explanations can
be put forward hypothetically; to decide which of these is the one
possible explanation is in his opinion impracticable in most cases, and
the attempt to do so leads to arbitrariness.

Only in a few points does Epicurus in his physics depart from
Democritus's atomic doctrine. The atoms and space become for him too
the principles, at which he arrives by starting from the given
phenomena. The corporeal, the only substantial thing, is partly a
composite, partly the individual, indivisible and unchangeable elements
of this composite. To the assumption of these elements one must
necessarily come, since one would otherwise have to hold
\opage{218}that the division of the composite would at last lead to
something that is not, which is unthinkable, since nothing comes from
what is not or is dissolved into what is not. Empty space, the only
incorporeal thing, is, as the condition of the motion of the atoms,
necessary in order to explain the phenomena. The atoms are by their
nature uncreated and imperishable; they differ from one another by the
general properties of the corporeal: size, shape and weight; the
properties of concrete things, e.g.\ color, smell, do not belong to
them. They are invisible, since, despite the difference in their size,
this size is nevertheless always so small that it cannot become an
object of immediate sensation. They are infinite in number; therefore
space and the world must be unlimited. Originally they move, in
accordance with their weight, downward in empty space and with equal
speed, not, as Democritus assumed, with a speed differing according to
size. In order to explain the combination among them, Epicurus assumes
--- as it were placing in the atoms the arbitrary self-determination he
wishes to preserve for the human individual --- that the atoms deviate
slightly from the perpendicular line. This gives rise to contacts,
aggregations and vortical motions, by which the worlds are formed. Of
worlds Epicurus, like Democritus, assumes a simultaneous infinite
multitude, which are separated by the intermundia and are as individual
divisions of \gk{τὸ ἄπειρον} [the infinite]. Our world comprises our
earth and all the heavenly bodies visible to us. The worlds, of which
some may by chance be like one another, are perishable, just as they
have come to be: the process of world-formation continues to infinity.
From the coming-to-be of the world and the formation of particular
things everything teleological is excluded: what is apparently
purposive, e.g.\ in the organic, is a purely accidental result of
effects that are necessary by nature. Among the natural forms that were
formed in a mechanical way, some were able to continue life and to
develop into higher forms (e.g.\ man), while the others soon perished
again. The ensoulment and divine nature of the heavenly bodies Epicurus
denies; their mo\opage{219}tions and the celestial phenomena have purely
mechanical causes. The existence of the gods Epicurus does not reject,
since the belief in them is grounded in a \gk{πρόληψις} that derives
from appearances of gods to men, but he rejects the notions of them in
popular religion, sprung from ignorance and fear, as impious and as the
source of the greatest evils for the human race. He ascribes to the
gods human shape because in theophanies they show themselves in such a
shape, and because it is the most beautiful. They must be thought of as
imperishable and blessed, and therefore a body of the finest atoms is
ascribed to them, and they are given a place outside the perishable
worlds, in the intermundia, where they live without concerning
themselves about the world, without labor, sorrows and passions, and
enjoy blessedness in the contemplation of their perfection. They are
the object of a free and pure worship, not for egoistic reasons, but
for the sake of their excellence. ---

In \emph{anthropology} the fundamental thought of physics is carried
through. The soul of man is, as standing in reciprocal action with the
body, itself corporeal. It is formed of the smoothest and roundest
atoms, and is described as a fiery or airy matter that is spread
through the whole body. The rational part, however, has its seat in the
breast. Its functions are sensation, notion, the movements of the will
and of the feelings; in the last instance life too depends on it.
Although it makes up one being with the irrational part, they can be
affected in different ways, and parts of the irrational soul can be
lost without the rational soul being disturbed. When the organism is
wholly dissolved, however, the light atoms of the soul are scattered at
once, and the soul does not survive the body.

In the conception of the relation between the rational and irrational
part of the soul, the dualism between the mental and the sensuous comes
to light, even though both are brought under the point of view of
matter.

Sensation Epicurus conceives as coming about in that corporeal images
(\gk{εἴδωλα} --- \gk{τύποι}) detach themselves from the
\opage{220}surface of things and through the separating medium reach
our senses. The notions to which no real object corresponds arise
partly through after-effects of images, partly through the
\gk{εἴδωλα} being mixed or disturbed before they reach us. Despite his
materialistic view, Epicurus holds firmly to \emph{the freedom of the
will}, because without it accountability would be canceled and man
subjected to a lamentable necessity. A fate (\gk{εἱμαρμένη}) there is
as little as a providence. A deeper inquiry into the essence of will and
freedom Epicurus did not, however, undertake; he contents himself with
the indeterminate notion of an accidental action, which in any case has
its prototype in the accidental motion of the atoms.

\divrule

e.\enspace \emph{Ethics.} --- Epicurus's ethics is a hedonism in which
Aristippus's doctrine comes forth in a modified and more reflective
shape. The highest good, \emph{happiness}, consists according to
Epicurus in \emph{pleasure}, toward which every being's natural desire
is directed, and which is the aim of every action; but pleasure he
conceives in a different way from Aristippus. For the latter the goal
was pleasure as the individual pleasure, as positive pleasure (pleasure
\gk{κατὰ κίνησιν}, i.e.\ that which arises through a motion), and as
bodily pleasure; for Epicurus the goal is pleasure as a total state, as
restful pleasure (\gk{καταστηματικὴ ἡδονή}), as freedom from pain, and
mental pleasure is higher than bodily. Positive pleasure we need only
when its absence would cause us pain. Mental pleasure, which Aristippus
had conceived as a weaker echo of the bodily, Epicurus conceives as the
higher, because, even though in the last instance it always stems from
something bodily, it is not, like the bodily, restricted to the present
moment, but through memory and hope relates itself also to the past and
the future, whereby the enjoyment of the moment is intensified. It also
has the power to bring momentary bodily displeasure into oblivion, as
\opage{221}Epicurus tells of himself that his bodily pain was
outweighed by the pleasure that the memory of his philosophical
discoveries afforded him. --- Since pleasure as a lasting state is the
goal, the task becomes to make a choice among the pleasures, which in
and for themselves are equally justified as natural, according to the
consideration of their consequences. A pleasure must be given up that
will have greater displeasure as its consequence; a displeasure is to
be chosen that can bring about a greater pleasure or remove greater
pains. There takes place, therefore, in the individual case a measuring
(\gk{συμμέτρησις}) of the immediate and mediate pleasure against the
immediate and mediate pain, and according to the result of this
measuring the decision is to be taken. As the condition of a lasting
and fresh enjoyment Epicurus sets frugality and limitation of the
desires. Only the natural (\gk{φυσικαί}), not the empty (\gk{κεναὶ})
desires are to be satisfied. What is essential is the soundness of the
body and the undisturbedness of the soul (\gk{ἡ τοῦ σῶματος ὑγίεια καὶ
ἡ τῆς ψυχῆς ἀταραξία}). In being able to make the right choice among
the pleasures, insight (\gk{φρόνησις}) manifests itself, which is the
highest virtue and the source of all other virtues. Without it a happy
life cannot be led, but when it is present, such a life will
necessarily follow along with it: the wise man, i.e.\ he who
understands how to strive after pleasure in the right way, has, through
rational insight, in every circumstance of life tranquillity and
contentment, since it frees him from fear and prejudices, gives him the
strength to moderate the desires and the affects, and raises him above
the contingencies of life. Virtue is thus the sure way, the sure
\emph{means} to the highest aim: happiness. The wise man's possession of
insight is an inalienable property; he can never become unwise. Whether
his life lasts a short or a long time makes no difference to the degree
of his happiness.

By Epicurus's conception of pleasure, which has its counterpart among
the later Cyrenaics, enjoyment is thus spiritualized, and so drawn out
of dependence on the external, that in reality only the name of pleasure
has remained. What is striven toward is the painless
undis\opage{222}turbedness of the individual life in the cultivated
mind's consciousness of independence of the external. The sensuous
content of enjoyment is pushed back; the real accent falls on the
mental determination: the rest in itself of infinite free
self-consciousness. Not sensuous enjoyment but mental consciousness, not
the given but the acquired, becomes the highest.

Since with Epicurus the individual stands as a mental atom, resting in
its abstract self-consciousness, it becomes estranged from and
indifferent to the forms of community in which the Greek had formerly
found the substantial content of his life: family life and the life of
the state. In the place of these forms of community there enters the
\emph{community of friendship}, grounded on a free choice. Friendship is
conceived in principle in an egoistic way, as a means of securing the
enjoyment of life; but the nobler conception of friendship as a
disinterested goodwill and as a natural league among the wise asserts
itself in the Epicureans' practice. In the Epicurean school there
developed, through the relation of friendship that bound its members
together, a spirit of humanity, by which a counterweight was formed
against the craving for enjoyment and the egoism that could not be
separated from the school's ethical principle. ---

Between the Stoic and the Epicurean ethics the opposition strikes the
eye at once, since the former conceives man from the side of his
universal reason and places the ethical task in unconditional
subordination to the rational world-law, while the latter conceives man
from the side of individuality, and places the ethical task in the
individual's satisfaction in painless undisturbedness by the external.
But beneath this opposition there nevertheless shows itself an essential
likeness, not only in the way in which the practical interest steps into
the foreground, but in the definitive determination of the ethical
ideal. Both systems seek the happiness of the subject, which is the
final goal of philosophy, in self-consciousness's rest in itself and
independence of the external, acquired through insight; and in both
systems, despite the unlike point of departure, this self-consciousness
freed by knowledge acquires an abstract universality, in which
\opage{223}there lies a point of attachment for an ethical universalism,
but which estranges it from the concrete forms of the life of the state
and of civic life. The Stoic and the Epicurean wise man are in the last
instance essentially alike: their world is the inner infinity of the
self, their happiness their rest in this infinity.

\divrule

\begin{center}
\textbf{5.\enspace The Skeptics and the Middle Academy.}
\end{center}

The same practical goal that the dogmatic systems strove for through
knowledge of the essence of things: freedom of mind, the subject's
rest in itself through independence of the external, skepticism
strives to reach through doubt, which cancels all certain objective
knowledge and makes impossible any dogmatic assertion about the worth of
the external. And while skepticism has the same goal as the
contemporary dogmatic systems it combats, it shows indirectly by its
arguments, by letting the certainty of knowledge fall with the certainty
of sensation, that it is influenced by the same materialistic tendency
as Stoicism and Epicureanism.

\divrule

a.\enspace The \emph{older skepticism,} which appears a little earlier
than Stoicism and Epicureanism, is represented by \emph{Pyrrho} of Elis
(360---270), the founder of this tendency, and by his disciple
\emph{Timon} (c.\ 325---c.\ 235). \emph{Pyrrho} seems to have been
influenced most directly by the Megarian philosophy and by Democritus's
doctrine, whose distinction between the original properties of things
and those arising through their relation to our sense organs may have
given him the essential skeptical motive. He wrote nothing; his doctrine
was brought into written form by Timon, but of the latter's writings we
have only reports.

When happiness is striven for, the nearest question becomes how we are
to comport ourselves toward things; but \opage{224}this leads back to
another: how things are in and for themselves. To this last question,
say the skeptics, we can give no answer, because the essence of things
hides itself from us, and the phenomenon is changing and inconstant.
Neither our senses nor our thinking can therefore decide anything
certain about things. But when this is the case, we dare state no
determination about them, and must not give our assent to any such
statement. For it does not express something that is in reality
(\gk{ἀληθείᾳ}, \gk{φύσει}), but an opinion, which at most can rest on
habit and convention. Every statement and its contradictory opposite
have equally valid grounds on their side: neither of them, therefore,
has justification. But when we [\dots] theoretical undecidedness
% printed "naar vi tilretiske Uafgjorthed": words evidently dropped (logged in odd-spots); sense supplied in brackets
(\gk{ἀῤῥεψία}), [we come] to the practical goal: undisturbedness of mind
(\gk{ἀταραξία}), apathy, which like a shadow accompanies that withholding
of judgment (\gk{ἐποχή}, \gk{ἀφασία}, \gk{ἀκαταληψία}). Nothing external
can disturb our tranquillity when we become conscious that we cannot
know anything certain about its essence. It then becomes an
\gk{ἀδιάφορον} [indifferent thing]. Only the calm disposition of the
mind acquires worth, and this disposition, which is one with virtue,
includes happiness.

The skeptical statements about the uncertainty of knowledge, which are
directed against dogmatism, are then at last, in order not to lead to a
new, negative dogmatism, applied also to themselves, whereby every
dogmatic assertion is to be removed and the theoretical suspension that
the skeptics strove for carried through. Those statements are not to
express an assertion but to state a condition of mind; they are not to
be a \gk{δόγμα}, but an \gk{ἐξομολόγησις}, a confession.

The skeptical tendency founded by Pyrrho, whose adherents designated
themselves as \gk{σκεπτικοί}, \gk{ζητητικοί}, \gk{ἀπορητικοί}, or
\gk{ἐφεκτικοί} [inquirers, seekers, doubters, or suspenders of
judgment], soon dies away again, and only after the lapse of several
centuries is skepticism renewed in a more thoroughgoing form, partly
with new motives grounded in the intervening development. In the
interval between the older and the younger \opage{225}skeptical school,
the so-called ``Middle Academy,'' in its polemic notably against the
dogmatism of the Stoics, represents the skeptical tendency in
philosophy, which gains strength from the deficient scientific grounding
of the new dogmatic systems and from the opposition between their
physical and ethical tenets.

\divrule

b.\enspace The founder of the \emph{Middle Academy, Arcesilaus}
(315---241), was the successor of Crates as head of the Academic school.
He combated the Stoic doctrine of \gk{φαντασία καταληπτικὴ} [cognitive
impression] as the criterion of truth, maintaining that the false notion
could have for the fool the same evidence and convincingness as the true
one for the wise man. And with the certainty of sensation he lets the
certainty of knowledge in general be canceled, since rational knowledge
rests on the former. The result, then, is that there is no knowledge at
all. Not even the fact that we know nothing can become an object of
knowledge. One must therefore refrain from every decision
(\gk{ἐπέχειν}), not give one's assent (\gk{συγκατάθεσις}) to any
positive statement. The highest we can reach is the probable
(\gk{τὸ εὔλογον}), and this must become decisive for our practical
conduct, in the choice between what we are to do or avoid.

Through \emph{Carneades} (214---129) Academic skepticism received a
more thoroughgoing form. His polemic was directed not only against the
doctrine of the criterion held by the dogmatic systems, notably
Stoicism, but with great acuteness against the tenets of the various
philosophical systems. Since his dialectic canceled every criterion of
truth, every positive decision became impossible, and dialectic gained
the right and the latitude to defend, with respect to every object and
with equal force, the contradictorily opposed views. The interest in the
practical, in ethical action, which a thoroughgoing \gk{ἐποχὴ} would make
impossible, led Carneades to develop more fully \opage{226}the doctrine
of the surrogate of truth, \emph{probability} (\gk{πιθανότης},
\gk{ἔμφασις}), and to determine its \emph{degrees}. The lowest degree of
probability belongs to the notion that makes an impression of truth but
stands isolated (it is merely \gk{πιθανή}); a higher degree to the one
that stands in connection with others and is not contradicted by them
(it is \gk{πιθανὴ καὶ ἀπερίσπαστος}); the highest degree to the one that
is not merely not contradicted, but is confirmed from all sides by the
notions with which it hangs together and whose probability is already
secured (it is \gk{πιθανὴ καὶ ἀπερίσπαστος καὶ περιωδευμένη}). To these
degrees of probability there correspond, since the relation of notions
\emph{to us} falls within our consciousness, the degrees of our
conviction, and by these in turn our practical choice is determined.
Only in what is indifferent can we content ourselves with the lowest
degree of probability; in what is more important a higher one is
required; in what is necessary, in what conditions happiness, the
highest. This doctrine of probability is the essential theoretical mark
of distinction between the Academics and the skeptics. In ethical
respects the skeptical Academics departed from the skeptics proper in
not, like these, setting ataraxia as the ethical goal. Carneades seems
to have set as the ethical goal virtue in conjunction with the things
that satisfy the original natural drives. After Carneades the Academy
returned to dogmatism and soon assumed, like the other contemporary
systems with the exception of Epicureanism, an eclectic character. ---

\divrule

c.\enspace The \emph{later Skepticism} is represented chiefly by
the renewer of this tendency, \emph{Aenesidemus} (1st cent.\ A.D.),
by \emph{Agrippa} and by \emph{Sextus Empiricus}, whose writings
(\gk{Πυρρώνειοι ὑποτυπώσεις} and 11 books against the Dogmatists) are
our most important sources for knowledge of ancient Skepticism.

\emph{Aenesidemus} of Cnossus in Crete appeared as a
teacher in Alexandria. From him derive the so-called 10
\opage{227}tropes (\gk{τρόποι}), i.e.\ ways of grounding Skepticism and of demonstrating the impossibility of objective knowledge. The \emph{first} emphasizes the dissimilarity between the various living beings, which
conditions different impressions of the same objects; --- the
\emph{second}, the mutual diversity of individual men in
bodily and mental respects, which entails the same
consequence; --- the \emph{third}, the diversity of our individual senses, their opposite reports about the same thing; --- the
\emph{fourth}, the influence of our changing mental and bodily states on perception; --- the \emph{fifth}, the influence of
distance and position on the way in which the object appears to us; --- the \emph{sixth}, the effect of the medium
through which the thing is perceived; --- the \emph{seventh}, the influence of the object's condition, of the composition of its parts, of quantity, on the way in which it shows itself to us; ---
the \emph{eighth}, the relativity in all the properties of the thing, which all
express, not its essence, but a relation to another thing,
and are at the same time changeable and accidental; --- the \emph{ninth}, the
different perception of things according as they are new or
strange to us, unaccustomed or accustomed; --- the \emph{tenth}, finally, the diversity in morals and customs, laws, religious notions,
philosophical opinions and the like, which shows that a universally valid
determination of the true and the right is not given. By these
tropes, which are all meant to demonstrate that the same thing
appears to us in opposite ways
% errata, p. 227: printed "opfattes af", corrected "viser sig for"
(\gk{τἀναντία περὶ τὸ αὐτὸ φαίνεσθαι}), Aenesidemus is said to have wished to pave the way for the
% errata, p. 227: printed "begrunde", corrected "bane Vei for"
Heraclitean doctrine that opposites are united in the same
thing (\gk{τἀναντία περὶ τὸ αὐτὸ ὑπάρχειν}); his Skepticism would
thus have been a polemical position in defense of a dogmatic view.

By \emph{Agrippa} the 10 tropes of Aenesidemus are reduced
to more general points of view, and, since the attack was directed essentially at the possibility of \emph{scientific} knowledge, 5 tropes were set up, among which the one that emphasized the diversity of opinions about the same objects was placed at the
\opage{228}head. The \emph{second} trope then emphasized the impossibility of
proof, on the ground that this must go on to infinity, since every ground of proof must itself in turn be proved. The
\emph{third} accentuated relativity, the object's different appearance according to the peculiarity of the perceiver and according to
its relation to other things. The \emph{fourth} showed that the Dogmatists,
in order to avoid going on into an infinity, stop at arbitrary
presuppositions. The \emph{fifth}, finally, pointed out the circle in their
proof, in that what is used as a ground of proof and what is
proved mutually presuppose one another (thus, notably,
the truth of sensation and of thinking). These tropes are reduced by younger Skeptics, with the same aim at the impossibility of
\emph{scientific} knowledge, still further, to
2, of which the \emph{first} emphasizes that the oppositions in views about everything sensible and thinkable show that nothing is
certain by itself, i.e.\ without proof; the \emph{second}, that such a
proof cannot be carried out, since every ground of proof would itself have to be proved, which would lead either to a circular proof or to a
regress to infinity.

\emph{Sextus Empiricus} (first half of the third cent.\
A.D.) gathers up the whole skeptical tradition, and sets out the skeptical arguments both against the formal conditions of knowledge and against
the material content of the sciences, with the same endeavor as in the earlier Skeptics to avoid a
negative dogmatism. In him Skepticism is directed more against
scientific knowledge and against the philosophical systems
than against experience and the pronouncements of sound common sense.
In the practical domain Sextus softens the standpoint of the older Skeptics.
He concedes that through experience,
which relies on \gk{τὸ φαινόμενον} [what appears], one can find certain rules for
practical procedure, and he concedes to habit
and convention a certain significance in those cases where no judgment based on experience is possible. On the whole he recognizes 4
norms for our action: immediate sensation and the
reflection arising from it, natural impulse and need, law
\opage{229}and convention, experience and art. Complete indifference
toward the external cannot be maintained when one has to act; therefore
the Skeptic is unavoidably led to seek such norms. And
since the Skeptic must concede that he too cannot withdraw from agreeable or disagreeable impressions, in general
from disturbances, the goal toward which he is to strive is
no longer determined as complete apathy, but as
a limitation of the excess of affect, a \emph{metriopathy},
which is posited as conditioned by this: that the unavoidable is not made
heavier by being determined as an evil, but that judgment about
it is withheld. ---

In ancient Skepticism the subject's withdrawal into itself,
which characterized Stoicism and Epicureanism, emerges in
its sharpest form. The Skeptic's ataraxia rests, not on
knowledge of the essence of things, but on the abandonment of all
knowledge, not on a positive conviction, but on indifference toward everything that stands as something unknown. It is characteristic
of ancient Skepticism that the abandonment of knowledge does not,
as in the Skeptics of modern times, leave behind a painful
sense of loss. But this very abandonment of knowledge is the sign of
the exhaustion of Greek thought, of its distrust of itself.
It feels itself powerless to reconcile the oppositions of existence, to find the truth independently; it therefore gives up
all knowledge in Skepticism in order to save the subject's happiness, or,
when the interest of action drives it to a theoretical judgment, it seeks
a basis for this in a choice of what is most probable among the various views. In Eclecticism there appears
a new symptom of the exhaustion of thought.

\divrule

\begin{center}
\textbf{6.\enspace Eclecticism.}
\end{center}

a.\enspace \emph{General character of Eclecticism.} --- Until
toward the beginning of the first pre-Christian century the various philosophical tendencies proceed unmixed side by side; then there begins a \opage{230}fusion through attempts at mediation and union, from which
only Epicureanism keeps itself untouched. The consciousness of the defects of the individual
systems, clarified by the mutual criticism of the various
schools, leads, since philosophical productivity
is exhausted, to an eclectic choosing of what is seemingly certain in the various
systems, which the individual school then
appropriates, while by an adaptation of its own doctrines it veils the difference from what is foreign. The
skeptical tendency, too, was bound in the end to lead to Eclecticism, inasmuch as,
after certain knowledge had been given up, the probable was, for the sake of
action, given a significance which, where the
practical need came still more decidedly to the fore, in turn
let it take the place of the true, and inasmuch as this probable could not
be sought by the skeptical consciousness in a single system,
but, as the subjective need directed, had to be sought in the
various systems, which for the Skeptic were equally justified. In particular,
from the successors of Carneades in the Academy there issued an eclectic
tendency, which relied on the necessity of definite
views as a basis for practical life. A great influence on the spread of Eclecticism was exerted by the transfer of Greek
philosophy to the \emph{Romans}, for whom on the whole
philosophy was not an end in itself, but a means to a training for practical life and to the grounding of a moral
and religious view of life. Therefore the interest in
systematic unity and strict consistency in philosophy receded before the interest in practical applicability.

Since Eclecticism demands that the individual select
the true in the various systems for his own use, it
in fact proceeds from the tacit presupposition that
the individual, prior to the scientific decision, bears within himself
a standard for distinguishing between the true and the false.
Eclecticism thus has as its background a faith in what
is immediately given in the subject's self-consciousness. It has
therein a motive to distinguish itself from the sensualism of the preceding systems, inasmuch as self-consciousness becomes a \opage{231}peculiar, \emph{original} source of knowledge. But Eclecticism can
thereby at the same time pave the way for those forms of Greek
philosophy which rely on the religious and on immediate divine revelations. What it finds in
the mind as an ``innate'' content of truth is not
comprehended in its origin from the active human mind,
but perceived as a gift of a higher power, and there is then
only one step to going back into self-consciousness merely in order
to receive in it divine revelations, a communication
of a truth which the human mind cannot itself find. On the
other hand, Eclecticism can also lead to a \emph{renewed
Skepticism}, since in its choosing there always lurks a residue of
doubt, which is held down only by the peremptory dictum of immediate
consciousness, and since its piecing together of the diverse
never avoids inconsistencies and contradictions. The renewed
Skepticism did indeed take as its essential motive the opposition
between the various systems. But this Skepticism then leads
in turn, indirectly, to the same goal as Eclecticism's faith
in the immediately given. When an eager need for knowledge of truth meets with doubt about the truth of all science,
the mind is forced more strongly toward the direction in which the source of knowledge is sought outside thought and science, partly in
the revelations of the Godhead within, partly in religious
tradition.

\divrule

b.\enspace \emph{The appearance of Eclecticism within the individual schools.} --- In the \emph{Academic} school Eclecticism emerges in pronounced form in \emph{Antiochus} of Ascalon,
the founder of the so-called fifth Academy. He maintains that
the Academics, the Peripatetics and the Stoics in reality
agree, and differ from one another only in words, or
at any rate in unessential points. But when he wishes to maintain this
agreement in the domains of the theory of knowledge, of metaphysics and of ethics, it is only by superficial \opage{232}minglings, by abandoning what is most essential in Platonic
philosophy: its doctrine of Ideas, and by inconsistencies, as when
in ethics he adopts the Stoics' doctrines about the wise man
and about apathy, but abandons their doctrine of the equality of all vices,
and in the concept of the highest good includes
the goods of the body, which are supposed to have value in and for themselves and to belong
to ``the happiest life'' (vita beatissima).

Among the Academics we may also name M. \emph{Tullius
Cicero} (106---43), insofar as he himself counts himself as belonging to this
school, on whose skeptical foundation his Eclecticism rests.
His Skepticism is, however, strongly limited, partly by his
conceding great value to probability in general, partly by his almost entirely excluding doubt from the ethical
and religious domain, which essentially interests him, while the
physical, which is more indifferent to him, is posited as the more uncertain.
Doubt is for him to be, first of all, a preparation for a positive conviction which, even if it has only the probable as its basis,
% errata, p. 232: printed "Sandsynlighedens Charakteer", corrected "det Sandsynlige til Grundlag"
nevertheless satisfies in practical respects. For the discovery of the probable the decisive thing is the \emph{immediate} testimony of the senses and of
self-consciousness, especially of the
latter. The immediate, \emph{inner} certainty, the innate feeling for truth, Cicero strongly emphasizes. It is a knowledge
which precedes all experience and science, and which, if
it developed undisturbed, would make all science superfluous.
Through it, which is like a commune judicium, we have consciousness of the fundamental concepts of right and morality (notiones
innatæ), of the Godhead and of the immortality of the soul.

In the domain of ethics Cicero polemicizes sharply against Epicureanism. The Peripatetics and the Academics he regards as
essentially in agreement; about the Stoics' relation to them
he expresses himself waveringly, now positing a real difference, now only a difference in expressions.
When he emphasizes the real difference, he often speaks with
% errata, p. 232: printed "han taler", corrected "han taler, naar han fremhæver den reale Forskjel, tidt"
enthusiasm of the Stoics' severity and ideality, and has misgivings about the Peripatetics' milder conception of the affects and about their determination of external goods as \opage{233}belonging to happiness. But he is in turn frightened back from Stoicism by the fact that its ideal does not fit into
actual life, that its distinctions are too abstract, and
that it does not take account of the \emph{whole} human nature;
and he is thereby led back again closer to the Peripatetics, and
even designates himself as one of them. He is conscious of his uncertainty, but finds excuses for it. His ethics has no original content; in general he is
everywhere dependent in his philosophy on his Greek predecessors.
Cicero's great significance lies in the fact that his writings have
been one of the most essential sources through which knowledge of Greek philosophy and interest in it were
spread among the Romans and later, in the Middle Ages
and modern times, among the peoples of the Occident. (The later
Academics will be discussed in the next §.)
In the \emph{Peripatetic} school, in which philosophical
productivity receded before a learned occupation with the
older philosophy and a careful interpretation of the master's
writings, Eclecticism emerges strongly in the treatise \emph{\gk{περὶ κόσμου}}, included among Aristotle's works, which was probably
composed about the beginning of the Christian era. With
the Aristotelian doctrine of the universe is combined the Stoic
view of the diffusion of the divine power through the
world as the world's governing law, and of the world's unity as
a \gk{συμπάθεια} [sympathy]. There are also some connections to Plato. ---
In \emph{Aristocles} of Messene (end of the 2nd cent.\
A.D.) an approach to Stoicism takes place in the transformation
of the Aristotelian doctrine of \gk{νοῦς παθητικὸς} and \gk{νοῦς ποιητικὸς} [active mind], inasmuch as \gk{νοῦς παθητικός}, the potential reason, is conceived
as the bodily mixture which offers such an organ
for the divine \gk{νοῦς} active in all bodies that the
activity of the divine \gk{νοῦς} can come to light in a
particular way, and an individual energy of thinking arise,
through the potential reason's being raised to an actual one. The
peculiarity of that mixture is posited, in Stoic fashion, as \opage{234}conditioned by the preponderance of fire. --- Even in Aristocles's pupil,
the famous interpreter of Aristotle, \emph{Alexander} of Aphrodisias (beg.\ of the 3rd cent.), Eclecticism is recognizable through
an approach to the Stoics, which betrays itself partly in the determination of universal concepts as merely subjective abstractions
from the real individual, partly in the conception, verging on materialism, of the soul's being conditioned by the body, partly in a
tendency, despite the distinction between the divine \gk{νοῦς} and the world being maintained, to transform the activity of the former
into a natural process.

In the \emph{Stoic} school an approach to the other
philosophical systems shows itself in \emph{Panaetius} (180---111 B.C.) and
his pupil \emph{Posidonius}. Panaetius appeals not only
to the masters of his school, but also to Plato and Aristotle,
and he strives to soften the harshness of Stoic ethics. Posidonius shows the same recognition of the great thinkers of the other schools,
and has a tendency to emphasize the agreement of the various philosophical systems in essentials. In anthropology he departs from the doctrine of the Stoics,
since he considers it self-contradictory that reason should
be the origin of the irrational and passionate, and therefore,
in Plato's manner, derives the affects from \gk{τὸ θυμοειδὲς} [the spirited part] and \gk{τὸ
ἐπιθυμητικόν} [the appetitive part], which are supposed to be determined by the constitution of the body. An opposition is thereby posited between the
rational and the irrational in the human being, and this
\emph{psychological dualism}, which gradually develops
more sharply, acquires significance as one of the moments that prepare the last forms of Greek philosophy.

Also in the later Stoics, in \emph{Seneca}, \emph{Epictetus} and \emph{Marcus Aurelius}, we find departures from the severity of Stoic doctrine; but what particularly emerges in them is that consciousness of man's lack of
moral strength, and that dualistic conception of the human being, which prepare the disappearance of Stoicism into Neoplatonism.

\divrule

\opage{235}
\begin{center}
\textbf{7.\enspace The Dualistic Systems.}
\end{center}

a.\enspace \emph{General character of the dualistic systems.} --- When the exhaustion of thought and doubt about the power of
human knowledge lead the mind to seek the source of the truth after which it strives outside thought
and outside science, in an immediate communication from the Godhead, in a revelation, through the religions or in the inner being of the individual, this change in the direction of the mind entails a changed determination of the divine and of
its relation to the world. Thought, which doubts itself, can
seek determinations for the divine neither in itself nor in natural existence, which, as
the correlate and opposite of thought, is something finite and therefore
untrue. It must
go beyond the oppositions, and intuit the divine
as the absolutely simple, which in its pure otherworldliness, in its
infinite exaltedness, is inaccessible to knowledge, but
is nevertheless always held fast as something ideal. Only in such an absolute
does the mind find the satisfaction for its need of happiness which neither pleasure nor Stoic apathy could give
it. This divine, which excludes from itself the fundamental determination of the natural, gets matter as its opposite,
and the dualistic distinction between the purely spiritual Godhead and matter is pronounced in principle. The human mind, which
seeks its satisfaction in the divine, therefore flees matter, disdains the world of sense, and seeks, through ascetic
liberation from everything that derives from what is opposed to the Godhead,
to draw near to the Godhead. But since the Godhead stands
as the transcendent, as unconditionally exalted above all contact with the world, intermediaries must be sought through which
the connection can be brought about, and these are found in intermediate beings
of divine kind, in divine emanations, in subordinate gods, in demons and the like. Through these the purified subject receives the revelation of the divine, and, having
opened itself to this, it can be led to a contact with
\opage{236}the Absolute itself in enthusiasm, in which the self blissfully loses itself in the abstract divine. ---

This peculiar turn in the development of Greek philosophy,
which emerges at about the same time as the preponderance of Eclecticism, and constantly gains greater power, is represented
chiefly by the Neopythagoreans and by Platonists who approach
them, but also within the other systems, notably in
Stoicism, we see its most essential distinguishing marks emerge.
And in the philosophical phenomena which, under the influence of Greek
philosophy, proceed from the religious foundation of Judaism, the same tendency is strongly pronounced, so that
the Jew Philo, the most important of the Jewish-Alexandrian philosophers, becomes in general its most important representative and
the one who most immediately prepares Neoplatonism.

Among the various peoples of antiquity there was present, at the time when
that tendency appeared, the same feeling of the weakness of the human mind, the same longing for divine assistance through a divine revelation. Similar philosophical phenomena could therefore arise independently of
one another among the various nations. But the contact between the peoples among whom the spiritual movement was most vigorous was nevertheless so strong, notably in Alexandria,
the real center for the friction and mixing of the
Greek and the Oriental, that an influence of the foreign
could not be avoided on either side. It was most
to be expected, given the superiority of Greek culture, that the philosophy which arose among the Jews would show itself essentially
influenced by Greek philosophy. And this is also the case with what we know of it. We see even in Philo, although
he takes his starting point in the Jewish religion and only
wishes to use philosophy as an auxiliary, that philosophy nevertheless
entirely transforms the positively religious, and not only determines
the scientific form of his system, but also furnishes by far
the greatest part of the individual concepts and doctrines.
Therefore we may bring this Jewish philosophy under
\opage{237}the Greek, just as we do with the philosophical attempts
which derive from Romans. --- But the Greek dualistic systems have, for their part, not remained without all influence from the
East, notably the Jewish. We may indeed derive all
their essential peculiarities from the development of the preceding Greek
philosophy, but the spiritual direction onto which
this led has evidently gained strength from the Jewish
religion's notions of a spiritual God infinitely exalted above the world, of this God's revelations, of the angels and
of prophecy and ecstasy, and from the ascetic tendency
which we see appear among Jewish sects in Palestine and Egypt.

\divrule

b.\enspace \emph{The Neopythagoreans.} --- Neopythagoreanism seems to
have arisen in Alexandria at the beginning of the first pre-Christian century. Cicero names P. \emph{Nigidius Figulus} (1st half of the 1st cent.) as the first who renewed the
Pythagorean philosophy; yet this is uncertain, and in any
case his significance seems to have been slight. Only
later do we find more prominent representatives of this
school, which spread especially widely in the Orient. Within it
a distinction can be drawn between a more \emph{practical-religious} tendency, which has its representative in \emph{Apollonius} of Tyana
(1st cent.\ A.D.), and a more \emph{speculative} one, which is represented by \emph{Moderatus} (contemporary with Apollonius) and \emph{Nicomachus} (middle of the 2nd cent.); yet the difference between
these two tendencies is only a difference of degree. The \emph{practical-religious} tendency conceives philosophy primarily as purity of
life and the right worship of God. Since all earthly things
are seen as defiled by their material existence, the
moral task is placed in asceticism, notably in abstention from the enjoyment of meat and wine and in celibacy, and in bodily purifications,
and all sacrifices to the Godhead are rejected. Through asceticism
man is to be freed from the bonds of the body and attain to a supernatural power and knowledge. The \emph{speculative} tendency has the
\opage{238}same practical goal, but lays greater stress on knowledge.
Taking its theoretical foundation from older philosophical
forms, notably from the Pythagorean doctrine of numbers as modified by Platonic elements,
it transforms what it has taken over, but covers the
discrepancy by the claim that it brings out the deeper
meaning which the older thinkers deliberately concealed under sensible
forms and signs. The doctrine of numbers is developed in a metaphysical-theological direction,
but the development soon loses itself in arbitrary and empty fantasizing.
Common to both tendencies is the determination of the principles. The
divine is conceived as purely spiritual, as exalted above every
contact with the world of sense, and as being an object of
knowledge only through its works. As its opposite is posited
matter, which is conceived more ethically than physically, and becomes the source of
evil, while all good comes from the Godhead. The two principles
are expressed by unity and the indefinite duality (\gk{δυὰς ἀόριστος}).
The numbers or Ideas, according to which the world is ordered, are thought of as
pre-existing in the divine understanding, and therefore not as inherent
in things as their substance, but as archetypes for them. As
intermediaries between the transcendent Godhead and finite existence
stand subordinate gods (such as the star-spirits) and demons.

\divrule

c.\enspace \emph{Pythagoreanizing Platonists.} --- Of these,
\emph{Plutarch} (c.\ 50---125 A.D.), \emph{Maximus} (middle of the 2nd
cent.) and \emph{Numenius} (in the second half of the 2nd cent.) are
the most important. The general tendency in them is essentially the
same as in the Neopythagoreans. They seek to conceive the divine in the
highest perfection, as raised above our knowledge, and designate it as
pure mind, as the One, the existent, the good, as reason. Matter, or the
evil world-soul, which becomes as it were a personification of the
possibility of evil lying in matter, becomes the opposite of the
Godhead, \opage{239}of the good. Here too the absoluteness of the
opposition prevents a direct action of God upon the world; therefore
the star-spirits and the daemons become active intermediaries, and in
Numenius a Demiurge, who essentially becomes identical with the
world-soul in the Platonic sense, is distinguished from the supreme God
as ``second God.''
Dualism penetrates into the conception of man's essence, and Numenius
ascribes to man a double soul: a rational and an irrational one, which
lie constantly in conflict with each other. The highest task for man
becomes to reach the state in which he can passively receive the
revelations of the Godhead.
This state is reached partly through the action of the good daemons,
partly through the soul's tranquility and its separation from the
sensible. In the inner revelation the soul is suddenly illuminated,
touches as by a leap the immaterial primal being, and in this contact
receives the initiation of truth.

\divrule

d.\enspace In the \emph{other philosophical schools} dualism and the
spiritual tendency connected with it emerge especially in the
\emph{Stoics}. The ethical dualism in the Stoics between reason (virtue)
and sensuousness gradually leads to an anthropological and metaphysical
dualism. \emph{Seneca} seeks the ground of moral weakness in the soul's
connection with the body, which becomes as it were a prison for it, and
he calls the released souls happy. \emph{Epictetus} posits the body
(\gk{σῶμα πήλινον}) and reason (\gk{λόγος}) as the two constituents of
the human being, sets the necessity to which the former is subject in
opposition to the freedom of the will, and calls man a soul that
carries a corpse. He says that the soul longs to leave the body and
return to its divine origin. \emph{Marcus Aurelius} stresses still more
strongly the opposition between the mind (\gk{νοῦς καὶ δαίμων}) and the
body (\gk{γῆ καὶ λύθρος}), and from the vital spirit (\gk{πνεῦμα}) he
distinguishes the incorporeal mind (\gk{νοῦς}), which is our proper
essence.
In him as well as in Epictetus there also appear, even if the old
doctrine is not expressly aban\opage{240}doned, intimations of the
metaphysical parallel to this anthropological dualism, when Epictetus
makes God's essence consist in reason and knowledge and has the
liberated souls return to him, and when M. Aurelius has God behold all
souls without their bodily wrapping, inasmuch as his reason immediately
touches its emanations.

Connected with dualism in these later Stoics is the consciousness
\emph{of man's moral impotence}, and the proud confidence of the old
Stoicism in the power of thought and will gives way to a surrender to
the Godhead, in whom the lacking insight and the lacking strength for
the good are sought. The moral is traced back to the religious, and in
the way in which, e.g.\ M. Aurelius depicts the soul's liberation by a
withdrawal from everything external, through which the external wholly
vanishes for it and every wish and every desire is cancelled, its
surrender to the will of the Godhead, and its relation to the daemon
within it, the concept of enthusiasm and ecstasy is already touched
upon.

\divrule

e.\enspace \emph{Philo.} --- As the representative of the
Jewish-Alexandrian philosophy, which in its peculiar determinateness
develops in Alexandria simultaneously with Greek Neopythagoreanism, we
single out \emph{Philo}, in whom this form of philosophy found its
highest development, and the concluding form of Greek philosophy,
Neoplatonism, its nearest preparation.

\emph{Philo} (born 25 B.C., still living toward the middle of the 1st
cent.\ A.D.) seeks in his philosophy to express the deeper content of
the Jewish sacred scriptures, which he means to find through an
\emph{allegorical} interpretation that penetrates through the letter to
the spiritual meaning hidden behind the wording. He is convinced that
Scripture, on account of its divine origin, conceals all truth within
itself, and that the truth in philosophy derives from the Greek
thinkers' having
\opage{241}known the Holy Scripture. In placing thoughts taken from
philosophy, especially from the Platonic and Stoic, into Scripture by
allegorical interpretation, he is assured that he is not placing
anything foreign into it, but only bringing out the deeper truth that
the spiritual man, through a divine inspiration, finds beneath the
garment of the letter.

In Philo's doctrine of \emph{God} the Jewish religion's notion of a
spiritual personality raised above everything natural forms the basis;
but through the more precise determinations that are to secure the
absoluteness of the divine, its elevation above all finite
determinateness, this notion is converted into the most abstract, most
universal concept. God is determined \emph{negatively} by the removal
of every property that belongs to the finite, not only of the
properties that express its imperfection: its materiality, its
divisibility and changeability, its boundedness and dependence; but
also of those that express its perfection: virtue, knowledge, the good,
the beautiful. Indeed, God is exalted even above these determinations
taken in and for themselves, not as properties of something finite: he
is higher than the very Idea of the good. His essence is
incomprehensible to us: we know only \emph{that} he is, not \emph{what}
he is, and the names we ascribe to him hold only in an improper sense.
Only \emph{being} is fit to be predicated of him: he is \gk{τὸ ὄν}
[that which is] or \gk{ὁ ὤν} [he who is]. But the need to conceive God
as the perfect actuality and as the author and governor of the world's
existence nevertheless brings Philo to give God more positive
determinations again, namely goodness and power, by which God is posited
as the highest cause of being and of all perfection in the finite. God
is present everywhere in the world with his power, but his essence is
not in space, not in the corporeal, to which he himself has given room
and place.

As a \emph{purely spiritual} being God might not touch impure matter,
which from the beginning stands as the opposite of the divine and
according to its essence is the non-existent (\gk{μὴ ὄν}), the passive,
the formless, the disordered,
\opage{242}the source of all imperfection in existence and of evil.
The formation of the world therefore takes place not immediately
through God, but through the \emph{Logos}, the highest divine power,
which encloses within itself the individual divine powers, the
individual \gk{λόγοι} or Ideas.
The Logos is with God as the divine thought, which is one with Wisdom
and encompasses all the Ideas as the \gk{τόπος} [place] of the world of
Ideas; it is diffused through the world as the world-forming divine
reason, as the active \gk{λόγος σπερματικός} [seminal reason], which
encompasses the individual \gk{λόγοι}, the active essences of things.
The Logos thus becomes an intermediate being between God and the world;
it is subordinate to the transcendent God, who does not step out of the
unity of his essence or communicate himself with it to the world; it is
not unoriginated like God, but neither has it originated in the same
way as finite things, and in relation to us it is a God. It is the
medium of God's revelation to the world and the world's representative
and advocate with God. In the conception of the Logos and of the
individual divine powers that the Logos comprehends there is an
indeterminacy in Philo, inasmuch as now the notion of them as distinct
personal beings, who surround the highest God as ministering angels,
now the notion of them as divine properties and powers comes to the
fore. Yet the conception, especially of the Logos, as an independent
hypostasis predominates.

The dualism in existence is mirrored in the \emph{human being} in the
sharp opposition between the rational soul, akin to the Logos,
(\gk{νοῦς} and \gk{πνεῦμα}) and the body together with the sensuousness
that has its source in the material. The ethical task must then become
chiefly this: ascetically to free oneself from all sensuousness, to root
out all sensuous pleasure and all affects, to win unconditional apathy.
But this ethical task nevertheless becomes only something provisional,
and it receives conditions that point beyond the ethical. The ethical
strength and the ethical insight must be given by God to the sinful
human mind, defiled by matter, and just as the strength for action, so
too the norm for
\opage{243}action must be sought in God; true virtue becomes: to
imitate God; the wisdom that lies at the basis of virtue receives only
one object: God, and the ground of wisdom is faith. Thus the way is
pointed beyond the ethical to a religious relation, which culminates in
immediate union with the divine. The whole practical life acquires
significance only as a preparatory exercise; it ranks unconditionally
below the contemplative; but within the latter, occupation with the
encyclical sciences and with logic and physics in turn ranks below
occupation with the science that concerns man's moral condition:
ethics. Ethics leads us to self-knowledge, but in self-knowledge we
become conscious of our nothingness and ignorance and give up every
claim to knowledge, and precisely this giving up of ourselves, which is
a surrender of our will and our thought to God, leads us to the
Infinite.
We are not content with the mediate, improper knowledge of God which
science gives in considering his works; we strive away from this
uncertain shadow-image to something higher, to immediate intuition, and
despite the transcendence of the divine such a beholding becomes
possible, by which the mind goes beyond the Ideas and the Logos,
intuits God in his pure unity, as the Logos intuits him, and becomes
one with God. But this beholding is reached only through a divine
revelation. Only when man, unconditionally without resistance,
surrenders himself to the divine working, goes as it were out of
himself, casts away the very consciousness of himself, and in ecstasy,
in prophetic rapture, sinks himself into the depth of the embracing
divine, does he find in union with the Absolute the highest bliss,
indescribable in words, for which the human mind yearns and which
nothing finite can give.

\divrule

\begin{center}
\textbf{8.\enspace Neoplatonism.}
\end{center}

The striving of the dualistic systems to conceive the divine in
absolute transcendence, and precisely in such an Absolute,
\opage{244}freed from all finitude of thought and of being, to secure
for the human mind, doubting itself, the source of the highest
happiness, finds its most consistent execution in Neoplatonism; but the
dualistic distinction in principle between God and matter is given up
here, and the divine, in being seen as the source of all existence, is
posited also as the author of matter, which becomes the last member in
a quantitatively diminishing emanation of power. The metaphysical
character of Neoplatonism is therefore that of monistic spiritualism,
and it forms the conclusion of Greek philosophy in that it seizes the
last of the possibilities that post-Aristotelian philosophy had at its
disposal, and in seizing it strives beyond its limit as Greek thought
and points toward a view of life that Greek thinking could approach but
not carry through. While it appears as a renewal of the idealistic
philosophy of the past, which, in its endeavor to restore it, it
essentially transforms, it absorbs into itself the various
philosophical systems, and becomes the representative of ancient
thinking over against the new, victorious world-religion, against which
it vainly seeks to maintain the faith of antiquity, sublimated by
philosophy. In this struggle, in which Christianity is supported by the
power of the state, Neoplatonism succumbs, but only after it has been
spiritually exhausted and transformed into fantastic scholasticism.
a.\enspace \emph{History of Neoplatonism.} --- The founder of the
Neoplatonic school is \emph{Ammonius} (Saccas), who lived in Alexandria
at the beginning of the third century after Christ. He left no
writings, and we do not know how much of the developed system that we
find in Plotinus goes back to him. \emph{Plotinus} (205---270) becomes
the most important representative of Neoplatonism; his writings
(collected by his disciple Porphyry in 6 Enneads) are our most
essential source of knowledge of this tendency. Through him
Neoplatonism was transferred to Rome, where it found wide diffusion, and
where Plotinus's disciple \emph{Porphyry} (233---304) also appeared as
a teacher. After Porphy\opage{245}ry the religious element gains the
upper hand in Neoplatonism, notably through the Syrian
\emph{Iamblichus} (d.\ under Constantine), for whom philosophy
essentially becomes a means to a speculative dogmatic grounding of
polytheism. In Iamblichus's school the theoretical interest recedes
entirely behind the religious, and philosophy becomes a guide to
attaining, through purifications and cultic acts, a supernatural power
that manifests itself in magic and theurgy. Only when Neoplatonism gets
its center in Athens is the scientific striving awakened to new life,
even though it now expresses itself only in renewed study and
interpretation of the old philosophers or in formal elaboration of the
system, and even though it cannot drive out the fantastic element.
Among the representatives of Neoplatonism in Athens the
Constantinopolitan \emph{Proclus} (411---485) occupies the foremost
place by his formal-dialectical skill and by his comprehensive knowledge
of the older philosophy, which he seeks to work together into his
system. It is chiefly through his and Plotinus's writings that
Neoplatonism comes to exercise its great influence on Greek patristics,
and on the thought of the Middle Ages and of the transitional period.
After Proclus the school languishes; only in the interpretation of the
great thinkers of the past does it still achieve anything significant,
and when Justinian in 529 by a peremptory decree closed the schools of
philosophy in Athens, the death sentence on Greek philosophy was only a
sentence on what had wholly outlived itself. In its yearning to reach
beyond the sphere of thought and to find something higher than thought,
Neoplatonism had ended, in its last representatives, in imagination
void of reality, and the dialectical virtuosity which forces the
figures of imagination into artificial and empty schemata of thought
without a foundation in reality loses itself, in seeking to reach the
Absolute, in the emptiness of the indeterminate and ends in a
\gk{ὑπεράγνοια}, in which all knowledge is annihilated.

\divrule

b.\enspace \emph{Plotinus.} --- In his doctrine of the divine primal being
Plotinus carries the conception of it as absolutely
\opage{246}raised above every determination of finitude through to its final consequence.
Whereas in the dualistic systems the divine had already been raised
above the sphere of thought, and had retained only the determination of
what is, which, through the opposition to matter, became equivalent to
the determination of a spiritual existent, Plotinus raises it above the
determination of being as well. --- The absolute must be \emph{absolute
unity}, because such a unity is the presupposition of all multiplicity.
Therefore its essence cannot be expressed by the concept that for
Aristotle was the highest: by \gk{νοῦς} [mind]; for in all thinking, even in
the divine, there is a duality, a distinction between the thinking and
the thought, which mutually presuppose and condition one another; the
divine must therefore be higher than \gk{νοῦς} and \gk{τό νοητόν} [the intelligible].
Nor can it be determined by will, life and activity, because these
determinations include a relation to something other than itself.
Indeed, even in the determination of being there is a plurality, since
all being is \emph{determinate} being, which points back to the pure
unity of the absolute as to its presupposition and source. It is then
beyond being (\gk{ἐπέκεινα τῆς οὐσίας}) as it is \gk{ἐπέκεινα τῆς νοήσεως} [beyond thinking].
The divine is thus raised above all positive predicates: it must be
determined negatively as infinite, without quality, or rather: as
raised above all quality. When it is determined by apparently positive
predicates, such as the First, the One, the Good, then by the First is
designated only a formal relation to what is derived, by unity only the
exclusion of all plurality, by the Good a relation to the world, which
has being through it and for which it is the Good, not a determination
of the absolute itself, which cannot be the Good for itself because it
cannot stand in need of itself. Only determinations of relation can be
stated in positive form of the absolute as the origin of the world; and
they are then to be understood thus, that they contain a statement about
the world rather than about the primal being, which as producing power
remains in itself and is untouched by the relation to what is produced.
In this absolute absence of determination the absolute
\opage{247}nevertheless shows itself here too to be thought as something
\emph{spiritual}, both by the determinateness of the emanations
(\gk{Νοῦς}, the first emanation, is the image of the primal ground) and by
the conditions for the union of the human mind with it.

The going-forth of the forms of existence from the divine One Plotinus
thinks of as taking place through a necessary radiation of power from
its infinite fullness, without any diminution or alteration of the
primal being, which during the outflow of power remains self-enclosed,
and does not communicate its \emph{essence} to what has come to be. The
One does not divide itself; for then it would cease to be the One; it
is everywhere with its undivided power, which bears what has radiated
forth, and yet it is not in any place. It is all only insofar as all
stems from it as its effect. That a manifold can proceed from it
without the unity originally containing a plurality is grounded in the
superabundant perfection of the unity, which, as it overflows, can let
something other, less perfect, come to be, which has not been as such
in the unity.

The series of what is grounded in the unity opens with \gk{Νοῦς}, which
is the image of the unity, by its essence stands nearest to it and,
though radiated forth, is in the \emph{unity}, borne and encompassed by
its power, as in general the lower is in the higher. The first
radiated likeness becomes \gk{Νοῦς} in that, following the impulse that
lies in everything that has come to be, it turns towards the
supersensible unity in order to grasp it, and is filled by it. In
\gk{Νοῦς} the imperfection which the effect must have in comparison with
its cause shows itself in this, that plurality and otherness
(\gk{ἑτερότης}) have already penetrated into it, since the knowing and
the known form a duality, at least in concept. And this plurality
receives its more determinate form in that the \emph{world of Ideas} is
enclosed in \gk{Νοῦς}; the Ideas are its substantial organic members,
whose multiplicity is absolutely pervaded by unity, and this unity of
the Ideas is \gk{Νοῦς} itself, since the difference between the
intelligible world of the Ideas (\gk{κόσμος νοητὸς}) and \gk{Νοῦς} is only
the conceptual difference between the
\opage{248}ideal existent as resting object of knowledge and as active
knowing. \gk{Νοῦς} thinks only the unity and itself, not what is lower.
The Ideas are the true existent (\gk{οὐσία}) and the truly living
(\gk{τὸ ὃ ἔστι ζῷον}), and \gk{Νοῦς} is thus the unity of thinking and
being. In the Ideas, as manifold and mutually different, there is
already a matter (a \gk{ὕλη}), but a supersensible one (a \gk{θεῖα ὕλη} [divine matter]),
which is unconditionally formed and ensouled by the higher.

From \gk{Νοῦς} there proceeds, as its radiation, the \emph{soul}, which
stands to it in the same way as \gk{Νοῦς} to the unity, and in which the
unity is present through it. In the soul the distance from the
perfection of the primal being comes forward more strongly; it is as it
were a double being, which unites within itself an ideal and
indivisible and a divisible element, and is drawn at once towards the
two worlds, the intelligible and the sensible, yet so that it stands
nearer to the former. It turns towards \gk{Νοῦς} as its source and
through it beholds the world of Ideas and the unity itself, but at the
same time it is drawn towards the material, which proceeds from it. As
it turns towards matter in order, as an active power, to communicate
itself to it and to form it after the archetypes of the world of Ideas,
the \emph{sensible world} arises, which despite its imperfect reality
shows itself, by its beauty and by the vital harmony of the whole, to be
a \gk{Κόσμος}, a likeness of the world of Ideas. The soul as unity is the
\emph{world-soul}: the individual souls proceed from it and are
comprehended by it. According to their essence the souls are immaterial
substances, separable from the body in \emph{all} their functions. They
pervade the body as light pervades the air; they are not in the body,
but the body is in them. As active they have a presence in the body,
are wholly in its whole and wholly in each of its parts. In and for
itself the soul is indivisible; only in relation to bodies is it
divided, inasmuch as these, being spatially determined, could not
receive it undivided. The soul forms the limit of the intelligible
world; beneath it lies that of the phenomenon.

\opage{249}The last stage in existence is occupied by \emph{matter},
which has proceeded from the soul in accordance with the active power
that the soul has from its source. Matter is, in relation to the
individual sensible things, the common substrate (\gk{τὸ ὑποκείμενον}),
which receives all forms. In relation to the ideal, whose opposite it
forms, it is the formless \gk{ἄπειρον} [unlimited], the dark, the non-existent
(\gk{μὴ ὄν}), and although, as determinable by form, it becomes an
intermediate being (\gk{μέσον ἀγαθοῦ καὶ κακοῦ} [intermediate between good and evil]),
it is, regarded in and for itself, an evil (\gk{κακόν}), the source of all
disturbance and imperfection in the \gk{Κόσμος}, the source of the
soul's defilement and finitude. Thus, through the \emph{quantitative}
decrease of perfection in the radiation of power, a \emph{qualitative}
opposite to the absolute is supposed to emerge, which by its negative
working shows itself as something independent and underived, in
\emph{dualistic} opposition to that which was to be the principle of
the unity of existence. This matter, which is the basis of the sensible
world, also has only the most abstract determination in common with
what in the world of the intelligible was called \gk{ὕλη}. They meet only
in being the substrate, but the \gk{ὕλη} of the sensible is only a
shadow-image of that of the intelligible, just as the sensible shape is
of the Idea. The \gk{ὕλη} of the intelligible is \gk{οὐσία}, that of the
sensible a non-existent. It receives only a relative being through the
active powers that, as \gk{λόγοι σπερματικοὶ} [seminal reasons], enter into it and
form it.

The fundamental view of existence determines the conception of the
\emph{human being} and of the task of human life. In the state of
pre-existence human souls had their home in the world of the
intelligible and there beheld \gk{Νοῦς} and the highest itself; but as
they, in accordance with their double nature, were magically drawn
downwards towards the corporeal in order to form it and be united with
it, something lower attached itself to the higher in them, the memory
of the divine was darkened in them, and as they wished to be
independent, they sank ever deeper. The task must therefore be this: to
turn away from the sensible, to free the soul so that it can again turn
towards its di\opage{250}vine origin and after death live a life in the
world of the intelligible, without falling prey to the transmigration
of souls that awaits the souls not perfectly purified. Every human
being has the possibility of accomplishing the task; for freedom is
preserved, and the essence is not disturbed, even if it is as it were
overgrown with sensuousness and guilt. The higher in us remains, even
in life in time, in the intelligible, and only the lower soul, which
through the soul's being drawn towards the body was attached to it,
reaches down into the sensible world. The way to the goal is virtue,
whose essence can be determined at once as an acting in accordance with
our essence, as an obeying of reason, and as a striving after likeness
to God. But in virtue there are different stages. The civic (ethical)
virtues, \gk{φρόνησις} [prudence], \gk{ἀνδρία} [courage], \gk{σωφροσύνη} [temperance] and
\gk{δικαιοσύνη} [justice], lead only to the right choice among what is in and
for itself something lower; they do not yet free from dependence on the
external, to which the agent stands in relation. The purifying virtues
free from all guilt through flight from sensuousness. The deifying
virtues, finally, not only free from guilt but make us divine. Their
essence is an acting in accordance with \emph{\gk{Νοῦς}}, a constant gazing
towards \gk{τὸ νοητόν} while the mind beholds itself, an immediate
knowledge of the content of the world of Ideas. Yet there is a still
higher goal: immediate contact with what stands above \gk{Νοῦς} itself,
with the divine undetermined unity. This can be reached only in an
ecstatic elevation above thought, even above its highest form, for which
the Ideas are the object; for even in this thinking there is still a
duality and a motion. The soul has in itself a higher faculty, which is
as it were its center, the expression of the unity in it, by which it
has the possibility of union with the highest, as it is suddenly
illuminated by this in self-absorption. This union is more than a
knowledge, more than a beholding; it is a contact, a fusion with the
divine unity itself, and in this contact the soul is freed from all
multiplicity (it is a \gk{ἁπλοῦν} [something simple]) and all motion; it is
ecstatically raised above self-\opage{251}consciousness and knowledge.
In it the soul finds the highest rest and the highest bliss. But as
long as we live, this state, which is reached by as it were leaping
over the intermediate members between the Godhead and the human soul,
the divine outflowings through which the way to the unity otherwise
goes, can, even in the wisest and most virtuous, be only rare and only
of short duration. The soul is as it were frightened back by the
absence of determination in the absolute; in order to possess it, it
strives to objectify it, and thus it steps out of the unity and turns
again towards the finite. ---

In Plotinus's system the distinctively formed Neoplatonic philosophy
appears in its most original and most surveyable shape. In the series
of stages, \gk{Νοῦς}, the soul (the world-soul), the sensible world,
matter, we have the familiar concepts from the systems in which Greek
philosophy culminated; but these concepts are as it were metamorphosed,
since the forms of existence they designate are conceived as effects
of the absolute raised above all determinations of thought, which in
its transcendence remains self-enclosed while through its radiation of
power it produces and bears existence.

To Plotinus's system Plato's philosophy, which he believes he is
reproducing in its genuine shape, furnishes the essential material; but
the fundamental concepts of that philosophy are fused with Aristotle's,
% errata, p. 251: printed "dets", corrected "dens"
whose philosophy is conceived as essentially in agreement with Plato's,
and the Aristotelian doctrine of the active forms is in turn brought
close to that of Stoicism, with the removal of the latter's
materialism. Of particular significance for the conception of the
transcendence of the absolute have been the dualistic systems,
especially those of Plutarch, Numenius and Philo. The last comes nearest
to Neoplatonism; yet he still differs from it in essential points: in
that, with regard to the conception of the divine, he is still
influenced by the Jewish religion's notion of the personality of God,
and will not yet set God above the concept of being; in that he
conceives matter in a dualistic opposition to the
di\opage{252}vine; and in that the communication of the primal being to
what has come to be, which in Neoplatonism is determined only
dynamically, as a communication of power, becomes something more with
him, since the intermediate beings seem to stand in a substantial
relation of unity to the primal being. Through the carrying-through of
transcendence and through its spiritualistic monism Neoplatonism
acquires its distinctive character, and the material taken over from
its predecessors is transformed and systematized from the independently
conceived principle.

\divrule

c.\enspace \emph{The later Neoplatonists.} --- It has already been
emphasized above that with these the practical and religious tendency
comes forward more strongly, and it is essentially this tendency,
proceeding from an ever stronger feeling of human weakness and
sinfulness, that leads to the simpler fundamental form of the Plotinian
system being articulated more richly. For thereby one seeks to make room
for all the innumerable divine figures of the popular religions, the
intermediate members between man and the primal ground, which are
ordered systematically and derived speculatively-dogmatically, and at
the same time to remove the primal being, with which they are to bring
us into connection, to a still greater height above thought and
reality. In \emph{Iamblichus}, above Plotinus's unity, which was still,
in the sense indicated above, designated as the Good, there is placed a
first unity which is higher than the Good, an ``absolutely ineffable
principle''. \gk{Νοῦς} divides into an intelligible (\gk{νοητὸς}) and an
intellectual (\gk{νοερὸς}) world, each of which is again divided according
to triads. The sphere of the soul is likewise expanded into a triad,
and the forms of the psychical within the \gk{Κόσμος} comprise a countless
host of gods and demons, which are ordered fantastically according to
Pythagorean mystical numbers, especially the number three and the
number seven and their multiples. All these divine figures are posited
as supra-rational, and are thereby raised above all objections of
thought.

In \emph{Proclus} the same striving comes forward to remove the primal
being to an infinite distance from thought. It is to be exalted above
all determinations, even above the determinations unity, the good and
cause, to be absolutely
\opage{253}unknowable, ``more ineffable than silence itself''. The forms
of existence are multiplied, not, as with Iamblichus, according to an
arbitrary schematism of numbers, but according to a triadic fundamental
form, which is supposed to be determined by the very law of the
world-process, but which in the execution nevertheless shows itself to
be only a logical form, by which, without regard to the real, abstract
presuppositions are spun out into their consequences and something
traditionally given is ordered according to an often remote analogy
with the moments of the concept. When a form of existence is produced,
three moments, according to Proclus, are to be distinguished in its
relation to its cause: \gk{μονὴ} [remaining], \gk{πρόοδος} [procession] and \gk{ἐπιστροφή} [reversion].
That which proceeds from the source is at once like and unlike its
origin, which in the producing remains in itself and preserves its
distinctive character. In virtue of the likeness, what is produced
remains in unity with its origin; in virtue of the unlikeness it
separates itself from it, but only in order to turn back to it again.
This motion repeats itself in all spheres, and the threefoldness in it
determines, as the moments are hypostatized, the articulation of the
forms of existence, whose perfection decreases with the distance from
the first cause, as multiplicity asserts itself ever more strongly in
relation to unity.

Nearest to the one primal being Proclus places, as its first effect, a
\emph{plurality of unities} (\gk{ἑνάδες}), which, though standing nearer
to or further from the primal being, are all absolutely simple, exalted
above the determinations of being, of life and of reason,
incomprehensible to everything lower, and despite their multiplicity
form an undisturbed unity, since they are wholly in one another. They
differ from the primal being in being \emph{determinate} unities, and
while the primal being is without any relation to the world, they exert
an effect on the world: in them is providence. They are gods in the
highest sense of the word. After the henads there then follows, as
proceeding from them and filling the place that \gk{Νοῦς} occupied in
Plotinus, the triad of the intelligible (\gk{τὸ νοητὸν}), the at once
intelligible and intellectual (\gk{τὸ νοητὸν ἅμα καὶ
νοερὸν}) and the intellectual (\gk{τὸ νοερόν}), in which triad the
\opage{254}unlimited becomes an element. The 3 members correspond to the
concepts of being, life and thinking. The first two members are each
again articulated into 3 triads, the third into 7 hebdomads, as the
first two members of its triad are expanded into new triads, and each
member of the hebdomad thus arising is once more expanded into
hebdomads. Room is thereby made for the incorporation of a host of
deities, which by allegorical interpretation are brought into
connection with the conceptual schemes of the triads and hebdomads. The
sphere of the psychical, too, is strongly articulated in order to make
room for a hierarchy of divine, demonic and human souls. For these last,
union with the primal being is to become possible through that faculty
of the soul which surpasses thought, through the \emph{supra-rational}
unity in it, which is pervaded by the divine power.

In the Neoplatonists who close the school, such as \emph{Damascius},
this supra-rational in us is emphasized more and more at the expense of
rational thought, and as the primal being in its infinite transcendence
withdraws from every determination of thought, and yet is to be held
fast by thought, it receives only determinations that are revoked in the
same moment, since even the negative statements, as opposites of the
positive ones, have a determinateness that makes them inapplicable.
Complaint is therefore made of the inadequacy of human language to
express the infinite, and thought, constantly destroying its own
determinations skeptically, loses itself in the infinite emptiness, to
which only imagination, in the intoxication of ecstasy, can give a
deceptive fullness. But that which this emptiness conceals is posited
as the highest reality, towards which the mind yearns; as thought
despairs of grasping it, it despairs of itself, and destroys itself as
philosophical thought. As it would go beyond itself, it sinks below
itself and becomes imagination without reality.

\divrule

\opage{255}
\begin{center}
{\Large\textbf{Chapter Five.}}
\end{center}

\divrule

\begin{center}
{\Large The Abiding Significance of Greek Philosophy; Its\\
Relation to the Thought of Later Ages.}
\end{center}

\divrule

\markboth{The Abiding Significance of Greek Philosophy}{The Abiding Significance of Greek Philosophy}
As Greek philosophy exhausts its possibilities and is dissolved as this
determinate historical main form of philosophy, it leaves behind an
\emph{abiding} result, which is taken up into the further development of
philosophical knowledge, and in that development keeps its life. In
that Greek thinking, as free thought, bound only by its own laws, has,
starting from the contemplation of nature and of the essence of the
human spirit, realized philosophy as a universal science, articulated
according to its essential main parts, and in that it has bent thought
back upon itself, to the knowledge of its own activity, it has acquired
an inalienable possession for the self-conscious human spirit and given
the abiding foundation for the realization of the theoretical aims of
spirit. And within each of the main parts of philosophy that Greek
thinking has formed, it has acquired for the future a fund of
fundamental concepts, indicated essential fundamental views, and posed
or prepared the decisive problems. \emph{Logic}, essentially as formal
logic, was developed by the Greeks in the form that it has kept right
down to Kant, and in the doctrine of categories the first incentive was
given to seek the fundamental determinations of thinking. In the domain
of the \emph{metaphysical} the concept of the Idea has been found: the
concept that expresses the primacy of the spiritual, and asserts an
eternal existent as governing what becomes. With keen dialectic the
fundamental thoughts of being and becoming, of Idea and matter (of form
and matter), of the active and the passive have been clarified, and a
series of the most decisive problems brought forward. In the domain of
the \emph{physical} the view of an organic whole of nature has been won,
and through the concept of the teleological and
\opage{256}of entelechy spirit is expressed as the immanent ensouling and
active principle in the continuous succession of natural forms. In the
domain of the \emph{psychological} the soul has been conceived as the
entelechy of the body, whereby the unity of the human being in its
various forms of development has been pointed out, and through the
theory of knowledge of Plato and Aristotle the self-active independence
of the mind in relation to the object of knowledge has been asserted.
In the domain of the \emph{ethical}, despite the difference from modern
ethics that comes strongly forward here, preparatory work has been done
by the development of the fundamental ethical concepts, by the
demonstration of the conditions of the ethical and of the ideal aim of
action, by the conception of the individual's connection with the
objective forms of human life, and, in the period after Aristotle, by
leading man back to the infinity of self-consciousness, and through it
to an ethical universalism, to a recognition of the humane as such. And
still more deeply than by Greek philosophy proper, the concept of modern
ethics of something unconditionally obligating, and of the free and
conscious self-determination of the ethical subject, was prepared by
Socrates. Finally, Greek philosophy, in Plato, Aristotle and the
Neoplatonists, has made essential contributions to the doctrine of
\emph{the beautiful}, and in the doctrine of \emph{the divine} the
determinations of Plato and Aristotle touch on the conception of God as
the infinite subjectivity and, ethically, as the one who is good.

In the historical development of philosophy the significance of Greek
philosophy has shown itself in this, that all the \emph{thought}-content
on which Christian antiquity and the Middle Ages and the transitional
period nourished themselves was drawn from it. And even if modern
philosophy, as it broke its own path, turned away from the thinking of
antiquity, it has nevertheless, in its unfolding, constantly returned to
it in order to be strengthened and enriched. For all times Greek
philosophy will stand as a model by the self-conscious freedom of its
thought, by its clarity, its plastic determinateness and honest
consistency, by its thinkers' unconditional devotion to objectivity and
to the laws of thought, and by the
\opage{257}confidence in the power of spirit that is the distinguishing
mark of the most eminent among them. And as Greek philosophy lies before
us, completed, it presents the picture of a lawful, undisturbed
development, governed by one fundamental thought, and led by its inner
law to the point where every possibility laid down in the original
disposition has been unfolded. Nowhere does the law of the historical
development of thinking come forward more clearly and purely.

The limitation of Greek philosophy has been shown by the critique of
the systems through which it attained its highest unfolding, and by this
critique it has become clear what Greek philosophy has realized of the
moments lying in the concept of philosophy (cf.\ the Introduction), and
what it has left as a task for the thinking of the subsequent age. It
has realized the determination of universality insofar as it has, as
free thought, encompassed reality, determined its highest principle and
given it as a whole the form of ideality, and it has led thought to make
itself, its essence and its fundamental relations its object. But,
bound by its inner limit, it has left to posterity the unfolding of the
world of thought from the principle, as demanded by the concept of
philosophy, through a dialectic that not only distinguishes the
opposites but unites them in a living unity, and the development of
that method of knowledge, the method of experience, which relates
itself to the real recognized in its self-validity. And in the
determination of the relation of the fundamental opposites it has not,
while it asserted the primacy of the ideal, been able to reach the
conception of the concept of \emph{spirit} as the ideal-real absolute
subjectivity, and by this concept to secure the unity of existence
against the dissolution into dualism that befell the Idea-determined
immediate unity of the world-artwork. It was not capable of this,
because its absolute, as the highest existent, even when it was
conceived as \gk{νόησις τῆς νοήσεως} [thinking of thinking], excluded negation,
opposition, from itself, and therefore could not, through negation,
close together with itself in the free infinity of absolute
subjectivity, but stood self-enclosed in its resting being, in a
\opage{258}separation proper to intuition, from the existence whose
life-principle it was posited to be, and from the human mind, which in
its highest form of reality was determined as one in essence with it.

These tasks, which Greek philosophy leaves unsolved, the presentation
of the subsequent main forms of philosophy will show to be, as it were,
divided among those forms as their \emph{essential} tasks.

\divrule

\begin{center}
{\Large\textbf{\emph{Appendix.}}}
\end{center}

\divrule

\begin{center}
{\Large Philosophy in Christian Antiquity and the Middle Ages.}
\end{center}

\divrule

\markboth{Philosophy in Christian Antiquity and the Middle Ages}{Philosophy in Christian Antiquity and the Middle Ages}
In the survey of the main historical forms of philosophy (p. 7) it was indicated how the Christian world-view and view of life, through its abstract, natureless spiritualism, made impossible the development of a Christian philosophy, that is of a \emph{whole of thought} independently unfolded from the principle of Christianity. We therefore do not speak of a Christian philosophy in antiquity and the Middle Ages, which, as a main historical form of philosophy, should stand beside Greek philosophy and the philosophy of modern times; but we speak of philosophy in Christian antiquity and the Middle Ages, since in this period we have a movement of thought influenced by Christianity, which, in calling forth the deepest movement in man's religious life, had also, in virtue of the unity of the human being, to set the life of thought in motion, however much the religious life it awakened turned away from thought. But this movement of thought fetched its forms of thought and its content of thought from the thinking of antiquity, and in the doctrinal edifices of Christian antiquity and the Middle Ages the philosophical element is a loan from the past, the properly philosophical movement only an echo of that of antiquity. The Christian natureless spirit lacked \opage{259}the capacity for an independent production of thought, and consistently the Christian consciousness could not even will to seek a real unfolding of thought, because its negative relation to the natural includes a negative relation to the natural capacities and powers of the human being, and so also to the thought that is borne by the natural and accords with nature, which it had to regard as lacking the possibility of the higher knowledge, which is that of faith, precisely because this thought must fall under the determinations of everything natural: nothingness and sinfulness. This negative relation to thought and to what is produced by thought had, however, since Christianity was to exist in the midst of the alien world, and in the struggle and the defense was not able to produce an independent content of thought out of its abstract negativity, to be transformed into a forced, factual recognition and adoption of what it denied in principle; and through the combination of the historically developed ancient forms of thought with the Christian there arose the phenomenon that has been designated as Christian philosophy.

That in this the philosophical element is a loan from the past is shown when we consider the movement of thought in the two divisions of time into which what has been called Christian philosophy naturally divides itself: the time of the Church Fathers and that of the Scholastics.

\emph{The time of the Church Fathers} is the time of the formation of dogma. In it the content of the Christian faith is unfolded into a circle of dogmas to which no essential member needed afterwards to be added; and in the course of polemic, first against what stood outside Christianity, then against what within Christianity was branded as heretical, the Fathers of the Church lay the foundation of a theological doctrinal edifice, and a systematizing of the circle of dogmas is begun on the basis of the main dogmas once fixed. The forms of thought used for the formation of the doctrine were taken, in the earliest time from the Stoics and the Alexandrian philosophy of religion, later from the renewed Platonism, whose spiritualistic tendency gave it so many points of contact with Christianity, and which towards the end of the time of the Church Fathers \opage{260}exercises an altogether overwhelming influence on Christian theology. Since the ancient forms of thought were employed, a difficulty had to arise in bringing the doctrines peculiar to Christianity, such as the doctrine of the God-man, of creation and of evil, into the alien forms without distortion. For the immediate religious consciousness of the believing Christian, for religious feeling and notion, the difference between the pagan and the Christian stood out immediately, and a distinction was drawn, with greater or less determinateness, between the \emph{content} of Christian faith and that of pagan knowledge. Yet the connection of the ancient Church with antiquity through nationality and language was so strong, and the men who, not only in the Eastern and Greek but also in the Western Church, represented the thought of Christian antiquity were, by their culture, by all the presuppositions of their mental life, so inwardly bound to the ancient, that a clear knowledge of the \emph{principle} of the opposition did not become possible, and in consequence clarity as regards the boundary between the domains had to remain incomplete. Essentially, one arrived only at setting a difference of degree between the more complete truth of Christianity and the limitation of philosophical knowledge, or at stating a difference in the practical motive and the practical consequences, not at stating the principle of two heterogeneous kinds of knowledge; and since, while urging the limitation of philosophical knowledge, one distinguished between philosophy as such and the individual philosophical systems, there lay in this a recognition of thinking, which had its consequences in the conception of the gnosis that is to develop out of faith. But since the difference of principle is not clarified, and the boundaries between the domains of faith and of thought therefore cannot be sharply maintained, the heterogeneous is involuntarily mixed. Even the most peculiar Christian doctrines, by being brought into the forms of ancient thought, were subjected to an alteration by which the peculiarity of the content of faith was impaired. The Christian \opage{261}concept of God was now led back to the Aristotelian determination of \gk{Νοῦς τέλειος} [perfect Mind], now the conception of the later philosophy had its effect in the notion of the corporeality of the divine. In the doctrine of creation the Greek view of an eternal world of Ideas and an eternal matter, or the Greek-Oriental view of an emanation of power, frequently asserted itself.
% errata, p. 261: printed "om", corrected "af"
In the doctrine of the God-man the ancient concept of the Logos forced its way in, as it had been developed by the Stoics and Philo, and especially in the Apologists we see the concepts of these thinkers applied immediately to the specifically Christian and transforming it into a speculative theorem about the divine reason immanent in the world. And even in the dogma of sin, despite the decisive divergence of the Christian religious consciousness from the ancient, the ancient metaphysical conception of evil as a non-being, as limitation and imperfection, has its effect, even in an Augustine.

During this unconscious metamorphosis of the Christian content of faith, the conscious relation to philosophy nevertheless remains this: that philosophy is posited as serving and subordinate to the supra-rational, revealed religious content. Taken as a whole, therefore, the doctrinal edifices of the Church Fathers, which essentially confine themselves to the religious problems and are poor in physical and ethical inquiries, are theological, not philosophical. There is not in them an unfolding of a total rational world-view and view of life, or a rational grounding of dogma. The world-view and view of life rest on suprarational presuppositions, and the rational grounding of dogma becomes illusory, since it is not the peculiarly Christian, the dogma in its specific form, that is maintained as a necessity of thought, but only the abstract metaphysical parallel to the dogma, such as the determination of the divine as the Absolute, as One and as Spirit, or the elements of thought taken over from the ancient, by which the Christian has been metamorphosed. Where the peculiarly Christian comes forward unaltered, it stands before thought either as something not comprehended, or as something formally derived from what is \opage{262}not comprehended. This stands out especially clearly from the time when the Catholic Church has given the fundamental dogmas their fixed form, and thinking now relates itself to this form in its striving after rational grounding. Earlier, while the central points of the religious view of life still had the mobility of life, they afforded easier points of attachment for that illusory grounding in thought, since the ancient content, to which a development of thought could attach itself, more easily slipped in under the notion of faith. ---

When in the \emph{Christian Middle Ages} a philosophical movement begins among the peoples of the West in Scholasticism, the dependence on antiquity that was grounded in the immediate connection of nationality has ceased, and new, vigorous races had replaced the nations of antiquity as bearers of the historical development. Among these new peoples the philosophical endeavor acquires a greater energy, and when the effort is made to give a Christian doctrinal edifice that comprehends the totality of faith and knowledge, and joins the content of faith and the content of thought into a whole, it is done with the clearer consciousness of the difference of principle that is the condition for a reconciliation, a real unity, being brought about. And yet here too we see the elements separate, the peculiarly Christian withdraw from knowledge, and the attempt to form a whole of thought miscarry.

If we designated the time of the Church Fathers as that of the formation of dogma, we may designate that of Scholasticism as the time of the working-up of dogma, of the formal appropriation of dogma. In the former division of time the Christian spirit had unfolded dogma into an essentially completed whole of objective truths; in the latter it has the unfolding of dogma behind it as a stage passed through, and the task now becomes this: to appropriate the objective truth for consciousness. As the means to this, philosophy is employed, which in this period too is taken over from antiquity, especially the philosophy of Plato and of Aristotle. Its forms are applied to the content of the Christian faith, which lay before it in the doctrine of the Church as an object of unshakable \opage{263}authority, and Scholasticism seeks, by the path of the reflecting and syllogistically demonstrating understanding, to justify the dogmas of the ecclesiastical system to the thinking consciousness, so far as this could be attempted at all from this standpoint. For the task could not, since the truth of the Church's faith stood as an unshakable presupposition, become this: to place itself above the object in order to let thinking freely produce its content; but only to place itself over against it in order to find the right interpretation of dogma, in order to understand the connection of the dogmas, in order to let the analyzing understanding give itself an account, not of the ratio of the whole, but of the rationes of the individual dogmas, and in order, with regard to the whole, to maintain in negative form the content of faith as the not-impossible.

The beginning of Scholasticism must be placed where the attempt at such a rationalizing of dogma was made with a definite consciousness of the task, and with a clear result of this striving. This first happened with \emph{Anselm of Canterbury} (1033---1109), who may therefore be regarded as the founder of this scholastic philosophizing. From him onwards the development of Scholasticism divides into 3 periods. The first seeks to determine the relation between faith and knowledge; it treats the main dogmas dialectically, and in comprehending the material of the Church's doctrine it seeks, starting from those dogmas, to gain a survey of it. The second period (from the 13th cent.\ to the beginning of the 14th) is the culminating time of Scholasticism. In it the whole circle of the dogmas is brought into a richly articulated system, and philosophy and theology come in it into as close a connection as is at all possible. In the third period this connection begins to loosen, at the same time as the systematizing power is lost. And as knowledge and faith fall apart from one another, through nominalism, which gives the universal an existence only in our notion, thinking and being fall apart from one another, and Scholasticism dissolves.

If we wish to form a clear notion of the nature of the combination of philosophy and dogma that we have in \opage{264}Scholasticism, this can be done partly by examining the successive development of the determination of the relation between the principles of the two spheres, knowledge and faith, partly by showing, through a consideration of the great scholastic systems, the relation between the heterogeneous components of these systems.

The determination of the relation between knowledge and faith shows a development that is essentially conditioned by the very endeavor to bring about a union, since this striving necessarily had to clarify consciousness about the relation.

In the important thinker who forms a transitional link between the time of the Church Fathers and Scholasticism, in \emph{Johannes Scotus Erigena} (born c.\ 810, died c.\ 877), the question of the relation of faith and knowledge had been answered by a postulate of the harmony between the activities of mind that spring from the same source, faith and reason, a harmony in which reason was to be the overreaching and primitive element. With the same confidence in knowledge that was peculiar to the old ecclesiastical gnosis and to the systems, strongly influenced by Neoplatonism, of the closing period of patristics, Scotus Erigena posits true philosophy as equal to true religion. Reason and authority both have their origin in the divine wisdom. But reason is what is given with creation itself; authority has proceeded from reason, and therefore reason has a priority over authority. In the time of Scholasticism proper, on the other hand, the development goes in the direction of subordinating the knowledge of reason to the knowledge of faith; and as the difference posited by the subordination is sharpened more and more by being led back to the incongruence between the natural and the supernatural (natural thinking and supernatural revelation), one arrives at last at an abstract opposition between faith and knowledge, between the truth of religion and that of science, so that the difference, at first merely quantitative, not only becomes a qualitative one but grows into an irreconcilable opposition; and as this development approaches its conclusion, one is no longer content to determine the relation between \opage{265}faith and knowledge as a relation between two forms of knowledge which in reality have only the name in common, but one leads them back to different functions of mind, faith being supposed to have its origin in us from an act of will; and one then lets theology, which rests on faith and whose aim is in the first instance a practical one, blessedness, become not a speculative doctrine but a practical doctrine, which is to grip and determine our will.

In the first period of Scholasticism \emph{Anselm}, with his credo ut intelligam [I believe in order that I may understand], posits faith as starting point and norm; but he demands a knowledge that takes the content of faith up into itself, and he ascribes to comprehended faith a higher worth than to uncomprehended faith. And \emph{Abelard}, with his intelligo ut credam [I understand in order that I may believe], gives knowledge a still greater significance as grounding faith, while \emph{Berengar} maintains reason as the higher than authority. In the great representatives of Scholasticism in its culminating time, \emph{Thomas Aquinas} and \emph{Duns Scotus}, a higher knowledge of faith is distinguished from the knowledge of reason not enlightened by grace; but Thomas maintains, alongside the necessity of faith, its agreement with reason, since reason too is a revelation of God; and Duns Scotus maintains, by his doctrine of the ``proportionality'' between opposites (e.g.\ between the acting and the suffering, form and matter), that what is supernaturally given is received naturally through self-active cooperation on the part of the rational subject, which possesses an infinite capacity for receptivity; yet he concedes that we cannot attain to comprehending faith, which he lets have its human origin in an \emph{act of will}. In the third period of Scholasticism, finally, in consequence of the absolute separation made between the supernatural and the natural, the difference between the higher knowledge of theology, which is given through the faith produced by a spiritual creation, and natural knowledge is sharpened to such a degree that secular science, which does not reach beyond the sphere of experience, is not even to be able to give a pointer to the divine truth, a preparation for our higher end. It is now declared, e.g.\ \opage{266}by \emph{William Occam}, that something can be true in theology that is false in philosophy, and conversely; and herewith Scholasticism has itself pronounced its judgment on the union that was its essential task, and on which it has labored with mighty powers and an enormous energy; it has itself separated the two elements and the two principles that it attempted to twine together.

To the same result as through the consideration of the relation of the principles we are led by the consideration of the most thoroughly worked-out scholastic systems, as they were formed in the flowering time of Scholasticism by men such as \emph{Albertus Magnus} (1193---1280), \emph{Thomas Aquinas} (1227---1274) and \emph{Duns Scotus} (\dag{} 1308). In these systems the philosophy of antiquity, especially that of Aristotle, is employed to the greatest extent, and the attempt to fuse it with Christianity is made with the greatest energy, and with an acuteness and a formal virtuosity of thought that is admirable. But even these powerful minds and incomparable master-builders of thought do not succeed in bringing the heterogeneous to unity, and they themselves pronounce the admission of this.

The task for these thinkers was to set up a Christian doctrinal edifice which in its systematic unity comprehended the totality of faith and knowledge. We now see the solution of this task attempted by them in the form that the ancient world-view, which was gathered up in the notion of the harmoniously ordered divine universe comprehending the whole of the forms, is, as it were, forced into the setting of the Christian. The intuition of the reality of the world in its fullness, which had to be taken from the philosophy for which the world was a reality, and which for antiquity was the whole, now becomes, joined with the Christian content of faith, a part, inserted into a larger whole whose beginning and end, which determine the character of the whole, belong to Christianity. Before the natural universe is set the nature-free, almighty, creating, triune God; after it the Christian heaven, the super\opage{267}natural kingdom of grace. But the beginning and the end were, by Scholasticism's own admission, something inaccessible to thought, something incomprehensible, hence in essence heterogeneous with what filled the interval between them. As regards creation, the production of the material world by the natureless, purely spiritual God, hence by a God in whose essence there is no point of attachment for the natural, is something that contradicts thought. What is posited as an effect of God is what is absolutely heterogeneous with his essence, and it is what, as opposed to his essence, would have to limit it, while, as the opposite of the spiritual, the only true, it would be something untrue whose existence was self-contradictory. An appeal to the divine omnipotence as the ground of explanation of creation is a help only for faith, but not for thought. In Scholasticism, as in our own time in the so-called speculative dogmatics, an attempt could be made to veil the incomprehensibility by attaching the content of nature more closely to God, e.g.\ by placing the ideas of nature in God's understanding, or by conceiving nature as having come to be through the communicability of God's essence in different degrees; but these hints, which would lead back to the view of antiquity, or at any rate away from the Christian view, since the notion of a communication of essence, which includes a necessary and eternal relation, stands in the sharpest opposition to the notion of an arbitrary creation in time and out of nothing, are not carried through. And Albertus Magnus expressly lets matter be created by God through an \emph{incomprehensible} act of will, just as Duns Scotus conceives the world as the product of the absolute, \emph{arbitrary} divine will; but by the determination of creation as something arbitrary its comprehensibility is expressly denied. And Thomas Aquinas denies the comprehensibility of God's very essence, because he is the absolutely infinite and the absolutely simple, while the dogma of the triune God is, from the culminating time of Scholasticism on, declared by all the Scholastics to be inaccessible to philosophy. But since the starting point of the system is thus posited as inaccessible to \opage{268}thought, it is thereby settled that from this presupposition no system could be developed that, as a whole of thought, bore the stamp of necessity determined by reason.

But what holds of the starting point of the systems holds also of their end point. The otherworldly kingdom of blessedness is the consummation of the kingdom of grace. But the working of grace is not a continuous unfolding of the natural capacities, but a miracle, by which these are in a supernatural way re-created, transformed into something new that is not even prepared for by natural development. Every working of grace therefore falls outside the sphere of thought, and the otherworldly heaven, in which the miracle of grace reveals itself, is posited as the absolutely incomprehensible. Thus, then, the beginning and the end of the systems, as objects of the ``higher knowledge'' of faith, which has only the name in common with the knowledge that belongs to knowing proper, withdraw from knowledge and show themselves heterogeneous with the material that was taken over from the world-view of antiquity, and the bringing about of a \emph{whole of thought} out of the different elements shows itself to be impossible.

There could then still only be a question of a \emph{homogeneous whole}, which might have been formed if what was taken over from antiquity had been so reshaped, and through the reshaping so \emph{assimilated} by the Christian, that the unity of a spiritualistic view of life, determined by the principle of the Christian that is inaccessible to thought, had been brought about. We indicated above that the character of the whole had to be determined by the beginning and the end; and when we consider the scholastic systems more closely, it also turns out, as indeed lies in the nature of the case, that what forms their beginning and end does not stand passive and indifferent beside what is enclosed by it, but that this is reshaped through its relation to what encloses it. Since God is posited as creator, the universe, which is transformed from an eternal cosmos into something created, is posited as a product of will, i.e.\ as that which lacks necessity and has no essentiality, as that which in itself is a nothing and is borne up only by the divine \opage{269}omnipotence. In other words: with creation, which is the first all-comprehending miracle, the miracle is posited as the point of view for the conception of nature. At every moment and at every point the miracle then breaks through the connection of nature and the law of nature, and thereby cancels law as such, and the very persistence, unaltered for a time, of the natural is only its preservation through the repetition of the miracle of creation. That which alone acts in nature is, for the eye of faith, the divine omnipotence; the supernatural is, as Duns Scotus expresses it, the concept and ground of nature, nature's proper essence. And in this nature, essenceless in and for itself, the Christian spirit has no rest; it is in it only as in something alien, which, when it acts as nature, is seen as evil, and which constantly, merely as existing in its nothingness, is felt as a hampering barrier and is wished annihilated. The direction of the spirit follows the supernatural working that proceeds from the almighty will, in order, in this existence, as miraculous working to break through the barriers of nature, as working of grace to break through the natural conditions of human life, and by the former as by the latter to anticipate the consummation that is accomplished when the rays of the supernatural gather into one light, in the kingdom of blessedness. In the notion of this kingdom, nature, whose nothingness was in the past posited for the whole through creation, and in the present is posited for the individual through the miracle, is in a new form posited as that which is a nothing for the spirit, and which therefore one day, when history is concluded and God's counsel accomplished, shall again cease, become nothing because it is from nothing. In the otherworldly existence all the barriers of nature, which are felt only as hampering bonds, are annihilated, and the natureless spirit lives, in the enjoying contemplation of the purely spiritual divine, its infinite blessed life.

Thus the ancient world-view, and the content of reality that was taken from it, is transformed at its ground by subordination to the Christian. The whole conception of the essence of nature, the whole theoretical and practical re\opage{270}lation of the spirit to it, is altered in principle. Nature, as created, has lost its essentiality; its laws are dissolved by the miracle; it can no longer be an object for theory, for the spirit seeks in it only the supernatural; it cannot be the foundation of the moral life, for the spirit flees nature, which, when it asserts itself as nature, is seen as evil; it does not, as something valid in itself, limit the arbitrariness of the spirit, for through imagination the abstract, natureless subjectivity gives itself an infinity for which no barrier exists, and the spirit rules, to all appearance, with unlimited power over impotent nature. And yet, despite this metamorphosis, there shows itself, and must show itself, a resistance on the part of the original ancient element, through which the apparently achieved unity becomes insecure and inconsistencies force their way in.

That a resistance and a reaction of the ancient element could not fail to appear lies already in this: that for man, who is a mental being only in being at the same time a natural being, the supranatural conception of existence always requires a naturalistic starting point; that the supernatural is supernatural only by rising above a natural that is presupposed, not explained by it; that man can live the purely spiritual life in its abstraction only by reaching it at every moment through turning away from the worldly, so that, where a spiritualistic standpoint is sought to be carried through, an independence must nevertheless indirectly be ascribed to the natural, as the opposite and condition of the supernatural. We accordingly find also in the Christian Middle Ages that the natural foundation by which the ancient view of life and world-view was borne asserts itself not only in the common consciousness but also in the philosophical. In the former it asserts itself in the demonic power that nature, denied and feared, continues, as an essential moment in the human, to exercise over consciousness. In the philosophical consciousness, in Scholasticism, its influence, through which the ancient world-view asserts itself, appears in the doctrine of the realization of form in the connection of nature, \opage{271}of the differences of degree in nature, and of the significance of the natural for the development of human life in intellectual and ethical respects. In all these determinations an essentiality is conceded to the natural precisely through the concepts of development and stage of development; for only within the essential can a difference of degree and an unfolding be given, not in the essenceless; only a nature valid in its essence can acquire significance for the development of the spiritual. The denied natural thus forces its way forward unconsciously; only through it does thought gain a foothold and the world-view a content; therefore it is unavoidable that it asserts itself. And even if its significance is cut off again, since form is to be seen as realized in an existence posited by a creation, whose continuity and determinateness of essence are cancelled by the miracle, and in which God's omnipotence is what immediately acts and is essential; and since the whole natural life, including the whole natural moral conduct, the ethical virtues, which were determined as the necessary presupposition of the theological, is posited only as a means for something supernatural, which does not realize itself through natural development, is not present in it, but enters from without, by a new miracle of creation, as a higher order complete in itself, as an alien eternal world over against the temporal; --- then it loses its significance only in order the next moment to wrest it back again. For only in the denial of the natural, in the flight from it, is the supranatural; and the denial of it and the flight from it, which presuppose that it was posited as a reality, exist only in that it is constantly posited anew as such a reality, certainly in order to be denied again, pushed away again, but in order then nevertheless, in an eternal cycle, to come back as what is uncomprehended in its origin, and therefore independent in its being. ---

The final result with regard to the scholastic systems is then this: that they seek in vain to unite heterogeneous principles and heterogeneous elements. A whole of thought they could not form; for everything specifically Christian is, as the content of supernatural revelation, expressly posited as in\opage{272}congruent with natural thinking. Only in that element of the ancient which asserts itself against the metamorphosis does thought have a temporary point of attachment; but this ancient element stands in the theological system as a foreign fragment; from the whole in which it finds a place through an unavoidable inconsistency, comprehending thought is excluded. The principle of the whole is the miracle, but the miracle is the annihilation of thought. Of a ``Christian philosophy'' there can thus be no question here, but only of a theological doctrinal edifice, which is nevertheless prevented, by the foreign element it cannot assimilate, from forming itself into a homogeneous whole. ---

It is true that in Scholasticism, since thought was taken into the service of the Church's doctrine, an enormous virtuosity of reflection developed; but the unfree thinking of Scholasticism had to become formalistic and devoid of reality, and unfit to unfold an independent philosophical world-view. Scholasticism, as we have seen, took its forms of thought, just as fully as patristics did, from the philosophy of antiquity. But since it took over forms from a view of life that was alien to it, it had involuntarily to separate these forms from the ancient view of life and take them up merely \emph{as} forms, in which it then moved by a reflection of the understanding that removed itself more and more from reality and its fullness of content. And since they were to be applied to the specifically Christian content, which stood above thought as an unshakable authority and proclaimed itself to be of a higher nature than thought, the nature of this content too had to prevent Scholasticism, in its use of them, from getting further than a formal thinking that had no hold in self-consciousness and in nature. The specifically Christian content is not dissolved in thought; for however long consequences of the incomprehensible presuppositions are spun out by means of the forms taken over, this still does not become a content of thought. Since the forms of thought become external and alien forms for the object of faith, the object cannot through them be appropriated by self-consciousness; at the most there can be brought about by them \opage{273}a sort of external connection between consciousness and the object, not a union of them in a living unity. The understanding, which reflectively uses the forms, which distinguishes, defines and proves syllogistically from commonly accepted presuppositions, does not comprehend; for only that is comprehended which has an attachment in self-consciousness and is assimilated by thought. Of a positive grounding of the content of faith by comprehending thought there is no question; it is only to be maintained negatively as the not-impossible.

Scholasticism is dissolved from within, as its heterogeneous elements separate from one another, each with its own truth, and theology itself leaves to philosophy the sphere of experience and of the natural. The thinking mind feels itself unsatisfied by its unfreedom (which remains in principle, even though the differences among the authorities and the unsettled shape of various dogmas give thought, as it were, some latitude), by its lack of a content of reality, by its alienness to self-consciousness. The religious spirit feels itself unsatisfied by its formalism of the understanding. And while in the religious domain mysticism gains strength through this feeling, there develops in the philosophical domain an independent philosophy, determined in its direction by the deficiencies of Scholasticism.

But the inner process of dissolution in Scholasticism and the development of a new thinking are essentially hastened by many influences coming from outside. Of essential significance in particular was the effect of the rise of the natural sciences, of the founding of new, freer universities, of the assertion of the nationalities, of the development of the vernacular literatures, of the unfolding of the forms of a secular moral life, of the beginning flowering of art, and finally of the revival of antiquity, which brought what was fermenting in the age to a powerful breakthrough. Through all the phenomena of the new movement of the spirit, which come forward at the same time as Scholasticism loses its support through the decline of the hierarchy during the schism of the papacy, there runs something common: the craving for freedom and the craving for reality, for \opage{274}an ideal fullness in the present. This craving led the spirit away from the authority of the Church and from the natureless transcendent spiritual; it led towards nature as the real content of thought, as the rational existing in the form of objective reality, as the foundation for a true and full concrete unfolding of human life; and it led towards self-consciousness as a source of truth. The mind, which in Scholasticism had bowed unfreely beneath the uncomprehended authority of dogma, which had busied itself in an external and formal manner with borrowed forms of thought, now seeks within itself, as self-conscious mind, to find the laws of thought, in order to be able to know itself as free, the first condition of all philosophy, and in order to be able to find itself again in what is thought as determined by reason. --- That the general direction of the age had to be essentially strengthened by the revival of acquaintance with classical antiquity and by the enthusiastic study of it lay in this: that this study led back to the intuition of an ideal human life, borne by nature and rich through its harmony with it. In this life the liberated self-consciousness found itself again and found a satisfaction for its natural craving. --- Through its altered direction the spirit had regained the condition for the development of a real philosophy (cf.\ p. 8); with the preponderance of this direction in thinking the dominion of Scholasticism ceases; with its carrying through in principle and consistently begins the independently formed \emph{modern philosophy}.

\dblrule


\chapter*{Corrections.}
\markboth{Corrections}{Corrections}

\opage{275}\relax
[The leaf of \emph{Rettelser} bound at the end of this part lists five
corrections (pp.\ 15, 63, 144 twice, 212); a further thirty for this
part, headed \emph{1ste Deel}, are printed on the \emph{Trykfeil og
Rettelser} leaf of the second part. Since they correct the Danish
wording, they are not translated here; both lists are transcribed in the
Danish editions. The English text above follows all of them at their
sites, where the Danish edition keeps the printed reading.]

\end{document}
